% Generated mechanically from frozen input inputs/p11/ja/stacks-morphisms.tex.
% Do not edit this generated file; edit the bound locale source instead.

% OK, start here.
%

\setcounter{chapter}{100}
\chapter{代数スタックの射}



\phantomsection
\label{stacks-morphisms-section-phantom}


\section{序論}
\label{stacks-morphisms-section-introduction}

\noindent
本章では、代数スタックの射のいくつかの型を導入する。
表現可能な対角射をもつ準分離代数スタックの場合の参考文献は
\cite{LM-B} である。

\medskip\noindent
目的は、代数空間の射の各型の定義を代数スタックの射へ拡張することである。
各場合はわずかに異なるため、すべてを個別に扱うのが最善と思われる。

\medskip\noindent
代数空間によって表現可能な代数スタックの射については、既に多くの型を
定義した。次を参照せよ：
Properties of Stacks,
Section \ref{stacks-properties-section-properties-morphisms}.
本章の対応する各場合について、一般の場合の定義がそこで与えた定義と
両立することを確認しなければならない。




\section{規約と用語の濫用}
\label{stacks-morphisms-section-conventions}

\noindent
引き続き、次で導入した規約と用語の濫用を用いる：
Properties of Stacks, Section \ref{stacks-properties-section-conventions}.




\section{対角射の性質}
\label{stacks-morphisms-section-diagonals}

\noindent
代数スタックの対角射は $\mathit{Isom}$-層と密接に関係する。次を参照せよ：
Algebraic Stacks, Lemma \ref{algebraic-lemma-representable-diagonal}.
代数スタックの第二の定義条件により、これらの $\mathit{Isom}$-層は常に
代数空間である。

\begin{lemma}
\label{stacks-morphisms-lemma-isom-locally-finite-type}
$\mathcal{X}$ を代数スタックとする。$T$ をスキームとし、
$x, y$ を $T$ 上の $\mathcal{X}$ のファイバー圏の対象とする。このとき射
$\mathit{Isom}_\mathcal{X}(x, y) \to T$ は局所有限型である。
\end{lemma}

\begin{proof}
次により、
Algebraic Stacks, Lemma \ref{algebraic-lemma-stack-presentation}
$\mathcal{X} = [U/R]$ と仮定してよい。ここで用いたのは代数空間における
ある滑らかなグルーポイドである。次により、
Descent on Spaces,
Lemma \ref{spaces-descent-lemma-descending-property-locally-finite-type}
この性質を $T$ 上 fppf 局所的に確認すれば十分である。したがって
$x, y$ は射 $x', y' : T \to U$ から来ると仮定してよい。次により、
Groupoids in Spaces,
Lemma \ref{spaces-groupoids-lemma-quotient-stack-morphisms}
この場合
$\mathit{Isom}_\mathcal{X}(x, y) = T \times_{(y', x'), U \times_S U} R$.
であることがわかる。よって $R \to U \times_S U$ が局所有限型であることを
示せば十分である。これは、合成 $s : R \to U \times_S U \to U$ が滑らか
（したがって局所有限型。次を参照せよ：
Morphisms of Spaces, Lemmas
\ref{spaces-morphisms-lemma-smooth-locally-finite-presentation} および
\ref{spaces-morphisms-lemma-finite-presentation-finite-type})
であることと、Morphisms of Spaces, Lemma
\ref{spaces-morphisms-lemma-permanence-finite-type} から従う。
\end{proof}

\begin{lemma}
\label{stacks-morphisms-lemma-isom-pseudo-torsor-aut}
$\mathcal{X}$ を代数スタックとする。$T$ をスキームとし、
$x, y$ を $T$ 上の $\mathcal{X}$ のファイバー圏の対象とする。このとき
\begin{enumerate}
\item $\mathit{Isom}_\mathcal{X}(y, y)$ は $T$ 上の群代数空間であり、
\item $\mathit{Isom}_\mathcal{X}(x, y)$ は $T$ 上
$\mathit{Isom}_\mathcal{X}(y, y)$ に対する擬トーサーである。
\end{enumerate}
\end{lemma}

\begin{proof}
次を参照せよ：
Groupoids in Spaces,
Definitions \ref{spaces-groupoids-definition-group-space} および
\ref{spaces-groupoids-definition-pseudo-torsor}.
補題は、$\mathcal{X}$ がグルーポイドのスタックであることから直ちに従う。
\end{proof}

\noindent
$f : \mathcal{X} \to \mathcal{Y}$ を代数スタックの射とする。
{\it $f$ の対角射}とは、射
$$
\Delta_f :
\mathcal{X}
\longrightarrow
\mathcal{X} \times_\mathcal{Y} \mathcal{X}
$$
をいう。任意の対角射は次の二つの性質をもつ。

\begin{lemma}
\label{stacks-morphisms-lemma-properties-diagonal}
\begin{slogan}
代数スタックの射の対角射は代数空間によって表現可能で、局所有限型である。
\end{slogan}
$f : \mathcal{X} \to \mathcal{Y}$ を代数スタックの射とする。このとき
\begin{enumerate}
\item $\Delta_f$ は代数空間によって表現可能であり、
\item $\Delta_f$ は局所有限型である。
\end{enumerate}
\end{lemma}

\begin{proof}
スキーム $T$ と射
$a : T \to \mathcal{X} \times_\mathcal{Y} \mathcal{X}$
を取る。ファイバー積の定義と $2$-米田の補題により、射 $a$ は三つ組
$a = (x, x', \alpha)$ で与えられる。ここで $x, x'$ は $T$ 上の
$\mathcal{X}$ の対象で、$\alpha : f(x) \to f(x')$ は $T$ 上の
$\mathcal{Y}$ のファイバー圏の射である。代数スタックの定義により、層
$\mathit{Isom}_\mathcal{X}(x, x')$ と
$\mathit{Isom}_\mathcal{Y}(f(x), f(x'))$ は $T$ 上の代数空間である。
この言葉では、$\alpha$ は射
$\mathit{Isom}_\mathcal{Y}(f(x), f(x')) \to T$ の切断を定める。
$T$ 上のスキーム $T' \to T$ に対し、
$\mathcal{X} \times_{\mathcal{X} \times_\mathcal{Y} \mathcal{X}, a} T$
の $T'$-値点は、$f$ による像が $\alpha|_{T'}$ となる同型
$x|_{T'} \to x'|_{T'}$ と同じものである。したがって
\begin{equation}
\label{stacks-morphisms-equation-diagonal}
\vcenter{
\xymatrix{
\mathcal{X} \times_{\mathcal{X} \times_\mathcal{Y} \mathcal{X}, a} T
\ar[d] \ar[r] &
\mathit{Isom}_\mathcal{X}(x, x') \ar[d] \\
T\ar[r]^-\alpha &
\mathit{Isom}_\mathcal{Y}(f(x), f(x'))
}
}
\end{equation}
は $T$ 上の層のファイバー積正方形である。特に
$\mathcal{X} \times_{\mathcal{X} \times_\mathcal{Y} \mathcal{X}, a} T$
は代数空間であり、補題の (1) が従う。

\medskip\noindent
第二の主張を証明するには、図式 (\ref{stacks-morphisms-equation-diagonal}) の左の垂直射が
局所有限型であることを示さなければならない。
Lemma \ref{stacks-morphisms-lemma-isom-locally-finite-type}
により、代数空間 $\mathit{Isom}_\mathcal{X}(x, x')$ は $T$ 上
局所有限型である。したがって図式 (\ref{stacks-morphisms-equation-diagonal}) の右の垂直射は
局所有限型である。次を参照せよ：
Morphisms of Spaces, Lemma \ref{spaces-morphisms-lemma-permanence-finite-type}.
結論は次から従う：
Morphisms of Spaces,
Lemma \ref{spaces-morphisms-lemma-base-change-finite-type}.
\end{proof}

\begin{lemma}
\label{stacks-morphisms-lemma-properties-diagonal-representable}
$f : \mathcal{X} \to \mathcal{Y}$ を代数空間によって表現可能な
代数スタックの射とする。このとき
\begin{enumerate}
\item $\Delta_f$ は（スキームによって）表現可能であり、
\item $\Delta_f$ は局所有限型であり、
\item $\Delta_f$ はモノ射であり、
\item $\Delta_f$ は分離であり、
\item $\Delta_f$ は局所準有限である。
\end{enumerate}
\end{lemma}

\begin{proof}
既に
Lemma \ref{stacks-morphisms-lemma-properties-diagonal}
で $\Delta_f$ が代数空間によって表現可能であることを見た。したがって
主張 (2) -- (5) には意味がある。次を参照せよ：
Properties of Stacks,
Section \ref{stacks-properties-section-properties-morphisms}.
また
Lemma \ref{stacks-morphisms-lemma-properties-diagonal}
により (2) が成り立つ。射
$T \to \mathcal{X} \times_\mathcal{Y} \mathcal{X}$ を取り、
図式 (\ref{stacks-morphisms-equation-diagonal}) を考える。次により、
Algebraic Stacks, Lemma
\ref{algebraic-lemma-criterion-map-representable-spaces-fibred-in-groupoids}
右の垂直射は層の写像として単射、すなわち代数空間のモノ射である。
したがって射
$T \times_{\mathcal{X} \times_\mathcal{Y} \mathcal{X}} \mathcal{X} \to T$
もモノ射である。よって (3) が成り立つ。既に
$T \times_{\mathcal{X} \times_\mathcal{Y} \mathcal{X}} \mathcal{X} \to T$
が局所有限型であることを知っている。したがって
Morphisms of Spaces, Lemma
\ref{spaces-morphisms-lemma-monomorphism-loc-finite-type-loc-quasi-finite}
により
$T \times_{\mathcal{X} \times_\mathcal{Y} \mathcal{X}} \mathcal{X} \to T$
は局所準有限かつ分離である。これで (4) と (5) が従う。最後に
Morphisms of Spaces, Proposition
\ref{spaces-morphisms-proposition-locally-quasi-finite-separated-over-scheme}
により
$T \times_{\mathcal{X} \times_\mathcal{Y} \mathcal{X}} \mathcal{X}$
はスキームである。これで (1) が従う。
\end{proof}

\begin{lemma}
\label{stacks-morphisms-lemma-representable-separated-diagonal-closed}
$f : \mathcal{X} \to \mathcal{Y}$ を代数空間によって表現可能な
代数スタックの射とする。このとき次は同値である。
\begin{enumerate}
\item $f$ は分離である。
\item $\Delta_f$ は閉埋入である。
\item $\Delta_f$ は固有である。
\item $\Delta_f$ は普遍閉である。
\end{enumerate}
\end{lemma}

\begin{proof}
「$f$ は分離である」「$\Delta_f$ は閉埋入である」
「$\Delta_f$ は普遍閉である」「$\Delta_f$ は固有である」という主張は、
次で定義した概念を指す：
Properties of Stacks,
Section \ref{stacks-properties-section-properties-morphisms}.
スキーム $V$ と全射で滑らかな射 $V \to \mathcal{Y}$ を選ぶ。
$U = \mathcal{X} \times_\mathcal{Y} V$ と置く。仮定によりこれは代数空間で、
射 $U \to \mathcal{X}$ は全射かつ滑らかである。次により、
Categories, Lemma \ref{categories-lemma-base-change-diagonal}
および
Properties of Stacks,
Lemma \ref{stacks-properties-lemma-check-property-covering}
その補題における任意の性質 $P$ に対し、$\Delta_f$ が $P$ をもつことと
$\Delta_{U/V} : U \to  U \times_V U$ が $P$ をもつことは同値である。
したがって (2), (3), (4) の同値性は
Morphisms of Spaces,
Lemma \ref{spaces-morphisms-lemma-separated-diagonal-proper}
を $U \to V$ に適用して従う。さらに (1) が成り立つならば、
$U \to V$ は分離で、$\Delta_{U/V}$ は閉埋入、すなわち (2) が成り立つ。
最後に (2) を仮定する。スキーム $T$ と射 $a : T \to \mathcal{Y}$ を取り、
$T' = \mathcal{X} \times_\mathcal{Y} T$ と置く。(1) を証明するには、
代数空間の射 $T' \to T$ が分離であることを示さなければならない。再び
Categories, Lemma \ref{categories-lemma-base-change-diagonal}
を用いると、$\Delta_{T'/T}$ は $\Delta_f$ の基底変換である。
したがって仮定 (2) により $\Delta_{T'/T}$ は閉埋入であり、
$T' \to T$ は所望どおり分離である。
\end{proof}

\begin{lemma}
\label{stacks-morphisms-lemma-representable-quasi-separated-diagonal-quasi-compact}
$f : \mathcal{X} \to \mathcal{Y}$ を代数空間によって表現可能な
代数スタックの射とする。このとき次は同値である。
\begin{enumerate}
\item $f$ は準分離である。
\item $\Delta_f$ は準コンパクトである。
\item $\Delta_f$ は有限型である。
\end{enumerate}
\end{lemma}

\begin{proof}
「$f$ は準分離である」「$\Delta_f$ は準コンパクトである」
「$\Delta_f$ は有限型である」という主張は、次で定義した概念を指す：
Properties of Stacks,
Section \ref{stacks-properties-section-properties-morphisms}.
(2) と (3) は同値であることに注意する。実際、
Lemma \ref{stacks-morphisms-lemma-properties-diagonal-representable}
（および
Algebraic Stacks, Lemma
\ref{algebraic-lemma-representable-transformations-property-implication}）
により $\Delta_f$ は局所有限型だからである。
スキーム $V$ と全射で滑らかな射 $V \to \mathcal{Y}$ を選ぶ。
$U = \mathcal{X} \times_\mathcal{Y} V$ と置く。仮定によりこれは代数空間で、
射 $U \to \mathcal{X}$ は全射かつ滑らかである。次により、
Categories, Lemma \ref{categories-lemma-base-change-diagonal}
および
Properties of Stacks,
Lemma \ref{stacks-properties-lemma-check-property-covering}
$\Delta_f$ が準コンパクトであることと
$\Delta_{U/V} : U \to  U \times_V U$ が準コンパクトであることは同値である。
(1) が成り立つならば、$U \to V$ は準分離であり、
$\Delta_{U/V}$ は準コンパクト、すなわち (2) が成り立つ。
(2) を仮定する。スキーム $T$ と射 $a : T \to \mathcal{Y}$ を取り、
$T' = \mathcal{X} \times_\mathcal{Y} T$ と置く。(1) を証明するには、
代数空間の射 $T' \to T$ が準分離であることを示さなければならない。再び
Categories, Lemma \ref{categories-lemma-base-change-diagonal}
を用いると、$\Delta_{T'/T}$ は $\Delta_f$ の基底変換である。
したがって仮定 (2) により $\Delta_{T'/T}$ は準コンパクトであり、
$T' \to T$ は所望どおり準分離である。
\end{proof}

\begin{lemma}
\label{stacks-morphisms-lemma-representable-locally-separated-diagonal-immersion}
$f : \mathcal{X} \to \mathcal{Y}$ を代数空間によって表現可能な
代数スタックの射とする。このとき次は同値である。
\begin{enumerate}
\item $f$ は局所分離である。
\item $\Delta_f$ は埋入である。
\end{enumerate}
\end{lemma}

\begin{proof}
「$f$ は局所分離である」「$\Delta_f$ は埋入である」という主張は、
次で定義した概念を指す：
Properties of Stacks,
Section \ref{stacks-properties-section-properties-morphisms}.
証明は省略する。ヒント：次の証明と同様に論じよ：
Lemmas \ref{stacks-morphisms-lemma-representable-separated-diagonal-closed} および
\ref{stacks-morphisms-lemma-representable-quasi-separated-diagonal-quasi-compact}.
\end{proof}






\section{分離公理}
\label{stacks-morphisms-section-separated}

\noindent
$\mathcal{X} = [U/R]$ を代数スタックの表示とする。このとき $S$ 上の
$\mathcal{X}$ の対角射の性質は、射 $j : R \to U \times_S U$ の性質である。
例えば $\mathcal{X} = [S/G]$ で、$G$ が $S$ 上の代数空間における
滑らかな群ならば、$j$ は構造射 $G \to S$ である。したがって対角射それ自体は
自動的には分離でない（スキームや代数空間の場合とは異なる）。
$[S/G]$ が $S$ 上準分離であるというなら、$G \to S$ は当然
準コンパクトでなければならない。しかし射 $G \to S$ が準分離であることも
要求せずに $[S/G]$ が $S$ 上準分離であるとは言いにくい。
言い換えると、対角射が準コンパクトであることだけを要求しても
「準分離代数スタック」という直観には十分一致せず、対角射それ自体も
準分離であることを要求すべきである。

\medskip\noindent
では「分離代数スタック」はどうであろうか。既に
Morphisms of Spaces,
Lemma \ref{spaces-morphisms-lemma-separated-diagonal-proper}
で、代数空間が分離であることと対角射が固有であることが同値だと見た。
これは通常、分離代数スタックの定義にも用いられる条件である。
上の例 $[S/G] \to S$ では、これは $G \to S$ が固有群スキームであることを
意味する。したがって $E$ が $k$ 上の楕円曲線ならば、形
$[\Spec(k)/E]$ の代数スタックは $k$ 上固有である
（将来の参照をここに挿入）。状況によっては、対角射が有限であると
仮定する方が自然かもしれない。

\begin{definition}
\label{stacks-morphisms-definition-separated}
$f : \mathcal{X} \to \mathcal{Y}$ を代数スタックの射とする。
\begin{enumerate}
\item $\Delta_f$ が非分岐であるとき、$f$ は {\it DM} であるという\footnote{
DM は Deligne--Mumford を表す。$f$ が DM ならば、任意のスキーム
$T$ と任意の射 $T \to \mathcal{Y}$ に対し、ファイバー積
$\mathcal{X}_T = \mathcal{X} \times_\mathcal{Y} T$
$T$ 上の代数スタックで、その対角射は非分岐、すなわち
$\mathcal{X}_T$ は DM である。したがって $\mathcal{X}_T$ は
Deligne--Mumford スタックである。Theorem \ref{stacks-morphisms-theorem-DM} を参照せよ。
言い換えると、DM 射とはその「ファイバー」が Deligne--Mumford スタック
となる射である。少なくともこれが用語の動機となることを期待する。}。
\item $\Delta_f$ が局所準有限であるとき、$f$ は {\it 準 DM} であるという
\footnote{$f$ が準 DM ならば、$\mathcal{X} \to \mathcal{Y}$ の
「ファイバー」$\mathcal{X}_T$ は準 DM である。代数スタック
$\mathcal{X}$ が準 DM であることと、スキーム $U$ および局所準有限な
有限表示の全射平坦射 $U \to \mathcal{X}$ が存在することは同値である。次を参照せよ：
Theorem \ref{stacks-morphisms-theorem-quasi-DM}.
スキームによるエタール被覆をもつことによって定義される
Deligne--Mumford 性との類似に注意する。}。
\item $\Delta_f$ が固有であるとき、$f$ は{\it 分離}であるという。
\item $\Delta_f$ が準コンパクトかつ準分離であるとき、
$f$ は{\it 準分離}であるという。
\end{enumerate}
\end{definition}

\noindent
この定義では、$\Delta_f$ が代数空間によって表現可能であることと、
次を用いている：
Properties of Stacks,
Section \ref{stacks-properties-section-properties-morphisms}
これにより $\Delta_f$ に条件を課すことに意味がある。
$f$ が代数空間によって表現可能な場合、これらの定義は既存の概念と
矛盾しない。次を参照せよ：
Lemmas \ref{stacks-morphisms-lemma-representable-quasi-separated-diagonal-quasi-compact} および
\ref{stacks-morphisms-lemma-representable-separated-diagonal-closed}.
高次対角射を見ることによってこれらの条件を特徴付ける興味深い方法がある。
次を参照せよ：
Lemma \ref{stacks-morphisms-lemma-definition-separated}.

\begin{definition}
\label{stacks-morphisms-definition-absolute-separated}
基底スキーム $S$ 上の代数スタック $\mathcal{X}$ を取り、
構造射を $p : \mathcal{X} \to S$ と表す。
\begin{enumerate}
\item $p : \mathcal{X} \to S$ が DM であるとき、
$\mathcal{X}$ は {\it $S$ 上 DM} であるという。
\item $p : \mathcal{X} \to S$ が準 DM であるとき、
$\mathcal{X}$ は {\it $S$ 上準 DM} であるという。
\item $p : \mathcal{X} \to S$ が分離であるとき、
$\mathcal{X}$ は {\it $S$ 上分離}であるという。
\item $p : \mathcal{X} \to S$ が準分離であるとき、
$\mathcal{X}$ は {\it $S$ 上準分離}であるという。
\item $\mathcal{X}$ が $\Spec(\mathbf{Z})$ 上
DM\footnote{Theorem \ref{stacks-morphisms-theorem-DM} により、これは
$\mathcal{X}$ が Deligne--Mumford スタックであることと同値である。}
であるとき、$\mathcal{X}$ は {\it DM} であるという。
\item $\mathcal{X}$ が $\Spec(\mathbf{Z})$ 上準 DM であるとき、
$\mathcal{X}$ は {\it 準 DM} であるという。
\item $\mathcal{X}$ が $\Spec(\mathbf{Z})$ 上分離であるとき、
$\mathcal{X}$ は{\it 分離}であるという。
\item $\mathcal{X}$ が $\Spec(\mathbf{Z})$ 上準分離であるとき、
$\mathcal{X}$ は{\it 準分離}であるという。
\end{enumerate}
最後の四つの定義では、次により $\mathcal{X}$ を
$\Spec(\mathbf{Z})$ 上の代数スタックとみなす：
Algebraic Stacks, Definition \ref{algebraic-definition-viewed-as}.
\end{definition}

\noindent
したがって各場合に、絶対的な概念と、与えられた基底スキームに相対的な概念が
ある（基底スキームへの言及は、次で導入した用語の濫用により通常省略する：
Properties of Stacks, Section \ref{stacks-properties-section-conventions}).
(1) $\Leftrightarrow$ (5) および (2) $\Leftrightarrow$ (6) を
Lemma \ref{stacks-morphisms-lemma-separated-implies-morphism-separated}.
で示す。以下では、これらの概念に関するいくつかの標準的結果を証明する。

\begin{lemma}
\label{stacks-morphisms-lemma-trivial-implications}
$f : \mathcal{X} \to \mathcal{Y}$ を代数スタックの射とする。
\begin{enumerate}
\item $f$ が分離ならば、$f$ は準分離である。
\item $f$ が DM ならば、$f$ は準 DM である。
\item $f$ が代数空間によって表現可能ならば、$f$ は DM である。
\end{enumerate}
\end{lemma}

\begin{proof}
(1) については、代数空間の固有射が準コンパクトかつ準分離であることに
注意する。次を参照せよ：
Morphisms of Spaces, Definition \ref{spaces-morphisms-definition-proper}.
(2) については、代数空間の非分岐射が局所準有限であることに注意する。
次を参照せよ：
Morphisms of Spaces, Lemma
\ref{spaces-morphisms-lemma-unramified-quasi-finite}.
最後に (3) は Lemma \ref{stacks-morphisms-lemma-properties-diagonal-representable} から従う。
\end{proof}

\begin{lemma}
\label{stacks-morphisms-lemma-base-change-separated}
次に列挙したすべての分離公理は基底変換で保たれる：
Definition \ref{stacks-morphisms-definition-separated}
\end{lemma}

\begin{proof}
$f : \mathcal{X} \to \mathcal{Y}$ と
$\mathcal{Y}' \to \mathcal{Y}$ を代数スタックの射とする。
$f' : \mathcal{Y}' \times_\mathcal{Y} \mathcal{X} \to \mathcal{Y}'$
を $\mathcal{Y}' \to \mathcal{Y}$ による $f$ の基底変換とする。
このとき $\Delta_{f'}$ は射
$\mathcal{X}' \times_{\mathcal{Y}'} \mathcal{X}' \to
\mathcal{X} \times_\mathcal{Y} \mathcal{X}$ による $\Delta_f$ の
基底変換である。次を参照せよ：
Categories, Lemma \ref{categories-lemma-base-change-diagonal}.
次の結果により、
Properties of Stacks,
Section \ref{stacks-properties-section-properties-morphisms}
で
Definition \ref{stacks-morphisms-definition-separated}
に用いた対角射の各性質は基底変換で保たれる。よって補題が従う。
\end{proof}

\begin{lemma}
\label{stacks-morphisms-lemma-check-separated-covering}
$f : \mathcal{X} \to \mathcal{Y}$ を代数スタックの射とする。
$W$ を代数空間とし、$W \to \mathcal{Y}$ を全射、平坦かつ
局所有限表示とする。基底変換
$W \times_\mathcal{Y} \mathcal{X} \to W$ が
Definition \ref{stacks-morphisms-definition-separated} の分離性質の一つをもつならば、
$f$ もその性質をもつ。
\end{lemma}

\begin{proof}
基底変換を $g : W \times_\mathcal{Y} \mathcal{X} \to W$ と表す。
このとき $\Delta_g$ は射
$q : W \times_\mathcal{Y} (\mathcal{X} \times_\mathcal{Y} \mathcal{X})
\to \mathcal{X} \times_\mathcal{Y} \mathcal{X}$ による $\Delta_f$ の
基底変換である。$q$ は $W \to \mathcal{Y}$ の基底変換だから、
$q$ は代数空間によって表現可能で、全射、平坦かつ局所有限表示である。
したがって結論は次から従う：
Properties of Stacks, Lemma
\ref{stacks-properties-lemma-check-property-weak-covering}.
\end{proof}

\begin{lemma}
\label{stacks-morphisms-lemma-change-of-base-separated}
スキーム $S$ を取る。$S$ 上準 DM、$S$ 上準分離、または $S$ 上分離である
という性質（次を参照せよ：
Definition \ref{stacks-morphisms-definition-absolute-separated})
）は基底スキームの変更で保たれる。次を参照せよ：
Algebraic Stacks, Definition \ref{algebraic-definition-change-of-base}.
\end{lemma}

\begin{proof}
Lemma \ref{stacks-morphisms-lemma-base-change-separated} から直ちに従う。
\end{proof}

\begin{lemma}
\label{stacks-morphisms-lemma-fibre-product-after-map}
$f : \mathcal{X} \to \mathcal{Z}$、$g : \mathcal{Y} \to \mathcal{Z}$、
および $\mathcal{Z} \to \mathcal{T}$ を代数スタックの射とする。誘導射
$i : \mathcal{X} \times_\mathcal{Z} \mathcal{Y} \to
\mathcal{X} \times_\mathcal{T} \mathcal{Y}$ を考える。このとき
\begin{enumerate}
\item $i$ は代数空間によって表現可能で、局所有限型である。
\item $\Delta_{\mathcal{Z}/\mathcal{T}}$ が準分離ならば、
$i$ は準分離である。
\item $\Delta_{\mathcal{Z}/\mathcal{T}}$ が分離ならば、
$i$ は分離である。
\item $\mathcal{Z} \to \mathcal{T}$ が DM ならば、
$i$ は非分岐である。
\item $\mathcal{Z} \to \mathcal{T}$ が準 DM ならば、
$i$ は局所準有限である。
\item $\mathcal{Z} \to \mathcal{T}$ が分離ならば、$i$ は固有である。
\item $\mathcal{Z} \to \mathcal{T}$ が準分離ならば、
$i$ は準コンパクトかつ準分離である。
\end{enumerate}
\end{lemma}

\begin{proof}
次の図式
$$
\xymatrix{
\mathcal{X} \times_\mathcal{Z} \mathcal{Y} \ar[r]_i \ar[d] &
\mathcal{X} \times_\mathcal{T} \mathcal{Y} \ar[d] \\
\mathcal{Z} \ar[r]^-{\Delta_{\mathcal{Z}/\mathcal{T}}} \ar[r] &
\mathcal{Z} \times_\mathcal{T} \mathcal{Z}
}
$$
は $2$-ファイバー積図式である。次を参照せよ：
Categories, Lemma \ref{categories-lemma-fibre-product-after-map}.
したがって $i$ は対角射 $\Delta_{\mathcal{Z}/\mathcal{T}}$ の基底変換である。
よって補題は
Lemma \ref{stacks-morphisms-lemma-properties-diagonal},
と、次の内容から従う：
Properties of Stacks,
Section \ref{stacks-properties-section-properties-morphisms}.
\end{proof}

\begin{lemma}
\label{stacks-morphisms-lemma-semi-diagonal}
$\mathcal{T}$ を代数スタックとする。
$g : \mathcal{X} \to \mathcal{Y}$ を $\mathcal{T}$ 上の代数スタックの射とし、
$g$ のグラフ $i : \mathcal{X} \to
\mathcal{X} \times_\mathcal{T} \mathcal{Y}$ を考える。このとき
\begin{enumerate}
\item $i$ は代数空間によって表現可能で、局所有限型である。
\item $\mathcal{Y} \to \mathcal{T}$ が DM ならば、$i$ は非分岐である。
\item $\mathcal{Y} \to \mathcal{T}$ が準 DM ならば、$i$ は局所準有限である。
\item $\mathcal{Y} \to \mathcal{T}$ が分離ならば、$i$ は固有である。
\item $\mathcal{Y} \to \mathcal{T}$ が準分離ならば、
$i$ は準コンパクトかつ準分離である。
\end{enumerate}
\end{lemma}

\begin{proof}
これは Lemma \ref{stacks-morphisms-lemma-fibre-product-after-map} を射
$\mathcal{X} = \mathcal{X} \times_\mathcal{Y} \mathcal{Y} \to
\mathcal{X} \times_\mathcal{T} \mathcal{Y}$ に適用した特別な場合である。
\end{proof}

\begin{lemma}
\label{stacks-morphisms-lemma-section-immersion}
$f : \mathcal{X} \to \mathcal{T}$ を代数スタックの射とする。
$s : \mathcal{T} \to \mathcal{X}$ を、$f \circ s$ が
$\text{id}_\mathcal{T}$ と $2$-同型であるような射とする。このとき
\begin{enumerate}
\item $s$ は代数空間によって表現可能で、局所有限型である。
\item $f$ が DM ならば、$s$ は非分岐である。
\item $f$ が準 DM ならば、$s$ は局所準有限である。
\item $f$ が分離ならば、$s$ は固有である。
\item $f$ が準分離ならば、$s$ は準コンパクトかつ準分離である。
\end{enumerate}
\end{lemma}

\begin{proof}
これは Lemma \ref{stacks-morphisms-lemma-semi-diagonal} を
$g = s$ および $\mathcal{Y} = \mathcal{T}$ に適用した特別な場合である。
この場合
$i : \mathcal{T} \to \mathcal{T} \times_\mathcal{T} \mathcal{X}$
は $s$ と $2$-同型である。
\end{proof}

\begin{lemma}
\label{stacks-morphisms-lemma-composition-separated}
次に列挙したすべての分離公理は射の合成で保たれる：
Definition \ref{stacks-morphisms-definition-separated}
\end{lemma}

\begin{proof}
$f : \mathcal{X} \to \mathcal{Y}$ と
$g : \mathcal{Y} \to \mathcal{Z}$ を、考えている公理が適用される
代数スタックの射とする。対角射 $\Delta_{\mathcal{X}/\mathcal{Z}}$ は合成
$$
\mathcal{X} \longrightarrow
\mathcal{X} \times_\mathcal{Y} \mathcal{X} \longrightarrow
\mathcal{X} \times_\mathcal{Z} \mathcal{X}.
$$
考えている分離公理は、対角射がある性質 $\mathcal{P}$ をもつことによって
定義される。上の
Lemma \ref{stacks-morphisms-lemma-fibre-product-after-map}
により、第二の射もこの性質をもつ。
性質 $\mathcal{P}$ をもつ、代数空間によって表現可能な射の合成も
性質 $\mathcal{P}$ をもつので、補題が従う。次の一般的議論を参照せよ：
Properties of Stacks,
Section \ref{stacks-properties-section-properties-morphisms}
および
Morphisms of Spaces, Lemmas
\ref{spaces-morphisms-lemma-composition-unramified},
\ref{spaces-morphisms-lemma-composition-quasi-finite},
\ref{spaces-morphisms-lemma-composition-proper},
\ref{spaces-morphisms-lemma-composition-quasi-compact}, および
\ref{spaces-morphisms-lemma-composition-separated}.
\end{proof}

\begin{lemma}
\label{stacks-morphisms-lemma-separated-over-separated}
$f : \mathcal{X} \to \mathcal{Y}$ を基底スキーム $S$ 上の
代数スタックの射とする。
\begin{enumerate}
\item $\mathcal{Y}$ が $S$ 上 DM で $f$ が DM ならば、
$\mathcal{X}$ は $S$ 上 DM である。
\item $\mathcal{Y}$ が $S$ 上準 DM で $f$ が準 DM ならば、
$\mathcal{X}$ は $S$ 上準 DM である。
\item $\mathcal{Y}$ が $S$ 上分離で $f$ が分離ならば、
$\mathcal{X}$ は $S$ 上分離である。
\item $\mathcal{Y}$ が $S$ 上準分離で $f$ が準分離ならば、
$\mathcal{X}$ は $S$ 上準分離である。
\item $\mathcal{Y}$ が DM で $f$ が DM ならば、$\mathcal{X}$ は DM である。
\item $\mathcal{Y}$ が準 DM で $f$ が準 DM ならば、
$\mathcal{X}$ は準 DM である。
\item $\mathcal{Y}$ が分離で $f$ が分離ならば、$\mathcal{X}$ は分離である。
\item $\mathcal{Y}$ が準分離で $f$ が準分離ならば、
$\mathcal{X}$ は準分離である。
\end{enumerate}
\end{lemma}

\begin{proof}
(1), (2), (3), (4) は
Lemma \ref{stacks-morphisms-lemma-composition-separated}
と
Definition \ref{stacks-morphisms-definition-absolute-separated}.
から直ちに従う。(5), (6), (7), (8) については、$\mathcal{X}$ と
$\mathcal{Y}$ を $\Spec(\mathbf{Z})$ 上の代数スタックとみなし、
Lemma \ref{stacks-morphisms-lemma-composition-separated}.
を適用する。詳細は省略する。
\end{proof}

\noindent
次の補題は代数空間に対する類似と少し異なる。比較のため、
Morphisms of Spaces, Lemma
\ref{spaces-morphisms-lemma-compose-after-separated} を参照せよ。

\begin{lemma}
\label{stacks-morphisms-lemma-compose-after-separated}
$f : \mathcal{X} \to \mathcal{Y}$ と
$g : \mathcal{Y} \to \mathcal{Z}$ を代数スタックの射とする。
\begin{enumerate}
\item $g \circ f$ が DM ならば、$f$ も DM である。
\item $g \circ f$ が準 DM ならば、$f$ も準 DM である。
\item $g \circ f$ が分離で $\Delta_g$ が分離ならば、
$f$ は分離である。
\item $g \circ f$ が準分離で $\Delta_g$ が準分離ならば、
$f$ は準分離である。
\end{enumerate}
\end{lemma}

\begin{proof}
対角射 $g \circ f$ の分解
$$
\mathcal{X} \to
\mathcal{X} \times_\mathcal{Y} \mathcal{X} \to
\mathcal{X} \times_\mathcal{Z} \mathcal{X}
$$
を考える。両射は代数空間によって表現可能である。次を参照せよ：
Lemmas \ref{stacks-morphisms-lemma-properties-diagonal} および
\ref{stacks-morphisms-lemma-fibre-product-after-map}.
したがって任意のスキーム $T$ と射
$T \to \mathcal{X} \times_\mathcal{Y} \mathcal{X}$
に対し、代数空間の射
$$
A = \mathcal{X} \times_{(\mathcal{X} \times_\mathcal{Z} \mathcal{X})} T
\longrightarrow
B = (\mathcal{X} \times_\mathcal{Y} \mathcal{X})
\times_{(\mathcal{X} \times_\mathcal{Z} \mathcal{X})} T
\longrightarrow
T
$$
$g \circ f$ が DM（それぞれ準 DM）ならば、合成 $A \to T$ は
非分岐（それぞれ局所準有限）である。したがって $A \to B$ は
非分岐（それぞれ局所準有限）である。次を参照せよ：
Morphisms of Spaces, Lemma
\ref{spaces-morphisms-lemma-permanence-unramified}
（それぞれ
Morphisms of Spaces,
Lemma \ref{spaces-morphisms-lemma-permanence-quasi-finite}）。
次に図式
$$
\xymatrix{
A \times_B T \ar[r] \ar[d] & T \ar[d] \\
A \ar[r] \ar[d] & B \ar[r] \ar[d] & T \ar[d] \\
\mathcal{X} \ar[r] &
\mathcal{X} \times_\mathcal{Y} \mathcal{X} \ar[r] &
\mathcal{X} \times_\mathcal{Z} \mathcal{X}
}
$$
を考える。すべての正方形は $2$-Cartesian であり、射 $T \to B$ は
最初に取った射 $T \to \mathcal{X} \times_\mathcal{Y} \mathcal{X}$
から得られる。射 $A \times_B T \to T$ は非分岐
（それぞれ局所準有限）である。次を参照せよ：Morphisms of Spaces, Lemma
\ref{spaces-morphisms-lemma-base-change-unramified}
（それぞれ
Morphisms of Spaces,
Lemma \ref{spaces-morphisms-lemma-base-change-quasi-finite}）。
$A \times_B T$ が
$\mathcal{X} \times_{\mathcal{X} \times_\mathcal{Y} \mathcal{X}} T$
に等しいことに注意すると、$f$ は DM（それぞれ準 DM）である。
これで (1) と (2) が従う。

\medskip\noindent
(4) の証明。$g \circ f$ は準分離で $\Delta_g$ は準分離であると仮定する。
対角射 $g \circ f$ の分解
$$
\mathcal{X} \to
\mathcal{X} \times_\mathcal{Y} \mathcal{X} \to
\mathcal{X} \times_\mathcal{Z} \mathcal{X}
$$
を考える。両射は代数空間によって表現可能で、第二の射は準分離である。
次を参照せよ：
Lemmas \ref{stacks-morphisms-lemma-properties-diagonal} および
\ref{stacks-morphisms-lemma-fibre-product-after-map}.
したがって任意のスキーム $T$ と射
$T \to \mathcal{X} \times_\mathcal{Y} \mathcal{X}$
に対し、代数空間の射
$$
A = \mathcal{X} \times_{(\mathcal{X} \times_\mathcal{Z} \mathcal{X})} T
\longrightarrow
B = (\mathcal{X} \times_\mathcal{Y} \mathcal{X})
\times_{(\mathcal{X} \times_\mathcal{Z} \mathcal{X})} T
\longrightarrow
T
$$
を得る。ここで $B \to T$ は準分離である。
$g \circ f$ は準分離と仮定したので、合成 $A \to T$ は
準コンパクトかつ準分離である。したがって $A \to B$ は準分離である。次を参照せよ：
Morphisms of Spaces,
Lemma \ref{spaces-morphisms-lemma-compose-after-separated}.
また $A \to B$ は準コンパクトである。次を参照せよ：
Morphisms of Spaces,
Lemma \ref{spaces-morphisms-lemma-quasi-compact-permanence}.
よって $f$ は準分離である。

\medskip\noindent
(3) の証明。$g \circ f$ は分離で $\Delta_g$ は分離であると仮定する。
対角射 $g \circ f$ の分解
$$
\mathcal{X} \to
\mathcal{X} \times_\mathcal{Y} \mathcal{X} \to
\mathcal{X} \times_\mathcal{Z} \mathcal{X}
$$
を考える。両射は代数空間によって表現可能で、第二の射は分離である。
次を参照せよ：
Lemmas \ref{stacks-morphisms-lemma-properties-diagonal} および
\ref{stacks-morphisms-lemma-fibre-product-after-map}.
したがって任意のスキーム $T$ と射
$T \to \mathcal{X} \times_\mathcal{Y} \mathcal{X}$
に対し、代数空間の射
$$
A = \mathcal{X} \times_{(\mathcal{X} \times_\mathcal{Z} \mathcal{X})} T
\longrightarrow
B = (\mathcal{X} \times_\mathcal{Y} \mathcal{X})
\times_{(\mathcal{X} \times_\mathcal{Z} \mathcal{X})} T
\longrightarrow
T
$$
を得る。ここで $B \to T$ は分離である。
$g \circ f$ は準分離と仮定したので、合成 $A \to T$ は固有である。
したがって $A \to B$ は固有である。次を参照せよ：
Morphisms of Spaces,
Lemma \ref{spaces-morphisms-lemma-universally-closed-permanence}
これは $f$ が分離であることを意味する。
\end{proof}

\begin{lemma}
\label{stacks-morphisms-lemma-separated-implies-morphism-separated}
基底スキーム $S$ 上の代数スタック $\mathcal{X}$ を取る。
\begin{enumerate}
\item
$\mathcal{X}$ は DM $\Leftrightarrow$
$\mathcal{X}$ は $S$ 上 DM。
\item
$\mathcal{X}$ は準 DM $\Leftrightarrow$
$\mathcal{X}$ は $S$ 上準 DM。
\item $\mathcal{X}$ が分離ならば、
$\mathcal{X}$ は $S$ 上分離である。
\item $\mathcal{X}$ が準分離ならば、
$\mathcal{X}$ は $S$ 上準分離である。
\end{enumerate}
$f : \mathcal{X} \to \mathcal{Y}$ を基底スキーム $S$ 上の
代数スタックの射とする。
\begin{enumerate}
\item[(5)] $\mathcal{X}$ が $S$ 上 DM ならば、$f$ は DM である。
\item[(6)] $\mathcal{X}$ が $S$ 上準 DM ならば、$f$ は準 DM である。
\item[(7)] $\mathcal{X}$ が $S$ 上分離で
$\Delta_{\mathcal{Y}/S}$ が分離ならば、$f$ は分離である。
\item[(8)] $\mathcal{X}$ が $S$ 上準分離で
$\Delta_{\mathcal{Y}/S}$ が準分離ならば、$f$ は準分離である。
\end{enumerate}
\end{lemma}

\begin{proof}
(5), (6), (7), (8) は
Lemma \ref{stacks-morphisms-lemma-compose-after-separated}
と
Spaces, Definition \ref{spaces-definition-separated}.
から直ちに従う。(3) と (4) を証明するには、$X$ と $Y$ を
$\Spec(\mathbf{Z})$ 上の代数スタックとみなし、
Lemma \ref{stacks-morphisms-lemma-compose-after-separated}.
を適用する。同様に (1) と (2) を証明するには、$\mathcal{X}$ を
$\Spec(\mathbf{Z})$ 上の代数スタックとみなし、射
$$
\mathcal{X} \longrightarrow
\mathcal{X} \times_S \mathcal{X} \longrightarrow
\mathcal{X} \times_{\Spec(\mathbf{Z})} \mathcal{X}
$$
を考える。両射は代数空間によって表現可能である。
第二の射は埋入 $\Delta_{S/\mathbf{Z}}$ の基底変換として、
非分岐かつ局所準有限である。したがって合成が非分岐
（それぞれ局所準有限）であることと、第一の射が非分岐
（それぞれ局所準有限）であることは同値である。次を参照せよ：
Morphisms of Spaces,
Lemmas \ref{spaces-morphisms-lemma-composition-unramified} および
\ref{spaces-morphisms-lemma-permanence-unramified}
（それぞれ Morphisms of Spaces,
Lemmas \ref{spaces-morphisms-lemma-composition-quasi-finite} および
\ref{spaces-morphisms-lemma-permanence-quasi-finite}）。
\end{proof}

\begin{lemma}
\label{stacks-morphisms-lemma-properties-covering-imply-diagonal}
$\mathcal{X}$ を代数スタックとする。$W$ を代数空間とし、
$f : W \to \mathcal{X}$ を全射、平坦かつ局所有限表示な射とする。
\begin{enumerate}
\item $f$ が非分岐（すなわちエタール、すなわち $\mathcal{X}$ が
Deligne--Mumford）ならば、$\mathcal{X}$ は DM である。
\item $f$ が局所準有限ならば、$\mathcal{X}$ は準 DM である。
\end{enumerate}
\end{lemma}

\begin{proof}
$f$ が非分岐ならば、次によりエタールであることに注意する：
Morphisms of Spaces, Lemma
\ref{spaces-morphisms-lemma-unramified-flat-lfp-etale}.
これが (1) の括弧内の記述を説明する。
$f$ は非分岐（それぞれ局所準有限）であると仮定する。
$\Delta_\mathcal{X} : \mathcal{X} \to \mathcal{X} \times \mathcal{X}$
が非分岐（それぞれ局所準有限）であることを示さなければならない。
$W \times W \to \mathcal{X} \times \mathcal{X}$ も全射、平坦かつ
局所有限表示であることに注意する。したがって
$$
W \times_{\mathcal{X} \times \mathcal{X}, \Delta_\mathcal{X}} \mathcal{X}
=
W \times_\mathcal{X} W
\longrightarrow
W \times W
$$
が非分岐（それぞれ局所準有限）であることを示せば十分である。次を参照せよ：
Properties of Stacks, Lemma
\ref{stacks-properties-lemma-check-property-covering}.
仮定により射 $\text{pr}_i : W \times_\mathcal{X} W \to W$ は
非分岐（それぞれ局所準有限）である。したがって表示した射は
非分岐（それぞれ局所準有限）である。次を参照せよ：
Morphisms of Spaces, Lemma
\ref{spaces-morphisms-lemma-permanence-unramified}
（それぞれ Morphisms of Spaces, Lemma
\ref{spaces-morphisms-lemma-permanence-quasi-finite}）。
\end{proof}

\begin{lemma}
\label{stacks-morphisms-lemma-monomorphism-separated}
代数スタックの単射は分離かつ DM である。
代数スタックの埋入についても同じことが成り立つ。
\end{lemma}

\begin{proof}
$f : \mathcal{X} \to \mathcal{Y}$ が代数スタックの単射ならば、
$\Delta_f$ は同型である。次を参照せよ：
Properties of Stacks, Lemma \ref{stacks-properties-lemma-monomorphism}.
代数空間の同型は固有かつ非分岐なので、$f$ は分離かつ DM である。
第二の主張は第一の主張から従う。実際、埋入は単射である。次を参照せよ：
Properties of Stacks,
Lemma \ref{stacks-properties-lemma-immersion-monomorphism}.
\end{proof}

\begin{lemma}
\label{stacks-morphisms-lemma-separation-properties-residual-gerbe}
$\mathcal{X}$ を代数スタックとし、$x \in |\mathcal{X}|$ とする。
$x$ における $\mathcal{X}$ の剰余ガーベ $\mathcal{Z}_x$ が存在すると仮定する。
$\mathcal{X}$ が DM（それぞれ準 DM、分離、準分離）ならば、
$\mathcal{Z}_x$ もそうである。
\end{lemma}

\begin{proof}
$\mathcal{Z}_x \to \mathcal{X}$ は単射であり、したがって
Lemma \ref{stacks-morphisms-lemma-monomorphism-separated}.
により DM かつ分離である。これに
Lemma \ref{stacks-morphisms-lemma-separated-over-separated}
を適用すればよい。
\end{proof}


















\section{慣性スタック}
\label{stacks-morphisms-section-inertia}

\noindent
グルーポイドをファイバーとするスタックの（相対）慣性スタックは
Stacks, Section \ref{stacks-section-the-inertia-stack}.
で定義されている。ファイバー圏の設定での実際の構成と、そのいくつかの性質は
Categories, Section \ref{categories-section-inertia}.
にある。

\begin{lemma}
\label{stacks-morphisms-lemma-inertia}
$\mathcal{X}$ を代数スタックとする。このとき慣性スタック
$\mathcal{I}_\mathcal{X}$ も代数スタックである。射
$$
\mathcal{I}_\mathcal{X} \longrightarrow \mathcal{X}
$$
は代数空間によって表現可能かつ局所有限型である。より一般に、
$f : \mathcal{X} \to \mathcal{Y}$ を代数スタックの射とする。このとき相対慣性
$\mathcal{I}_{\mathcal{X}/\mathcal{Y}}$ は代数スタックであり、射
$$
\mathcal{I}_{\mathcal{X}/\mathcal{Y}} \longrightarrow \mathcal{X}
$$
は代数空間によって表現可能かつ局所有限型である。
\end{lemma}

\begin{proof}
Categories, Lemma \ref{categories-lemma-inertia-fibred-category}
により同値
$$
\mathcal{I}_\mathcal{X} \to
\mathcal{X} \times_{\Delta, \mathcal{X} \times_S \mathcal{X}, \Delta}
\mathcal{X}
\quad\text{and}\quad
\mathcal{I}_{\mathcal{X}/\mathcal{Y}} \to
\mathcal{X}
\times_{\Delta, \mathcal{X} \times_\mathcal{Y} \mathcal{X}, \Delta}
\mathcal{X}
$$
があり、慣性スタックが代数スタックであることが分かる。
$T$ 上の $\mathcal{X}$ のファイバー圏の対象 $x$ が与える射
$T \to \mathcal{X}$ を取る。このとき $2$-ファイバー積図式
$$
\xymatrix{
\mathit{Isom}_\mathcal{X}(x, x) \ar[d] \ar[r] &
\mathcal{I}_\mathcal{X} \ar[d] \\
T \ar[r]^x & \mathcal{X}
}
$$
を得る。これは $\mathcal{I}_\mathcal{X}$ の定義から直ちに従う。
$\mathit{Isom}_\mathcal{X}(x, x)$ は常に $T$ 上局所有限型な代数空間なので（次を参照せよ：
Lemma \ref{stacks-morphisms-lemma-isom-locally-finite-type})
）、$\mathcal{I}_\mathcal{X} \to \mathcal{X}$ は代数空間によって表現可能かつ
局所有限型である。最後に、相対慣性について
$$
\vcenter{
\xymatrix{
\mathit{Isom}_\mathcal{X}(x, x) \ar[d] &
K \ar[l] \ar[d] \ar[r] &
\mathcal{I}_{\mathcal{X}/\mathcal{Y}} \ar[d] \\
\mathit{Isom}_\mathcal{Y}(f(x), f(x)) &
T \ar[l]_-e \ar[r]^x & \mathcal{X}
}
}
$$
を得て、両正方形は $2$-ファイバー積である。これは
Categories, Lemma \ref{categories-lemma-relative-inertia-as-fibre-product}.
から従う。左の縦射は $T$ 上局所有限型な代数空間の射であり、
したがって局所有限型である。次を参照せよ：
Morphisms of Spaces,
Lemma \ref{spaces-morphisms-lemma-permanence-finite-type}.
ゆえに $K$ は代数空間であり、$K \to T$ は局所有限型である。
これで相対慣性に関する主張が証明された。
\end{proof}

\begin{remark}
\label{stacks-morphisms-remark-inertia-is-group-in-spaces}
$f : \mathcal{X} \to \mathcal{Y}$ を代数スタックの射とする。
Properties of Stacks, Remark \ref{stacks-properties-remark-representable-over}
で見たように、終域を $\mathcal{X}$ とする代数空間によって表現可能な射
$\mathcal{Z} \to \mathcal{X}$ の $2$-圏は圏をなす。
この圏において $\mathcal{X}/\mathcal{Y}$ の慣性スタックは
{\it 群対象}である。対象
$\mathcal{I}_{\mathcal{X}/\mathcal{Y}}$
は単に対 $(x, \alpha)$ であり、ここで $x$ は $\mathcal{X}$ の対象、
$\alpha$ は $x$ が属する $\mathcal{X}$ のファイバー圏における $x$ の自己同型で、
$f(\alpha) = \text{id}$ を満たすことを思い出そう。合成
$$
c :
\mathcal{I}_{\mathcal{X}/\mathcal{Y}}
\times_\mathcal{X} \mathcal{I}_{\mathcal{X}/\mathcal{Y}}
\longrightarrow
\mathcal{I}_{\mathcal{X}/\mathcal{Y}}
$$
は対象上の規則
$$
((x, \alpha), (x', \alpha'), \beta) \mapsto
(x, \alpha \circ \beta^{-1} \circ \alpha' \circ \beta)
$$
で与えられる。ファイバー積の定義により $\beta : x \to x'$ は
ファイバー圏における同型なので、これは意味をもつ。単位元
$e : \mathcal{X} \to \mathcal{I}_{\mathcal{X}/\mathcal{Y}}$ は関手
$x \mapsto (x, \text{id}_x)$ で与えられる。群対象の公理が成り立つことの証明は省略する。
\end{remark}

\noindent
$f : \mathcal{X} \to \mathcal{Y}$ を代数スタックの射とし、
$\mathcal{I}_{\mathcal{X}/\mathcal{Y}}$ をその慣性スタックとする。
$T$ をスキーム、$x$ を $T$ 上の $\mathcal{X}$ の対象とする。
$y = f(x)$ と置く。Lemma \ref{stacks-morphisms-lemma-inertia} の証明で見たように、
任意のスキーム $T$ と $T$ 上の $\mathcal{X}$ の対象 $x$ に対して、
群の層の完全列
\begin{equation}
\label{stacks-morphisms-equation-exact-sequence-isom}
0 \to
\mathit{Isom}_{\mathcal{X}/\mathcal{Y}}(x, x) \to
\mathit{Isom}_\mathcal{X}(x, x) \to
\mathit{Isom}_\mathcal{Y}(y, y)
\end{equation}
がある。第二項と第三項の群構造は Lemma \ref{stacks-morphisms-lemma-isom-pseudo-torsor-aut}
で定義したものであり、この列が第一項に意味を与える。また、標準的な
Cartesian 図式
$$
\xymatrix{
\mathit{Isom}_{\mathcal{X}/\mathcal{Y}}(x, x) \ar[d] \ar[r] &
\mathcal{I}_{\mathcal{X}/\mathcal{Y}} \ar[d] \\
T \ar[r]^x & \mathcal{X}
}
$$
がある。実際、Remark \ref{stacks-morphisms-remark-inertia-is-group-in-spaces} で論じた
$\mathcal{I}_{\mathcal{X}/\mathcal{Y}}$ 上の群構造は、
$\mathit{Isom}_{\mathcal{X}/\mathcal{Y}}(x, x)$ 上の群構造を誘導する。
これにより層
$\mathit{Isom}_{\mathcal{X}/\mathcal{Y}}(x, x)$
を、代数空間から $\mathcal{X}$ への射に対しても定義できる。これを次の定義で定式化する。

\begin{definition}
\label{stacks-morphisms-definition-isom}
$f : \mathcal{X} \to \mathcal{Y}$ を代数スタックの射とし、
$Z$ を代数空間とする。
\begin{enumerate}
\item $x : Z \to \mathcal{X}$ を射とする。次のように置く：
$$
\mathit{Isom}_{\mathcal{X}/\mathcal{Y}}(x, x) =
Z \times_{x, \mathcal{X}} \mathcal{I}_{\mathcal{X}/\mathcal{Y}}
$$
Remark \ref{stacks-morphisms-remark-inertia-is-group-in-spaces} で論じた合成則を引き戻すことにより、
これに $Z$ 上の群代数空間の構造を入れる。
$\mathit{Isom}_{\mathcal{X}/\mathcal{Y}}(x, x)$ を
{\it $x$ の相対自己同型の層}と呼ぶことがある。
\item $x_1, x_2 : Z \to \mathcal{X}$ を射とし、
$y_i = f \circ x_i$ と置く。$\alpha : y_1 \to y_2$ を $2$-射とする。
このとき $\alpha$ は射
$\Delta^\alpha : Z \to Z \times_{y_1, \mathcal{Y}, y_2} Z$ を定め、次のように置く：
$$
\mathit{Isom}_{\mathcal{X}/\mathcal{Y}}^\alpha(x_1, x_2) =
(Z \times_{x_1, \mathcal{X}, x_2} Z)
\times_{Z \times_{y_1, \mathcal{Y}, y_2} Z, \Delta^\alpha} Z.
$$
$\mathit{Isom}_{\mathcal{X}/\mathcal{Y}}^\alpha(x_1, x_2)$
を{\it $x_1$ から $x_2$ への相対同型の層}と呼ぶことがある。
\end{enumerate}
もし $\mathcal{Y} = \Spec(\mathbf{Z})$、またはより一般に $\mathcal{Y}$ が
代数空間ならば、記法
 $\mathit{Isom}_\mathcal{X}(x, x)$ および $\mathit{Isom}_\mathcal{X}(x_1, x_2)$
を用い、{\it $x$ の自己同型の層}および
{\it $x_1$ から $x_2$ への同型の層}という用語を用いる。
\end{definition}

\begin{lemma}
\label{stacks-morphisms-lemma-isom-pseudo-torsor-aut-over-space}
$f : \mathcal{X} \to \mathcal{Y}$ を代数スタックの射とする。
$Z$ を代数空間とし、$x_i : Z \to \mathcal{X}$（$i = 1, 2$）を射とする。このとき
\begin{enumerate}
\item $\mathit{Isom}_{\mathcal{X}/\mathcal{Y}}(x_2, x_2)$
は $Z$ 上の群代数空間である。
\item 群の完全列
$$
0 \to \mathit{Isom}_{\mathcal{X}/\mathcal{Y}}(x_2, x_2)
\to \mathit{Isom}_\mathcal{X}(x_2, x_2)
\to \mathit{Isom}_\mathcal{Y}(f \circ x_2, f \circ x_2)
$$
がある。
\item 代数空間の写像
$
\mathit{Isom}_\mathcal{X}(x_1, x_2)
\to \mathit{Isom}_\mathcal{Y}(f \circ x_1, f \circ x_2)
$
があり、任意の $2$-射 $\alpha : f \circ x_1 \to f \circ x_2$ に対して
Cartesian 図式
$$
\xymatrix{
\mathit{Isom}_{\mathcal{X}/\mathcal{Y}}^\alpha(x_1, x_2) \ar[d] \ar[r] &
Z \ar[d]^\alpha \\
\mathit{Isom}_\mathcal{X}(x_1, x_2) \ar[r] &
\mathit{Isom}_\mathcal{Y}(f \circ x_1, f \circ x_2)
}
$$
を得る。
\item 任意の $2$-射 $\alpha : f \circ x_1 \to f \circ x_2$ に対して、
代数空間 $\mathit{Isom}_{\mathcal{X}/\mathcal{Y}}^\alpha(x_1, x_2)$ は
$Z$ 上の $\mathit{Isom}_{\mathcal{X}/\mathcal{Y}}(x_2, x_2)$ に対する擬トーサーである。
\end{enumerate}
\end{lemma}

\begin{proof}
(1) は Definition \ref{stacks-morphisms-definition-isom} から従う。
(2) は $Z$ 上エタール局所的に、完全列 (\ref{stacks-morphisms-equation-exact-sequence-isom})
から得られる。(3) は定義を展開すれば分かる。
$Z$ 上エタール位相について局所的に、(4) は
Lemma \ref{stacks-morphisms-lemma-isom-pseudo-torsor-aut} の (2) に帰着する。
\end{proof}

\begin{lemma}
\label{stacks-morphisms-lemma-cartesian-square-inertia}
$\pi : \mathcal{X} \to \mathcal{Y}$ と
$f : \mathcal{Y}' \to \mathcal{Y}$ を代数スタックの射とする。
$\mathcal{X}' = \mathcal{X} \times_\mathcal{Y} \mathcal{Y}'$ と置く。
このとき図式
$$
\xymatrix{
\mathcal{I}_{\mathcal{X}'/\mathcal{Y}'} \ar[r]
\ar[d]_{
\text{Categories, Equation}\ (\ref{categories-equation-functorial})
} &
\mathcal{X}' \ar[r]_{\pi'} \ar[d] & \mathcal{Y}' \ar[d]^f \\
\mathcal{I}_{\mathcal{X}/\mathcal{Y}} \ar[r] &
\mathcal{X} \ar[r]^\pi & \mathcal{Y}
}
$$
の両正方形はファイバー積図式である。
\end{lemma}

\begin{proof}
慣性スタック $\mathcal{I}_{\mathcal{X}'/\mathcal{Y}'}$ は、$x'$ が
$\mathcal{X}'$ の対象であり、
$\alpha'$ が $\pi'(\alpha') = \text{id}$ を満たす $x'$ の自己同型であるような
対 $(x', \alpha')$ の圏として定義される。次を参照せよ：
Categories, Section \ref{categories-section-inertia}.
$x'$ がスキーム $U$ 上にあり、$\mathcal{X}$ の対象 $x$ に写ると仮定する。
Categories, Lemma \ref{categories-lemma-2-product-categories-over-C}
における $2$-ファイバー積の構成により、
$x' = (U, x, y', \beta)$ であることが分かる。ここで $y'$ は $U$ 上の
$\mathcal{Y}'$ の対象であり、$\beta$ は $U$ 上の $\mathcal{Y}$ の
ファイバー圏における同型 $\beta : \pi(x) \to f(y')$ である。
$2$-ファイバー積の構成そのものにより、自己同型 $\alpha'$ は対
$(\alpha, \gamma)$ であり、ここで $\alpha$ は $U$ 上の $x$ の自己同型、
$\gamma$ は $U$ 上の $y'$ の自己同型であり、$\alpha$ と $\gamma$ は
$\beta$ を介して両立する。条件 $\pi'(\alpha') = \text{id}$ は
$\gamma = \text{id}$ を意味し、このとき $\alpha, \beta, \gamma$ が両立する
という条件はちょうど $\pi(\alpha) = \text{id}$ という条件、すなわち
$(x, \alpha)$ が $\mathcal{I}_{\mathcal{X}/\mathcal{Y}}$ の対象であるという条件である。
以上から左の正方形はファイバー積図式である（細部は省略する）。
\end{proof}

\begin{lemma}
\label{stacks-morphisms-lemma-monomorphism-cartesian-square-inertia}
$f : \mathcal{X} \to \mathcal{Y}$ を代数スタックの単射とする。このとき図式
$$
\xymatrix{
\mathcal{I}_\mathcal{X} \ar[r] \ar[d] &
\mathcal{X} \ar[d] \\
\mathcal{I}_\mathcal{Y} \ar[r] &
\mathcal{Y}
}
$$
はファイバー積図式である。
\end{lemma}

\begin{proof}
$f$ が充満忠実であること（次を参照せよ：
Properties of Stacks, Lemma \ref{stacks-properties-lemma-monomorphism})
と、次における慣性の定義から直ちに従う：
Categories, Section \ref{categories-section-inertia}.
すなわち、スキーム $T$ 上の $\mathcal{I}_\mathcal{X}$ の対象とは、
$T$ 上の $\mathcal{X}$ の対象 $x$ と、$T$ 上の $\mathcal{X}$ の
ファイバー圏における射 $\alpha : x \to x$ とからなる対 $(x, \alpha)$ にほかならない。
$f$ は充満忠実なので、$\alpha$ を与えることは $T$ 上の $\mathcal{Y}$ の
ファイバー圏における射 $\beta : f(x) \to f(x)$ を与えることと同じである。
したがって、$T$ 上の $\mathcal{I}_\mathcal{X}$ の対象を三つ組
$((y, \beta), x, \gamma)$ とみなせる。ここで $y$ は $T$ 上の
$\mathcal{Y}$ の対象、$\beta : y \to y$ は $\mathcal{Y}_T$ における射、
$\gamma : y \to f(x)$ は $T$ 上の同型である。すなわちこれは
$T$ 上の $\mathcal{I}_\mathcal{Y} \times_\mathcal{Y} \mathcal{X}$ の対象である。
\end{proof}

\begin{lemma}
\label{stacks-morphisms-lemma-presentation-inertia}
$\mathcal{X}$ を代数スタックとし、$[U/R] \to \mathcal{X}$ を表示とする。
$G/U$ をグルーポイド $(U, R, s, t, c)$ に付随する安定化群代数空間とする。このとき
$$
\xymatrix{
G \ar[d] \ar[r] & U \ar[d] \\
\mathcal{I}_\mathcal{X} \ar[r] & \mathcal{X}
}
$$
はファイバー積図式である。
\end{lemma}

\begin{proof}
Groupoids in Spaces, Lemma \ref{spaces-groupoids-lemma-2-cartesian-inertia}.
から直ちに従う。
\end{proof}









\section{高次対角射}
\label{stacks-morphisms-section-higher-diagonals}

\noindent
$f : \mathcal{X} \to \mathcal{Y}$ を代数スタックの射とする。
この状況では対角射
$$
\Delta_f : \mathcal{X} \to \mathcal{X} \times_\mathcal{Y} \mathcal{X}
$$
だけでなく、対角射の対角射、すなわち射
$$
\Delta_{\Delta_f} :
\mathcal{X}
\longrightarrow
\mathcal{X} \times_{(\mathcal{X} \times_\mathcal{Y} \mathcal{X})} \mathcal{X}
$$
も考えることができる。このため、次の用語を用いることがある。
$\Delta_{f, 0} = f$ を{\it 第零対角射}、
$\Delta_{f, 1} = \Delta_f$ を{\it 第一対角射}、
$\Delta_{f, 2} = \Delta_{\Delta_f}$ を{\it 第二対角射}と呼ぶ。
$\Delta_{f, 1}$ は代数空間によって表現可能かつ局所有限型であることに注意する。次を参照せよ：
Lemma \ref{stacks-morphisms-lemma-properties-diagonal}.
したがって $\Delta_{f, 2}$ は表現可能な単射であり、局所有限型、
分離かつ局所準有限である。次を参照せよ：
Lemma \ref{stacks-morphisms-lemma-properties-diagonal-representable}.

\medskip\noindent
第二対角射は相対慣性スタックを用いて記述できる。すなわち、ファイバー積
$\mathcal{X}
\times_{(\mathcal{X} \times_\mathcal{Y} \mathcal{X})} \mathcal{X}$
は次により相対慣性スタック $\mathcal{I}_{\mathcal{X}/\mathcal{Y}}$ と同値である：
Categories, Lemma \ref{categories-lemma-inertia-fibred-category}.
さらに、この同定のもとで第二対角射は相対慣性スタックの{\it 単位切断}
$$
\Delta_{f, 2} = e : \mathcal{X} \to \mathcal{I}_{\mathcal{X}/\mathcal{Y}}
$$
となる。代数空間のグルーポイドに起こること
(Groupoids in Spaces, Lemma \ref{spaces-groupoids-lemma-diagonal})
との類比により、次の同値を得る。

\begin{lemma}
\label{stacks-morphisms-lemma-diagonal-diagonal}
$f : \mathcal{X} \to \mathcal{Y}$ を代数スタックの射とする。
\begin{enumerate}
\item
次は同値である：
\begin{enumerate}
\item $\mathcal{I}_{\mathcal{X}/\mathcal{Y}} \to \mathcal{X}$
は分離である。
\item $\Delta_{f, 1} = \Delta_f :
\mathcal{X} \to \mathcal{X} \times_\mathcal{Y} \mathcal{X}$
は分離である。
\item $\Delta_{f, 2} = e :
\mathcal{X} \to \mathcal{I}_{\mathcal{X}/\mathcal{Y}}$
は閉埋入である。
\end{enumerate}
\item
次は同値である：
\begin{enumerate}
\item $\mathcal{I}_{\mathcal{X}/\mathcal{Y}} \to \mathcal{X}$
は準分離である。
\item $\Delta_{f, 1} = \Delta_f :
\mathcal{X} \to \mathcal{X} \times_\mathcal{Y} \mathcal{X}$
は準分離である。
\item $\Delta_{f, 2} = e :
\mathcal{X} \to \mathcal{I}_{\mathcal{X}/\mathcal{Y}}$
は準コンパクトである。
\end{enumerate}
\item
次は同値である：
\begin{enumerate}
\item $\mathcal{I}_{\mathcal{X}/\mathcal{Y}} \to \mathcal{X}$
は局所分離である。
\item $\Delta_{f, 1} = \Delta_f :
\mathcal{X} \to \mathcal{X} \times_\mathcal{Y} \mathcal{X}$
は局所分離である。
\item $\Delta_{f, 2} = e :
\mathcal{X} \to \mathcal{I}_{\mathcal{X}/\mathcal{Y}}$
は埋入である。
\end{enumerate}
\item
次は同値である：
\begin{enumerate}
\item $\mathcal{I}_{\mathcal{X}/\mathcal{Y}} \to \mathcal{X}$
は非分岐である。
\item $f$ は DM である。
\end{enumerate}
\item
次は同値である：
\begin{enumerate}
\item $\mathcal{I}_{\mathcal{X}/\mathcal{Y}} \to \mathcal{X}$
は局所準有限である。
\item $f$ は準 DM である。
\end{enumerate}
\end{enumerate}
\end{lemma}

\begin{proof}
(1), (2), (3) の証明。代数空間 $U$ と全射な滑らかな射
$U \to \mathcal{X}$ を選ぶ。このとき
$G = U \times_\mathcal{X} \mathcal{I}_{\mathcal{X}/\mathcal{Y}}$
は $U$ 上の代数空間である（Lemma \ref{stacks-morphisms-lemma-inertia}）。
実際、Remark \ref{stacks-morphisms-remark-inertia-is-group-in-spaces} で構成した相対慣性上の
群則により、$G$ は $U$ 上の群代数空間である。さらに
$G \to \mathcal{I}_{\mathcal{X}/\mathcal{Y}}$ は
$U \to \mathcal{X}$ の基底変換として全射かつ滑らかである。最後に、
$e : \mathcal{X} \to \mathcal{I}_{\mathcal{X}/\mathcal{Y}}$
の $G \to \mathcal{I}_{\mathcal{X}/\mathcal{Y}}$ による基底変換は、
$G/U$ の単位元 $U \to G$ である。したがって (a) と (c) の同値は
Groupoids in Spaces, Lemma
\ref{spaces-groupoids-lemma-group-scheme-separated}.
から従う。$\Delta_{f, 2}$ は $\Delta_f$ の対角射なので、
定義により (b) $\Leftrightarrow$ (c) である。

\medskip\noindent
(4) と (5) の証明。(4)(b) は $\Delta_f$ が非分岐であることを意味し、
(5)(b) は $\Delta_f$ が局所準有限であることを意味することを思い出そう。
スキーム $Z$ と射
$a : Z \to \mathcal{X} \times_\mathcal{Y} \mathcal{X}$.
を選ぶ。このとき $a = (x_1, x_2, \alpha)$ であり、ここで
$x_i : Z \to \mathcal{X}$、また
$\alpha : f \circ x_1  \to f \circ x_2$ は $2$-射である。図式
$$
\vcenter{
\xymatrix{
\mathit{Isom}_{\mathcal{X}/\mathcal{Y}}^\alpha(x_1, x_2)
\ar[d] \ar[r] &
Z \ar[d] \\
\mathcal{X} \ar[r]^{\Delta_f} &
\mathcal{X} \times_\mathcal{Y} \mathcal{X}
}
}
\quad\text{and}\quad
\vcenter{
\xymatrix{
\mathit{Isom}_{\mathcal{X}/\mathcal{Y}}(x_2, x_2)
\ar[d] \ar[r] &
Z \ar[d]^{x_2} \\
\mathcal{I}_{\mathcal{X}/\mathcal{Y}} \ar[r] &
\mathcal{X}
}
}
$$
はいずれも Cartesian 図式である。
Lemma \ref{stacks-morphisms-lemma-isom-pseudo-torsor-aut-over-space} により、代数空間
$\mathit{Isom}_{\mathcal{X}/\mathcal{Y}}^\alpha(x_1, x_2)$ は
$Z$ 上の $\mathit{Isom}_{\mathcal{X}/\mathcal{Y}}(x_2, x_2)$ に対する
擬トーサーである。したがって (4) と (5) の同値は
Groupoids in Spaces, Lemma
\ref{spaces-groupoids-lemma-pseudo-torsor-implications}.
から従う。
\end{proof}

\begin{lemma}
\label{stacks-morphisms-lemma-second-diagonal}
$f : \mathcal{X} \to \mathcal{Y}$ を代数スタックの射とする。次は同値である：
\begin{enumerate}
\item 射 $f$ は代数空間によって表現可能である。
\item $f$ の第二対角射は同型である。
\item 群スタック $ \mathcal{I}_{\mathcal{X}/\mathcal{Y}}$ は
$\mathcal X$ 上自明である。
\item スキーム $T$ と射 $x : T \to \mathcal{X}$ に対して、
$\mathit{Isom}_\mathcal{X}(x, x) \to
\mathit{Isom}_\mathcal{Y}(f(x), f(x))$ の核は自明である。
\end{enumerate}
\end{lemma}

\begin{proof}
まず (1) と (2) の同値を証明する。すなわち、$f$ が代数空間によって
表現可能であることと $f$ が忠実であることは同値である。次を参照せよ：
Algebraic Stacks,
Lemma \ref{algebraic-lemma-characterize-representable-by-algebraic-spaces}.
一方、$f$ が忠実であることと、スキーム $T$ 上の $\mathcal{X}$ の
任意の対象 $x$ に対して関手 $f$ が単射
$\mathit{Isom}_\mathcal{X}(x, x) \to
\mathit{Isom}_\mathcal{Y}(f(x), f(x))$,
を誘導することは同値である。これは核 $K$ が自明であることと同値であり、
さらに任意の $x : T \to \mathcal{X}$ に対して $e : T \to K$ が同型であることと同値である。
上で論じたように $K = T \times_{x, \mathcal{X}} \mathcal{I}_{\mathcal{X}/\mathcal{Y}}$
なので、(1) と (2) の同値が証明された。(2) と (3) の同値を証明するには、
上の議論により、群スタックが自明であることとその単位切断が同型であることが
同値であると注意すれば十分である。最後に (3) と (4) の同値は定義から従う。
実際、Lemma \ref{stacks-morphisms-lemma-inertia} の証明で、(4) の核が $T$ 上のファイバー積
$T \times_{x, \mathcal{X}} \mathcal{I}_{\mathcal{X}/\mathcal{Y}}$ に対応することを見た。
\end{proof}

\noindent
この補題から、代数スタックの射について次の階層が得られる。

\begin{lemma}
\label{stacks-morphisms-lemma-hierarchy}
代数スタックの射 $f : \mathcal{X} \to \mathcal{Y}$ について、
\begin{enumerate}
\item 単射であることと $\Delta_{f, 1}$ が同型であることは同値である。
\item 代数空間によって表現可能であることと $\Delta_{f, 1}$ が単射であることは同値である。
\end{enumerate}
さらに、第二対角射 $\Delta_{f, 2}$ は常に単射である。
\end{lemma}

\begin{proof}
Properties of Stacks, Lemma
\ref{stacks-properties-lemma-monomorphism}
より、代数スタックの射が単射であることと、その対角射がスタックの同型であることは
同値であることを思い出そう。したがって Lemma \ref{stacks-morphisms-lemma-second-diagonal} は、
対角射が単射ならば、その射が代数空間によって表現可能である、と言い換えられる。
特に、Lemma \ref{stacks-morphisms-lemma-properties-diagonal-representable} の条件 (3) は
実は必要十分条件であることが分かる。すなわち、代数スタックの射が
代数空間によって表現可能であることと、その対角射が単射であることは同値である。
\end{proof}

\begin{lemma}
\label{stacks-morphisms-lemma-first-diagonal-separated-second-diagonal-closed}
$f : \mathcal{X} \to \mathcal{Y}$ を代数スタックの射とする。このとき
\begin{enumerate}
\item $\Delta_{f, 1}$ は分離 $\Leftrightarrow$
$\Delta_{f, 2}$ は閉埋入 $\Leftrightarrow$
$\Delta_{f, 2}$ は固有 $\Leftrightarrow$
$\Delta_{f, 2}$ は普遍閉である。
\item $\Delta_{f, 1}$ は準分離 $\Leftrightarrow$
$\Delta_{f, 2}$ は有限型 $\Leftrightarrow$ $\Delta_{f, 2}$ は準コンパクトである。
\item $\Delta_{f, 1}$ は局所分離 $\Leftrightarrow$
$\Delta_{f, 2}$ は埋入である。
\end{enumerate}
\end{lemma}

\begin{proof}
$\Delta_{f, 1}$ に
Lemmas \ref{stacks-morphisms-lemma-representable-separated-diagonal-closed},
\ref{stacks-morphisms-lemma-representable-quasi-separated-diagonal-quasi-compact}, および
\ref{stacks-morphisms-lemma-representable-locally-separated-diagonal-immersion}
を適用すれば従う。
\end{proof}

\noindent
次の補題は少し巧妙であり、これらの条件を高次代数スタックへ一般化することを示唆するかもしれない。

\begin{lemma}
\label{stacks-morphisms-lemma-definition-separated}
$f : \mathcal{X} \to \mathcal{Y}$ を代数スタックの射とする。このとき
\begin{enumerate}
\item $f$ が分離であることと、$\Delta_{f, 1}$ および $\Delta_{f, 2}$ が
普遍閉であることは同値である。
\item $f$ が準分離であることと、$\Delta_{f, 1}$ および $\Delta_{f, 2}$ が
準コンパクトであることは同値である。
\item $f$ が準 DM であることと、$\Delta_{f, 1}$ および $\Delta_{f, 2}$ が
局所準有限であることは同値である。
\item $f$ が DM であることと、$\Delta_{f, 1}$ および $\Delta_{f, 2}$ が
非分岐であることは同値である。
\end{enumerate}
\end{lemma}

\begin{proof}
(1) の証明。$\Delta_{f, 2}$ と $\Delta_{f, 1}$ は普遍閉であると仮定する。
このとき $\Delta_{f, 1}$ は次により分離かつ普遍閉である：
Lemma \ref{stacks-morphisms-lemma-first-diagonal-separated-second-diagonal-closed}.
Morphisms of Spaces,
Lemma \ref{spaces-morphisms-lemma-universally-closed-quasi-compact}
および
Algebraic Stacks,
Lemma \ref{algebraic-lemma-representable-transformations-property-implication}
により $\Delta_{f, 1}$ は準コンパクトである。したがってこれは
準コンパクト、分離、普遍閉かつ局所有限型である（次による：
Lemma \ref{stacks-morphisms-lemma-properties-diagonal})
）ので固有である。これで (1) の ``$\Leftarrow$'' が証明された。
逆向きの含意の証明は省略する。

\medskip\noindent
(2) の証明。これは
Lemma \ref{stacks-morphisms-lemma-first-diagonal-separated-second-diagonal-closed}.
から直ちに従う。

\medskip\noindent
(3) の証明。$\Delta_f = \Delta_{f, 1}$ に
Lemma \ref{stacks-morphisms-lemma-properties-diagonal-representable}
を適用すれば、$\Delta_{f, 2}$ は常に局所準有限であることから従う。

\medskip\noindent
(4) の証明。$\Delta_f = \Delta_{f, 1}$ に
Lemma \ref{stacks-morphisms-lemma-properties-diagonal-representable}
を適用すると、$\Delta_{f, 2}$ が局所有限型かつ単射であることが分かるので、
$\Delta_{f, 2}$ は常に非分岐である。次を参照せよ：
More on Morphisms of Spaces,
Lemma \ref{spaces-more-morphisms-lemma-universally-injective-unramified}.
\end{proof}

\begin{lemma}
\label{stacks-morphisms-lemma-separated-implies-isom}
$f : \mathcal{X} \to \mathcal{Y}$ を代数スタックの分離
（それぞれ準分離、準 DM、DM）射とする。このとき
\begin{enumerate}
\item 代数空間 $T_i$（$i = 1, 2$）と射
$x_i : T_i \to \mathcal{X}$ が与えられ、$y_i = f \circ x_i$ と置くと、射
$$
T_1 \times_{x_1, \mathcal{X}, x_2} T_2 \longrightarrow
T_1 \times_{y_1, \mathcal{Y}, y_2} T_2
$$
は固有（それぞれ準コンパクトかつ準分離、局所準有限、非分岐）である。
\item 代数空間 $T$ と射 $x_i : T \to \mathcal{X}$（$i = 1, 2$）が与えられ、
$y_i = f \circ x_i$ と置くと、射
$$
\mathit{Isom}_\mathcal{X}(x_1, x_2) \longrightarrow
\mathit{Isom}_\mathcal{Y}(y_1, y_2)
$$
は固有（それぞれ準コンパクトかつ準分離、局所準有限、非分岐）である。
\end{enumerate}
\end{lemma}

\begin{proof}
(1) の証明。図式
$$
\xymatrix{
T_1 \times_{x_1, \mathcal{X}, x_2} T_2 \ar[d] \ar[r] &
T_1 \times_{y_1, \mathcal{Y}, y_2} T_2 \ar[d] \\
\mathcal{X} \ar[r] & \mathcal{X} \times_\mathcal{Y} \mathcal{X}
}
$$
が Cartesian であることに注意する。したがって、$f$ が分離
（それぞれ準分離、準 DM、DM）であることと、対角射が固有
（それぞれ準コンパクトかつ準分離、局所準有限、非分岐）であることが
同値であることから従う。

\medskip\noindent
(2) の証明。これは
$$
\mathit{Isom}_\mathcal{X}(x_1, x_2) =
(T \times_{x_1, \mathcal{X}, x_2} T) \times_{T \times T, \Delta_T} T
$$
であり、したがって (2) の射が (1) の射の基底変換だからである。
\end{proof}








\section{準コンパクトな射}
\label{stacks-morphisms-section-quasi-compact}

\noindent
$f$ を代数空間によって表現可能な代数スタックの射とする。
Properties of Stacks, Section
\ref{stacks-properties-section-properties-morphisms}
で $f$ が準コンパクトであることの意味を定義した。別の特徴づけを与える。

\begin{lemma}
\label{stacks-morphisms-lemma-characterize-representable-quasi-compact}
$f : \mathcal{X} \to \mathcal{Y}$ を代数空間によって表現可能な
代数スタックの射とする。次は同値である：
\begin{enumerate}
\item $f$ は準コンパクトである（Properties of Stacks,
Section \ref{stacks-properties-section-properties-morphisms} の意味で）。
\item 任意の準コンパクトな代数スタック $\mathcal{Z}$ と任意の射
$\mathcal{Z} \to \mathcal{Y}$ に対して、代数スタック
$\mathcal{Z} \times_\mathcal{Y} \mathcal{X}$ は準コンパクトである。
\end{enumerate}
\end{lemma}

\begin{proof}
(1) を仮定し、$\mathcal{Z}$ が準コンパクトであるような代数スタックの射
$\mathcal{Z} \to \mathcal{Y}$ を取る。
Properties of Stacks,
Lemma \ref{stacks-properties-lemma-quasi-compact-stack}
により、準コンパクトなスキーム $U$ と全射な滑らかな射
$U \to \mathcal{Z}$ が存在する。$f$ は代数空間によって表現可能かつ
準コンパクトなので、定義により
$U \times_\mathcal{Y} \mathcal{X}$ は代数空間であり、
$U \times_\mathcal{Y} \mathcal{X} \to U$ は準コンパクトである。
したがって $U \times_\mathcal{Y} \mathcal{X}$ は準コンパクトな代数空間である。射
$U \times_\mathcal{Y} \mathcal{X} \to
\mathcal{Z} \times_\mathcal{Y} \mathcal{X}$
は滑らかな全射 $U \to \mathcal{Z}$ の基底変換として、滑らかかつ全射である。
したがって、もう一度
Properties of Stacks,
Lemma \ref{stacks-properties-lemma-quasi-compact-stack}
を適用すれば、$\mathcal{Z} \times_\mathcal{Y} \mathcal{X}$ は準コンパクトである。

\medskip\noindent
(2) を仮定する。$Z$ をスキームとし、$Z \to \mathcal{Y}$ を射とする。
代数空間の射 $p : Z \times_\mathcal{Y} \mathcal{X} \to Z$ が
準コンパクトであることを示さなければならない。
$U \subset Z$ をアフィン開集合とする。このとき
$p^{-1}(U) = U \times_\mathcal{Y} Z$
であり、代数空間 $U \times_\mathcal{Y} Z$ は仮定 (2) により
準コンパクトである。したがって $p$ は準コンパクトである。次を参照せよ：
Morphisms of Spaces, Lemma \ref{spaces-morphisms-lemma-quasi-compact-local}.
\end{proof}

\noindent
以上から次の定義が動機づけられる。

\begin{definition}
\label{stacks-morphisms-definition-quasi-compact}
$f : \mathcal{X} \to \mathcal{Y}$ を代数スタックの射とする。
任意の準コンパクトな代数スタック $\mathcal{Z}$ と射
$\mathcal{Z} \to \mathcal{Y}$ に対して、ファイバー積
$\mathcal{Z} \times_\mathcal{Y} \mathcal{X}$ が準コンパクトであるとき、
$f$ は{\it 準コンパクト}であるという。
\end{definition}

\noindent
上の Lemma \ref{stacks-morphisms-lemma-characterize-representable-quasi-compact} により、
これは代数空間によって表現可能な代数スタックの射に対する既存の概念と一致する。
特に、この概念は代数スタック間およびスキーム間の射について既に定義した概念と一致する。

\begin{lemma}
\label{stacks-morphisms-lemma-base-change-quasi-compact}
代数スタックの準コンパクトな射を代数スタックの任意の射で基底変換したものは
準コンパクトである。
\end{lemma}

\begin{proof}
省略する。
\end{proof}

\begin{lemma}
\label{stacks-morphisms-lemma-composition-quasi-compact}
代数スタックの準コンパクトな二つの射の合成は準コンパクトである。
\end{lemma}

\begin{proof}
省略する。
\end{proof}

\begin{lemma}
\label{stacks-morphisms-lemma-closed-immersion-quasi-compact}
代数スタックの閉埋入は準コンパクトである。
\end{lemma}

\begin{proof}
埋入は常に表現可能であるという事実と、代数空間の閉埋入についての対応する事実から従う。
\end{proof}

\begin{lemma}
\label{stacks-morphisms-lemma-surjection-from-quasi-compact}
代数スタックの射の $2$-可換図式
$$
\xymatrix{
\mathcal{X} \ar[rr]_f \ar[rd]_p & &
\mathcal{Y} \ar[dl]^q \\
& \mathcal{Z}
}
$$
を取る。$f$ が全射で $p$ が準コンパクトならば、$q$ は準コンパクトである。
\end{lemma}

\begin{proof}
$\mathcal{T}$ を準コンパクトな代数スタックとし、
$\mathcal{T} \to \mathcal{Z}$ を射とする。
Properties of Stacks,
Lemma \ref{stacks-properties-lemma-base-change-surjective}
により、射
$\mathcal{T} \times_\mathcal{Z} \mathcal{X} \to
\mathcal{T} \times_\mathcal{Z} \mathcal{Y}$
は全射であり、仮定により $\mathcal{T} \times_\mathcal{Z} \mathcal{X}$ は
準コンパクトである。したがって
$\mathcal{T} \times_\mathcal{Z} \mathcal{Y}$ は次により準コンパクトである：
Properties of Stacks, Lemma \ref{stacks-properties-lemma-quasi-compact-stack}.
\end{proof}

\begin{lemma}
\label{stacks-morphisms-lemma-quasi-compact-permanence}
$f : \mathcal{X} \to \mathcal{Y}$ と
$g : \mathcal{Y} \to \mathcal{Z}$ を代数スタックの射とする。
$g \circ f$ が準コンパクトで $g$ が準分離ならば、$f$ は準コンパクトである。
\end{lemma}

\begin{proof}
$f$ は合成
$(1, f) : \mathcal{X} \to \mathcal{X} \times_\mathcal{Z} \mathcal{Y} \to
\mathcal{Y}$.
に等しい。第一の写像は、準分離な射
$\mathcal{X} \times_\mathcal{Z} \mathcal{Y} \to \mathcal{X}$
の切断なので、次により準コンパクトである：
Lemma \ref{stacks-morphisms-lemma-section-immersion}
（これは $g$ の基底変換である。次を参照せよ：
Lemma \ref{stacks-morphisms-lemma-base-change-separated}).
第二の写像は $f$ の基底変換なので準コンパクトである。次を参照せよ：
Lemma \ref{stacks-morphisms-lemma-base-change-quasi-compact}.
また、準コンパクトな射の合成は準コンパクトである。次を参照せよ：
Lemma \ref{stacks-morphisms-lemma-composition-quasi-compact}.
\end{proof}

\begin{lemma}
\label{stacks-morphisms-lemma-quasi-compact-quasi-separated-permanence}
$f : \mathcal{X} \to \mathcal{Y}$ を代数スタックの射とする。
\begin{enumerate}
\item $\mathcal{X}$ が準コンパクトで $\mathcal{Y}$ が準分離ならば、
$f$ は準コンパクトである。
\item $\mathcal{X}$ が準コンパクトかつ準分離で、$\mathcal{Y}$ が準分離ならば、
$f$ は準コンパクトかつ準分離である。
\item 準コンパクトかつ準分離な代数スタックのファイバー積は、
準コンパクトかつ準分離である。
\end{enumerate}
\end{lemma}

\begin{proof}
(1) は
Lemma \ref{stacks-morphisms-lemma-quasi-compact-permanence}.
から従う。(2) は (1) と
Lemma \ref{stacks-morphisms-lemma-compose-after-separated}.
から従う。(3) について、$\mathcal{X} \to \mathcal{Y}$ と
$\mathcal{Z} \to \mathcal{Y}$ を準コンパクトかつ準分離な代数スタックの射とする。
このとき $\mathcal{X} \times_\mathcal{Y} \mathcal{Z} \to \mathcal{Z}$ は、
$\mathcal{X} \to \mathcal{Y}$ の基底変換として準コンパクトかつ準分離である。
ここで (2) と Lemmas \ref{stacks-morphisms-lemma-base-change-quasi-compact} および
\ref{stacks-morphisms-lemma-base-change-separated}.
を用いた。したがって $\mathcal{X} \times_\mathcal{Y} \mathcal{Z}$ は、
$\mathcal{Z}$ 上準コンパクトかつ準分離な代数スタックとして、
準コンパクトかつ準分離である。次を参照せよ：
Lemmas \ref{stacks-morphisms-lemma-separated-over-separated} および
\ref{stacks-morphisms-lemma-composition-quasi-compact}.
\end{proof}

\begin{lemma}
\label{stacks-morphisms-lemma-reach-points-scheme-theoretic-image}
$f : \mathcal{X} \to \mathcal{Y}$ を代数スタックの準コンパクトな射とする。
$y \in |\mathcal{Y}|$ を $|f|$ の像の閉包に属する点とする。
分数体 $K$ をもつ付値環 $A$ と可換図式
$$
\xymatrix{
\Spec(K) \ar[r] \ar[d] & \mathcal{X} \ar[d] \\
\Spec(A) \ar[r] & \mathcal{Y}
}
$$
が存在し、$\Spec(A)$ の閉点は $y$ に写る。
\end{lemma}

\begin{proof}
アフィンスキーム $V$、点 $v \in V$、および $v$ を $y$ に写す滑らかな射
$V \to \mathcal{Y}$ を選ぶ。基底変換図式
$$
\xymatrix{
V \times_\mathcal{Y} \mathcal{X} \ar[r] \ar[d]_g & \mathcal{X} \ar[d]^f \\
V \ar[r] & \mathcal{Y}
}
$$
を考える。射 $|V \times_\mathcal{Y} \mathcal{X}| \to
|V| \times_{|\mathcal{Y}|} |\mathcal{X}|$ は全射であることを思い出そう
(Properties of Stacks, Lemma \ref{stacks-properties-lemma-points-cartesian}).
$|V| \to |\mathcal{Y}|$ は開写像なので
(Properties of Stacks, Lemma \ref{stacks-properties-lemma-topology-points})
、$v$ は $|g|$ の像の閉包に属する。したがって、準コンパクトな射 $g$ について
補題を証明すれば十分である（Lemma \ref{stacks-morphisms-lemma-base-change-quasi-compact}）。
これを次の段落で行う。

\medskip\noindent
$\mathcal{Y} = Y$ がアフィンスキームであると仮定する。このとき $f$ は
準コンパクトなので $\mathcal{X}$ は準コンパクトである
(Definition \ref{stacks-morphisms-definition-quasi-compact}).
アフィンスキーム $W$ と全射な滑らかな射 $W \to \mathcal{X}$ を選ぶ。
このとき $|f|$ の像は $W \to Y$ の像である。
Morphisms, Lemma \ref{morphisms-lemma-reach-points-scheme-theoretic-image}
により図式
$$
\xymatrix{
\Spec(K) \ar[r] \ar[d] & W \ar[d] \ar[r] & \mathcal{X} \ar[d] \\
\Spec(A) \ar[r] & Y \ar[r] & Y
}
$$
を選べて、$\Spec(A)$ の閉点は $y$ に写る。
$W \to \mathcal{X}$ と合成すれば求める図式を得る。
\end{proof}

\begin{lemma}
\label{stacks-morphisms-lemma-check-quasi-compact-covering}
$f : \mathcal{X} \to \mathcal{Y}$ を代数スタックの射とする。
$W$ を代数空間とし、$W \to \mathcal{Y}$ を全射、平坦かつ局所有限表示な射とする。
基底変換 $W \times_\mathcal{Y} \mathcal{X} \to W$ が準コンパクトならば、
$f$ は準コンパクトである。
\end{lemma}

\begin{proof}
$W \times_\mathcal{Y} \mathcal{X} \to W$ は準コンパクトであると仮定する。
$\mathcal{Z}$ が準コンパクトな代数スタックであるような射
$\mathcal{Z} \to \mathcal{Y}$ を取る。スキーム $U$ と全射な滑らかな射
$U \to W \times_\mathcal{Y} \mathcal{Z}$.
を選ぶ。$U \to \mathcal{Z}$ は平坦、全射かつ局所有限表示であり、
$\mathcal{Z}$ は準コンパクトなので、$U' \to \mathcal{Z}$ が全射となる
準コンパクトな開部分スキーム $U' \subset U$ を取れる。このとき
$U' \times_\mathcal{Y} \mathcal{X} =
U' \times_W (W \times_\mathcal{Y} \mathcal{X})$
は仮定により準コンパクトで、
$\mathcal{Z} \times_\mathcal{Y} \mathcal{X}$ へ全射する。
したがって $\mathcal{Z} \times_\mathcal{Y} \mathcal{X}$ は準コンパクトである。
\end{proof}











\section{Noether 代数スタック}
\label{stacks-morphisms-section-noetherian}

\noindent
局所 Noether 代数スタックは既に
Properties of Stacks, Section \ref{stacks-properties-section-types-properties}.
で定義した。

\begin{definition}
\label{stacks-morphisms-definition-noetherian}
$\mathcal{X}$ を代数スタックとする。$\mathcal{X}$ が準コンパクト、準分離かつ
局所 Noether であるとき、$\mathcal{X}$ は{\it Noether} であるという。
\end{definition}

\noindent
Noether 代数スタック $\mathcal{X}$ は準コンパクトかつ局所 Noether であるだけでなく、
準分離でもあることに注意する。
Section \ref{stacks-morphisms-section-higher-diagonals}
の用語で、$p : \mathcal{X} \to \Spec(\mathbf{Z})$ を「絶対」構造射
（すなわち、$\mathcal{X}$ を $\mathbf{Z}$ 上の代数スタックとみなした構造射）と書けば、
$$
\mathcal{X}\text{ Noetherian}
\Leftrightarrow
\mathcal{X}\text{ locally Noetherian and }
\Delta_{p, 0}, \Delta_{p, 1}, \Delta_{p, 2}
\text{ quasi-compact}.
$$
これは後に、Noether 代数スタック上有限型な代数スタックが自動的には
Noether にならないことを意味する。

\begin{lemma}
\label{stacks-morphisms-lemma-locally-closed-in-noetherian}
$j : \mathcal{X} \to \mathcal{Y}$ を代数スタックの埋入とする。
\begin{enumerate}
\item $\mathcal{Y}$ が局所 Noether ならば、$\mathcal{X}$ は局所 Noether であり、
$j$ は準コンパクトである。
\item $\mathcal{Y}$ が Noether ならば、$\mathcal{X}$ は Noether である。
\end{enumerate}
\end{lemma}

\begin{proof}
スキーム $V$ と全射な滑らかな射 $V \to \mathcal{Y}$ を選ぶ。
このとき $U = \mathcal{X} \times_\mathcal{Y} V$ はスキームであり、
$V \to U$ は埋入である。次を参照せよ：
Properties of Stacks, Definition \ref{stacks-properties-definition-immersion}.
$\mathcal{Y}$ が局所 Noether であることと $V$ が局所 Noether であることは
同値であることを思い出そう。この場合 $U$ も局所 Noether であり
(Morphisms, Lemmas \ref{morphisms-lemma-immersion-locally-finite-type} および
\ref{morphisms-lemma-finite-type-noetherian})、また $U \to V$ は準コンパクトである
(Properties, Lemma \ref{properties-lemma-immersion-into-noetherian}).
したがって $j$ は準コンパクトであり
(Lemma \ref{stacks-morphisms-lemma-check-quasi-compact-covering})
、$\mathcal{X}$ は局所 Noether である。最後に $\mathcal{Y}$ が Noether ならば、
以上から $\mathcal{X}$ は準コンパクトかつ局所 Noether である。
証明を終えるため、$j$ は分離であり、$\mathcal{Y}$ は準分離なので、
$\mathcal{X}$ は次により準分離であることに注意する：
Lemma \ref{stacks-morphisms-lemma-separated-over-separated}.
\end{proof}

\begin{lemma}
\label{stacks-morphisms-lemma-Noetherian-topology}
$\mathcal{X}$ を代数スタックとする。
\begin{enumerate}
\item $\mathcal{X}$ が局所 Noether ならば、$|\mathcal{X}|$ は
局所 Noether 位相空間である。
\item $\mathcal{X}$ が準コンパクトかつ局所 Noether ならば、
$|\mathcal{X}|$ は Noether 位相空間である。
\end{enumerate}
\end{lemma}

\begin{proof}
$\mathcal{X}$ は局所 Noether であると仮定する。スキーム $U$ と
全射な滑らかな射 $U \to \mathcal{X}$ を選ぶ。$\mathcal{X}$ は局所 Noether
なので、$U$ も局所 Noether である。
Properties, Lemma \ref{properties-lemma-Noetherian-topology}
により、これは $|U|$ が局所 Noether 位相空間であることを意味する。
$|U| \to |\mathcal{X}|$ は開かつ全射なので、次により
$|\mathcal{X}|$ は局所 Noether である：
Topology, Lemma \ref{topology-lemma-image-Noetherian}.
これで (1) が証明された。$\mathcal{X}$ が準コンパクトかつ局所 Noether ならば、
$|\mathcal{X}|$ は準コンパクトかつ局所 Noether である。したがって
$|\mathcal{X}|$ は次により Noether である：Topology, Lemma
\ref{topology-lemma-quasi-compact-locally-Noetherian-Noetherian}.
\end{proof}

\begin{lemma}
\label{stacks-morphisms-lemma-locally-Noetherian-quasi-sober}
$\mathcal{X}$ を局所 Noether 代数スタックとする。このとき
$|\mathcal{X}|$ は準 sober である（Topology, Definition
\ref{topology-definition-generic-point}).
\end{lemma}

\begin{proof}
任意の既約閉部分集合 $T \subset |\mathcal{X}|$ が生成点をもつことを示さなければならない。
アフィンスキーム $U$ と滑らかな射 $f : U \to \mathcal{X}$ を、
$f^{-1}(T) \subset |U|$ が空でないように選ぶ。$U$ は Noether なので、
閉部分集合 $f^{-1}(T)$ は有限個の既約成分をもつ
(Topology, Lemma \ref{topology-lemma-Noetherian}).
$f^{-1}(T) = Z_1 \cup \ldots \cup Z_n$ を既約成分分解とする。
$f$ は開写像なので、$f|_{f^{-1}(T)} : f^{-1}(T) \to T$ の像は
$T$ の空でない開部分集合を含む。$T$ は既約なので、これは
$f(f^{-1}(T))$ が稠密であることを意味する。$T$ は既約なので、ある $i$ に対して
$f(Z_i)$ は稠密である。したがって、$\xi_i \in Z_i$ が生成点ならば、
$f(\xi_i)$ は $T$ の生成点である。
\end{proof}





\section{アフィン射}
\label{stacks-morphisms-section-affine}

\noindent
代数スタックのアフィン射を次のように定義する。

\begin{definition}
\label{stacks-morphisms-definition-affine}
代数スタックの射が表現可能であり、かつ
Properties of Stacks, Section
\ref{stacks-properties-section-properties-morphisms}.
の意味でアフィンであるとき、その射は{\it アフィン}であるという。
\end{definition}

\noindent
代数スタックのアフィン射を、代数空間によって表現可能であり、かつ
Properties of Stacks, Section
\ref{stacks-properties-section-properties-morphisms}.
の意味でアフィンである代数スタックの射と考える方が、ここではやや便利である。
（Stacks project で単に「表現可能」というときは、既定ではスキームによって
表現可能であることを思い出そう。）これは今定義した概念と明らかに同値なので、
以下では断りなくこの特徴づけを用いる。この概念についていくつか簡単な補題を証明する。

\begin{lemma}
\label{stacks-morphisms-lemma-base-change-affine}
$\mathcal{X} \to \mathcal{Y}$ を代数スタックの射とし、
$\mathcal{Z} \to \mathcal{Y}$ を代数スタックのアフィン射とする。
このとき $\mathcal{Z} \times_\mathcal{Y} \mathcal{X} \to \mathcal{X}$ は
代数スタックのアフィン射である。
\end{lemma}

\begin{proof}
Properties of Stacks, Section
\ref{stacks-properties-section-properties-morphisms}.
の議論から従う。
\end{proof}

\begin{lemma}
\label{stacks-morphisms-lemma-composition-affine}
代数スタックのアフィン射の合成はアフィンである。
\end{lemma}

\begin{proof}
Properties of Stacks, Section
\ref{stacks-properties-section-properties-morphisms}
の議論と
Morphisms of Spaces, Lemma \ref{spaces-morphisms-lemma-composition-affine}.
から従う。
\end{proof}

\begin{lemma}
\label{stacks-morphisms-lemma-affine-permanence}
代数スタックの射の可換図式
$$
\xymatrix{
\mathcal{X} \ar[rr]_f \ar[rd]_a & & \mathcal{Y} \ar[dl]^b \\
& \mathcal{Z}
}
$$
を取る。$a$ と $\Delta_b$ がアフィンならば、$f$ はアフィンである。
\end{lemma}

\begin{proof}
射 $a$ の基底変換
$\text{pr}_2 : \mathcal{X} \times_\mathcal{Z} \mathcal{Y} \to \mathcal{Y}$
は Lemma \ref{stacks-morphisms-lemma-base-change-affine} によりアフィンである。射
$(1, f) : \mathcal{X} \to \mathcal{X} \times_\mathcal{Z} \mathcal{Y}$
は射
$\Delta_b : \mathcal{Y} \to \mathcal{Y} \times_\mathcal{Z} \mathcal{Y}$
の、射 $\mathcal{X} \times_\mathcal{Z} \mathcal{Y} \to
\mathcal{Y} \times_\mathcal{Z} \mathcal{Y}$ による基底変換である
（Categories, Section \ref{categories-section-2-fibre-products} の内容を参照せよ）。
したがって
Lemma \ref{stacks-morphisms-lemma-base-change-affine}.
によりアフィンである。アフィン射の合成 $f = \text{pr}_2 \circ (1, f)$ は
Lemma \ref{stacks-morphisms-lemma-composition-affine} によりアフィンであり、証明が終わる。
\end{proof}





\section{整射と有限射}
\label{stacks-morphisms-section-integral}

\noindent
代数スタックの整射と有限射を次のように定義する。

\begin{definition}
\label{stacks-morphisms-definition-integral}
$f : \mathcal{X} \to \mathcal{Y}$ を代数スタックの射とする。
\begin{enumerate}
\item $f$ が表現可能であり、かつ Properties of Stacks, Section
\ref{stacks-properties-section-properties-morphisms}.
の意味で整であるとき、$f$ は{\it 整}であるという。
\item $f$ が表現可能であり、かつ Properties of Stacks, Section
\ref{stacks-properties-section-properties-morphisms}.
の意味で有限であるとき、$f$ は{\it 有限}であるという。
\end{enumerate}
\end{definition}

\noindent
代数スタックの整射（それぞれ有限射）を、代数空間によって表現可能であり、かつ
Properties of Stacks, Section
\ref{stacks-properties-section-properties-morphisms}.
の意味で整（それぞれ有限）である代数スタックの射と考える方が、ここではやや便利である。
（Stacks project で単に「表現可能」というときは、既定ではスキームによって
表現可能であることを思い出そう。）これは今定義した概念と明らかに同値なので、
以下では断りなくこの特徴づけを用いる。この概念についていくつか簡単な補題を証明する。

\begin{lemma}
\label{stacks-morphisms-lemma-base-change-integral}
$\mathcal{X} \to \mathcal{Y}$ を代数スタックの射とし、
$\mathcal{Z} \to \mathcal{Y}$ を代数スタックの整（または有限）射とする。このとき
$\mathcal{Z} \times_\mathcal{Y} \mathcal{X} \to \mathcal{X}$
は代数スタックの整（または有限）射である。
\end{lemma}

\begin{proof}
Properties of Stacks, Section
\ref{stacks-properties-section-properties-morphisms}.
の議論から従う。
\end{proof}

\begin{lemma}
\label{stacks-morphisms-lemma-composition-integral}
代数スタックの整射（それぞれ有限射）の合成は整（それぞれ有限）である。
\end{lemma}

\begin{proof}
Properties of Stacks, Section
\ref{stacks-properties-section-properties-morphisms}
の議論と
Morphisms of Spaces, Lemma \ref{spaces-morphisms-lemma-composition-integral}.
から従う。
\end{proof}




\section{開射}
\label{stacks-morphisms-section-open}

\noindent
$f$ を代数空間によって表現可能な代数スタックの射とする。
Properties of Stacks, Section
\ref{stacks-properties-section-properties-morphisms}
で $f$ が普遍開であることの意味を定義した。別の特徴づけを与える。

\begin{lemma}
\label{stacks-morphisms-lemma-characterize-representable-universally-open}
$f : \mathcal{X} \to \mathcal{Y}$ を代数空間によって表現可能な
代数スタックの射とする。次は同値である：
\begin{enumerate}
\item $f$ は普遍開である（Properties of Stacks,
Section \ref{stacks-properties-section-properties-morphisms} の意味で）。
\item 代数スタックの任意の射 $\mathcal{Z} \to \mathcal{Y}$ に対して、
位相空間の写像 $|\mathcal{Z} \times_\mathcal{Y} \mathcal{X}| \to |\mathcal{Z}|$
は開写像である。
\end{enumerate}
\end{lemma}

\begin{proof}
(1) を仮定し、$\mathcal{Z} \to \mathcal{Y}$ を (2) の射とする。
スキーム $V$ と全射な滑らかな射 $V \to \mathcal{Z}$ を選ぶ。
仮定により代数空間の射 $V \times_\mathcal{Y} \mathcal{X} \to V$ は普遍開であり、
特に写像 $|V \times_\mathcal{Y} \mathcal{X}| \to |V|$ は開写像である。
Properties of Stacks, Section \ref{stacks-properties-section-points}
により、可換図式
$$
\xymatrix{
|V \times_\mathcal{Y} \mathcal{X}| \ar[r] \ar[d] &
|\mathcal{Z} \times_\mathcal{Y} \mathcal{X}| \ar[d] \\
|V| \ar[r] & |\mathcal{Z}|
}
$$
の横の矢印は開かつ全射であり、さらに
$$
|V \times_\mathcal{Y} \mathcal{X}| \longrightarrow
|V| \times_{|\mathcal{Z}|} |\mathcal{Z} \times_\mathcal{Y} \mathcal{X}|
$$
は全射である。したがって左の縦射が開であることから右の縦射も開である。
これで (2) が証明された。含意 (2) $\Rightarrow$ (1) は定義から従う。
\end{proof}

\noindent
したがって、次の自然な定義を用いることができる。

\begin{definition}
\label{stacks-morphisms-definition-open}
$f : \mathcal{X} \to \mathcal{Y}$ を代数スタックの射とする。
\begin{enumerate}
\item 位相空間の写像 $|\mathcal{X}| \to |\mathcal{Y}|$ が開写像であるとき、
$f$ は{\it 開}であるという。
\item 代数スタックの任意の射 $\mathcal{Z} \to \mathcal{Y}$ に対して、
位相空間の写像
$$
|\mathcal{Z} \times_\mathcal{Y} \mathcal{X}| \to |\mathcal{Z}|
$$
が開写像である、すなわち基底変換
$\mathcal{Z} \times_\mathcal{Y} \mathcal{X} \to \mathcal{Z}$ が開であるとき、
$f$ は{\it 普遍開}であるという。
\end{enumerate}
\end{definition}

\begin{lemma}
\label{stacks-morphisms-lemma-base-change-universally-open}
代数スタックの普遍開な射を代数スタックの任意の射で基底変換したものは普遍開である。
\end{lemma}

\begin{proof}
定義から直ちに従う。
\end{proof}

\begin{lemma}
\label{stacks-morphisms-lemma-composition-universally-open}
代数スタックの二つの（普遍）開射の合成は（普遍）開である。
\end{lemma}

\begin{proof}
省略する。
\end{proof}







\section{商写像的な射}
\label{stacks-morphisms-section-submersive}

\noindent
$f$ を代数空間によって表現可能な代数スタックの射とする。
Properties of Stacks, Section
\ref{stacks-properties-section-properties-morphisms}
で $f$ が普遍的に商写像的であることの意味を定義した。別の特徴づけを与える。

\begin{lemma}
\label{stacks-morphisms-lemma-characterize-representable-universally-submersive}
$f : \mathcal{X} \to \mathcal{Y}$ を代数空間によって表現可能な
代数スタックの射とする。次は同値である：
\begin{enumerate}
\item $f$ は普遍的に商写像的である（Properties of Stacks,
Section \ref{stacks-properties-section-properties-morphisms} の意味で）。
\item 代数スタックの任意の射 $\mathcal{Z} \to \mathcal{Y}$ に対して、
位相空間の写像 $|\mathcal{Z} \times_\mathcal{Y} \mathcal{X}| \to |\mathcal{Z}|$
は商写像である。
\end{enumerate}
\end{lemma}

\begin{proof}
(1) を仮定し、$\mathcal{Z} \to \mathcal{Y}$ を (2) の射とする。
スキーム $V$ と全射な滑らかな射 $V \to \mathcal{Z}$ を選ぶ。
仮定により代数空間の射 $V \times_\mathcal{Y} \mathcal{X} \to V$ は
普遍的に商写像的であり、特に写像 $|V \times_\mathcal{Y} \mathcal{X}| \to |V|$
は商写像である。
Properties of Stacks, Section \ref{stacks-properties-section-points}
により、可換図式
$$
\xymatrix{
|V \times_\mathcal{Y} \mathcal{X}| \ar[r] \ar[d] &
|\mathcal{Z} \times_\mathcal{Y} \mathcal{X}| \ar[d] \\
|V| \ar[r] & |\mathcal{Z}|
}
$$
の横の矢印は開かつ全射であり、さらに
$$
|V \times_\mathcal{Y} \mathcal{X}| \longrightarrow
|V| \times_{|\mathcal{Z}|} |\mathcal{Z} \times_\mathcal{Y} \mathcal{X}|
$$
は全射である。したがって左の縦射が商写像であることから、右の縦射も商写像である。
これで (2) が証明された。含意 (2) $\Rightarrow$ (1) は定義から従う。
\end{proof}

\noindent
したがって、次の自然な定義を用いることができる。

\begin{definition}
\label{stacks-morphisms-definition-submersive}
$f : \mathcal{X} \to \mathcal{Y}$ を代数スタックの射とする。
\begin{enumerate}
\item 連続写像 $|\mathcal{X}| \to |\mathcal{Y}|$ が商写像であるとき、
$f$ は{\it 商写像的}であるという\footnote{これは微分多様体の沈め込みの概念とは
大きく異なる。}。次を参照せよ：
Topology, Definition \ref{topology-definition-submersive}.
\item 代数スタックの任意の射 $\mathcal{Y}' \to \mathcal{Y}$ に対して、基底変換
$\mathcal{Y}' \times_\mathcal{Y} \mathcal{X} \to \mathcal{Y}'$ が商写像的であるとき、
$f$ は{\it 普遍的に商写像的}であるという。
\end{enumerate}
\end{definition}

\noindent
商写像的な射は特に全射であることに注意する。

\begin{lemma}
\label{stacks-morphisms-lemma-base-change-universally-submersive}
代数スタックの普遍的に商写像的な射を代数スタックの任意の射で基底変換したものは、
普遍的に商写像的である。
\end{lemma}

\begin{proof}
定義から直ちに従う。
\end{proof}

\begin{lemma}
\label{stacks-morphisms-lemma-composition-universally-submersive}
代数スタックの二つの（普遍的に）商写像的な射の合成は（普遍的に）商写像的である。
\end{lemma}

\begin{proof}
省略する。
\end{proof}












\section{普遍閉射}
\label{stacks-morphisms-section-universally-closed}

\noindent
$f$ を代数空間によって表現可能な代数スタックの射とする。
Properties of Stacks, Section
\ref{stacks-properties-section-properties-morphisms}
で $f$ が普遍閉であることの意味を定義した。別の特徴づけを与える。

\begin{lemma}
\label{stacks-morphisms-lemma-characterize-representable-universally-closed}
$f : \mathcal{X} \to \mathcal{Y}$ を代数空間によって表現可能な
代数スタックの射とする。次は同値である：
\begin{enumerate}
\item $f$ は普遍閉である（Properties of Stacks,
Section \ref{stacks-properties-section-properties-morphisms} の意味で）。
\item 代数スタックの任意の射 $\mathcal{Z} \to \mathcal{Y}$ に対して、
位相空間の写像 $|\mathcal{Z} \times_\mathcal{Y} \mathcal{X}| \to |\mathcal{Z}|$
は閉写像である。
\end{enumerate}
\end{lemma}

\begin{proof}
(1) を仮定し、$\mathcal{Z} \to \mathcal{Y}$ を (2) の射とする。
スキーム $V$ と全射な滑らかな射 $V \to \mathcal{Z}$ を選ぶ。
仮定により代数空間の射 $V \times_\mathcal{Y} \mathcal{X} \to V$ は普遍閉であり、
特に写像 $|V \times_\mathcal{Y} \mathcal{X}| \to |V|$ は閉写像である。
Properties of Stacks, Section \ref{stacks-properties-section-points}
により、可換図式
$$
\xymatrix{
|V \times_\mathcal{Y} \mathcal{X}| \ar[r] \ar[d] &
|\mathcal{Z} \times_\mathcal{Y} \mathcal{X}| \ar[d] \\
|V| \ar[r] & |\mathcal{Z}|
}
$$
の横の矢印は開かつ全射であり、さらに
$$
|V \times_\mathcal{Y} \mathcal{X}| \longrightarrow
|V| \times_{|\mathcal{Z}|} |\mathcal{Z} \times_\mathcal{Y} \mathcal{X}|
$$
は全射である。したがって左の縦射が閉であることから、右の縦射も閉である。
これで (2) が証明された。含意 (2) $\Rightarrow$ (1) は定義から従う。
\end{proof}

\noindent
したがって、次の自然な定義を用いることができる。

\begin{definition}
\label{stacks-morphisms-definition-closed}
$f : \mathcal{X} \to \mathcal{Y}$ を代数スタックの射とする。
\begin{enumerate}
\item 位相空間の写像 $|\mathcal{X}| \to |\mathcal{Y}|$ が閉写像であるとき、
$f$ は{\it 閉}であるという。
\item 代数スタックの任意の射 $\mathcal{Z} \to \mathcal{Y}$ に対して、
位相空間の写像
$$
|\mathcal{Z} \times_\mathcal{Y} \mathcal{X}| \to |\mathcal{Z}|
$$
が閉写像である、すなわち基底変換
$\mathcal{Z} \times_\mathcal{Y} \mathcal{X} \to \mathcal{Z}$ が閉であるとき、
$f$ は{\it 普遍閉}であるという。
\end{enumerate}
\end{definition}

\begin{lemma}
\label{stacks-morphisms-lemma-base-change-universally-closed}
代数スタックの普遍閉射を代数スタックの任意の射で基底変換したものは普遍閉である。
\end{lemma}

\begin{proof}
定義から直ちに従う。
\end{proof}

\begin{lemma}
\label{stacks-morphisms-lemma-composition-universally-closed}
代数スタックの二つの（普遍）閉射の合成は（普遍）閉である。
\end{lemma}

\begin{proof}
省略する。
\end{proof}

\begin{lemma}
\label{stacks-morphisms-lemma-universally-closed-local}
$f : \mathcal{X} \to \mathcal{Y}$ を代数スタックの射とする。次は同値である：
\begin{enumerate}
\item $f$ は普遍閉である。
\item 任意のスキーム $Z$ と任意の射 $Z \to \mathcal{Y}$ に対して、射影
$|Z \times_\mathcal{Y} \mathcal{X}| \to |Z|$ は閉写像である。
\item 任意のアフィンスキーム $Z$ と任意の射 $Z \to \mathcal{Y}$ に対して、射影
$|Z \times_\mathcal{Y} \mathcal{X}| \to |Z|$ は閉写像である。
\item 代数空間 $V$ と全射な滑らかな射 $V \to \mathcal{Y}$ が存在し、
$V \times_\mathcal{Y} \mathcal{X} \to V$ は代数スタックの普遍閉射である。
\end{enumerate}
\end{lemma}

\begin{proof}
(1) が (2) を含意すること、および (2) が (3) を含意することの証明は省略する。

\medskip\noindent
(3) を仮定する。全射な滑らかな射 $V \to \mathcal{Y}$ を選ぶ。
$V \times_\mathcal{Y} \mathcal{X} \to V$ が代数スタックの普遍閉射であることを示す。
$\mathcal{Z} \to V$ を代数スタックから $V$ への射とする。
$W = \coprod W_i$ がアフィンスキームの非交和であるような全射な滑らかな射
$W \to \mathcal{Z}$ を取る。このとき可換図式
$$
\xymatrix{
\coprod_i |W_i \times_\mathcal{Y} \mathcal{X}| \ar@{=}[r] \ar[d] &
|W \times_\mathcal{Y} \mathcal{X}| \ar[r] \ar[d] &
|\mathcal{Z} \times_\mathcal{Y} \mathcal{X}| \ar[d] \ar@{=}[r] &
|\mathcal{Z} \times_V (V \times_\mathcal{Y} \mathcal{X})| \ar[ld] \\
\coprod |W_i| \ar@{=}[r] &
|W| \ar[r] &
|\mathcal{Z}|
}
$$
がある。南東向きの矢印が閉であることを示さなければならない。
中央の横の矢印は全射かつ開である
(Properties of Stacks, Lemma \ref{stacks-properties-lemma-topology-points}).
仮定 (3) と $W_i$ がアフィンであることから、左の縦射は閉である。
したがって右の縦射は閉である。

\medskip\noindent
(4) を仮定する。$f$ が普遍閉であることを示す。
$\mathcal{Z} \to \mathcal{Y}$ を代数スタックの射とする。図式
$$
\xymatrix{
|(V \times_\mathcal{Y} \mathcal{Z})
\times_V (V \times_\mathcal{Y} \mathcal{X})| \ar@{=}[r] \ar[rd] &
|V \times_\mathcal{Y} \mathcal{X}| \ar[r] \ar[d] &
|Z \times_\mathcal{Y} \mathcal{X}| \ar[d] \\
 &
|V \times_\mathcal{Y} \mathcal{Z}| \ar[r] &
|\mathcal{Z}|
}
$$
を考える。南西向きの矢印は仮定により閉である。横の矢印は、対応する
代数スタックの射が全射かつ滑らかなので全射かつ開である（上の参照を見よ）。
したがって右の縦射は閉である。
\end{proof}










\section{普遍単射な射}
\label{stacks-morphisms-section-universally-injective}

\noindent
$f$ を代数空間によって表現可能な代数スタックの射とする。
Properties of Stacks, Section
\ref{stacks-properties-section-properties-morphisms}
で $f$ が普遍単射であることの意味を定義した。別の特徴づけを与える。

\begin{lemma}
\label{stacks-morphisms-lemma-characterize-representable-universally-injective}
$f : \mathcal{X} \to \mathcal{Y}$ を代数空間によって表現可能な
代数スタックの射とする。次は同値である：
\begin{enumerate}
\item $f$ は普遍単射である（Properties of Stacks,
Section \ref{stacks-properties-section-properties-morphisms} の意味で）。
\item 代数スタックの任意の射 $\mathcal{Z} \to \mathcal{Y}$ に対して、写像
$|\mathcal{Z} \times_\mathcal{Y} \mathcal{X}| \to |\mathcal{Z}|$ は単射である。
\end{enumerate}
\end{lemma}

\begin{proof}
(1) を仮定し、$\mathcal{Z} \to \mathcal{Y}$ を (2) の射とする。
スキーム $V$ と全射な滑らかな射 $V \to \mathcal{Z}$ を選ぶ。
仮定により代数空間の射 $V \times_\mathcal{Y} \mathcal{X} \to V$ は普遍単射であり、
特に写像 $|V \times_\mathcal{Y} \mathcal{X}| \to |V|$ は単射である。
Properties of Stacks, Section \ref{stacks-properties-section-points}
により、可換図式
$$
\xymatrix{
|V \times_\mathcal{Y} \mathcal{X}| \ar[r] \ar[d] &
|\mathcal{Z} \times_\mathcal{Y} \mathcal{X}| \ar[d] \\
|V| \ar[r] & |\mathcal{Z}|
}
$$
の横の矢印は開かつ全射であり、さらに
$$
|V \times_\mathcal{Y} \mathcal{X}| \longrightarrow
|V| \times_{|\mathcal{Z}|} |\mathcal{Z} \times_\mathcal{Y} \mathcal{X}|
$$
は全射である。したがって左の縦射が単射であることから、右の縦射も単射である。
これで (2) が証明された。含意 (2) $\Rightarrow$ (1) は定義から従う。
\end{proof}

\noindent
したがって、次の自然な定義を用いることができる。

\begin{definition}
\label{stacks-morphisms-definition-universally-injective}
$f : \mathcal{X} \to \mathcal{Y}$ を代数スタックの射とする。
代数スタックの任意の射 $\mathcal{Z} \to \mathcal{Y}$ に対して写像
$$
|\mathcal{Z} \times_\mathcal{Y} \mathcal{X}| \to |\mathcal{Z}|
$$
が単射であるとき、$f$ は{\it 普遍単射}であるという。
\end{definition}

\begin{lemma}
\label{stacks-morphisms-lemma-base-change-universally-injective}
代数スタックの普遍単射な射を代数スタックの任意の射で基底変換したものは普遍単射である。
\end{lemma}

\begin{proof}
定義から直ちに従う。
\end{proof}

\begin{lemma}
\label{stacks-morphisms-lemma-composition-universally-injective}
代数スタックの二つの普遍単射な射の合成は普遍単射である。
\end{lemma}

\begin{proof}
省略する。
\end{proof}

\begin{lemma}
\label{stacks-morphisms-lemma-universally-injective}
$f : \mathcal{X} \to \mathcal{Y}$ を代数スタックの射とする。次は同値である：
\begin{enumerate}
\item $f$ は普遍単射である。
\item $\Delta : \mathcal{X} \to \mathcal{X} \times_\mathcal{Y} \mathcal{X}$
は全射である。
\item 代数閉体について、$x_1, x_2 : \Spec(k) \to \mathcal{X}$ と
$2$-射 $\beta : f \circ x_1 \to f \circ x_2$ に対して、
$\beta = \text{id}_f \star \alpha$ を満たす $2$-射
$\alpha : x_1 \to x_2$ が存在する。
\end{enumerate}
\end{lemma}

\begin{proof}
(1) $\Rightarrow$ (2)。$f$ が普遍単射ならば、第一射影
$|\mathcal{X} \times_\mathcal{Y} \mathcal{X}| \to |\mathcal{X}|$
は単射であり、したがって $|\Delta|$ は全射である。

\medskip\noindent
(2) $\Rightarrow$ (1)。$\Delta$ は全射であると仮定する。このとき
$\Delta$ の任意の基底変換は全射である（Properties of Stacks, Section
\ref{stacks-properties-section-surjective}).
$f$ の基底変換の対角射は $\Delta$ の基底変換なので、
$|\mathcal{X}| \to |\mathcal{Y}|$ が単射であることを示せば十分である。
そうでないならば、Properties of Stacks, Lemma
\ref{stacks-properties-lemma-points-cartesian}
により第一射影
$|\mathcal{X} \times_\mathcal{Y} \mathcal{X}| \to |\mathcal{X}|$
は単射でない。これはもちろん $|\Delta|$ が全射でないことを意味する。

\medskip\noindent
(3) $\Rightarrow$ (2)。
$t \in |\mathcal{X} \times_\mathcal{Y} \mathcal{X}|$.
を取る。このとき $t$ は、$k$ を代数閉体とする射
$t : \Spec(k) \to \mathcal{X} \times_\mathcal{Y} \mathcal{X}$
で表せる。$2$-ファイバー積の構成により、$t$ は
$(x_1, x_2, \beta)$ で表せる。ここで $x_1, x_2 : \Spec(k) \to \mathcal{X}$、
$\beta : f \circ x_1 \to f \circ x_2$ は $2$-射である。
このとき (3) により、$\beta$ に写る $2$-射 $\alpha : x_1 \to x_2$ が存在する。
これはちょうど $\Delta(x_1) = (x_1, x_1, \text{id})$ が $t$ と同型であることを
意味する。したがって (2) が成り立つ。

\medskip\noindent
(2) $\Rightarrow$ (3)。$k$ を代数閉体とし、
$x_1, x_2 : \Spec(k) \to \mathcal{X}$ を射とする。
$\beta : f \circ x_1 \to f \circ x_2$ を $2$-射とする。
前の段落と同様に、射
$t = (x_1, x_2, \beta) :
\Spec(k) \to \mathcal{X} \times_\mathcal{Y} \mathcal{X}$.
を得る。Lemma \ref{stacks-morphisms-lemma-properties-diagonal} により
$$
T = \mathcal{X}
\times_{\Delta, \mathcal{X} \times_\mathcal{Y} \mathcal{X}, t} \Spec(k)
$$
は $\Spec(k)$ 上局所有限型な代数空間である。条件 (2) により $T$ は空でない。
$k$ は代数閉なので、$T$ には $k$-点がある。定義を展開すると、これは
$\beta = \text{id}_f \star \alpha$ を満たす射
$\alpha : x_1 \to x_2$ が $\Mor(\Spec(k), \mathcal{X})$ に存在することを意味する。
\end{proof}

\begin{lemma}
\label{stacks-morphisms-lemma-universally-injective-point}
$f : \mathcal{X} \to \mathcal{Y}$ を代数スタックの普遍単射な射とする。
$k$ を代数閉体とし、$y : \Spec(k) \to \mathcal{Y}$ を射とする。
$y$ が $|\mathcal{X}| \to |\mathcal{Y}|$ の像に属するならば、
$y = f \circ x$ を満たす射 $x : \Spec(k) \to \mathcal{X}$ が存在する。
\end{lemma}

\begin{proof}
まず、この補題は自明ではないことに注意する。$y$ が $|f|$ の像に属するという仮定は、
$k$ を拡大体で置き換えれば $y$ を $\mathcal{X}$ への射に持ち上げられる、ということしか
意味しないからである。補題を証明するため、$f$ を $y$ で基底変換してよい。
したがって、空でない代数スタック $\mathcal{X}$ と普遍単射な射
$\mathcal{X} \to \Spec(k)$ があり、$\mathcal{X}$ の $k$-値点を求める場合に帰着する。
$\mathcal{X}$ をその被約化で置き換えてよい。体 $k'$ と全射、平坦かつ局所有限型な射
$\Spec(k') \to \mathcal{X}$ を選べる。次を参照せよ：
Properties of Stacks, Lemma \ref{stacks-properties-lemma-unique-point}.
$\mathcal{X} \to \Spec(k)$ は普遍単射なので、
$$
\Spec(k') \times_\mathcal{X} \Spec(k') \to \Spec(k' \otimes_k k')
$$
は全射
$\Delta : \mathcal{X} \to \mathcal{X} \times_{\Spec(k)} \mathcal{X}$
の基底変換として全射である（Lemma \ref{stacks-morphisms-lemma-universally-injective}）。
$k$ は代数閉なので、$k' \otimes_k k'$ は整域である
(Algebra, Lemma
\ref{algebra-lemma-geometrically-integral-any-integral-base-change}).
$\xi \in \Spec(k') \times_\mathcal{X} \Spec(k')$ を
$\Spec(k' \otimes_k k')$ の生成点に写る点とする。
$U$ を、$\xi$ を含む $\Spec(k') \times_\mathcal{X} \Spec(k')$ の連結成分上の
被約誘導閉部分スキーム構造とする。このとき二つの射影
$U \to \Spec(k')$ は局所有限型である。実際、射影
$\Spec(k') \times_\mathcal{X} \Spec(k') \to \Spec(k')$
が射 $\Spec(k') \to \mathcal{X}$ の基底変換としてそうであった。
Varieties, Proposition \ref{varieties-proposition-unique-base-field}
を適用すると、$\Gamma(U, \mathcal{O}_U)$ における $k'$ の二つの像の
整閉包は等しいことが分かる。$\kappa(\xi)$ で考えると、これは
$\lambda \otimes 1$ の形の任意の元が部分体
$$
1 \otimes k' \subset 
(\text{fraction field of }k' \otimes_k k') \subset
\kappa(\xi).
$$
上代数的従属であることを意味する。$k$ は代数閉なので、これは
$k' = k$ のときにしか起こらず、証明が完了する。
\end{proof}

\begin{lemma}
\label{stacks-morphisms-lemma-universally-injective-local}
$f : \mathcal{X} \to \mathcal{Y}$ を代数スタックの射とする。次は同値である：
\begin{enumerate}
\item $f$ は普遍単射である。
\item 任意のアフィンスキーム $Z$ と任意の射 $Z \to \mathcal{Y}$ に対して、射
$Z \times_\mathcal{Y} \mathcal{X} \to Z$ は普遍単射である。
\item ここにさらに追加する。
\end{enumerate}
\end{lemma}

\begin{proof}
(1) $\Rightarrow$ (2) は直ちに従う。(2) が成り立つと仮定する。
$\Delta_f : \mathcal{X} \to \mathcal{X} \times_\mathcal{Y} \mathcal{X}$
が全射であることを示す。そうすれば Lemma \ref{stacks-morphisms-lemma-universally-injective} により
(1) が従う。アフィンスキーム $V$ と滑らかな射 $V \to \mathcal{Y}$ を取る。
$g : V \times_\mathcal{Y} \mathcal{X} \to V$
は (2) により普遍単射なので、$\Delta_g$ は全射である。
一方、$\Delta_g$ は $\Delta_f$ の滑らかな射 $V \to \mathcal{Y}$ による基底変換である。
このような射 $V \to \mathcal{Y}$ の族は共同で全射なので、$\Delta_f$ は全射である。
\end{proof}

\begin{lemma}
\label{stacks-morphisms-lemma-check-universally-injective-covering}
$f : \mathcal{X} \to \mathcal{Y}$ を代数スタックの射とする。
$W$ を代数空間とし、$W \to \mathcal{Y}$ を全射、平坦かつ局所有限表示な射とする。
基底変換 $W \times_\mathcal{Y} \mathcal{X} \to W$ が普遍単射ならば、
$f$ は普遍単射である。
\end{lemma}

\begin{proof}
射
$g : W \times_\mathcal{Y} \mathcal{X} \to W$
の対角射 $\Delta_g$ は $\Delta_f$ の $W \to \mathcal{Y}$ による基底変換であることに注意する。
したがって、$\Delta_g$ が全射ならば、Properties of Stacks, Lemma
\ref{stacks-properties-lemma-check-property-covering}.
により $\Delta_f$ も全射である。よって補題は
Lemma \ref{stacks-morphisms-lemma-universally-injective} の特徴づけ (2) から従う。
\end{proof}








\section{普遍同相射}
\label{stacks-morphisms-section-universal-homeomorphisms}

\noindent
$f$ を代数空間によって表現可能な代数スタックの射とする。
Properties of Stacks, Section
\ref{stacks-properties-section-properties-morphisms}
で $f$ が普遍同相射であることの意味を定義した。別の特徴づけを与える。

\begin{lemma}
\label{stacks-morphisms-lemma-characterize-representable-universal-homeomorphism}
$f : \mathcal{X} \to \mathcal{Y}$ を代数空間によって表現可能な
代数スタックの射とする。次は同値である：
\begin{enumerate}
\item $f$ は普遍同相射である（Properties of Stacks,
Section \ref{stacks-properties-section-properties-morphisms}）。
\item 代数スタックの任意の射 $\mathcal{Z} \to \mathcal{Y}$ に対して、位相空間の写像
$|\mathcal{Z} \times_\mathcal{Y} \mathcal{X}| \to |\mathcal{Z}|$ は同相写像である。
\end{enumerate}
\end{lemma}

\begin{proof}
(1) を仮定し、$\mathcal{Z} \to \mathcal{Y}$ を (2) の射とする。
スキーム $V$ と全射な滑らかな射 $V \to \mathcal{Z}$ を選ぶ。
仮定により代数空間の射 $V \times_\mathcal{Y} \mathcal{X} \to V$ は普遍同相射であり、
特に写像 $|V \times_\mathcal{Y} \mathcal{X}| \to |V|$ は同相写像である。
Properties of Stacks, Section \ref{stacks-properties-section-points}
により、可換図式
$$
\xymatrix{
|V \times_\mathcal{Y} \mathcal{X}| \ar[r] \ar[d] &
|\mathcal{Z} \times_\mathcal{Y} \mathcal{X}| \ar[d] \\
|V| \ar[r] & |\mathcal{Z}|
}
$$
の横の矢印は開かつ全射であり、さらに
$$
|V \times_\mathcal{Y} \mathcal{X}| \longrightarrow
|V| \times_{|\mathcal{Z}|} |\mathcal{Z} \times_\mathcal{Y} \mathcal{X}|
$$
は全射である。したがって左の縦射が同相写像であることから、右の縦射も同相写像である。
これで (2) が証明された。含意 (2) $\Rightarrow$ (1) は定義から従う。
\end{proof}

\noindent
したがって、次の自然な定義を用いることができる。

\begin{definition}
\label{stacks-morphisms-definition-universal-homeomorphism}
$f : \mathcal{X} \to \mathcal{Y}$ を代数スタックの射とする。
代数スタックの任意の射 $\mathcal{Z} \to \mathcal{Y}$ に対して、位相空間の写像
$$
|\mathcal{Z} \times_\mathcal{Y} \mathcal{X}| \to |\mathcal{Z}|
$$
が同相写像であるとき、$f$ は{\it 普遍同相射}であるという。
\end{definition}

\begin{lemma}
\label{stacks-morphisms-lemma-base-change-universal-homeomorphism}
代数スタックの普遍同相射を代数スタックの任意の射で基底変換したものは普遍同相射である。
\end{lemma}

\begin{proof}
定義から直ちに従う。
\end{proof}

\begin{lemma}
\label{stacks-morphisms-lemma-composition-universal-homeomorphism}
代数スタックの二つの普遍同相射の合成は普遍同相射である。
\end{lemma}

\begin{proof}
省略する。
\end{proof}

\begin{lemma}
\label{stacks-morphisms-lemma-check-universal-homeomorphism-covering}
$f : \mathcal{X} \to \mathcal{Y}$ を代数スタックの射とする。
$W$ を代数空間とし、$W \to \mathcal{Y}$ を全射、平坦かつ局所有限表示な射とする。
基底変換 $W \times_\mathcal{Y} \mathcal{X} \to W$ が普遍同相射ならば、
$f$ は普遍同相射である。
\end{lemma}

\begin{proof}
$g : W \times_\mathcal{Y} \mathcal{X} \to W$ は普遍同相射であると仮定する。
このとき $g$ は普遍単射であり、したがって $f$ は次により普遍単射である：
Lemma \ref{stacks-morphisms-lemma-check-universally-injective-covering}.
一方、$\mathcal{Z}$ を代数スタックとし、$\mathcal{Z} \to \mathcal{Y}$ を射とする。
スキーム $U$ と全射な滑らかな射 $U \to W \times_\mathcal{Y} \mathcal{Z}$ を選ぶ。
図式
$$
\xymatrix{
W \times_\mathcal{Y} \mathcal{X} \ar[d]^g &
U \times_\mathcal{Y} \mathcal{X} \ar[d] \ar[l] \ar[r] &
\mathcal{Z} \times_\mathcal{Y} \mathcal{X} \ar[d] \\
W &
U \ar[l] \ar[r] &
\mathcal{Z}
}
$$
を考える。$g$ に関する仮定により、中央の縦射は位相空間上の同相写像を誘導する。
射 $U \to \mathcal{Z}$ と
$U \times_\mathcal{Y} \mathcal{X} \to
\mathcal{Z} \times_\mathcal{Y} \mathcal{X}$
は全射、平坦かつ局所有限表示なので、位相空間上の開写像を誘導する。
したがって
$|\mathcal{Z} \times_\mathcal{Y} \mathcal{X}| \to |\mathcal{Z}|$
は開写像である。全射性の証明は容易なので省略する。
\end{proof}










\section{始域と終域について滑らか局所的な射の型}
\label{stacks-morphisms-section-local-source-target}

\noindent
代数空間の射の性質で、{\it 始域と終域について滑らか局所的}なものが与えられたとする。次を参照せよ：
Descent on Spaces,
Definition \ref{spaces-descent-definition-local-source-target}
これを用いて代数スタックの射の対応する性質を定義できる。すなわち、次の補題の
同値な条件のいずれかを課せばよい。

\begin{lemma}
\label{stacks-morphisms-lemma-local-source-target}
$\mathcal{P}$ を、始域と終域について滑らか局所的な代数空間の射の性質とする。
$f : \mathcal{X} \to \mathcal{Y}$ を代数スタックの射とする。可換図式
$$
\xymatrix{
U \ar[d]_a \ar[r]_h & V \ar[d]^b \\
\mathcal{X} \ar[r]^f & \mathcal{Y}
}
$$
を考える。ここで $U$ と $V$ は代数空間で、縦射は滑らかである。次は同値である：
\begin{enumerate}
\item 上のような任意の図式で、さらに $U \to \mathcal{X} \times_\mathcal{Y} V$
が滑らかならば、射 $h$ は性質 $\mathcal{P}$ をもつ。
\item 上のようなある図式で、$a : U \to \mathcal{X}$ が全射かつ
射 $h$ が性質 $\mathcal{P}$ をもつ。
\end{enumerate}
$\mathcal{X}$ と $\mathcal{Y}$ が代数空間によって表現可能ならば、これは
（代数空間の射としての）$f$ が性質 $\mathcal{P}$ をもつこととも同値である。
$\mathcal{P}$ がさらに任意の基底変換で保たれ、基底上 fppf 局所的ならば、
代数空間によって表現可能な射 $f$ について、これは
Properties of Stacks,
Section \ref{stacks-properties-section-properties-morphisms}.
の意味で $f$ が性質 $\mathcal{P}$ をもつこととも同値である。
\end{lemma}

\begin{proof}
含意 (1) $\Rightarrow$ (2) を証明する。代数空間 $V$ と全射な滑らかな射
$V \to \mathcal{Y}$ を選ぶ。代数空間 $U$ と全射な滑らかな射
$U \to \mathcal{X} \times_\mathcal{Y} V$ を選ぶ。
$U \to \mathcal{X}$ も全射かつ滑らかであることに注意する。実際、これは基底変換
$\mathcal{X} \times_\mathcal{Y} V \to \mathcal{X}$ と、選んだ写像
$U \to \mathcal{X} \times_\mathcal{Y} V$ の合成である。したがって (1) のような
図式を得る。よって (1) が成り立つならば $h : U \to V$ は性質 $\mathcal{P}$ をもち、
$U \to \mathcal{X}$ は全射なので (2) が成り立つ。

\medskip\noindent
逆に (2) が成り立つと仮定し、$U, V, a, b, h$ を (2) のものとする。
次に $U', V', a', b', h'$ を (1) の任意の図式とする。図示すると
$$
\xymatrix{
U \ar[d] \ar[r]_h & V \ar[d] \\
\mathcal{X} \ar[r]^f & \mathcal{Y}
}
\quad\quad
\xymatrix{
U' \ar[d] \ar[r]_{h'} & V' \ar[d] \\
\mathcal{X} \ar[r]^f & \mathcal{Y}
}
$$
である。(2) が (1) を含意することを示すには、$h'$ が $\mathcal{P}$ をもつことを
証明しなければならない。そのために、代数空間の可換図式
$$
\xymatrix{
U \ar[dd]^h &
U \times_\mathcal{X} U' \ar[d] \ar[l] \ar@/^6ex/[dd]^{(h, h')} \ar[r] &
U' \ar[dd]^{h'} \\
& U \times_\mathcal{Y} V' \ar[lu] \ar[d] & \\
V &
V \times_\mathcal{Y} V' \ar[l] \ar[r] &
V'
}
$$
を考える。横の矢印は、滑らかな射 $V \to \mathcal{Y}$、
$V' \to \mathcal{Y}$、$U \to \mathcal{X}$、$U' \to \mathcal{X}$ の
基底変換として滑らかであることに注意する。また
$$
\xymatrix{
U \times_\mathcal{X} U' \ar[d] \ar[r] & U' \ar[d] \\
U \times_\mathcal{Y} V' \ar[r] & \mathcal{X} \times_\mathcal{Y} V'
}
$$
は Cartesian であり、$U', V', a', b', h'$ は (1) のものなので、左の縦射は滑らかである。
$\mathcal{P}$ は次により終域について滑らか局所的なので、
Descent on Spaces, Lemma
\ref{spaces-descent-lemma-local-source-target-implies} part (2)
基底変換 $U \times_\mathcal{Y} V' \to V \times_\mathcal{Y} V'$ は
$\mathcal{P}$ をもつ。$\mathcal{P}$ は次により始域について滑らか局所的なので、
Descent on Spaces, Lemma
\ref{spaces-descent-lemma-local-source-target-implies} part (1)
滑らかな射 $U \times_\mathcal{X} U' \to U \times_\mathcal{Y} V'$ を前合成して、
$(h, h')$ が $\mathcal{P}$ をもつと結論できる。
$V \times_\mathcal{Y} V' \to V'$ は滑らかなので、次により
$U \times_\mathcal{X} U' \to V'$ は $\mathcal{P}$ をもつ：
Descent on Spaces, Lemma
\ref{spaces-descent-lemma-local-source-target-implies} part (3).
最後に、$U \times_X U' \to U'$ は全射かつ滑らかであり、$\mathcal{P}$ は
始域について滑らか局所的なので（同じ補題）、$h'$ は $\mathcal{P}$ をもつ。
これで (1) と (2) の同値の証明が終わる。

\medskip\noindent
$\mathcal{X}$ と $\mathcal{Y}$ が表現可能ならば、
Descent on Spaces,
Lemma \ref{spaces-descent-lemma-local-source-target-characterize}
を適用でき、(1) と (2) は $f$ が $\mathcal{P}$ をもつことと同値である。

\medskip\noindent
最後に、$f$ は表現可能であり、$U, V, a, b, h$ は補題の (2) のもので、
$\mathcal{P}$ は任意の基底変換で保たれると仮定する。任意のスキーム $Z$ と
射 $Z \to \mathcal{Y}$ に対して、基底変換
$Z \times_\mathcal{Y} \mathcal{X} \to Z$
が性質 $\mathcal{P}$ をもつことを示さなければならない。図式
$$
\xymatrix{
Z \times_\mathcal{Y} U \ar[d] \ar[r] &
Z \times_\mathcal{Y} V \ar[d] \\
Z \times_\mathcal{Y} \mathcal{X} \ar[r] &
Z
}
$$
を考える。上の横射は $h$ の基底変換なので性質 $\mathcal{P}$ をもつ。
左の縦射は滑らかかつ全射で、右の縦射は滑らかである。したがって
Descent on Spaces,
Lemma \ref{spaces-descent-lemma-local-source-target-characterize}
を適用すれば、$Z \times_\mathcal{Y} \mathcal{X} \to Z$ は性質 $\mathcal{P}$ をもつ。
\end{proof}

\begin{definition}
\label{stacks-morphisms-definition-P}
$\mathcal{P}$ を、始域と終域について滑らか局所的な代数空間の射の性質とする。
Lemma \ref{stacks-morphisms-lemma-local-source-target} の同値な条件が成り立つとき、
代数スタックの射 $f : \mathcal{X} \to \mathcal{Y}$ は
{\it 性質 $\mathcal{P}$ をもつ}という。
\end{definition}

\begin{remark}
\label{stacks-morphisms-remark-composition}
$\mathcal{P}$ を、始域と終域について滑らか局所的で合成について安定な
代数空間の射の性質とする。このとき Definition \ref{stacks-morphisms-definition-P} で定義した
代数スタックの射の性質は合成について安定である。実際、
$f : \mathcal{X} \to \mathcal{Y}$ と $g : \mathcal{Y} \to \mathcal{Z}$ を、
性質 $\mathcal{P}$ をもつ代数スタックの射とする。代数空間 $W$ と全射な滑らかな射
$W \to \mathcal{Z}$ を選ぶ。代数空間 $V$ と全射な滑らかな射
$V \to \mathcal{Y} \times_\mathcal{Z} W$ を選ぶ。最後に代数空間 $U$ と
全射な滑らかな射 $U \to \mathcal{X} \times_\mathcal{Y} V$ を選ぶ。
すると定義により射 $V \to W$ と $U \to V$ は性質 $\mathcal{P}$ をもつ。
$\mathcal{P}$ は合成について安定と仮定したので、$U \to W$ は性質 $\mathcal{P}$ をもつ。
したがって再び定義により、$g \circ f : \mathcal{X} \to \mathcal{Z}$ は
性質 $\mathcal{P}$ をもつ。
\end{remark}

\begin{remark}
\label{stacks-morphisms-remark-base-change}
$\mathcal{P}$ を、始域と終域について滑らか局所的で基底変換について安定な
代数空間の射の性質とする。このとき Definition \ref{stacks-morphisms-definition-P} で定義した
代数スタックの射の性質は基底変換について安定である。実際、
$f : \mathcal{X} \to \mathcal{Y}$ と $g : \mathcal{Y}' \to \mathcal{Y}$ を
代数スタックの射とし、$f$ は性質 $\mathcal{P}$ をもつと仮定する。
代数空間 $V$ と全射な滑らかな射 $V \to \mathcal{Y}$ を選ぶ。
代数空間 $U$ と全射な滑らかな射 $U \to \mathcal{X} \times_\mathcal{Y} V$ を選ぶ。
最後に代数空間 $V'$ と全射な滑らかな射
$V' \to \mathcal{Y}' \times_\mathcal{Y} V$ を選ぶ。定義により射
$U \to V$ は性質 $\mathcal{P}$ をもつ。$\mathcal{P}$ は基底変換について
安定と仮定したので、$V' \times_V U \to V'$ は性質 $\mathcal{P}$ をもつ。
図式
$$
\xymatrix{
V' \times_V U \ar[r] \ar[d] &
\mathcal{Y}' \times_\mathcal{Y} \mathcal{X} \ar[r] \ar[d] &
\mathcal{X} \ar[d] \\
V' \ar[r] & \mathcal{Y}' \ar[r] & \mathcal{Y}
}
$$
を考えると、左上の横射は滑らかかつ全射であることが分かる。したがって定義により、
射影 $\mathcal{Y}' \times_\mathcal{Y} \mathcal{X} \to \mathcal{Y}'$ は
性質 $\mathcal{P}$ をもつ。
\end{remark}

\begin{remark}
\label{stacks-morphisms-remark-implication}
$\mathcal{P}, \mathcal{P}'$ を、始域と終域について滑らか局所的な
代数空間の射の性質とする。代数空間の射について
$\mathcal{P} \Rightarrow \mathcal{P}'$ であると仮定する。このとき
$\mathcal{P}$ と $\mathcal{P}'$ を用いて Definition \ref{stacks-morphisms-definition-P} で
定義した代数スタックの射の性質についても
$\mathcal{P} \Rightarrow \mathcal{P}'$ である。これは定義から明らかである。
\end{remark}









\section{有限型の射}
\label{stacks-morphisms-section-finite-type}

\noindent
代数空間の射が「局所有限型」であるという性質は、始域と終域について
滑らか局所的である。Descent on Spaces, Remark
\ref{spaces-descent-remark-list-local-source-target} を参照せよ。
また、この性質は基底変換について安定であり、終域について fpqc 局所的である。
Morphisms of Spaces,
Lemma \ref{spaces-morphisms-lemma-base-change-finite-type} および
Descent on Spaces,
Lemma \ref{spaces-descent-lemma-descending-property-locally-finite-type}
を参照せよ。したがって上の Lemma \ref{stacks-morphisms-lemma-local-source-target} により、
代数空間の射が局所有限型であることを次のように定義できる。この定義は、
射が代数空間によって表現可能な場合には、Properties of Stacks,
Section \ref{stacks-properties-section-properties-morphisms}
で定義された既存の概念と一致する。

\begin{definition}
\label{stacks-morphisms-definition-locally-finite-type}
代数スタックの射 $f : \mathcal{X} \to \mathcal{Y}$ をとる。
\begin{enumerate}
\item $\mathcal{P} = \text{locally of finite type}$ としたときに
Lemma \ref{stacks-morphisms-lemma-local-source-target} の同値な条件が成り立つならば、
$f$ は{\it 局所有限型}であるという。
\item $f$ が局所有限型かつ準コンパクトであるならば、
この射は{\it 有限型}であるという。
\end{enumerate}
\end{definition}

\begin{lemma}
\label{stacks-morphisms-lemma-composition-finite-type}
有限型の射の合成は有限型である。
局所有限型についても同じことが成り立つ。
\end{lemma}

\begin{proof}
Remark \ref{stacks-morphisms-remark-composition} と
Morphisms of Spaces, Lemma
\ref{spaces-morphisms-lemma-composition-finite-type} を組み合わせればよい。
\end{proof}

\begin{lemma}
\label{stacks-morphisms-lemma-base-change-finite-type}
有限型の射の基底変換は有限型である。
局所有限型についても同じことが成り立つ。
\end{lemma}

\begin{proof}
Remark \ref{stacks-morphisms-remark-base-change} と
Morphisms of Spaces, Lemma
\ref{spaces-morphisms-lemma-base-change-finite-type} を組み合わせればよい。
\end{proof}

\begin{lemma}
\label{stacks-morphisms-lemma-immersion-locally-finite-type}
埋入は局所有限型である。
\end{lemma}

\begin{proof}
Remark \ref{stacks-morphisms-remark-implication} と Morphisms of Spaces,
Lemma \ref{spaces-morphisms-lemma-immersion-locally-finite-type}
を組み合わせればよい。
\end{proof}

\begin{lemma}
\label{stacks-morphisms-lemma-locally-finite-type-locally-noetherian}
代数スタックの射 $f : \mathcal{X} \to \mathcal{Y}$ をとる。
$f$ が局所有限型で $\mathcal{Y}$ が局所 Noether ならば、
$\mathcal{X}$ は局所 Noether である。
\end{lemma}

\begin{proof}
可換図式
$$
\xymatrix{
U \ar[d] \ar[r] & V \ar[d] \\
\mathcal{X} \ar[r] & \mathcal{Y}
}
$$
をとる。ここで $U$, $V$ はスキームであり、$V \to \mathcal{Y}$ は
全射かつ滑らか、$U \to V \times_\mathcal{Y} \mathcal{X}$ は
全射かつ滑らかである。このとき $U \to V$ は局所有限型である。
$\mathcal{Y}$ が局所 Noether ならば、$V$ は局所 Noether である。
Morphisms, Lemma \ref{morphisms-lemma-finite-type-noetherian} により
$U$ は局所 Noether であり、これは $\mathcal{X}$ が局所 Noether であることを意味する。
\end{proof}

\noindent
次の二つの補題は、局所有限表示な代数スタックの射を論じた後で改良される。

\begin{lemma}
\label{stacks-morphisms-lemma-check-finite-type-covering}
代数スタックの射 $f : \mathcal{X} \to \mathcal{Y}$ をとる。
$W$ を代数空間とし、$W \to \mathcal{Y}$ を全射、平坦、かつ局所有限表示な
射とする。基底変換 $W \times_\mathcal{Y} \mathcal{X} \to W$ が
局所有限型ならば、$f$ は局所有限型である。
\end{lemma}

\begin{proof}
代数空間 $V$ と全射な滑らかな射 $V \to \mathcal{Y}$ を選ぶ。
代数空間 $U$ と全射な滑らかな射
$U \to V \times_\mathcal{Y} \mathcal{X}$ を選ぶ。
$U \to V$ が局所有限表示であることを示せばよい。
ここですべてを $W \to \mathcal{Y}$ によって基底変換し、
$U' = W \times_\mathcal{Y} U$,
$V' = W \times_\mathcal{Y} V$,
$\mathcal{X}' = W \times_\mathcal{Y} \mathcal{X}$,
$\mathcal{Y}' = W \times_\mathcal{Y} \mathcal{Y} = W$
とおく。基底変換により、$U' \to V' \times_{\mathcal{Y}'} \mathcal{X}'$
は依然として滑らかである。したがって、代数スタックの局所有限型の射の定義と
$\mathcal{X}' \to \mathcal{Y}'$ が局所有限型であるという仮定から、
$U' \to V'$ は局所有限型である。一方、$V' \to V$ は
$W \to \mathcal{Y}$ の基底変換として全射、平坦、かつ局所有限表示である。
ゆえに Descent on Spaces, Lemma
\ref{spaces-descent-lemma-descending-property-locally-finite-type}
により $U \to V$ は局所有限型であり、これで結論を得る。
\end{proof}

\begin{lemma}
\label{stacks-morphisms-lemma-check-finite-type-precompose}
代数スタックの射 $\mathcal{X} \to \mathcal{Y} \to \mathcal{Z}$ をとる。
$\mathcal{X} \to \mathcal{Z}$ は局所有限型であり、
$\mathcal{X} \to \mathcal{Y}$ は代数空間によって表現可能で、全射、平坦、かつ
局所有限表示であると仮定する。このとき
$\mathcal{Y} \to \mathcal{Z}$ は局所有限型である。
\end{lemma}

\begin{proof}
代数空間 $W$ と全射な滑らかな射 $W \to \mathcal{Z}$ を選ぶ。
代数空間 $V$ と全射な滑らかな射
$V \to W \times_\mathcal{Z} \mathcal{Y}$ を選ぶ。
$U = V \times_\mathcal{Y} \mathcal{X}$ とおくと、これは代数空間である。
$U \to V$ は全射、平坦、かつ局所有限表示であり、
$U \to W$ は局所有限型であることが分かっている。
したがって補題は代数空間の射の場合に帰着する。この場合は
Descent on Spaces, Lemma
\ref{spaces-descent-lemma-locally-finite-type-fppf-local-source}
である。
\end{proof}

\begin{lemma}
\label{stacks-morphisms-lemma-finite-type-permanence}
代数スタックの射 $f : \mathcal{X} \to \mathcal{Y}$ と
$g : \mathcal{Y} \to \mathcal{Z}$ をとる。
$g \circ f : \mathcal{X} \to \mathcal{Z}$ が局所有限型ならば、
$f : \mathcal{X} \to \mathcal{Y}$ は局所有限型である。
\end{lemma}

\begin{proof}
図式
$$
\xymatrix{
U \ar[r] \ar[d] & V \ar[r] \ar[d] & W \ar[d] \\
\mathcal{X} \ar[r] & \mathcal{Y} \ar[r] & \mathcal{Z}
}
$$
をとることができる。ここで $U$, $V$, $W$ はスキームであり、縦の射
$W \to \mathcal{Z}$ は全射かつ滑らか、射
$V \to \mathcal{Y} \times_\mathcal{Z} W$ は全射かつ滑らか、さらに射
$U \to \mathcal{X} \times_\mathcal{Y} V$ も全射かつ滑らかである。
このとき $U \to \mathcal{X} \times_\mathcal{Z} V$ も全射かつ滑らかである
（全射な滑らかな射と、そのような射の基底変換との合成だからである）。
定義により $U \to W$ は局所有限型である。したがって
Morphisms, Lemma \ref{morphisms-lemma-permanence-finite-type}
により $U \to V$ は局所有限型であり、これは定義により
$\mathcal{X} \to \mathcal{Y}$ が局所有限型であることを意味する。
\end{proof}










\section{有限型点}
\label{stacks-morphisms-section-points-finite-type}

\noindent
$\mathcal{X}$ を代数スタックとする。有限型点 $x \in |\mathcal{X}|$ とは、
局所有限型の射 $\Spec(k) \to \mathcal{X}$ によって表される点である。
代数空間および代数スタックにおいて、有限型点は閉点の適切な代替となる。
例えば、有限型点は常に「十分多く」存在する。

\begin{lemma}
\label{stacks-morphisms-lemma-point-finite-type}
$\mathcal{X}$ を代数スタックとし、$x \in |\mathcal{X}|$ とする。
次の条件は同値である。
\begin{enumerate}
\item 局所有限型で $x$ を表す射 $\Spec(k) \to \mathcal{X}$ が存在する。
\item スキーム $U$、閉点 $u \in U$、および
$\varphi(u) = x$ を満たす滑らかな射 $\varphi : U \to \mathcal{X}$ が存在する。
\end{enumerate}
\end{lemma}

\begin{proof}
(2) のような $u \in U$ と $U \to \mathcal{X}$ をとる。このとき
$\Spec(\kappa(u)) \to U$ は有限型であり、$U \to \mathcal{X}$ は
表現可能かつ局所有限型である（Morphisms of Spaces,
Lemmas \ref{spaces-morphisms-lemma-etale-locally-finite-presentation} および
\ref{spaces-morphisms-lemma-finite-presentation-finite-type} による）。
したがって Lemma \ref{stacks-morphisms-lemma-composition-finite-type} により (1) が成り立つ。

\medskip\noindent
逆に、$\Spec(k) \to \mathcal{X}$ が局所有限型で $x$ を表すと仮定する。
$U$ をスキームとし、$U \to \mathcal{X}$ を全射な滑らかな射とする。
仮定により $U \times_\mathcal{X} \Spec(k) \to U$ は局所有限型の
代数空間の射である。$U \times_\mathcal{X} \Spec(k)$ の有限型点 $v$ を選ぶ
（少なくとも一つ存在する。Morphisms of Spaces,
Lemma \ref{spaces-morphisms-lemma-identify-finite-type-points} を参照せよ）。
Morphisms of Spaces,
Lemma \ref{spaces-morphisms-lemma-finite-type-points-morphism} により、
$v$ の像 $u \in U$ は $U$ の有限型点である。したがって
Morphisms, Lemma \ref{morphisms-lemma-identify-finite-type-points} により、
$U$ を縮小した後で $u$ が $U$ の閉点であると仮定できる。すなわち (2) が成り立つ。
\end{proof}

\begin{definition}
\label{stacks-morphisms-definition-finite-type-point}
$\mathcal{X}$ を代数スタックとする。点 $x \in |\mathcal{X}|$ が
Lemma \ref{stacks-morphisms-lemma-point-finite-type} の同値な条件を満たすとき、
この点を{\it 有限型点}\footnote{「局所有限型点」と呼ぶ方が正確かもしれないので、
これは若干の用語の濫用である。}という。
$\mathcal{X}$ の有限型点全体の集合を $\mathcal{X}_{\text{ft-pts}}$ と書く。
\end{definition}

\noindent
有限型点の集合は次のように記述できる。

\begin{lemma}
\label{stacks-morphisms-lemma-identify-finite-type-points}
$\mathcal{X}$ を代数スタックとする。このとき
$$
\mathcal{X}_{\text{ft-pts}} =
\bigcup\nolimits_{\varphi : U \to \mathcal{X}\text{ smooth}} |\varphi|(U_0)
$$
である。ここで $U_0$ は $U$ の閉点全体の集合である。
このとき $U$ は $\mathcal{X}$ 上滑らかなすべてのスキームを動いてもよく、
$\mathcal{X}$ 上滑らかなすべてのアフィンスキームを動いてもよい。
\end{lemma}

\begin{proof}
Lemma \ref{stacks-morphisms-lemma-point-finite-type} から直ちに従う。
\end{proof}

\begin{lemma}
\label{stacks-morphisms-lemma-finite-type-points-morphism}
代数スタックの射 $f : \mathcal{X} \to \mathcal{Y}$ をとる。
$f$ が局所有限型ならば
$f(\mathcal{X}_{\text{ft-pts}}) \subset \mathcal{Y}_{\text{ft-pts}}$.
である。
\end{lemma}

\begin{proof}
$x \in \mathcal{X}_{\text{ft-pts}}$ をとる。$x$ を局所有限型の射
$x : \Spec(k) \to \mathcal{X}$ によって表す。
Lemma \ref{stacks-morphisms-lemma-composition-finite-type} により $f \circ x$ は局所有限型である。
したがって $f(x) \in \mathcal{Y}_{\text{ft-pts}}$ である。
\end{proof}

\begin{lemma}
\label{stacks-morphisms-lemma-finite-type-points-surjective-morphism}
代数スタックの射 $f : \mathcal{X} \to \mathcal{Y}$ をとる。
$f$ が局所有限型かつ全射ならば
$f(\mathcal{X}_{\text{ft-pts}}) = \mathcal{Y}_{\text{ft-pts}}$.
である。
\end{lemma}

\begin{proof}
Lemma \ref{stacks-morphisms-lemma-finite-type-points-morphism} により
$f(\mathcal{X}_{\text{ft-pts}}) \subset \mathcal{Y}_{\text{ft-pts}}$
である。$y \in |\mathcal{Y}|$ を有限型点とし、局所有限型の射
$\Spec(k) \to \mathcal{Y}$ によって $y$ を表す。
$f$ は全射なので、代数スタック
$\mathcal{X}_k = \Spec(k) \times_\mathcal{Y} \mathcal{X}$ は空でない。
したがって Lemma \ref{stacks-morphisms-lemma-identify-finite-type-points} により、
これは有限型点 $x \in |\mathcal{X}_k|$ をもつ。
ここで $\mathcal{X}_k \to \mathcal{X}$ は
$\Spec(k) \to \mathcal{Y}$ の基底変換として局所有限型である
（Lemma \ref{stacks-morphisms-lemma-base-change-finite-type}）。
ゆえに Lemma \ref{stacks-morphisms-lemma-finite-type-points-morphism} により、
$\mathcal{X}$ における $x$ の像は有限型点であり、構成により $y$ に写る。
\end{proof}

\begin{lemma}
\label{stacks-morphisms-lemma-enough-finite-type-points}
$\mathcal{X}$ を代数スタックとする。任意の局所閉部分集合
$T \subset |\mathcal{X}|$ に対して
$$
T \not = \emptyset
\Rightarrow
T \cap \mathcal{X}_{\text{ft-pts}} \not = \emptyset.
$$
である。特に、任意の閉部分集合 $T \subset |\mathcal{X}|$ に対して
$T \cap \mathcal{X}_{\text{ft-pts}}$ は $T$ で稠密である。
\end{lemma}

\begin{proof}
$i : \mathcal{Z} \to \mathcal{X}$ を $T$ 上の被約誘導部分スタック構造とする。
Properties of Stacks, Remark
\ref{stacks-properties-remark-stack-structure-locally-closed-subset}
を参照せよ。埋入は局所有限型である。Lemma
\ref{stacks-morphisms-lemma-immersion-locally-finite-type} を参照せよ。したがって
Lemma \ref{stacks-morphisms-lemma-finite-type-points-morphism} により
$\mathcal{Z}_{\text{ft-pts}} \subset \mathcal{X}_{\text{ft-pts}} \cap T$.
である。最後に、$\mathcal{Z}$ への滑らかな射をもつ空でない任意の
アフィンスキーム $U$ は少なくとも一つの閉点をもつ。したがって
Lemma \ref{stacks-morphisms-lemma-identify-finite-type-points} により、
$\mathcal{Z}$ は少なくとも一つの有限型点をもつ。これで補題が従う。
\end{proof}

\noindent
ここで、代数スタック上の有限型点について、もう一つのより技術的な特徴付けを与える。
これは特に、$x$ が有限型点ならば $\mathcal{X}$ の $x$ における
剰余ガーベが存在することを示す。

\begin{lemma}
\label{stacks-morphisms-lemma-point-finite-type-monomorphism}
$\mathcal{X}$ を代数スタックとし、$x \in |\mathcal{X}|$ とする。
次の条件は同値である。
\begin{enumerate}
\item $x$ は有限型点である。
\item 基礎位相空間 $|\mathcal{Z}|$ が一点集合である代数スタック
$\mathcal{Z}$ と、$\{x\} = |f|(|\mathcal{Z}|)$ を満たす局所有限型の射
$f : \mathcal{Z} \to \mathcal{X}$ が存在する。
\item $\mathcal{X}$ の $x$ における剰余ガーベ $\mathcal{Z}_x$ が存在し、
包含射 $\mathcal{Z}_x \to \mathcal{X}$ は局所有限型である。
\end{enumerate}
\end{lemma}

\begin{proof}
（この段落に現れる射はすべて代数空間によって表現可能なので、
Properties of Stacks,
Section \ref{stacks-properties-section-properties-morphisms}
の規約と結果を適用できる。）
$x$ が有限型点であると仮定する。アフィンスキーム $U$、閉点
$u \in U$、および $\varphi(u) = x$ を満たす滑らかな射
$\varphi : U \to \mathcal{X}$ を選ぶ。Lemma
\ref{stacks-morphisms-lemma-identify-finite-type-points} を参照せよ。
通常どおり $u = \Spec(\kappa(u))$ とおく。
$R = u \times_\mathcal{X} u$ とおくと、代数空間における亜群
$(u, R, s, t, c)$ を得る。Algebraic Stacks, Lemma
\ref{algebraic-lemma-map-space-into-stack} を参照せよ。
射影 $R \to u$ は合成
$$
R = u \times_\mathcal{X} u \to
u \times_\mathcal{X} U \to
u \times_\mathcal{X} X = u
$$
である。ここで最初の矢印は有限型（スキームの閉埋入 $u \to U$ の基底変換）であり、
第二の矢印は滑らか（滑らかな射 $U \to \mathcal{X}$ の基底変換）である。
したがって $s, t : R \to u$ は局所有限型である
（合成であることと Morphisms of Spaces, Lemma
\ref{spaces-morphisms-lemma-composition-finite-type} を参照せよ）。
$u$ は体のスペクトルなので、Morphisms of Spaces, Lemma
\ref{spaces-morphisms-lemma-noetherian-finite-type-finite-presentation}
により $s, t$ は平坦かつ局所有限表示である。Criteria for Representability,
Theorem \ref{criteria-theorem-flat-groupoid-gives-algebraic-stack} により
$\mathcal{Z} = [u/R]$ は代数スタックである。
Algebraic Stacks, Lemma \ref{algebraic-lemma-map-space-into-stack} により、
標準的な射
$$
f : \mathcal{Z} \longrightarrow \mathcal{X}
$$
を得て、これは充満忠実である。したがってこの射は代数空間によって表現可能であり
（Algebraic Stacks, Lemma
\ref{algebraic-lemma-characterize-representable-by-algebraic-spaces}）、
単射である（Properties of Stacks, Lemma
\ref{stacks-properties-lemma-monomorphism}）。
ゆえに $\mathcal{X}$ の $x$ における剰余ガーベ
$\mathcal{Z}_x \subset \mathcal{X}$ が存在し、$f$ は同値
$\mathcal{Z} \to \mathcal{Z}_x$ を経由する。Properties of Stacks,
Lemma \ref{stacks-properties-lemma-residual-gerbe-unique} を参照せよ。
構成により図式
$$
\xymatrix{
u \ar[d] \ar[r] & U \ar[d] \\
\mathcal{Z} \ar[r]^f & \mathcal{X}
}
$$
は可換である。Criteria for Representability,
Lemma \ref{criteria-lemma-flat-quotient-flat-presentation} により、
左の縦射は全射、平坦、かつ局所有限表示である。図式
$$
\xymatrix{
u \times_\mathcal{X} U \ar[d] \ar[r] &
\mathcal{Z} \times_\mathcal{X} U \ar[r] \ar[d] & U \ar[d] \\
u \ar[r] & \mathcal{Z} \ar[r]^f & \mathcal{X}
}
$$
を考える。$u \to \mathcal{X}$ は局所有限型なので、その基底変換
$u \times_\mathcal{X} U \to U$ は局所有限型である。また、
$u \times_\mathcal{X} U \to \mathcal{Z} \times_\mathcal{X} U$ は
$u \to \mathcal{Z}$ の基底変換として全射、平坦、かつ局所有限表示である。
したがって
$\{u \times_\mathcal{X} U \to \mathcal{Z} \times_\mathcal{X} U\}$
は代数空間の fppf 被覆であり、Descent on Spaces, Lemma
\ref{spaces-descent-lemma-locally-finite-presentation-fppf-local-source}.
により $\mathcal{Z} \times_\mathcal{X} U \to U$ は局所有限型である。
定義により、これは $f$ が局所有限型であることを意味する
（実際、縦射 $\mathcal{Z} \times_\mathcal{X} U \to \mathcal{Z}$ は
$U \to \mathcal{X}$ の基底変換として滑らかであり、$\mathcal{Z}$ は一点しか
もたないので全射である）。$\mathcal{Z} = \mathcal{Z}_x$ なので (3) が成り立つ。

\medskip\noindent
(3) が (2) を含意することは明らかである。(2) が成り立つならば、
Lemma \ref{stacks-morphisms-lemma-finite-type-points-morphism} と、
$\mathcal{Z}_{\text{ft-pts}}$ が空でないこと、すなわち $\mathcal{Z}$ の唯一の点が
$\mathcal{Z}$ の有限型点であることを示す
Lemma \ref{stacks-morphisms-lemma-enough-finite-type-points} により、$x$ は
$\mathcal{X}$ の有限型点である。
\end{proof}







\section{自己同型群}
\label{stacks-morphisms-section-automorphism-groups}

\noindent
$\mathcal{X}$ を代数スタックとする。$x \in |\mathcal{X}|$ が
$x : \Spec(k) \to \mathcal{X}$ に対応するとする。この状況ではしばしば
「$x$ の自己同型群代数空間を $G_x/k$ とする」という言い方を用いる。
これは単に
$$
G_x = \mathit{Isom}_\mathcal{X}(x, x) =
\Spec(k) \times_\mathcal{X} \mathcal{I}_\mathcal{X}
$$
が $x$ の自己同型の群代数空間であることを意味する。これは
$\Spec(k)$ 上の群代数空間である。$k'/k$ が体の拡大ならば、誘導された射
$x' : \Spec(k') \to \mathcal{X}$ の自己同型群代数空間は、
$G_x$ の $\Spec(k')$ への基底変換である。

\begin{lemma}
\label{stacks-morphisms-lemma-automorphism-group-scheme}
上の状況で、次のいずれかが成り立つならば $G_x$ はスキームである。
\begin{enumerate}
\item $\Delta : \mathcal{X} \to \mathcal{X} \times \mathcal{X}$ が準分離的である。
\item $\Delta : \mathcal{X} \to \mathcal{X} \times \mathcal{X}$ が局所分離的である。
\item $\mathcal{X}$ が準 DM である。
\item $\mathcal{I}_\mathcal{X} \to \mathcal{X}$ が準分離的である。
\item $\mathcal{I}_\mathcal{X} \to \mathcal{X}$ が局所分離的である。
\item $\mathcal{I}_\mathcal{X} \to \mathcal{X}$ が局所準有限である。
\end{enumerate}
\end{lemma}

\begin{proof}
Lemma \ref{stacks-morphisms-lemma-diagonal-diagonal} により (1) $\Rightarrow$ (4)、(2)
$\Rightarrow$ (5)、(3) $\Rightarrow$ (6) であることに注意する。
(4) の場合 $G_x$ は準分離的な代数空間であり、(5) の場合 $G_x$ は
局所分離的な代数空間であることが分かる。いずれの場合も $G_x$ は良い代数空間である
(Decent Spaces, Section \ref{decent-spaces-section-reasonable-decent} および
Lemma \ref{decent-spaces-lemma-locally-separated-decent})。
したがって More on Groupoids in Spaces, Lemma
\ref{spaces-more-groupoids-lemma-group-scheme-over-field-separated} により
$G_x$ は分離的であり、More on Groupoids in Spaces, Proposition
\ref{spaces-more-groupoids-proposition-group-space-scheme-over-field} により
$G_x$ はスキームである。(6) の場合 $G_x \to \Spec(k)$ は局所準有限である。
したがって Spaces over Fields, Lemma
\ref{spaces-over-fields-lemma-locally-quasi-finite-over-field} により $G_x$ はスキームである。
\end{proof}

\begin{lemma}
\label{stacks-morphisms-lemma-property-automorphism-groups}
代数スタック $\mathcal{X}$ と点 $x \in |\mathcal{X}|$ をとる。
$P$ を体上の代数空間の性質で、基礎体の拡大に関して不変なものとする。例えば
$P(X/k) = X \to \Spec(k)\text{ is finite}$ とする。
次の条件は同値である。
\begin{enumerate}
\item $x$ の同値類に属するある射 $x : \Spec(k) \to \mathcal{X}$ について、
自己同型群代数空間 $G_x/k$ が $P$ をもつ。
\item $x$ の同値類に属する任意の射 $x : \Spec(k) \to \mathcal{X}$ について、
自己同型群代数空間 $G_x/k$ が $P$ をもつ。
\end{enumerate}
\end{lemma}

\begin{proof}
証明は省略する。
\end{proof}

\begin{remark}
\label{stacks-morphisms-remark-property-automorphism-groups}
$P$ を体上の代数空間の性質で、基礎体の拡大に関して不変なものとする。
代数スタック $\mathcal{X}$ と $x \in |\mathcal{X}|$ に対し、
Lemma \ref{stacks-morphisms-lemma-property-automorphism-groups} の同値な条件が成り立つならば、
$\mathcal{X}$ の $x$ における自己同型群は $P$ をもつという。
例えば、$x : \Spec(k) \to \mathcal{X}$ が $x$ の代表であるとき常に
$G_x \to \Spec(k)$ が有限ならば、{\it $\mathcal{X}$ の $x$ における
自己同型群は有限}であるという。滑らか、固有などについても同様である。
（ここでは明らかに用語を濫用しているが、混乱や不正確さは生じないと考える。）
\end{remark}

\begin{lemma}
\label{stacks-morphisms-lemma-iso-automorphism-groups}
代数スタックの射 $f : \mathcal{X} \to \mathcal{Y}$ と点
$x \in |\mathcal{X}|$ をとる。次の条件は同値である。
\begin{enumerate}
\item $x$ の同値類に属するある射 $x : \Spec(k) \to \mathcal{X}$ について
$y = f \circ x$ とおくと、自己同型群代数空間の射
$G_x \to G_y$ は同型である。
\item $x$ の同値類に属する任意の射 $x : \Spec(k) \to \mathcal{X}$ について
$y = f \circ x$ とおくと、自己同型群代数空間の射
$G_x \to G_y$ は同型である。
\end{enumerate}
\end{lemma}

\begin{proof}
これは同型であることが終域について fpqc 局所的であることから従う。
Descent on Spaces, Lemma
\ref{spaces-descent-lemma-descending-property-isomorphism} を参照せよ。
すなわち $k'/k$ を体の拡大とし、合成
$x' : \Spec(k') \to \mathcal{X}$ を考え、$y' = f \circ x'$ とおく。
このとき $G_{x'} \to G_{y'}$ は $G_x \to G_y$ の
$\Spec(k') \to \Spec(k)$ による基底変換である。したがって
$G_x \to G_y$ が同型であることと $G_{x'} \to G_{y'}$ が同型であることは同値である。
ゆえに、ある一つについて成り立つならば、この性質は同値類全体に伝播する。
\end{proof}

\begin{remark}
\label{stacks-morphisms-remark-identify-automorphism-groups}
代数スタックの射 $f : \mathcal{X} \to \mathcal{Y}$ と点
$x \in |\mathcal{X}|$ をとる。Lemma
\ref{stacks-morphisms-lemma-iso-automorphism-groups} の同値な条件が $f$ と $x$ について
成り立つことを示すため、文献では{\it $f$ は $x$ で安定化群を保つ}
または {\it $f$ は $x$ で固定点を反映する} という用語を用いる。
我々は{\it $f$ は $x$ と $f(x)$ における自己同型群の間の同型を誘導する}
と言う方を好む。
\end{remark}





\section{代数スタックの表示と性質}
\label{stacks-morphisms-section-presentations-properties}

\noindent
$(U, R, s, t, c)$ を代数空間における亜群とする。
$s, t : R \to U$ が平坦かつ局所有限表示ならば、商スタック
$[U/R]$ は代数スタックである。Criteria for Representability, Theorem
\ref{criteria-theorem-flat-groupoid-gives-algebraic-stack} を参照せよ。
本節では、$(U, R, s, t, c)$ のどの性質が代数スタック $[U/R]$ に
どのような性質をもたらすかを調べる。

\begin{lemma}
\label{stacks-morphisms-lemma-properties-diagonal-from-presentation}
$(U, R, s, t, c)$ を、$s, t : R \to U$ が平坦かつ局所有限表示である
代数空間における亜群とする。上で述べた代数スタック
$\mathcal{X} = [U/R]$ を考える。
\begin{enumerate}
\item $R \to U \times U$ が分離的ならば、$\Delta_\mathcal{X}$ は分離的である。
\item $U$, $R$ が分離的ならば、$\Delta_\mathcal{X}$ は分離的である。
\item $R \to U \times U$ が局所準有限ならば、$\mathcal{X}$ は準 DM である。
\item $s, t : R \to U$ が局所準有限ならば、$\mathcal{X}$ は準 DM である。
\item $R \to U \times U$ が固有ならば、$\mathcal{X}$ は分離的である。
\item $s, t : R \to U$ が固有で $U$ が分離的ならば、$\mathcal{X}$ は分離的である。
\item ここにさらに追加する。
\end{enumerate}
\end{lemma}

\begin{proof}
射 $U \to \mathcal{X}$ は全射、平坦、かつ局所有限表示である。
Criteria for Representability, Lemma
\ref{criteria-lemma-flat-quotient-flat-presentation} による。
したがって $U \times U \to \mathcal{X} \times \mathcal{X}$ についても同じである。
次の Cartesian 図式を得る。
$$
\xymatrix{
R  = U \times_\mathcal{X} U \ar[r] \ar[d] & U \times U \ar[d] \\
\mathcal{X} \ar[r] & \mathcal{X} \times \mathcal{X}
}
$$
(Groupoids in Spaces, Lemma
\ref{spaces-groupoids-lemma-quotient-stack-2-cartesian} を参照せよ。)
したがって $\Delta_\mathcal{X}$ が Properties of Stacks, Section
\ref{stacks-properties-section-properties-morphisms} に列挙された性質の一つを
もつことと、射 $R \to U \times U$ がそれをもつことは同値である。
Properties of Stacks, Lemma
\ref{stacks-properties-lemma-check-property-covering} を参照せよ。
これにより (1), (3), (5) が成り立つ理由が分かる。(2) の条件は
$R \to U \times U$ が分離的であることを含意するので、(2) は (1) から従う。
(4) の条件は (3) の条件を含意するので、(4) は (3) から従う。
(6) の条件は Morphisms of Spaces, Lemma
\ref{spaces-morphisms-lemma-universally-closed-permanence} により
(5) の条件を含意するので、(6) は (5) から従う。
\end{proof}

\begin{lemma}
\label{stacks-morphisms-lemma-points-presentation}
$(U, R, s, t, c)$ を、$s, t : R \to U$ が平坦かつ局所有限表示である
代数空間における亜群とする。上で述べた代数スタック
$\mathcal{X} = [U/R]$ を考える。このとき
$|R| \to |U| \times |U|$ の像は同値関係であり、$|\mathcal{X}|$ は
この同値関係による $|U|$ の商である。
\end{lemma}

\begin{proof}
誘導された射 $p : U \to \mathcal{X}$ は全射、平坦、かつ局所有限表示である。
Criteria for Representability, Lemma
\ref{criteria-lemma-flat-quotient-flat-presentation} を参照せよ。
したがって Properties of Stacks, Lemma
\ref{stacks-properties-lemma-characterize-surjective} により
$|U| \to |\mathcal{X}|$ は全射である。$R = U \times_\mathcal{X} U$
であることに注意する。Groupoids in Spaces, Lemma
\ref{spaces-groupoids-lemma-quotient-stack-2-cartesian} を参照せよ。
したがって Properties of Stacks, Lemma
\ref{stacks-properties-lemma-points-cartesian} により射
$$
|R| \longrightarrow |U| \times_{|\mathcal{X}|} |U|
$$
は全射である。ゆえに $|R| \to |U| \times |U|$ の像は、
組 $(u_1, u_2) \in |U| \times |U|$ のうち、$u_1$ と $u_2$ の
$|\mathcal{X}|$ における像が同じであるものの集合に正確に一致する。
この二つの主張を合わせれば補題の結論を得る。
\end{proof}






\section{代数スタックの特殊な表示}
\label{stacks-morphisms-section-presentations}

\noindent
本節では二つの重要な定理を証明する。第一は、準 DM スタック
$\mathcal{X}$ が、$s, t : R \to U$ が局所準有限（さらに平坦かつ局所有限表示）
である形 $\mathcal{X} = [U/R]$ のスタックとして特徴付けられることである。
第二は、DM 代数スタックが Deligne--Mumford であるという主張である。

\medskip\noindent
次の補題は、表示の「切片」がなお代数スタック上平坦であるための判定条件を与える。

\begin{lemma}
\label{stacks-morphisms-lemma-slice}
$\mathcal{X}$ を代数スタックとする。次の Cartesian 図式を考える。
$$
\xymatrix{
U \ar[d] & F \ar[l]^p \ar[d] \\
\mathcal{X} & \Spec(k) \ar[l]
}
$$
ここで $U$ は代数空間、$k$ は体であり、$U \to \mathcal{X}$ は平坦かつ
局所有限表示とする。$f_1, \ldots, f_r \in \Gamma(U, \mathcal{O}_U)$
および $z \in |F|$ を、$f_1, \ldots, f_r$ が局所環
$\mathcal{O}_{F, \overline{z}}$ における正則列へ写るようにとる。
$p(z)$ を含む開部分空間で $U$ を置き換えた後、射
$$
V(f_1, \ldots, f_r) \longrightarrow \mathcal{X}
$$
は平坦かつ局所有限表示である。
\end{lemma}

\begin{proof}
スキーム $W$ と全射な滑らかな射 $W \to \mathcal{X}$ を選ぶ。
体の拡大 $k'/k$ と射 $w : \Spec(k') \to W$ を選び、
$\Spec(k') \to W \to \mathcal{X}$ が
$\Spec(k') \to \Spec(k) \to \mathcal{X}$ と $2$-同型になるようにする。
$W \to \mathcal{X}$ は全射なので、これは可能である。可換図式
$$
\xymatrix{
U \ar[d] &
U \times_\mathcal{X} W \ar[l]^-{\text{pr}_0} \ar[d] &
F' \ar[l]^-{p'} \ar[d] \\
\mathcal{X} &
W \ar[l] &
\Spec(k') \ar[l]
}
$$
二つの正方形は Cartesian である。$w$ の選び方により
$F' = F \times_{\Spec(k)} \Spec(k')$ である。したがって $F' \to F$ は全射であり、
$z$ に写る点 $z' \in |F'|$ を選べる。$F' \to F$ は平坦なので
$\mathcal{O}_{F, \overline{z}} \to \mathcal{O}_{F', \overline{z}'}$ は平坦である。
Morphisms of Spaces, Lemma
\ref{spaces-morphisms-lemma-flat-at-point-etale-local-rings} を参照せよ。
したがって $f_1, \ldots, f_r$ は
$\mathcal{O}_{F', \overline{z}'}$ における正則列へ写る。
Algebra, Lemma \ref{algebra-lemma-flat-increases-depth} を参照せよ。
$U \times_\mathcal{X} W \to W$ は平坦かつ局所有限表示の代数空間の射である。
したがって More on Morphisms of Spaces, Lemma
\ref{spaces-more-morphisms-lemma-slice} により、$p'(z')$ を含む
$U \times_\mathcal{X} W$ の開部分空間 $U'$ であって、交わり
$U' \cap (V(f_1, \ldots, f_r) \times_\mathcal{X} W)$ が $W$ 上平坦かつ
局所有限表示となるものが存在する。$\text{pr}_0(U')$ は $p(z)$ を含む
$U$ の開部分空間である。$\text{pr}_0$ は滑らか、したがって開だからである。
ここで
$U' \cap (V(f_1, \ldots, f_r) \times_\mathcal{X} W) \to \mathcal{X}$
は合成
$$
U' \cap (V(f_1, \ldots, f_r) \times_\mathcal{X} W) \to W \to \mathcal{X}.
$$
であるから平坦かつ局所有限表示である。したがって Properties of Stacks,
Lemma \ref{stacks-properties-lemma-check-property-after-precomposing} により
$\text{pr}_0(U') \cap V(f_1, \ldots, f_r) \to \mathcal{X}$ も平坦かつ
局所有限表示であり、所望の結論を得る。
\end{proof}

\begin{lemma}
\label{stacks-morphisms-lemma-quasi-finite-at-point}
$\mathcal{X}$ を代数スタックとする。次の Cartesian 図式を考える。
$$
\xymatrix{
U \ar[d] & F \ar[l]^p \ar[d] \\
\mathcal{X} & \Spec(k) \ar[l]
}
$$
ここで $U$ は代数空間、$k$ は体であり、$U \to \mathcal{X}$ は局所有限型とする。
$\dim_z(F) = 0$ を満たす $z \in |F|$ をとる。$p(z)$ を含む開部分空間で
$U$ を置き換えた後、射
$$
U \longrightarrow \mathcal{X}
$$
は局所準有限である。
\end{lemma}

\begin{proof}
射 $f : U \to \mathcal{X}$ は局所有限型なので、制限
$f|_{W(f)} : W(f) \to \mathcal{X}$ が局所準有限となる最大の開集合
$W(f) \subset U$ が存在する。Properties of Stacks, Remark
\ref{stacks-properties-remark-local-source-apply}
(\ref{stacks-properties-item-loc-quasi-finite}) を参照せよ。
したがって、$p(z)$ が $W(f)$ の点であることを示せばよい。
さらに、上記の Remark は、代数空間によって表現可能な射による任意の
基底変換と $W(f)$ の形成が可換することも示す。ゆえに、$F \to \Spec(k)$ が
$z$ において局所準有限であることを示せば十分である。これは直ちに
Morphisms of Spaces, Lemma
\ref{spaces-morphisms-lemma-locally-quasi-finite-rel-dimension-0} から従う。
\end{proof}

\noindent
準 DM スタックは、スキームによる局所準有限な「被覆」をもつ。

\begin{theorem}
\label{stacks-morphisms-theorem-quasi-DM}
$\mathcal{X}$ を代数スタックとする。次の条件は同値である。
\begin{enumerate}
\item $\mathcal{X}$ は準 DM である。
\item スキーム $W$ と、全射、平坦、局所有限表示、かつ局所準有限な射
$W \to \mathcal{X}$ が存在する。
\end{enumerate}
\end{theorem}

\begin{proof}
(2) $\Rightarrow$ (1) は Lemma \ref{stacks-morphisms-lemma-properties-covering-imply-diagonal} である。
 (1) を仮定する。$x \in |\mathcal{X}|$ を有限型点とする。
$x$ の近傍で「機能する」$\mathcal{X}$ 上のスキームを構成する。
証明の最後に、これらすべての非交和をとって結論を得る。

\medskip\noindent
$U$ をアフィンスキーム、$U \to \mathcal{X}$ を滑らかな射とし、
$u \in U$ を $x$ に写る閉点とする。Lemma
\ref{stacks-morphisms-lemma-point-finite-type} を参照せよ。通常どおり
$u = \Spec(\kappa(u))$ と書く。次の可換図式を考える。
$$
\xymatrix{
u \ar[d] & R \ar[l] \ar[d] \\
U \ar[d] & F \ar[d] \ar[l]^p \\
\mathcal{X} & u \ar[l]
}
$$
二つの正方形はいずれもファイバー積正方形であり、特に
$R = u \times_\mathcal{X} u$ である。Lemma
\ref{stacks-morphisms-lemma-point-finite-type-monomorphism} の証明で見たように、
$(u, R, s, t, c)$ は $s, t$ が局所有限型である代数空間における亜群である。
$G \to u$ を安定化群代数空間とする（
Groupoids in Spaces, Definition
\ref{spaces-groupoids-definition-stabilizer-groupoid} を参照せよ）。
次に注意する：
$$
G = R \times_{(u \times u)} u =
(u \times_\mathcal{X} u) \times_{(u \times u)} u =
\mathcal{X} \times_{\mathcal{X} \times \mathcal{X}} u.
$$
 $\mathcal{X}$ は準 DM なので $G$ は $u$ 上局所準有限である。また
More on Groupoids in Spaces, Lemma
\ref{spaces-more-groupoids-lemma-groupoid-on-field-dimension-equal-stabilizer}
により $\dim(R) = 0$ である。

\medskip\noindent
$e : u \to R$ を亜群の単位射とする。合成 $u \to R \to u$ は
$u$ の恒等射に等しい。$u \subset U$ は閉部分スキームなので
$R \subset F$ は閉部分空間である。したがって $e$ を $F$ の点とみなすこともでき、
これを $z$ と呼ぶ。エタール局所環の射
$$
\mathcal{O}_{U, u}
\xrightarrow{p^\sharp}
\mathcal{O}_{F, \overline{z}}
\longrightarrow
\mathcal{O}_{R, \overline{e}}
$$
第一段落の結果により $\mathcal{O}_{R, \overline{e}}$ の次元は $0$ である。
一方、$R$ は $\mathfrak m_u$ によって $F$ 内に切り出されるので、第二の射の核は
$p^\sharp(\mathfrak m_u)\mathcal{O}_{F, \overline{z}}$ である。したがって
$$
\mathfrak m_{\overline{z}} =
\sqrt{p^\sharp(\mathfrak m_u)\mathcal{O}_{F, \overline{z}}}
$$
一方、$U \to \mathcal{X}$ は滑らかなので $F \to u$ は代数空間の滑らかな射である。
これは $F$ が正則代数空間であることを意味する
(Spaces over Fields, Lemma \ref{spaces-over-fields-lemma-smooth-regular})。
したがって $\mathcal{O}_{F, \overline{z}}$ は正則局所環である
(Properties of Spaces, Lemma \ref{spaces-properties-lemma-regular})。
正則局所環は Cohen--Macaulay であることに注意する
(Algebra, Lemma \ref{algebra-lemma-regular-ring-CM})。
$d = \dim(\mathcal{O}_{F, \overline{z}})$ とおく。Algebra, Lemma
\ref{algebra-lemma-find-sequence-image-regular} により、
$f_1, \ldots, f_d \in \mathcal{O}_{U, u}$ を選び、その像
$\varphi(f_1), \ldots, \varphi(f_d)$ が
$\mathcal{O}_{F, \overline{z}}$ における正則列となるものを選べる。Lemma
\ref{stacks-morphisms-lemma-slice} により $U$ を縮小した後、
$Z = V(f_1, \ldots, f_d) \to \mathcal{X}$ は平坦かつ局所有限表示であると仮定できる。
構成により $F_Z = Z \times_\mathcal{X} u$ は
$F = U \times_\mathcal{X} u$ の閉部分空間であり、$z$ はこの閉部分空間の点であり、
さらに
$$
\dim(\mathcal{O}_{F_Z, \overline{z}}) = 0.
$$
Morphisms of Spaces, Lemma
\ref{spaces-morphisms-lemma-dimension-fibre-at-a-point} により、
$u$ に対する $z$ の超越次数が zero なので $\dim_z(F_Z) = 0$ である。
したがって Lemma \ref{stacks-morphisms-lemma-quasi-finite-at-point} により、
必要なら $U$ をさらに縮小した後で射 $Z \to \mathcal{X}$ は局所準有限である。

\medskip\noindent
これで、$\mathcal{X}$ の任意の有限型点 $x$ に対し、局所準有限、平坦、かつ
局所有限表示な射 $f_x : Z_x \to \mathcal{X}$ で $x$ が $|f_x|$ の像に入るものが
存在する。$W = \coprod_x Z_x$ および $f = \coprod f_x$ とおく。
すると $f$ は平坦、局所有限表示、かつ局所準有限である。特に $|f|$ の像は開である。
Properties of Stacks, Lemma \ref{stacks-properties-lemma-topology-points} を参照せよ。
構成によりその像は $\mathcal{X}$ のすべての有限型点を含む。したがって
Lemma \ref{stacks-morphisms-lemma-enough-finite-type-points} および Properties of Stacks, Lemma
\ref{stacks-properties-lemma-characterize-surjective} により $f$ は全射である。
\end{proof}

\begin{lemma}
\label{stacks-morphisms-lemma-DM-residual-gerbe}
$\mathcal{Z}$ を DM、局所 Noether、被約な代数スタックとし、
$|\mathcal{Z}|$ は一点集合とする。このとき、体 $k$ と全射なエタール射
$\Spec(k) \to \mathcal{Z}$ が存在する。
\end{lemma}

\begin{proof}
Properties of Stacks, Lemma
\ref{stacks-properties-lemma-unique-point-better} により、体 $k$ と全射、平坦、
かつ局所有限表示な射 $\Spec(k) \to \mathcal{Z}$ が存在する。
$U = \Spec(k)$、$R = U \times_\mathcal{Z} U$ とおくと、代数空間における亜群
$(U, R, s, t, c)$ を得る。Algebraic Stacks, Lemma
\ref{algebraic-lemma-criterion-map-representable-spaces-fibred-in-groupoids}.
Algebraic Stacks, Remark \ref{algebraic-remark-flat-fp-presentation} により
同値
$$
f_{can} : [U/R] \longrightarrow \mathcal{Z}
$$
射影 $s, t : R \to U$ は局所有限表示である。
$\mathcal{Z}$ は DM なので、安定化群代数空間
$$
G = U \times_{U \times U} R = U \times_{U \times U} (U \times_\mathcal{Z} U) =
U \times_{\mathcal{Z} \times \mathcal{Z}, \Delta_\mathcal{Z}} \mathcal{Z}
$$
は $U$ 上非分岐である。特に $\dim(G) = 0$ であり、
More on Groupoids in Spaces, Lemma
\ref{spaces-more-groupoids-lemma-groupoid-on-field-dimension-equal-stabilizer}
により $\dim(R) = 0$ である。したがって $R$ はスキームである。
次を参照せよ：
Spaces over Fields, Lemma
\ref{spaces-over-fields-lemma-locally-finite-type-dim-zero}.
Varieties, Lemma \ref{varieties-lemma-algebraic-scheme-dim-0} により、
$R$（および $G$）は、$s$ または $t$ のいずれによっても $k$ 上有限な
Artin 局所環のスペクトルの非交和である。$P = \Spec(A) \subset R$ を、
亜群スキーム $(U, R, s, t, c)$ の単位元 $e$ を基礎点とする開かつ閉な
部分スキームとする。$s \circ e = t \circ e = \text{id}_{\Spec(k)}$ なので、
$A$ は剰余体が $s^\sharp : k \to A$ または $t^\sharp : k \to A$ のいずれか
によって $k$ と同一視される Artin 局所環である。$s, t : \Spec(A) \to \Spec(k)$
は有限である（上で参照した補題による）。$G \to \Spec(k)$ は非分岐なので
$$
G \cap P = P \times_{U \times U} U = \Spec(A \otimes_{k \otimes k} k)
$$
は $k$ 上非分岐である。一方 $A \otimes_{k \otimes k} k$ は $A$ の商なので局所環であり、
$k$ へ全射する。よって $A \otimes_{k \otimes k} k = k$ である。
したがって $P \to U \times U$ は（$P$ は剰余体 $k$ をもつ一点しかもたないので）
普遍単射であり、（上の唯一の像点上のファイバーの計算により）非分岐であり、
（$s, t$ が有限型なので）有限型である。ゆえに単射である（\'Etale Morphisms, Lemma
\ref{etale-lemma-universally-injective-unramified} を参照せよ）。
したがって $s|_P, t|_P : P \to U$ は有限平坦な同値関係を定める。
よって
Groupoids, Proposition \ref{groupoids-proposition-finite-flat-equivalence}
を適用して $U/P$ が存在し、それがスキーム $\overline{U}$ であることを得る。
さらに $U \to \overline{U}$ は有限局所自由であり、
$P = U \times_{\overline{U}} U$ である。実際、
$\overline{U} = \Spec(k_0)$ であり、$k_0 \subset k$ は $R$-不変関数の環である。
$k$ は体なので、定義
Groupoids, Equation (\ref{groupoids-equation-invariants})
から $k_0$ は体である。

\medskip\noindent
次の射
\begin{equation}
\label{stacks-morphisms-equation-etale-covering}
\Spec(k_0) = \overline{U} = U/P \to [U/R] = \mathcal{Z}
\end{equation}
は所望の全射エタール射である。Properties of Stacks, Lemma
\ref{stacks-properties-lemma-flat-cover-by-field} により、この射は全射である。
したがって (\ref{stacks-morphisms-equation-etale-covering}) がエタールであることを示せば十分である
\footnote{この事実について読者自身の証明を見つけることを勧める。実際、この議論は
Bootstrap, Theorem \ref{bootstrap-theorem-final-bootstrap} の証明の最後の議論と
多くの共通点をもつので、どこかで独立の補題にすべきだろう。}。
エタール性を直接証明する代わりに、まず Bootstrap, Lemma
\ref{bootstrap-lemma-divide-subgroupoid} を適用して、亜群スキーム
$(\overline{U}, \overline{R}, \overline{s}, \overline{t}, \overline{c})$
が存在し、$(U, R, s, t, c)$ が
$(\overline{U}, \overline{R}, \overline{s}, \overline{t}, \overline{c})$
の商射 $U \to \overline{U}$ による制限となることが分かる。
（上の補題の仮定は、$j : R \to U \times U$ が分離的かつ局所準有限であるという
主張を除いてすべて確認済みであり、これは $R$ が $k$ 上分離的で局所準有限な
スキームであることから従う。）$U \to \overline{U}$ は有限局所自由なので、
$[U/R] \to [\overline{U}/\overline{R}]$ は同値である。次を参照せよ：
Groupoids in Spaces,
\ref{spaces-groupoids-lemma-quotient-stack-restrict-equivalence}.

\medskip\noindent
注意すべきは、$s, t$ が $U \to \overline{U}$ による
$\overline{s}, \overline{t}$ の基底変換であることである。
$\{U \to \overline{U}\}$ は fppf 被覆なので、
$\overline{s}, \overline{t}$ は平坦、局所有限表示、かつ局所準有限である。
次を参照せよ：
Descent, Lemmas \ref{descent-lemma-descending-property-flat},
\ref{descent-lemma-descending-property-locally-finite-presentation}, および
\ref{descent-lemma-descending-property-quasi-finite}.
可換図式
$$
\xymatrix{
U \times_{\overline{U}} U \ar@{=}[r] \ar[rd] & P \ar[r] \ar[d] & R \ar[d] \\
& \overline{U} \ar[r]^{\overline{e}} & \overline{R}
}
$$
制限についての一般事実により、外側の四つの角は Cartesian 図式をなす。
等式により内側の正方形も Cartesian である。$P$ は $R$ の開部分なので、
$\overline{e}$ は開埋入である。次を参照せよ：
Descent, Lemma \ref{descent-lemma-descending-property-open-immersion}.

\medskip\noindent
もちろん、$\overline{e}$ が開埋入で、$\overline{s}, \overline{t}$ が平坦かつ
局所有限表示ならば、射 $\overline{t}, \overline{s}$ はエタールである。
例えば次を適用すればよい：
More on Groupoids, Lemma \ref{more-groupoids-lemma-sheaf-differentials}
これは $\Omega_{\overline{R}/\overline{U}} = 0$ が
$\overline{s}, \overline{t} : \overline{R} \to \overline{U}$ の非分岐性を含意することを示す
（
Morphisms, Lemma \ref{morphisms-lemma-unramified-omega-zero}）。
さらに Morphisms, Lemma
\ref{morphisms-lemma-flat-unramified-etale} により
$\overline{s}, \overline{t}$ はエタールである。したがって
$\mathcal{Z} = [\overline{U}/\overline{R}]$ は代数スタック $\mathcal{Z}$ の
エタール表示であり、$\overline{U} \to \mathcal{Z}$ はエタールである。次を参照せよ：
Properties of Stacks, Lemma
\ref{stacks-properties-lemma-check-property-covering}.
\end{proof}

\begin{lemma}
\label{stacks-morphisms-lemma-etale-at-point}
$\mathcal{X}$ を代数スタックとする。次の Cartesian 図式を考える。
$$
\xymatrix{
U \ar[d] & F \ar[l]^p \ar[d] \\
\mathcal{X} & \Spec(k) \ar[l]
}
$$
ここで $U$ は代数空間、$k$ は体であり、$U \to \mathcal{X}$ は平坦かつ
局所有限表示とする。$F \to \Spec(k)$ が $z$ で非分岐となる
$z \in |F|$ をとる。$p(z)$ を含む開部分空間で $U$ を置き換えた後、射
$$
U \longrightarrow \mathcal{X}
$$
はエタールである。
\end{lemma}

\begin{proof}
射 $f : U \to \mathcal{X}$ は平坦かつ局所有限表示なので、制限
$f|_{W(f)} : W(f) \to \mathcal{X}$ がエタールとなる最大の開集合
$W(f) \subset U$ が存在する。Properties of Stacks, Remark
\ref{stacks-properties-remark-local-source-apply}
(\ref{stacks-properties-item-etale}) を参照せよ。
したがって $p(z)$ が $W(f)$ の点であることを示せばよい。
さらに、上記の Remark は代数空間によって表現可能な射による任意の基底変換と
$W(f)$ の形成が可換することも示す。ゆえに $F \to \Spec(k)$ が
$z$ でエタールであることを示せば十分である。これは $U \to \mathcal{X}$ の
基底変換として平坦かつ局所有限表示であり、仮定により
$F \to \Spec(k)$ は $z$ で非分岐なので、次から従う：
Morphisms of Spaces, Lemma
\ref{spaces-morphisms-lemma-unramified-flat-lfp-etale}。
\end{proof}

\noindent
DM スタックは Deligne--Mumford スタックである。

\begin{theorem}
\label{stacks-morphisms-theorem-DM}
$\mathcal{X}$ を代数スタックとする。次の条件は同値である。
\begin{enumerate}
\item $\mathcal{X}$ は DM である。
\item $\mathcal{X}$ は Deligne--Mumford である。
\item スキーム $W$ と全射なエタール射 $W \to \mathcal{X}$ が存在する。
\end{enumerate}
\end{theorem}

\begin{proof}
 (3) は (2) の定義であることを思い出そう。Algebraic Stacks, Definition
\ref{algebraic-definition-deligne-mumford} を参照せよ。(3) $\Rightarrow$ (1) は
Lemma \ref{stacks-morphisms-lemma-properties-covering-imply-diagonal} である。
(1) を仮定する。$x \in |\mathcal{X}|$ を有限型点とする。
$x$ の近傍で「機能する」$\mathcal{X}$ 上のスキームを構成する。
証明の最後に、これらすべての非交和をとって結論を得る。

\medskip\noindent
Lemma \ref{stacks-morphisms-lemma-point-finite-type-monomorphism} により、$x$ における
$\mathcal{X}$ の剰余ガーベ $\mathcal{Z}_x$ が存在し、
$\mathcal{Z}_x \to \mathcal{X}$ は局所有限型である。Lemma
\ref{stacks-morphisms-lemma-separation-properties-residual-gerbe} により、代数スタック
$\mathcal{Z}_x$ は DM である。Lemma \ref{stacks-morphisms-lemma-DM-residual-gerbe} により、
体 $k$ と全射エタール射 $z : \Spec(k) \to \mathcal{Z}_x$ が存在する。
特に、合成 $x : \Spec(k) \to \mathcal{X}$ は局所有限型である
（Morphisms of Spaces, Lemmas
\ref{spaces-morphisms-lemma-composition-finite-type} および
\ref{spaces-morphisms-lemma-etale-locally-finite-type} による）。

\medskip\noindent
$x$ が $|U| \to |\mathcal{X}|$ の像に入るようなスキーム $U$ と
滑らかな射 $U \to \mathcal{X}$ を選ぶ。次のファイバー積正方形を考える。
$$
\xymatrix{
U \ar[d] & F  \ar[l] \ar[d] \\
\mathcal{X} & \Spec(k) \ar[l]_-x
}
$$
すなわち $F = U \times_{\mathcal{X}, x} \Spec(k)$ である。
Properties of Stacks, Lemma \ref{stacks-properties-lemma-points-cartesian} により
$F$ は空でない。$\mathcal{Z}_x \to \mathcal{X}$ は単射なので
$$
\Spec(k) \times_{z, \mathcal{Z}_x, z} \Spec(k)
=
\Spec(k) \times_{x, \mathcal{X}, x} \Spec(k)
$$
 $z$ の構成により、これは $\Spec(k)$ へのエタールな射影をもつ。
等式
$$
F \times_U F =
(\Spec(k) \times_\mathcal{X} \Spec(k))
\times_{\Spec(k)} F
$$
が成り立つので、射影 $F \times_U F \to F$ もエタールである。したがって
$\Delta_{F/U} : F \to F \times_U F$ はエタールである（
Morphisms of Spaces, Lemma \ref{spaces-morphisms-lemma-etale-permanence}).
Morphisms of Spaces, Lemma
\ref{spaces-morphisms-lemma-etale-universally-injective-open} により
$\Delta_{F/U}$ は開埋入である。さらに Morphisms of Spaces, Lemma
\ref{spaces-morphisms-lemma-diagonal-unramified-morphism} により
$F \to U$ は非分岐である。

\medskip\noindent
空でないアフィン・スキーム $V$ とエタール射 $V \to F$ を取る。
（$F$ そのものを直接扱えばこれは避けられるが、こうする方が何が起きているかを
説明しやすい。）図のようになる。
$$
\xymatrix{
U \ar[d] & F  \ar[l] \ar[d] & V \ar[l] \ar[ld] \\
\mathcal{X} & \Spec(k) \ar[l]_-x
}
$$
このとき $V \to \Spec(k)$ はスキームの滑らかな射であり、$V \to U$ は
スキームの非分岐射である（
Morphisms of Spaces, Lemmas
\ref{spaces-morphisms-lemma-composition-smooth} および
\ref{spaces-morphisms-lemma-composition-unramified}）。
$k \subset \kappa(v)$ が有限分離となる閉点 $v \in V$ を取る（
Varieties, Lemma \ref{varieties-lemma-smooth-separable-closed-points-dense}）。
$u \in U$ をその像点とする。局所環
$\mathcal{O}_{V, v}$ は正則であり（
Varieties, Lemma \ref{varieties-lemma-smooth-regular}）、局所環準同型
$$
\varphi : \mathcal{O}_{U, u} \longrightarrow \mathcal{O}_{V, v}
$$
は $V \to U$ から来るもので、$\varphi(\mathfrak m_u)\mathcal{O}_{V, v} = \mathfrak m_v$
を満たす（Morphisms, Lemma \ref{morphisms-lemma-unramified-at-point}）。
したがって
$f_1, \ldots, f_d \in \mathcal{O}_{U, u}$ を取って、その像
$\varphi(f_1), \ldots, \varphi(f_d)$ が
$\mathfrak m_v/\mathfrak m_v^2$ の $\kappa(v)$ 上の基底となるようにできる。
$\mathcal{O}_{V, v}$ は正則局所環なので、これは
$\varphi(f_1), \ldots, \varphi(f_d)$ が
$\mathcal{O}_{V, v}$ の正則列をなすことを意味する（
Algebra, Lemma \ref{algebra-lemma-regular-ring-CM}）。
$U$ を $u$ の開近傍で置き換えることで
$f_1, \ldots, f_d \in \Gamma(U, \mathcal{O}_U)$ と仮定してよい。
さらに $U$ を $u$ の、場合によってはもっと小さい開近傍で置き換えることで、
$V(f_1, \ldots, f_d) \to \mathcal{X}$ が平坦かつ局所有限表示であると
仮定してよい（Lemma \ref{stacks-morphisms-lemma-slice}）。
構成により
$$
V(f_1, \ldots, f_d) \times_\mathcal{X} \Spec(k)
\longleftarrow
V(f_1, \ldots, f_d) \times_U V
$$
はエタールであり、$V(f_1, \ldots, f_d) \times_U V$ は
$f_1|_V, \ldots, f_d|_V$ で切り出される V の閉部分スキーム
$T \subset V$ である。したがって構成により $v \in T$ であり、
$$
\mathcal{O}_{T, v} =
\mathcal{O}_{V, v}/(\varphi(f_1), \ldots, \varphi(f_d)) = \kappa(v)
$$
は $k$ の有限分離拡大である。これから $T \to \Spec(k)$ は $v$ で
非分岐である（Morphisms, Lemma \ref{morphisms-lemma-unramified-at-point}）。
代数空間の非分岐射の定義により、これは
$V(f_1, \ldots, f_d) \times_\mathcal{X} \Spec(k) \to \Spec(k)$
が
$V(f_1, \ldots, f_d) \times_\mathcal{X} \Spec(k)$
における $v$ の像で非分岐であることを意味する。
これを
Lemma \ref{stacks-morphisms-lemma-etale-at-point}
により、$U$ を $u$ のさらに別の開近傍へ縮小すれば
$V(f_1, \ldots, f_d) \to \mathcal{X}$ がエタールとなることが分かる。

\medskip\noindent
以上から、$\mathcal{X}$ の任意の有限型点 $x$ に対して、
$x$ が $|f_x|$ の像に入るエタール射
$f_x : W_x \to \mathcal{X}$ が存在する。$W = \coprod_x W_x$、
$f = \coprod f_x$ と置く。このとき $f$ はエタールである。
特に $|f|$ の像は開である（
Properties of Stacks, Lemma \ref{stacks-properties-lemma-topology-points}）。
構成によりこの像は $\mathcal{X}$ のすべての有限型点を含むので、
$f$ は
Lemma \ref{stacks-morphisms-lemma-enough-finite-type-points}（および
Properties of Stacks, Lemma
\ref{stacks-properties-lemma-characterize-surjective}）により全射である。
\end{proof}

\noindent
これは、代数スタックの DM 射の「ファイバー」が
Deligne--Mumford であることを示す有用な系である。

\begin{lemma}
\label{stacks-morphisms-lemma-DM}
代数スタックの DM 射 $f : \mathcal{X} \to \mathcal{Y}$ を取る。このとき
\begin{enumerate}
\item 任意の DM 代数スタック $\mathcal{Z}$ と射
$\mathcal{Z} \to \mathcal{Y}$ に対し、スキーム U と全射エタール射
$U \to \mathcal{X} \times_\mathcal{Y} \mathcal{Z}$ が存在する。
\item 任意の代数空間 $Z$ と射 $Z \to \mathcal{Y}$ に対し、スキーム U と
全射エタール射 $U \to \mathcal{X} \times_\mathcal{Y} Z$ が存在する。
\end{enumerate}
\end{lemma}

\begin{proof}
（1）の証明。$f$ は DM なので、Lemma
\ref{stacks-morphisms-lemma-base-change-separated} により基底変換
$\mathcal{X} \times_\mathcal{Y} \mathcal{Z} \to \mathcal{Z}$ は DM である。
$\mathcal{Z}$ が DM であることから、Lemma
\ref{stacks-morphisms-lemma-separated-over-separated} により
$\mathcal{X} \times_\mathcal{Y} \mathcal{Z}$ も DM である。したがって
Theorem \ref{stacks-morphisms-theorem-DM} により、スキーム $U$ と全射エタール射
$U \to \mathcal{X} \times_\mathcal{Y} \mathcal{Z}$ が存在する。
（2）は（1）の特別な場合である。というのも、代数空間は
（代数スタックとみなせば）Lemma \ref{stacks-morphisms-lemma-trivial-implications} により DM
だからである。
\end{proof}






\section{Deligne--Mumford 軌跡}
\label{stacks-morphisms-section-DM}

\noindent
任意の代数スタックは、それ自体が Deligne--Mumford スタックである最大の
開部分スタックをもつ。このことはほぼ明らかだが、以下では証明も記す。
もちろんこの部分スタックは空でもよく、例えば
$X = [\Spec(\mathbf{Z})/\mathbf{G}_{m, \mathbf{Z}}]$ の場合がそうである。
以下では DM 軌跡の点を特徴づける。

\begin{lemma}
\label{stacks-morphisms-lemma-open-DM-locus}
代数スタック $\mathcal{X}$ を取る。開部分スタック
$$
\mathcal{X}'' \subset \mathcal{X}' \subset \mathcal{X}
$$
が存在して、$\mathcal{X}''$ は DM、$\mathcal{X}'$ は準DM であり、
これらはそれぞれこの性質をもつ最大の開部分スタックである。
\end{lemma}

\begin{proof}
ここで実際に主張しているのは、$\mathcal{U} \subset \mathcal{X}$ と
$\mathcal{V} \subset \mathcal{X}$ が DM である開部分スタックならば、
$|\mathcal{W}| = |\mathcal{U}| \cup |\mathcal{V}|$ を満たす開部分スタック
$\mathcal{W} \subset \mathcal{X}$ も DM であるということである。
（準DM についても同様である。）これは少しごまかしだが、これを証明するために
Theorem \ref{stacks-morphisms-theorem-DM} を用いる。
この定理により、スキーム $U$, $V$ と全射エタール射
$U \to \mathcal{U}$, $V \to \mathcal{V}$ を選べる。
するともちろん $U \amalg V \to \mathcal{W}$ は全射かつエタールである。
準DM の場合も Theorem \ref{stacks-morphisms-theorem-quasi-DM} を用いてまったく同じ方法で証明できる。
\end{proof}

\begin{lemma}
\label{stacks-morphisms-lemma-points-DM-locus}
代数スタック $\mathcal{X}$ を取る。$x \in |\mathcal{X}|$ が
$x : \Spec(k) \to \mathcal{X}$ に対応するとする。$G_x/k$ を
$x$ の自己同型群代数空間とする。このとき
\begin{enumerate}
\item $x$ が $\mathcal{X}$ の DM 軌跡に入ることと
$G_x \to \Spec(k)$ が非分岐であることは同値であり、
\item $x$ が $\mathcal{X}$ の準DM 軌跡に入ることと
$G_x \to \Spec(k)$ が局所準有限であることは同値である。
\end{enumerate}
\end{lemma}

\begin{proof}
（2）の証明。スキーム $U$ と全射滑らかな射
$U \to \mathcal{X}$ を選ぶ。ファイバー積
$$
\xymatrix{
G \ar[r] \ar[d] & \mathcal{I}_\mathcal{X} \ar[d] \\
U \ar[r] & \mathcal{X}
}
$$
 $G$ は $U \to \mathcal{X}$ の自己同型群代数空間であることを思い出そう。
Groupoids in Spaces, Lemma
\ref{spaces-groupoids-lemma-open-over-which-unramified-or-lqf}
により、$G_{U'} \to U'$ が局所準有限となる最大の開部分スキーム
$U' \subset U$ が存在する。さらに $U'$ の構成は任意の基底変換と可換である。
特に $R = U \times_\mathcal{X} U$ における $U'$ の二つの逆像は
$R$ の同じ開部分空間である（結局、$R \to \mathcal{X}$ の二つの射は同型であり、
したがって自己同型群空間も同型だからである）。
よって $U'$ は、次による開部分スタック
$\mathcal{X}' \subset \mathcal{X}$ の逆像である：
Properties of Stacks, Lemma
\ref{stacks-properties-lemma-substacks-presentation}
そして次の Cartesian 図式を得る。
$$
\xymatrix{
G_{U'} \ar[r] \ar[d] & \mathcal{I}_{\mathcal{X}'} \ar[d] \\
U' \ar[r] & \mathcal{X}'
}
$$
したがって射 $\mathcal{I}_{\mathcal{X}'} \to \mathcal{X}'$ は局所準有限であり、
Lemma \ref{stacks-morphisms-lemma-diagonal-diagonal} の (5) により
$\mathcal{X}'$ は準DM である。これに対し、$\mathcal{W} \subset \mathcal{X}$ が
準DM である開部分スタックならば、$\mathcal{W}$ の逆像
$W \subset U$ は $U'$ に含まれなければならない。これは $U'$ の構成により、
$\mathcal{I}_\mathcal{W} =
\mathcal{W} \times_\mathcal{X} \mathcal{I}_\mathcal{X}$
が $\mathcal{W}$ 上局所準有限だからである。したがって
$\mathcal{X}'$ は準DM 軌跡である。
最後に体拡大 $K/k$ と $2$-可換図式
$$
\xymatrix{
\Spec(K) \ar[r] \ar[d] & \Spec(k) \ar[d]^x \\
U \ar[r] & \mathcal{X}
}
$$
を選ぶ。このとき同型
$G_x \times_{\Spec(k)} \Spec(K) \cong G \times_U \Spec(K)$
なる $K$ 上の群代数空間の同型を得る。したがって $G_x$ が $k$ 上局所準有限である
ことと、$\Spec(K) \to U$ が $U'$ に写ることは同値である
（$U'$ の構成との可換性と
Groupoids in Spaces, Lemma
\ref{spaces-groupoids-lemma-open-over-which-unramified-or-lqf}
を $\Spec(K) \to \Spec(k)$ と $G_x$ に適用すれば分かる）。
これで (2) の証明が終わる。(1) の証明もまったく同じである。
\end{proof}









\section{局所準有限射}
\label{stacks-morphisms-section-locally-quasi-finite}

\noindent
代数空間の射の性質「局所準有限」は始域と終域について滑らか局所的
ではない。したがって、代数スタックの局所準有限射を定義するために
Section \ref{stacks-morphisms-section-local-source-target} の内容を使うことはできない。
代数空間によって表現可能な代数スタックの射について局所準有限が何を意味するかは
すでに知っている。次を参照せよ：
Properties of Stacks, Section
\ref{stacks-properties-section-properties-morphisms}.
一般の射に適した条件を見つけるため、次の観察をする。

\begin{lemma}
\label{stacks-morphisms-lemma-representable-by-spaces-quasi-finite}
代数スタックの射 $f : \mathcal{X} \to \mathcal{Y}$ を取る。
$f$ は代数空間によって表現可能であると仮定する。以下は同値である。
\begin{enumerate}
\item $f$ が局所準有限である（Properties of Stacks,
Section \ref{stacks-properties-section-properties-morphisms} の意味である）。
\item $f$ は局所有限型であり、任意の体 $k$ と射
$\Spec(k) \to \mathcal{Y}$ に対して空間
$|\Spec(k) \times_\mathcal{Y} \mathcal{X}|$ が離散である。
\end{enumerate}
\end{lemma}

\begin{proof}
（1）を仮定する。この場合、代数空間の射
$\mathcal{X}_k \to \Spec(k)$ は $f$ の基底変換なので局所準有限である。
したがって $|\mathcal{X}_k|$ は
Morphisms of Spaces, Lemma \ref{spaces-morphisms-lemma-locally-quasi-finite}
により離散である。
逆に（2）を仮定する。スキーム $V$ からの全射滑らかな射
$V \to \mathcal{Y}$ を取る。代数空間の射
$V \times_\mathcal{Y} \mathcal{X} \to V$ が局所準有限であることを示せば
十分である。次を参照せよ：
Properties of Stacks, Lemma
\ref{stacks-properties-lemma-check-property-covering}.
仮定により射 $V \times_\mathcal{Y} \mathcal{X} \to V$ は局所有限型である。
体 $k$ と任意の射 $\Spec(k) \to V$ に対して
$$
\Spec(k) \times_V (V \times_\mathcal{Y} \mathcal{X}) =
\Spec(k) \times_\mathcal{Y} \mathcal{X}
$$
は仮定により点の空間が離散である。したがって
$V \times_\mathcal{Y} \mathcal{X} \to V$ は
Morphisms of Spaces, Lemma \ref{spaces-morphisms-lemma-locally-quasi-finite}.
により局所準有限である。
\end{proof}

\noindent
代数空間によって表現可能な代数スタックの射は準DM である。次を参照せよ：
Lemma \ref{stacks-morphisms-lemma-trivial-implications}.
上の補題と合わせると、次の定義は代数空間によって表現可能な射の場合に
すでにある概念と矛盾しないことが分かる。

\begin{definition}
\label{stacks-morphisms-definition-locally-quasi-finite}
代数スタックの射 $f : \mathcal{X} \to \mathcal{Y}$ を取る。
$f$ が {\it 局所準有限} であるとは、$f$ が準DM かつ局所有限型であり、
さらに体 $k$ と任意の射 $\Spec(k) \to \mathcal{Y}$ に対して空間
$|\mathcal{X}_k|$ が離散であることをいう。
\end{definition}

\noindent
射 $f$ を準DM とする条件は自然である。例えば体 $k$ と射
$\pi : [\Spec(k)/\mathbf{G}_m] \to \Spec(k)$
を考える。この射はファイバーが一点で局所有限型である。後で見るように、
この射は相対次元 $-1$ の滑らかな射であり、局所準有限射には相対次元 $0$ を
もたせたい。
また切断 $\Spec(k) \to [\Spec(k)/\mathbf{G}_m]$ は離散ファイバーをもたないので、
局所準有限ではないことに注意する。さらに局所準有限射について次の永続性を
もたせたい：$f : \mathcal{X} \to \mathcal{X}'$ が代数スタック
$\mathcal{Y}$ 上局所準有限な代数スタックの射ならば、$f$ は局所準有限である
（実際には少し強い主張が成り立つ。次を参照せよ：
Lemma \ref{stacks-morphisms-lemma-locally-quasi-finite-permanence}).

\medskip\noindent
上の定義の別の根拠は、以下の
Lemma \ref{stacks-morphisms-lemma-characterize-locally-quasi-finite}
である。この補題は、$\mathcal{X}$ と $\mathcal{Y}$ の適切な
「表示」または「被覆」の存在により局所準有限性を特徴づける。

\begin{lemma}
\label{stacks-morphisms-lemma-base-change-locally-quasi-finite}
局所準有限射の基底変換は局所準有限である。
\end{lemma}

\begin{proof}
これは準DM 射については
Lemma \ref{stacks-morphisms-lemma-base-change-separated}
で、局所有限型射については
Lemma \ref{stacks-morphisms-lemma-base-change-finite-type}.
で見た。ファイバーに関する条件が基底変換に遺伝することは明らかである。
\end{proof}

\begin{lemma}
\label{stacks-morphisms-lemma-locally-quasi-finite-over-field}
局所準有限射 $\mathcal{X} \to \Spec(k)$ を取る。ここで
$\mathcal{X}$ は代数スタック、$k$ は体とする。
$V$ がスキームである局所準有限射 $f : V \to \mathcal{X}$ を取る。
このとき $V \to \Spec(k)$ は局所準有限である。
\end{lemma}

\begin{proof}
Lemma \ref{stacks-morphisms-lemma-composition-finite-type}
により $V \to \Spec(k)$ は局所有限型である。
矛盾を得るため $V \to \Spec(k)$ が局所準有限でないと仮定する。
すると $V$ の点の非自明な特殊化 $v \leadsto v'$ が存在する（
Morphisms, Lemma \ref{morphisms-lemma-quasi-finite-at-point-characterize}）。
特に
$\text{trdeg}_k(\kappa(v)) > \text{trdeg}_k(\kappa(v'))$ である（
Morphisms, Lemma \ref{morphisms-lemma-dimension-fibre-specialization}）。
$|\mathcal{X}|$ は離散なので $|f|(v) = |f|(v')$ である。
$R = V \times_\mathcal{X} V$ と置く。$R$ は代数空間であり、
射影 $s, t : R \to V$ は $V \to \mathcal{X}$ の基底変換として局所準有限である
（これは代数空間によって表現可能なので、次の議論から従う：
Properties of Stacks, Section
\ref{stacks-properties-section-properties-morphisms}).
Properties of Stacks, Lemma \ref{stacks-properties-lemma-points-cartesian}
により $s(r) = v$, $t(r) = v'$ を満たす $r \in |R|$ が存在する。
Morphisms of Spaces, Lemma \ref{spaces-morphisms-lemma-compare-tr-deg}
により、$v/k$ の超越次数は $r/k$ の超越次数に等しく、
さらに $v'/k$ の超越次数にも等しい。これは矛盾であり、補題が証明された。
\end{proof}

\begin{lemma}
\label{stacks-morphisms-lemma-composition-locally-quasi-finite}
局所準有限射の合成は局所準有限である。
\end{lemma}

\begin{proof}
これは準DM 射については
Lemma \ref{stacks-morphisms-lemma-composition-separated}
で、局所有限型射については
Lemma \ref{stacks-morphisms-lemma-composition-finite-type}.
で見た。$\mathcal{X} \to \mathcal{Y}$ と $\mathcal{Y} \to \mathcal{Z}$ を
局所準有限とする。$k$ を体とし、
$\Spec(k) \to \mathcal{Z}$ を射とする。
$|\mathcal{X}_k|$ が離散であることを示せば十分である。
Lemma \ref{stacks-morphisms-lemma-base-change-locally-quasi-finite}
により射 $\mathcal{X}_k \to \mathcal{Y}_k$ と
$\mathcal{Y}_k \to \Spec(k)$ は局所準有限である。
特に $\mathcal{Y}_k$ は準DM 代数スタックである。次を参照せよ：
Lemma \ref{stacks-morphisms-lemma-separated-implies-morphism-separated}.
Theorem \ref{stacks-morphisms-theorem-quasi-DM}
により、スキーム $V$ と全射、平坦、局所有限表示、局所準有限な射
$V \to \mathcal{Y}_k$ を取れる。
Lemma \ref{stacks-morphisms-lemma-locally-quasi-finite-over-field}
により $V$ は $k$ 上局所準有限であり、特に $|V|$ は離散である。射
$V \times_{\mathcal{Y}_k} \mathcal{X}_k \to \mathcal{X}_k$ は全射、平坦、局所有限表示
なので、
$|V \times_{\mathcal{Y}_k} \mathcal{X}_k| \to |\mathcal{X}_k|$ は全射かつ開である。
したがって
$|V \times_{\mathcal{Y}_k} \mathcal{X}_k|$ が離散であることを示せば十分である。
$V$ は剰余体 $k_i$ をもつアルティン局所 $k$-代数 $A_i$ のスペクトルの
非交和であることに注意する（
Varieties, Lemma \ref{varieties-lemma-algebraic-scheme-dim-0}）。
したがって各
$$
|\Spec(A_i) \times_{\mathcal{Y}_k} \mathcal{X}_k| =
|\Spec(k_i) \times_{\mathcal{Y}_k} \mathcal{X}_k| =
|\Spec(k_i) \times_\mathcal{Y} \mathcal{X}|
$$
が離散であることを示せばよい。これは
$\mathcal{X} \to \mathcal{Y}$ が局所準有限であるという仮定から従う。
\end{proof}

\noindent
局所準有限射を被覆で特徴づける前に、まず準DM 射について行う。

\begin{lemma}
\label{stacks-morphisms-lemma-characterize-quasi-DM}
代数スタックの射 $f : \mathcal{X} \to \mathcal{Y}$ を取る。
以下は同値である。
\begin{enumerate}
\item $f$ は準DM である。
\item $V$ が代数空間である任意の射 $V \to \mathcal{Y}$ に対し、
全射、平坦、局所有限表示、局所準有限な射
$U \to \mathcal{X} \times_\mathcal{Y} V$ が存在する。ただし
$U$ は代数空間である。
\item 代数空間 $U$, $V$ と射 $V \to \mathcal{Y}$ で全射、平坦、局所有限表示なもの、
および射 $U \to \mathcal{X} \times_\mathcal{Y} V$ で全射、平坦、局所有限表示かつ
局所準有限なものが存在する。
\end{enumerate}
\end{lemma}

\begin{proof}
（2）$\Rightarrow$（3）は明らかである。

\medskip\noindent
（1）を仮定し、（2）のような $V \to \mathcal{Y}$ を取る。このとき
$\mathcal{X} \times_\mathcal{Y} V \to V$ は準DM である。次を参照せよ：
Lemma \ref{stacks-morphisms-lemma-base-change-separated}.
Lemma \ref{stacks-morphisms-lemma-trivial-implications}
により代数空間 $V$ は DM、したがって準DM である。ゆえに
$\mathcal{X} \times_\mathcal{Y} V$ は
Lemma \ref{stacks-morphisms-lemma-separated-over-separated}.
により準DM である。したがって
Theorem \ref{stacks-morphisms-theorem-quasi-DM}
を適用して、（2）のような射
$U \to \mathcal{X} \times_\mathcal{Y} V$ を得る。

\medskip\noindent
（3）を仮定する。（3）のような $V \to \mathcal{Y}$ と
$U \to \mathcal{X} \times_\mathcal{Y} V$ を取る。
$f$ が準DM であることを示すには、
$\mathcal{X} \times_\mathcal{Y} V \to V$ が準DM であることを示せば十分である。次を参照せよ：
Lemma \ref{stacks-morphisms-lemma-check-separated-covering}.
Lemma \ref{stacks-morphisms-lemma-properties-covering-imply-diagonal}
により $\mathcal{X} \times_\mathcal{Y} V$ は準DM である。
したがって $\mathcal{X} \times_\mathcal{Y} V \to V$ は
Lemma \ref{stacks-morphisms-lemma-separated-implies-morphism-separated}
により準DM であり、（1）が成り立つ。これで補題の証明が終わる。
\end{proof}

\begin{lemma}
\label{stacks-morphisms-lemma-characterize-locally-quasi-finite}
代数スタックの射 $f : \mathcal{X} \to \mathcal{Y}$ を取る。
以下は同値である。
\begin{enumerate}
\item $f$ は局所準有限である。
\item $f$ は準DM であり、$V$ が代数空間である任意の射
$V \to \mathcal{Y}$ と、$U$ が代数空間である任意の局所準有限射
$U \to \mathcal{X} \times_\mathcal{Y} V$ に対して、射 $U \to V$ は局所準有限である。
\item 代数空間 $V$ からの任意の射 $V \to \mathcal{Y}$ に対し、
全射、平坦、局所有限表示かつ局所準有限な射
$U \to \mathcal{X} \times_\mathcal{Y} V$ が存在し、$U$ は代数空間で、
$U \to V$ は局所準有限である。
\item 代数空間 $U$, $V$、全射、平坦、局所有限表示な射
$V \to \mathcal{Y}$、および全射、平坦、局所有限表示かつ局所準有限な射
$U \to \mathcal{X} \times_\mathcal{Y} V$ が存在し、$U \to V$ は局所準有限である。
\end{enumerate}
\end{lemma}

\begin{proof}
（1）を仮定する。仮定により $f$ は準DM である。（2）のような
$V \to \mathcal{Y}$ と $U \to \mathcal{X} \times_\mathcal{Y} V$ を取る。
Lemma \ref{stacks-morphisms-lemma-composition-locally-quasi-finite}
により合成
$U \to \mathcal{X} \times_\mathcal{Y} V \to V$ は局所準有限である。
したがって（1）から（2）が従う。

\medskip\noindent
（2）を仮定する。（3）のような $V \to \mathcal{Y}$ を取る。
Lemma \ref{stacks-morphisms-lemma-characterize-quasi-DM}
により、代数空間 $U$ と全射、平坦、局所有限表示、局所準有限な射
$U \to \mathcal{X} \times_\mathcal{Y} V$ を取れる。（2）により合成
$U \to V$ は局所準有限である。したがって（2）から（3）が従う。

\medskip\noindent
（3）から（4）が従うことは明らかである。

\medskip\noindent
（4）を仮定する。（1）が成り立つことを示せば証明は終わる。
Lemma \ref{stacks-morphisms-lemma-characterize-quasi-DM}
により $f$ は準DM である。$f$ が局所有限型であることを示すには、
$g : \mathcal{X} \times_\mathcal{Y} V \to V$ が局所有限型であることを示せば
十分である。次を参照せよ：
Lemma \ref{stacks-morphisms-lemma-check-finite-type-covering}.
さらに $g$ と
$h : U \to \mathcal{X} \times_\mathcal{Y} V$ の合成が局所有限型であることを
確認すれば十分である。次を参照せよ：
Lemma \ref{stacks-morphisms-lemma-check-finite-type-precompose}.
仮定により $g \circ h : U \to V$ は局所準有限なのでこれは成り立ち、
$f$ は局所有限型である。
最後に $k$ を体とし、$\Spec(k) \to \mathcal{Y}$ を射とする。このとき
$V \times_\mathcal{Y} \Spec(k)$ は $k$ 上局所有限表示の空でない代数空間である。
したがって有限拡大 $k'/k$ と射 $\Spec(k') \to V$ を取って
次の図式が可換となる：
$$
\xymatrix{
\Spec(k') \ar[r] \ar[d] & V \ar[d] \\
\Spec(k) \ar[r] & \mathcal{Y}
}
$$
は可換である（詳細は省略）。このとき $\mathcal{X}_{k'} \to \mathcal{X}_k$ は
（スキームによって）表現可能で、全射かつ有限局所自由である。特に
$|\mathcal{X}_{k'}| \to |\mathcal{X}_k|$ は全射かつ開である。
したがって $|\mathcal{X}_{k'}|$ が離散であることを示せば十分である。
次の等式により
$$
U \times_V \Spec(k') =
U \times_{\mathcal{X} \times_\mathcal{Y} V} \mathcal{X}_{k'}
$$
から、$U \times_V \Spec(k') \to \mathcal{X}_{k'}$ は
（$U \to \mathcal{X} \times_\mathcal{Y} V$ の基底変換として）全射、平坦、
局所有限表示である。したがって
$|U \times_V \Spec(k')| \to |\mathcal{X}_{k'}|$ は全射かつ開である。
ゆえに $|U \times_V \Spec(k')|$ が離散であることを示せば十分である。
これは $U \to V$ が局所準有限であることから従う（上の定義によっても、
あるいは代数空間の射に対する元の定義と
Morphisms of Spaces, Lemma \ref{spaces-morphisms-lemma-locally-quasi-finite}
によってもよい）。
\end{proof}

\begin{lemma}
\label{stacks-morphisms-lemma-locally-quasi-finite-permanence}
代数スタックの射 $\mathcal{X} \to \mathcal{Y} \to \mathcal{Z}$ を取る。
$\mathcal{X} \to \mathcal{Z}$ は局所準有限、$\mathcal{Y} \to \mathcal{Z}$ は準DM
と仮定する。このとき $\mathcal{X} \to \mathcal{Y}$ は局所準有限である。
\end{lemma}

\begin{proof}
射 $\mathcal{X} \to \mathcal{Y}$ を次の合成として書く。
$$
\mathcal{X} \longrightarrow
\mathcal{X} \times_\mathcal{Z} \mathcal{Y} \longrightarrow
\mathcal{Y}
$$
第二の矢印は $\mathcal{X} \to \mathcal{Z}$ の基底変換なので局所準有限である。次を参照せよ：
Lemma \ref{stacks-morphisms-lemma-base-change-locally-quasi-finite}.
第一の矢印は
Lemma \ref{stacks-morphisms-lemma-semi-diagonal}
により局所準有限である。$\mathcal{Y} \to \mathcal{Z}$ が準DM だからである。
したがって $\mathcal{X} \to \mathcal{Y}$ は
Lemma \ref{stacks-morphisms-lemma-composition-locally-quasi-finite}.
により局所準有限である。
\end{proof}

















\section{準有限射}
\label{stacks-morphisms-section-quasi-finite}

\noindent
代数スタックの局所準有限射は Section \ref{stacks-morphisms-section-locally-quasi-finite} で、
準コンパクト射は Section \ref{stacks-morphisms-section-quasi-compact} で定義した。
代数空間の射は定義により、局所準有限かつ準コンパクトであることと準有限である
ことが同値である（Morphisms of Spaces, Definition
\ref{spaces-morphisms-definition-locally-quasi-finite}),
したがって代数スタックの射が準有限であることを次のように定義できる。
これは、代数空間によって表現可能な射について
Properties of Stacks,
Section \ref{stacks-properties-section-properties-morphisms}
で既に定義されている概念と一致する。

\begin{definition}
\label{stacks-morphisms-definition-quasi-finite}
\begin{reference}
\cite{rydh_approx}
\end{reference}
代数スタックの射 $f : \mathcal{X} \to \mathcal{Y}$ を取る。
$f$ が {\it 準有限} であるとは、$f$ が局所準有限
（Definition \ref{stacks-morphisms-definition-locally-quasi-finite}）かつ準コンパクト
（Definition \ref{stacks-morphisms-definition-quasi-compact}）であることをいう。
\end{definition}

\begin{lemma}
\label{stacks-morphisms-lemma-composition-quasi-finite}
準有限射の合成は準有限である。
\end{lemma}

\begin{proof}
次を組み合わせればよい：
Lemmas \ref{stacks-morphisms-lemma-composition-locally-quasi-finite} および
\ref{stacks-morphisms-lemma-composition-quasi-compact}.
\end{proof}

\begin{lemma}
\label{stacks-morphisms-lemma-base-change-quasi-finite}
準有限射の基底変換は準有限である。
\end{lemma}

\begin{proof}
次を組み合わせればよい：
Lemmas \ref{stacks-morphisms-lemma-base-change-locally-quasi-finite} および
\ref{stacks-morphisms-lemma-base-change-quasi-compact}.
\end{proof}

\begin{lemma}
\label{stacks-morphisms-lemma-quasi-finite-permanence}
代数スタックの射 $f : \mathcal{X} \to \mathcal{Y}$ と
$g : \mathcal{Y} \to \mathcal{Z}$ を取る。$g \circ f$ が準有限で、
$g$ が準分離かつ準DM ならば、$f$ は準有限である。
\end{lemma}

\begin{proof}
次を組み合わせればよい：
Lemmas \ref{stacks-morphisms-lemma-locally-quasi-finite-permanence} および
\ref{stacks-morphisms-lemma-quasi-compact-permanence}.
\end{proof}












\section{平坦射}
\label{stacks-morphisms-section-flat}

\noindent
代数空間の射の性質「平坦」は始域と終域について滑らか局所的である。次を参照せよ：
Descent on Spaces, Remark \ref{spaces-descent-remark-list-local-source-target}.
また基底変換で保たれ、終域について fpqc 局所的である。次を参照せよ：
Morphisms of Spaces,
Lemma \ref{spaces-morphisms-lemma-base-change-flat}
および
Descent on Spaces, Lemma
\ref{spaces-descent-lemma-descending-property-flat}.
したがって上の
Lemma \ref{stacks-morphisms-lemma-local-source-target}
により、代数スタックの射が平坦であることを次のように定義できる。
これは、代数空間によって表現可能な射について
Properties of Stacks,
Section \ref{stacks-properties-section-properties-morphisms}
で既に定義されている概念と一致する。

\begin{definition}
\label{stacks-morphisms-definition-flat}
代数スタックの射 $f : \mathcal{X} \to \mathcal{Y}$ を取る。
$f$ が {\it 平坦} であるとは、次の同値条件が
Lemma \ref{stacks-morphisms-lemma-local-source-target}
を $\mathcal{P} = \text{flat}$ として満たすことをいう。
\end{definition}

\begin{lemma}
\label{stacks-morphisms-lemma-composition-flat}
平坦射の合成は平坦である。
\end{lemma}

\begin{proof}
次を組み合わせればよい：
Remark \ref{stacks-morphisms-remark-composition}
と
Morphisms of Spaces, Lemma
\ref{spaces-morphisms-lemma-composition-flat}.
\end{proof}

\begin{lemma}
\label{stacks-morphisms-lemma-base-change-flat}
平坦射の基底変換は平坦である。
\end{lemma}

\begin{proof}
次を組み合わせればよい：
Remark \ref{stacks-morphisms-remark-base-change}
と
Morphisms of Spaces, Lemma
\ref{spaces-morphisms-lemma-base-change-flat}.
\end{proof}

\begin{lemma}
\label{stacks-morphisms-lemma-descent-flat}
代数スタックの射 $f : \mathcal{X} \to \mathcal{Y}$ を取る。
$\mathcal{Z} \to \mathcal{Y}$ を代数スタックの全射平坦射とする。基底変換
$\mathcal{Z} \times_\mathcal{Y} \mathcal{X} \to \mathcal{Z}$
が平坦ならば、$f$ は平坦である。
\end{lemma}

\begin{proof}
代数空間 $W$ と全射滑らかな射 $W \to \mathcal{Z}$ を選ぶ。
$W \to \mathcal{Z}$ は全射かつ平坦であり（Morphisms of Spaces,
Lemma \ref{spaces-morphisms-lemma-smooth-flat}）、したがって
$W \to \mathcal{Y}$ は全射かつ平坦である（次により：
Properties of Stacks, Lemma
\ref{stacks-properties-lemma-composition-surjective}
および
Lemma \ref{stacks-morphisms-lemma-composition-flat}).
 $W \to \mathcal{Z}$ による
$\mathcal{Z} \times_\mathcal{Y} \mathcal{X} \to \mathcal{Z}$
の基底変換は平坦射である（Lemma \ref{stacks-morphisms-lemma-base-change-flat}）から、
$\mathcal{Z}$ を $W$ で置き換えてよい。

\medskip\noindent
代数空間 $V$ と全射滑らかな射 $V \to \mathcal{Y}$ を選ぶ。代数空間 $U$ と
全射滑らかな射 $U \to V \times_\mathcal{Y} \mathcal{X}$ を選ぶ。
$U \to V$ が平坦であることを示さなければならない。ここで
$W \to \mathcal{Y}$ によってすべてを基底変換する。すなわち
$U' = W \times_\mathcal{Y} U$、
$V' = W \times_\mathcal{Y} V$、
$\mathcal{X}' = W \times_\mathcal{Y} \mathcal{X}$、
$\mathcal{Y}' = W \times_\mathcal{Y} \mathcal{Y} = W$ と置く。
すると基底変換により $U' \to V' \times_{\mathcal{Y}'} \mathcal{X}'$ は
依然として滑らかである。したがって代数スタックの平坦射の定義と
$\mathcal{X}' \to \mathcal{Y}'$ が平坦であるという仮定から、
$U' \to V'$ は平坦である。さらに
$V' \to V$ は $W \to \mathcal{Y}$ の基底変換として全射なので、
$U \to V$ は次により平坦である：
Morphisms of Spaces, Lemma
\ref{spaces-morphisms-lemma-base-change-module-flat} (2)
これで証明が終わる。
\end{proof}

\begin{lemma}
\label{stacks-morphisms-lemma-flat-permanence}
代数スタックの射 $\mathcal{X} \to \mathcal{Y} \to \mathcal{Z}$ を取る。
$\mathcal{X} \to \mathcal{Z}$ が平坦で、$\mathcal{X} \to \mathcal{Y}$ が全射かつ
平坦ならば、$\mathcal{Y} \to \mathcal{Z}$ は平坦である。
\end{lemma}

\begin{proof}
代数空間 $W$ と全射滑らかな射 $W \to \mathcal{Z}$ を選ぶ。代数空間 $V$ と
全射滑らかな射 $V \to W \times_\mathcal{Z} \mathcal{Y}$ を選ぶ。代数空間 $U$ と
全射滑らかな射 $U \to V \times_\mathcal{Y} \mathcal{X}$ を選ぶ。
$U \to V$ と $U \to W$ は平坦である。また $\mathcal{X} \to \mathcal{Y}$ は全射
なので、$U \to V$ は全射である（全射の合成として）。したがってこの補題は
代数空間の射の場合に帰着する。代数空間の射の場合は
Morphisms of Spaces, Lemma
\ref{spaces-morphisms-lemma-flat-permanence}.
\end{proof}

\begin{lemma}
\label{stacks-morphisms-lemma-lift-valuation-ring-through-flat-morphism}
代数スタックの平坦射 $f : \mathcal{X} \to \mathcal{Y}$ を取る。
$A$ が付値環である射 $\Spec(A) \to \mathcal{Y}$ を取る。
$\Spec(A)$ の閉点の像が、$|\mathcal{X|} \to |\mathcal{Y}|$ の像に入る
$|\mathcal{Y}|$ の点ならば、可換図式
$$
\xymatrix{
\Spec(A') \ar[r] \ar[d] & \mathcal{X} \ar[d] \\
\Spec(A) \ar[r] & \mathcal{Y}
}
$$
が存在する。ここで $A \to A'$ は付値環の拡大である
（More on Algebra, Definition
\ref{more-algebra-definition-extension-valuation-rings}).
\end{lemma}

\begin{proof}
基底変換 $\mathcal{X}_A \to \Spec(A)$ は平坦であり
（Lemma \ref{stacks-morphisms-lemma-base-change-flat}）、$\Spec(A)$ の閉点は
$|\mathcal{X}_A| \to |\Spec(A)|$ の像に入る
（Properties of Stacks, Lemma \ref{stacks-properties-lemma-points-cartesian}）。
したがって $\mathcal{Y} = \Spec(A)$ と仮定してよい。$U$ がスキームである
全射滑らかな射 $U \to \mathcal{X}$ を取る。すると Morphisms of Spaces, Lemma
\ref{spaces-morphisms-lemma-lift-valuation-ring-through-flat-morphism}
を射 $U \to \Spec(A)$ に適用して結論を得る。
\end{proof}







\section{点での平坦性}
\label{stacks-morphisms-section-flat-at-point}

\noindent
代数スタックの射が点で所与の性質をもつとは何を意味するかを述べるために、
一般的な道具立てをまだ整備する必要がある。ここでは次の補題で十分である。

\begin{lemma}
\label{stacks-morphisms-lemma-flat-at-point}
代数スタックの射 $f : \mathcal{X} \to \mathcal{Y}$ を取る。
$x \in |\mathcal{X}|$ とする。可換図式
$$
\vcenter{
\xymatrix{
U \ar[d]_a \ar[r]_h & V \ar[d]^b \\
\mathcal{X} \ar[r]^f & \mathcal{Y}
}
}
\quad\text{with points}
\vcenter{
\xymatrix{
u \in |U| \ar[d] \\
x \in |\mathcal{X}|
}
}
$$
ここで $U$ と $V$ は代数空間、$b$ は平坦であり、
$(a, h) : U \to \mathcal{X} \times_\mathcal{Y} V$
は平坦である。以下は同値である。
\begin{enumerate}
\item 上のような一つの図式で $h$ が $u$ で平坦である。
\item 上のようなすべての図式で $h$ が $u$ で平坦である。
\end{enumerate}
\end{lemma}

\begin{proof}
補題のような二つ目の図式 $U', V', u', a', b', h'$ が与えられているとする。
すると次の図式を考えることができる：
$$
\xymatrix{
U \ar[d] & U \times_\mathcal{X} U' \ar[l] \ar[d] \ar[r] & U' \ar[d] \\
V & V \times_\mathcal{Y} V' \ar[l] \ar[r] & V'
}
$$
Properties of Stacks, Lemma \ref{stacks-properties-lemma-points-cartesian}
により、$u$ および $u'$ に写る点 $u'' \in |U \times_\mathcal{X} U'|$ が存在する。
$h$ が $u$ で平坦なら、基底変換
$U \times_V (V \times_\mathcal{Y} V') \to V \times_\mathcal{Y} V'$
は $u$ の上の任意の点で平坦である。これは
Morphisms of Spaces, Lemma \ref{spaces-morphisms-lemma-base-change-module-flat}.
一方、射
$$
U \times_\mathcal{X} U' \to
U \times_\mathcal{X} (\mathcal{X} \times_\mathcal{Y} V') =
U \times_\mathcal{Y} V' =
U \times_V (V \times_\mathcal{Y} V')
$$
は $(a', h')$ の基底変換として平坦である。これは Lemma
\ref{stacks-morphisms-lemma-base-change-flat} による。
これと
Morphisms of Spaces, Lemma \ref{spaces-morphisms-lemma-composition-module-flat}
を用いて合成すると、$U \times_\mathcal{X} U' \to V \times_\mathcal{Y} V'$
は $u''$ で平坦であることが分かる。次に平坦射
$V \times_\mathcal{Y} V' \to V'$
との合成により、$U \times_\mathcal{X} U' \to V'$ は $u''$ で平坦である
ことが分かる。最後に、$U \times_\mathcal{X} U' \to U'$ が $u''$ で平坦であり、
$u''$ が $u'$ に写るので、
Morphisms of Spaces, Lemma \ref{spaces-morphisms-lemma-flat-permanence}.
により $U' \to V'$ は $u'$ で平坦であると結論する。
\end{proof}

\begin{definition}
\label{stacks-morphisms-definition-flat-at-point}
代数スタックの射 $f : \mathcal{X} \to \mathcal{Y}$ を取る。
$x \in |\mathcal{X}|$ とする。Lemma \ref{stacks-morphisms-lemma-flat-at-point} の
同値な条件が成り立つとき、$f$ は {\it $x$ で平坦} であるという。
\end{definition}





\section{有限表示の射}
\label{stacks-morphisms-section-finite-presentation}

\noindent
代数空間の射が「局所有限表示」であるという性質は始域・終域上滑らか局所的である。
次を参照せよ：
Descent on Spaces, Remark \ref{spaces-descent-remark-list-local-source-target}.
また、この性質は基底変換で保たれ、終域上 fpqc 局所的である。次を参照せよ：
Morphisms of Spaces,
Lemma \ref{spaces-morphisms-lemma-base-change-finite-presentation}
および
Descent on Spaces, Lemma
\ref{spaces-descent-lemma-descending-property-locally-finite-presentation}.
したがって、上の Lemma \ref{stacks-morphisms-lemma-local-source-target} により、
代数スタックの射が局所有限表示であることを次のように定義できる。
射が代数空間で表現可能な場合、この定義は既に
Properties of Stacks,
Section \ref{stacks-properties-section-properties-morphisms}
で定義されている概念と一致する。

\begin{definition}
\label{stacks-morphisms-definition-locally-finite-presentation}
代数スタックの射 $f : \mathcal{X} \to \mathcal{Y}$ を取る。
\begin{enumerate}
\item
Lemma \ref{stacks-morphisms-lemma-local-source-target}
の同値な条件が
$\mathcal{P} = \text{locally of finite presentation}$ として成り立つとき、
$f$ は {\it 局所有限表示} であるという。
\item $f$ が局所有限表示、準コンパクト、かつ準分離であるとき、
この射は {\it 有限表示} であるという。
\end{enumerate}
\end{definition}

\noindent
有限表示の射は、単に局所有限表示かつ準コンパクトな射なのでは
{\bf ない}ことに注意せよ。

\begin{lemma}
\label{stacks-morphisms-lemma-composition-finite-presentation}
有限表示の射の合成は有限表示である。
局所有限表示の射についても同じことが成り立つ。
\end{lemma}

\begin{proof}
次を組み合わせればよい：
Remark \ref{stacks-morphisms-remark-composition} と
Morphisms of Spaces, Lemma
\ref{spaces-morphisms-lemma-composition-finite-presentation}.
\end{proof}

\begin{lemma}
\label{stacks-morphisms-lemma-base-change-finite-presentation}
有限表示の射の基底変換は有限表示である。
局所有限表示の射についても同じことが成り立つ。
\end{lemma}

\begin{proof}
次を組み合わせればよい：
Remark \ref{stacks-morphisms-remark-base-change} と
Morphisms of Spaces, Lemma
\ref{spaces-morphisms-lemma-base-change-finite-presentation}.
\end{proof}

\begin{lemma}
\label{stacks-morphisms-lemma-finite-presentation-finite-type}
局所有限表示の射は局所有限型である。
有限表示の射は有限型である。
\end{lemma}

\begin{proof}
次を組み合わせればよい：
Remark \ref{stacks-morphisms-remark-implication} と
Morphisms of Spaces, Lemma
\ref{spaces-morphisms-lemma-finite-presentation-finite-type}.
\end{proof}

\begin{lemma}
\label{stacks-morphisms-lemma-noetherian-finite-type-finite-presentation}
代数スタックの射 $f : \mathcal{X} \to \mathcal{Y}$ を取る。
\begin{enumerate}
\item $\mathcal{Y}$ が局所 Noether で $f$ が局所有限型ならば、
$f$ は局所有限表示である。
\item $\mathcal{Y}$ が局所 Noether で $f$ が有限型かつ準分離ならば、
$f$ は有限表示である。
\end{enumerate}
\end{lemma}

\begin{proof}
 $f : \mathcal{X} \to \mathcal{Y}$ が局所有限型で、
$\mathcal{Y}$ が局所 Noether であると仮定する。
これは、
Lemma \ref{stacks-morphisms-lemma-local-source-target}
のような図式で、$h$ が局所有限型かつ垂直射 $a$ が全射であるものが存在する
ことを意味する。
Morphisms of Spaces, Lemma
\ref{spaces-morphisms-lemma-noetherian-finite-type-finite-presentation}
により、$h$ は局所有限表示である。
したがって定義から $\mathcal{X} \to \mathcal{Y}$ は局所有限表示である。
これで (1) が証明された。
$f$ が有限型かつ準分離ならば、これは準コンパクトかつ準分離でもあるから、
(2) が直ちに従う。
\end{proof}

\begin{lemma}
\label{stacks-morphisms-lemma-finite-presentation-permanence}
代数スタックの射 $f : \mathcal{X} \to \mathcal{Y}$ および
$g : \mathcal{Y} \to \mathcal{Z}$ を取る。
$g \circ f$ が局所有限表示で $g$ が局所有限型ならば、
$f$ は局所有限表示である。
\end{lemma}

\begin{proof}
代数空間 $W$ と全射滑らかな射 $W \to \mathcal{Z}$ を選ぶ。
代数空間 $V$ と全射滑らかな射
$V \to \mathcal{Y} \times_\mathcal{Z} W$ を選ぶ。
代数空間 $U$ と全射滑らかな射
$U \to \mathcal{X} \times_\mathcal{Y} V$ を選ぶ。
射 $U \to V \to W$ に
Morphisms of Spaces, Lemma
\ref{spaces-morphisms-lemma-finite-presentation-permanence}
を適用すれば補題が従う。
\end{proof}

\begin{lemma}
\label{stacks-morphisms-lemma-diagonal-morphism-finite-type}
対角射
$\Delta : \mathcal{X} \to \mathcal{X} \times_\mathcal{Y} \mathcal{X}$.
をもつ代数スタックの射 $f : \mathcal{X} \to \mathcal{Y}$ を取る。
$f$ が局所有限型ならば、$\Delta$ は局所有限表示である。
$f$ が準分離かつ局所有限型ならば、$\Delta$ は有限表示である。
\end{lemma}

\begin{proof}
 $\Delta$ は（第二射影を介して）$\mathcal{X}$ 上の射であることに注意する。
$f$ が局所有限型ならば、$\mathcal{X}$ は $\mathcal{X}$ 上有限表示であり、
$\text{pr}_2 : \mathcal{X} \times_\mathcal{Y} \mathcal{X} \to \mathcal{X}$
は Lemma \ref{stacks-morphisms-lemma-base-change-finite-type} により局所有限型である。
したがって最初の主張は Lemma \ref{stacks-morphisms-lemma-finite-presentation-permanence}
から従う。第二の主張は最初の主張と定義から従う（定義により、
$f$ が準分離であるとは $\Delta_f$ が準コンパクトかつ準分離であることを
意味するためである）。
\end{proof}

\begin{lemma}
\label{stacks-morphisms-lemma-open-immersion-locally-finite-presentation}
開埋入は局所有限表示である。
\end{lemma}

\begin{proof}
Properties of Stacks, Definition
\ref{stacks-properties-definition-immersion}
を考慮すれば、これは
Morphisms of Spaces,
Lemma \ref{spaces-morphisms-lemma-open-immersion-locally-finite-presentation}.
から従う。
\end{proof}

\begin{lemma}
\label{stacks-morphisms-lemma-check-property-after-fppf-base-change}
$P$ を、終域上 fppf 局所的で任意の基底変換により保たれる
代数空間の射の性質とする。
$f : \mathcal{X} \to \mathcal{Y}$ を代数空間で表現可能な
代数スタックの射とする。
$\mathcal{Z} \to \mathcal{Y}$ を全射、平坦、かつ局所有限表示な
代数スタックの射とする。
$\mathcal{W} = \mathcal{Z} \times_\mathcal{Y} \mathcal{X}$ と置く。このとき
$$
(f\text{ has }P) \Leftrightarrow
(\text{the projection }\mathcal{W} \to \mathcal{Z}\text{ has }P).
$$
この主張の意味については次を参照せよ：
Properties of Stacks, Section
\ref{stacks-properties-section-properties-morphisms}.
\end{lemma}

\begin{proof}
代数空間 $W$ と全射、平坦、かつ局所有限表示な射
$W \to \mathcal{Z}$ を選ぶ。
Properties of Stacks, Lemma
\ref{stacks-properties-lemma-composition-surjective}
および Lemmas \ref{stacks-morphisms-lemma-composition-flat} と
\ref{stacks-morphisms-lemma-composition-finite-presentation}
により、合成 $W \to \mathcal{Y}$ も全射、平坦、かつ局所有限表示である。
次のように記す：
$V = W \times_\mathcal{Z} \mathcal{W} = V \times_\mathcal{Y} \mathcal{X}$.
Properties of Stacks, Lemma
\ref{stacks-properties-lemma-check-property-covering}
により、$f$ が $\mathcal{P}$ をもつことと $V \to W$ が
その性質をもつことは同値であり、
$\mathcal{W} \to \mathcal{Z}$ が $\mathcal{P}$ をもつことと
$V \to W$ がその性質をもつことも同値である。よって補題が従う。
\end{proof}

\begin{lemma}
\label{stacks-morphisms-lemma-descent-property}
$\mathcal{P}$ を、始域・終域上滑らか局所的かつ終域上 fppf 局所的な
代数空間の射の性質とする。
$f : \mathcal{X} \to \mathcal{Y}$ を代数スタックの射とする。
$\mathcal{Z} \to \mathcal{Y}$ を全射、平坦、かつ局所有限表示な
代数スタックの射とする。基底変換
$\mathcal{Z} \times_\mathcal{Y} \mathcal{X} \to \mathcal{Z}$
が $\mathcal{P}$ をもつならば、$f$ は $\mathcal{P}$ をもつ。
\end{lemma}

\begin{proof}
$\mathcal{Z} \times_\mathcal{Y} \mathcal{X} \to \mathcal{Z}$ が
$\mathcal{P}$ をもつと仮定する。代数空間 $W$ と全射滑らかな射
$W \to \mathcal{Z}$ を選ぶ。次に注意する：
$W \times_\mathcal{Z} \mathcal{Z} \times_\mathcal{Y} \mathcal{X} =
W \times_\mathcal{Y} \mathcal{X}$.
したがって、$\mathcal{Z} \times_\mathcal{Y} \mathcal{X} \to \mathcal{Z}$
が $\mathcal{P}$ をもつことの定義そのもの（Definition \ref{stacks-morphisms-definition-P}
および Lemma \ref{stacks-morphisms-lemma-local-source-target} を参照）から、
$W \times_\mathcal{Y} \mathcal{X} \to W$ は $\mathcal{P}$ をもつ。
一方、$W \to \mathcal{Z}$ は全射、平坦、かつ局所有限表示である
（Morphisms of Spaces, Lemmas
\ref{spaces-morphisms-lemma-smooth-flat} および
\ref{spaces-morphisms-lemma-smooth-locally-finite-presentation})
ので、$W \to \mathcal{Y}$ も全射、平坦、かつ局所有限表示である（次による：
Properties of Stacks, Lemma
\ref{stacks-properties-lemma-composition-surjective}
ならびに
Lemmas \ref{stacks-morphisms-lemma-composition-flat} と
\ref{stacks-morphisms-lemma-composition-finite-presentation}).
したがって $\mathcal{Z}$ を $W$ で置き換えてよい。

\medskip\noindent
代数空間 $V$ と全射滑らかな射 $V \to \mathcal{Y}$ を選ぶ。
代数空間 $U$ と全射滑らかな射
$U \to V \times_\mathcal{Y} \mathcal{X}$ を選ぶ。
$U \to V$ が $\mathcal{P}$ をもつことを示さなければならない。
ここですべてを $W \to \mathcal{Y}$ により基底変換する。次のように置く：
$U' = W \times_\mathcal{Y} U$,
$V' = W \times_\mathcal{Y} V$,
$\mathcal{X}' = W \times_\mathcal{Y} \mathcal{X}$,
および $\mathcal{Y}' = W \times_\mathcal{Y} \mathcal{Y} = W$。
基底変換により、$U' \to V' \times_{\mathcal{Y}'} \mathcal{X}'$ は
依然として滑らかである。したがって、$\mathcal{X}' \to
\mathcal{Y}' = W$ が $\mathcal{P}$ をもつことの定義に用いられる
Lemma \ref{stacks-morphisms-lemma-local-source-target} により、$U' \to V'$ は
$\mathcal{P}$ をもつ。さらに $V' \to V$ は
$W \to \mathcal{Y}$ の基底変換として全射、平坦、かつ局所有限表示であり、
$\mathcal{P}$ は終域上 fppf 位相で局所的なので、
$U \to V$ は $\mathcal{P}$ をもつ。
\end{proof}

\begin{lemma}
\label{stacks-morphisms-lemma-descent-finite-presentation}
$f : \mathcal{X} \to \mathcal{Y}$ を代数スタックの射とする。
$\mathcal{Z} \to \mathcal{Y}$ を全射、平坦、かつ局所有限表示な
代数スタックの射とする。基底変換
$\mathcal{Z} \times_\mathcal{Y} \mathcal{X} \to \mathcal{Z}$
が局所有限表示ならば、$f$ は局所有限表示である。
\end{lemma}

\begin{proof}
「局所有限表示」という性質は Lemma \ref{stacks-morphisms-lemma-descent-property}
の条件を満たす。始域・終域上滑らか局所的であることは本節の冒頭で見た。
終域上 fppf 局所的であることは
Descent on Spaces, Lemma
\ref{spaces-descent-lemma-descending-property-locally-finite-presentation}.
である。
\end{proof}

\begin{lemma}
\label{stacks-morphisms-lemma-flat-finite-presentation-permanence}
$\mathcal{X} \to \mathcal{Y} \to \mathcal{Z}$ を代数スタックの射とする。
$\mathcal{X} \to \mathcal{Z}$ が局所有限表示であり、
$\mathcal{X} \to \mathcal{Y}$ が全射、平坦、かつ局所有限表示ならば、
$\mathcal{Y} \to \mathcal{Z}$ は局所有限表示である。
\end{lemma}

\begin{proof}
代数空間 $W$ と全射滑らかな射 $W \to \mathcal{Z}$ を選ぶ。
代数空間 $V$ と全射滑らかな射
$V \to W \times_\mathcal{Z} \mathcal{Y}$ を選ぶ。
代数空間 $U$ と全射滑らかな射
$U \to V \times_\mathcal{Y} \mathcal{X}$ を選ぶ。
$U \to V$ は平坦かつ局所有限表示であり、$U \to W$ は局所有限表示である。
また、$\mathcal{X} \to \mathcal{Y}$ が全射なので、$U \to V$ は
（全射の合成として）全射である。
したがって補題は代数空間の射の場合に帰着する。
代数空間の射の場合は次である：
Descent on Spaces, Lemma
\ref{spaces-descent-lemma-locally-finite-presentation-fppf-local-source}.
\end{proof}

\begin{lemma}
\label{stacks-morphisms-lemma-surjective-flat-locally-finite-presentation}
$f : \mathcal{X} \to \mathcal{Y}$ を全射、平坦、かつ局所有限表示な
代数スタックの射とする。このとき任意のスキーム $U$ と
$U$ 上の $\mathcal{Y}$ の対象 $y$ に対して、fppf 被覆
$\{U_i \to U\}$ と $U_i$ 上の $\mathcal{X}$ の対象 $x_i$ で、
$\mathcal{Y}_{U_i}$ において $f(x_i) \cong y|_{U_i}$ を満たすものが存在する。
\end{lemma}

\begin{proof}
$y$ を射 $U \to \mathcal{Y}$ とみなしてよい。
Properties of Stacks, Lemma
\ref{stacks-properties-lemma-base-change-surjective}
および
Lemmas \ref{stacks-morphisms-lemma-base-change-finite-presentation} と
\ref{stacks-morphisms-lemma-base-change-flat}
により、$\mathcal{X} \times_\mathcal{Y} U \to U$ は
全射、平坦、かつ局所有限表示である。
$V$ をスキームとし、$V \to \mathcal{X} \times_\mathcal{Y} U$ を
滑らかかつ全射とする。このとき
$V \to \mathcal{X} \times_\mathcal{Y} U$ も
全射、平坦、かつ局所有限表示である（次を参照：
Morphisms of Spaces, Lemmas
\ref{spaces-morphisms-lemma-smooth-flat} および
\ref{spaces-morphisms-lemma-smooth-locally-finite-presentation}).
したがって $V \to U$ も全射、平坦、かつ局所有限表示である。次を参照せよ：
Properties of Stacks, Lemma
\ref{stacks-properties-lemma-composition-surjective}
および
Lemmas \ref{stacks-morphisms-lemma-composition-finite-presentation} と
\ref{stacks-morphisms-lemma-composition-flat}.
ゆえに $\{V \to U\}$ は求める fppf 被覆であり、
$x : V \to \mathcal{X}$ は求める対象である。
\end{proof}

\begin{lemma}
\label{stacks-morphisms-lemma-surjective-family-flat-locally-finite-presentation}
$f_j : \mathcal{X}_j \to \mathcal{X}$, $j \in J$ を、各々が平坦かつ
局所有限表示であり、合わせて全射、すなわち
$|\mathcal{X}| = \bigcup |f_j|(|\mathcal{X}_j|)$.
を満たす代数スタックの射の族とする。このとき任意のスキーム $U$ と
$U$ 上の $\mathcal{X}$ の対象 $x$ に対して、fppf 被覆
$\{U_i \to U\}_{i \in I}$、写像 $a : I \to J$、および
$U_i$ 上の $\mathcal{X}_{a(i)}$ の対象 $x_i$ で、
$\mathcal{X}_{U_i}$ において $f_{a(i)}(x_i) \cong y|_{U_i}$ を
満たすものが存在する。
\end{lemma}

\begin{proof}
射 $\coprod_{j \in J} \mathcal{X}_j \to \mathcal{X}$ に
Lemma \ref{stacks-morphisms-lemma-surjective-flat-locally-finite-presentation}
を適用する。（ここには、我々の設定に起因する軽微な集合論上の問題があるが、
これは無視する。）最後に、射
$x_i : U_i \to \coprod_{j \in J} \mathcal{X}_j$ は非交和分解
$U_i = \coprod U_{i, j}$ と射 $U_{i, j} \to \mathcal{X}_j$ により
与えられることに注意する。すると fppf 被覆 $\{U_{i, j} \to U\}$ と
射 $U_{i, j} \to \mathcal{X}_j$ が求めるものである。
\end{proof}

\begin{lemma}
\label{stacks-morphisms-lemma-fppf-open}
$f : \mathcal{X} \to \mathcal{Y}$ を平坦かつ局所有限表示とする。
このとき $|f| : |\mathcal{X}| \to |\mathcal{Y}|$ は開写像である。
\end{lemma}

\begin{proof}
スキーム $V$ と滑らかで全射な射 $V \to \mathcal{Y}$ を選ぶ。
スキーム $U$ と滑らかで全射な射
$U \to V \times_\mathcal{Y} \mathcal{X}$ を選ぶ。
仮定により、スキームの射 $U \to V$ は平坦かつ局所有限表示である。
したがって
Morphisms, Lemma \ref{morphisms-lemma-fppf-open}.
により $U \to V$ は開写像である。$|\mathcal{Y}|$ 上の位相の構成から、
写像 $|V| \to |\mathcal{Y}|$ は開写像である。
写像 $|U| \to |\mathcal{X}|$ は全射である。
初等的な位相の議論により、これらの事実から結果が従う。
\end{proof}

\begin{lemma}
\label{stacks-morphisms-lemma-descent-quasi-compact}
$f : \mathcal{X} \to \mathcal{Y}$ を代数スタックの射とする。
$\mathcal{Z} \to \mathcal{Y}$ を全射、平坦、かつ局所有限表示な
代数スタックの射とする。基底変換
$\mathcal{Z} \times_\mathcal{Y} \mathcal{X} \to \mathcal{Z}$
が準コンパクトならば、$f$ は準コンパクトである。
\end{lemma}

\begin{proof}
$\mathcal{Y}'$ が準コンパクトであるような
$\mathcal{Y}' \to \mathcal{Y}$ が与えられたとき、
$\mathcal{Y}' \times_\mathcal{Y} \mathcal{X}$ が準コンパクトであることを
示さなければならない。
$\mathcal{Z}' = \mathcal{Z} \times_\mathcal{Y} \mathcal{Y}'$ と記す。
Lemma \ref{stacks-morphisms-lemma-fppf-open} により
$|\mathcal{Z}'| \to |\mathcal{Y}'|$ は開写像である。
したがって、$\mathcal{Y}'$ の上へ写る準コンパクトな開部分スタック
$\mathcal{W} \subset \mathcal{Z}'$ を見つけることができる。
$\mathcal{Z} \times_\mathcal{Y} \mathcal{X} \to \mathcal{Z}$
は準コンパクトなので、
$$
\mathcal{W} \times_\mathcal{Z} \mathcal{Z} \times_\mathcal{Y} \mathcal{X} =
\mathcal{W} \times_\mathcal{Y} \mathcal{X}
$$
は準コンパクトである。また、写像
$\mathcal{W} \times_\mathcal{Y} \mathcal{X} \to
\mathcal{Y}' \times_\mathcal{Y} \mathcal{X}$
は全射なので結論を得る。細部は省略する。
\end{proof}

\begin{lemma}
\label{stacks-morphisms-lemma-check-separated-on-ui-cover}
$f : \mathcal{X} \to \mathcal{Y}$、
$g : \mathcal{Y} \to \mathcal{Z}$ を合成可能な代数スタックの射とし、
その合成を $h = g \circ f : \mathcal{X} \to \mathcal{Z}$ とする。
$f$ が全射、平坦、局所有限表示、かつ普遍的単射であり、
$h$ が分離ならば、$g$ は分離である。
\end{lemma}

\begin{proof}
次の図式を考える：
$$
\xymatrix{
\mathcal{X} \ar[r]_\Delta \ar[rd] &
\mathcal{X} \times_\mathcal{Y} \mathcal{X} \ar[r] \ar[d] &
\mathcal{X} \times_\mathcal{Z} \mathcal{X} \ar[d] \\
& \mathcal{Y} \ar[r] & \mathcal{Y} \times_\mathcal{Z} \mathcal{Y}
}
$$
正方形は Cartesian である。下の水平射が固有であることを
示さなければならない。この射が代数空間で表現可能かつ局所有限型である
ことは既に分かっている（Lemma \ref{stacks-morphisms-lemma-properties-diagonal}）。
右の垂直射は全射、平坦、かつ局所有限表示なので、
右上の水平射が固有であることを示せば十分である
（Lemma \ref{stacks-morphisms-lemma-check-property-after-fppf-base-change}）。
$h$ は分離なので、上の水平射の合成は固有である。

\medskip\noindent
$f$ は普遍的単射なので、$\Delta$ は全射である
（Lemma \ref{stacks-morphisms-lemma-universally-injective}）。
$\Delta$ と射影
$\mathcal{X} \times_\mathcal{Y} \mathcal{X} \to \mathcal{X}$
との合成は恒等射だから、$\Delta$ は普遍的閉である。
Morphisms of Spaces, Lemma
\ref{spaces-morphisms-lemma-image-universally-closed-separated}
により、$\mathcal{X} \times_\mathcal{Y} \mathcal{X} \to
\mathcal{X} \times_\mathcal{Z} \mathcal{X}$
は分離であると結論する。実際、
$\mathcal{X} \to \mathcal{X} \times_\mathcal{Z} \mathcal{X}$ は分離である。
ここでは、代数空間の射の性質間の含意を、代数空間で表現可能な
代数スタックの射の性質間の同じ含意へ移せることを用いた。
これは Properties of Stacks, Section
\ref{stacks-properties-section-properties-morphisms}.
で論じられている。最後に同じ原理と
$\mathcal{X} \times_\mathcal{Y} \mathcal{X} \to
\mathcal{X} \times_\mathcal{Z} \mathcal{X}$ が固有であることを
Morphisms of Spaces, Lemma
\ref{spaces-morphisms-lemma-image-proper-is-proper}.
から結論する。
\end{proof}














\section{ガーベ}
\label{stacks-morphisms-section-gerbes}

\noindent
代数スタックの重要な型の一つは $[B/G]$ という形のスタックである。
ここで $B$ は代数空間、$G$ は $B$ 上平坦かつ局所有限表示な群代数空間であり、
$B$ 上自明に作用する。次を参照せよ：
Criteria for Representability, Lemma \ref{criteria-lemma-BG-algebraic}.
代数スタックが fppf 位相で局所的にこの形をもつとき、ガーベであることが分かる。
次を参照せよ：
Lemma \ref{stacks-morphisms-lemma-local-structure-gerbe}.
本節ではこの概念と対応する相対的概念を簡潔に論じる。

\begin{definition}
\label{stacks-morphisms-definition-gerbe}
$f : \mathcal{X} \to \mathcal{Y}$ を代数スタックの射とする。
$\mathcal{X}$ が $(\Sch/S)_{fppf}$ 上の亜群のスタックとして
$\mathcal{Y}$ 上のガーベであるとき、
$\mathcal{X}$ は $\mathcal{Y}$ {\it 上のガーベ} であるという。次を参照せよ：
Stacks, Definition \ref{stacks-definition-gerbe-over-stack-in-groupoids}.
ある代数空間 $X$ への射 $\mathcal{X} \to X$ で
$\mathcal{X}$ を $X$ 上のガーベにするものが存在するとき、
代数スタック $\mathcal{X}$ は {\it ガーベ} であるという。
\end{definition}

\noindent
$\mathcal{X}$ が $\mathcal{Y}$ 上のガーベであるという条件は、
与えられた代数スタックの基礎にある位相と圏論だけを用いて定義される。
しかし後で見るように、この条件には幾何学的帰結がある。例えば、
$\mathcal{X} \to \mathcal{Y}$ は全射、平坦、かつ局所有限表示となる。
次を参照せよ：
Lemma \ref{stacks-morphisms-lemma-gerbe-fppf}.
絶対的概念の解釈はより微妙である。というのも、最初は $X$ が一意に
定まることが明らかでないからである。実際には一意に定まる。

\begin{lemma}
\label{stacks-morphisms-lemma-gerbe-over-iso-classes}
$\mathcal{X}$ を代数スタックとする。$\mathcal{X}$ がガーベならば、前層
$$
(\Sch/S)_{fppf}^{opp} \to \textit{Sets}, \quad
U \mapsto \Ob(\mathcal{X}_U)/\!\!\cong
$$
の層化は代数空間であり、$\mathcal{X}$ はその上のガーベである。
\end{lemma}

\begin{proof}
（この証明では Section \ref{stacks-morphisms-section-conventions} で導入した
言葉の濫用が大いに役立つ。）
$X$ を代数空間とし、$\mathcal{X}$ を $X$ 上のガーベにする射
$\pi : \mathcal{X} \to X$ を選ぶ。
$X$ が補題に表示された前層 $\mathcal{F}$ の層化であることを示せば十分である。
写像 $c : \mathcal{F} \to X$ が存在することは明らかである。
Stacks, Lemma \ref{stacks-lemma-when-gerbe}
の性質 (2)(a) および (2)(b) を用いて、写像
$c^\# : \mathcal{F}^\# \to X$ が全射かつ単射、したがって同型であることを
示す。次を参照せよ：
Sites, Lemma \ref{sites-lemma-mono-epi-sheaves}.
全射性：$T$ をスキーム、$f : T \to X$ を射とする。
性質 (2)(a) により、fppf 被覆 $\{h_i : T_i \to T\}$ と射
$x_i : T_i \to \mathcal{X}$ で、$f \circ h_i$ が
$\pi \circ x_i$ に対応するものが存在する。したがって
$f|_{T_i}$ は $c$ の像に入る。
単射性：$T$ をスキームとし、$c \circ x = c \circ y$ を満たす射
$x, y : T \to \mathcal{X}$ を取る。(2)(b) により、被覆
$\{T_i \to T\}$ とファイバー圏 $\mathcal{X}_{T_i}$ における射
$x|_{T_i} \to y|_{T_i}$ を見つけることができる。
したがって制限 $x|_{T_i}, y|_{T_i}$ は $\mathcal{F}(T_i)$ で等しい。
これは望みどおり、$x, y$ が $T$ 上の $\mathcal{F}^\#$ の同じ切断を
与えることを示す。
\end{proof}

\begin{lemma}
\label{stacks-morphisms-lemma-base-change-gerbe}
次の図式を
$$
\xymatrix{
\mathcal{X}' \ar[r] \ar[d] & \mathcal{X} \ar[d] \\
\mathcal{Y}' \ar[r] & \mathcal{Y}
}
$$
を代数スタックのファイバー積とする。
$\mathcal{X}$ が $\mathcal{Y}$ 上のガーベならば、
$\mathcal{X}'$ は $\mathcal{Y}'$ 上のガーベである。
\end{lemma}

\begin{proof}
定義と
Stacks, Lemma \ref{stacks-lemma-base-change-gerbe}.
から直ちに従う。
\end{proof}

\begin{lemma}
\label{stacks-morphisms-lemma-composition-gerbe}
$\mathcal{X} \to \mathcal{Y}$ および $\mathcal{Y} \to \mathcal{Z}$
を代数スタックの射とする。
$\mathcal{X}$ が $\mathcal{Y}$ 上のガーベで、
$\mathcal{Y}$ が $\mathcal{Z}$ 上のガーベならば、
$\mathcal{X}$ は $\mathcal{Z}$ 上のガーベである。
\end{lemma}

\begin{proof}
次から直ちに従う：
Stacks, Lemma \ref{stacks-lemma-composition-gerbe}.
\end{proof}

\begin{lemma}
\label{stacks-morphisms-lemma-gerbe-descent}
次の図式を
$$
\xymatrix{
\mathcal{X}' \ar[r] \ar[d] & \mathcal{X} \ar[d] \\
\mathcal{Y}' \ar[r] & \mathcal{Y}
}
$$
を代数スタックのファイバー積とする。
$\mathcal{Y}' \to \mathcal{Y}$ が全射、平坦、かつ局所有限表示であり、
$\mathcal{X}'$ が $\mathcal{Y}'$ 上のガーベならば、
$\mathcal{X}$ は $\mathcal{Y}$ 上のガーベである。
\end{lemma}

\begin{proof}
次から直ちに従う：
Lemma \ref{stacks-morphisms-lemma-surjective-flat-locally-finite-presentation}
および
Stacks, Lemma \ref{stacks-lemma-gerbe-descent}.
\end{proof}

\begin{lemma}
\label{stacks-morphisms-lemma-gerbe-with-section}
$\pi : \mathcal{X} \to U$ を代数スタックから代数空間への射とし、
$x : U \to \mathcal{X}$ を $\pi$ の切断とする。
$G = \mathit{Isom}_\mathcal{X}(x, x)$ と置く。次を参照せよ：
Definition \ref{stacks-morphisms-definition-isom}.
$\mathcal{X}$ が $U$ 上のガーベならば、
\begin{enumerate}
\item 亜群のスタックの標準的な同値
$$
x_{can} : [U/G] \longrightarrow \mathcal{X}.
$$
が存在する。ここで $[U/G]$ は $G$ の $U$ への自明な作用に対する商スタックである。
\item $G \to U$ は平坦かつ局所有限表示である。
\item $U \to \mathcal{X}$ は全射、平坦、かつ局所有限表示である。
\end{enumerate}
\end{lemma}

\begin{proof}
$R = U \times_{x, \mathcal{X}, x} U$ と置く。
射 $R \to U \times U$ は $U \times_U U = U$ を経由するので、
対角射 $\Delta_U : U \to U \times U$ を経由する。したがって
\begin{align*}
G & = \mathit{Isom}_\mathcal{X}(x, x) \\
& = U \times_{x, \mathcal{X}} \mathcal{I}_\mathcal{X} \\
& = U \times_{x, \mathcal{X}}
(\mathcal{X}
\times_{\Delta, \mathcal{X} \times_S \mathcal{X}, \Delta}
\mathcal{X}) \\
& = (U \times_{x, \mathcal{X}, x} U) \times_{U \times U, \Delta_U} U \\
& = R \times_{U \times U, \Delta_U} U \\
& = R
\end{align*}
により $R = G$ である。第四の等号には
Categories, Lemma \ref{categories-lemma-diagonal-2}.
を用いる。$t, s : R \to U$ を射影とする。
Algebraic Stacks, Lemma \ref{algebraic-lemma-map-space-into-stack}
で $R$ 上に構成された合成則 $c : R \times_{s, U, t} R \to R$ は
$G$ 上の群則と一致する（証明は省略する）。したがって
Algebraic Stacks, Lemma \ref{algebraic-lemma-map-space-into-stack}
により、$(\Sch/S)_{fppf}$ 上の亜群のスタックの標準的な
充満忠実 $1$-射
$$
x_{can} : [U/G] \longrightarrow \mathcal{X}
$$
を得る。これが同値であることを示すには、本質的に全射であることを
示せば十分である。そのためには、スキーム $T$ 上の $\mathcal{X}$ の
任意の対象が、射 $T \to U$ を介して $x$ から fppf 局所的に得られる
ことを示せば十分である。次を参照せよ：
Stacks, Lemma \ref{stacks-lemma-characterize-essentially-surjective-when-ff}.
しかし、これは $\pi$ が $\mathcal{X}$ を $U$ 上のガーベにするという
条件から従う。次の性質 (2)(a) を参照せよ：
Stacks, Lemma \ref{stacks-lemma-when-gerbe}.

\medskip\noindent
Criteria for Representability, Lemma \ref{criteria-lemma-BG-algebraic}
により、$G \to U$ は平坦かつ局所有限表示である。
最後に、
Criteria for Representability, Lemma
\ref{criteria-lemma-flat-quotient-flat-presentation}.
により、$U \to \mathcal{X}$ は全射、平坦、かつ局所有限表示である。
\end{proof}

\begin{lemma}
\label{stacks-morphisms-lemma-local-structure-gerbe}
$\pi : \mathcal{X} \to \mathcal{Y}$ を代数スタックの射とする。
以下は同値である：
\begin{enumerate}
\item $\mathcal{X}$ は $\mathcal{Y}$ 上のガーベである。
\item 代数空間 $U$、$U$ 上平坦かつ局所有限表示な群代数空間 $G$、
および全射、平坦、かつ局所有限表示な射 $U \to \mathcal{Y}$ で、
$U$ 上 $\mathcal{X} \times_\mathcal{Y} U \cong [U/G]$ を満たすものが存在する。
\end{enumerate}
\end{lemma}

\begin{proof}
(2) を仮定する。
Lemma \ref{stacks-morphisms-lemma-gerbe-descent}
により、(1) を示すには $[U/G]$ が $U$ 上のガーベであることを示せば十分である。
これは次から直ちに従う：
Groupoids in Spaces, Lemma \ref{spaces-groupoids-lemma-group-quotient-gerbe}.

\medskip\noindent
(1) を仮定する。$\pi$ の任意の基底変換はガーベである。次を参照せよ：
Lemma \ref{stacks-morphisms-lemma-base-change-gerbe}.
まずスキーム $V$ と全射滑らかな射 $V \to \mathcal{Y}$ を選ぶ。
したがって $\pi : \mathcal{X} \to V$ がスキーム上のガーベであると
仮定してよい。これは、ファイバー圏 $\mathcal{X}_{V_i}$ が空でないような
fppf 被覆 $\{V_i \to V\}$ が存在することを意味する。次を参照せよ：
Stacks, Lemma \ref{stacks-lemma-when-gerbe} (2)(a).
$U = \coprod V_i \to V$ は全射、平坦、かつ局所有限表示であることに
注意する。したがって $V$ を $U$ で置き換え、
$\pi : \mathcal{X} \to U$ がスキーム $U$ 上のガーベであり、
$U$ 上の $\mathcal{X}$ の対象 $x$ が存在すると仮定してよい。
Lemma \ref{stacks-morphisms-lemma-gerbe-with-section}
により、ある $U$ 上平坦かつ局所有限表示な群代数空間 $G$ に対して、
$U$ 上 $\mathcal{X} = [U/G]$ である。
\end{proof}

\begin{lemma}
\label{stacks-morphisms-lemma-gerbe-fppf}
$\pi : \mathcal{X} \to \mathcal{Y}$ を代数スタックの射とする。
$\mathcal{X}$ が $\mathcal{Y}$ 上のガーベならば、
$\pi$ は全射、平坦、かつ局所有限表示である。
\end{lemma}

\begin{proof}
Properties of Stacks, Lemma
\ref{stacks-properties-lemma-descent-surjective}
ならびに
Lemmas \ref{stacks-morphisms-lemma-descent-flat} と
\ref{stacks-morphisms-lemma-descent-finite-presentation}
により、$\pi$ を全射、平坦、かつ局所有限表示な射
$\mathcal{Y}' \to \mathcal{Y}$ による基底変換で置き換えた後に
補題を示せば十分である。
Lemma \ref{stacks-morphisms-lemma-local-structure-gerbe}
により、$\mathcal{Y} = U$ は代数空間で、$U$ 上
$\mathcal{X} = [U/G]$ であると仮定してよい。
このとき $U \to [U/G]$ は全射、平坦、かつ局所有限表示である。
次を参照せよ：
Lemma \ref{stacks-morphisms-lemma-gerbe-with-section}.
これと
Properties of Stacks,
Lemma \ref{stacks-properties-lemma-surjective-permanence}
ならびに
Lemmas \ref{stacks-morphisms-lemma-flat-permanence} と
\ref{stacks-morphisms-lemma-flat-finite-presentation-permanence}.
により、$\pi$ は全射、平坦、かつ局所有限表示である。
\end{proof}

\begin{proposition}
\label{stacks-morphisms-proposition-when-gerbe}
$\mathcal{X}$ を代数スタックとする。以下は同値である：
\begin{enumerate}
\item $\mathcal{X}$ はガーベである。
\item $\mathcal{I}_\mathcal{X} \to \mathcal{X}$ は平坦かつ局所有限表示である。
\end{enumerate}
\end{proposition}

\begin{proof}
(1) を仮定する。代数空間 $X$ への射 $\mathcal{X} \to X$ で、
$\mathcal{X}$ を $X$ 上のガーベにするものを選ぶ。
$X' \to X$ を全射、平坦、かつ局所有限表示な射とし、
$\mathcal{X}' = X' \times_X \mathcal{X}$ と置く。
Lemma \ref{stacks-morphisms-lemma-base-change-gerbe}.
により、$\mathcal{X}'$ は $X'$ 上のガーベであることに注意する。
次の二つの正方形はいずれもファイバー積である：
$$
\xymatrix{
\mathcal{I}_{\mathcal{X}'} \ar[r] \ar[d] &
\mathcal{X}' \ar[r] \ar[d] & X' \ar[d] \\
\mathcal{I}_\mathcal{X} \ar[r] &
\mathcal{X} \ar[r] & X
}
$$
次を参照せよ：
Lemma \ref{stacks-morphisms-lemma-cartesian-square-inertia}.
したがって、
Lemmas \ref{stacks-morphisms-lemma-descent-flat} と
\ref{stacks-morphisms-lemma-descent-finite-presentation}.
により、$\mathcal{I}_\mathcal{X} \to \mathcal{X}$ が平坦かつ
局所有限表示であることを示すには、このような基底変換の後に示せば十分である。
そこで
Lemma \ref{stacks-morphisms-lemma-local-structure-gerbe}
を適用し、$\mathcal{X} = [U/G]$ と仮定できる。
Lemma \ref{stacks-morphisms-lemma-gerbe-with-section}
により、$G$ は $U$ 上平坦かつ局所有限表示であり、
$x : U \to [U/G]$ は全射、平坦、かつ局所有限表示である。
さらに、$x$ による $\mathcal{I}_\mathcal{X}$ の引き戻しは $G$ である。
再び降下、すなわち
Lemmas \ref{stacks-morphisms-lemma-descent-flat} と
\ref{stacks-morphisms-lemma-descent-finite-presentation}.
により、(2) が成り立つと結論する。

\medskip\noindent
逆に (2) を仮定する。滑らかな表示 $\mathcal{X} = [U/R]$ を選ぶ。
次を参照せよ：
Algebraic Stacks, Section \ref{algebraic-section-stack-to-presentation}.
$(U, R, s, t, c, e, i)$ を亜群とし、その安定化群代数空間を
$G \to U$ と記す。次を参照せよ：
Groupoids in Spaces, Definition
\ref{spaces-groupoids-definition-stabilizer-groupoid}.
Lemma \ref{stacks-morphisms-lemma-presentation-inertia}
により、$G \to U$ は $\mathcal{I}_\mathcal{X} \to \mathcal{X}$ の
基底変換として平坦かつ局所有限表示である。次を参照せよ：
Lemmas \ref{stacks-morphisms-lemma-base-change-flat} と
\ref{stacks-morphisms-lemma-base-change-finite-presentation}.
次の $G$ の $R$ への作用を考える：
$$
a : G \times_{U, t} R \to R, \quad (g, r) \mapsto c(g, r)
$$
$R$ は亜群なので、この作用は任意のスキーム $T$ に対し
$T$-値点上自由である。したがって $R' = R/G$ は代数空間であり、
Bootstrap, Lemma \ref{bootstrap-lemma-quotient-free-action}.
により、商射 $\pi : R \to R'$ は全射、平坦、かつ局所有限表示である。
射影 $s, t : R \to U$ は $G$-不変なので、
$s = s' \circ \pi$ および $t = t' \circ \pi$ を満たす射
$s' , t' : R' \to U$ を得る。
$s, t : R \to U$ は平坦かつ局所有限表示なので、
$s', t'$ も平坦かつ局所有限表示である。次を参照せよ：
Morphisms of Spaces, Lemmas
\ref{spaces-morphisms-lemma-flat-permanence} および
Descent on Spaces, Lemma
\ref{spaces-descent-lemma-locally-finite-presentation-fppf-local-source}.
射
$$
j' = (t', s') : R' \longrightarrow U \times U.
$$
を考える。これは単射であると主張する。実際、$T$ をスキームとし、
$a, b : T \to R'$ を $U \times U$ で同じ像をもつ射とする。
商 $R' = R/G$ の定義により、fppf 被覆 $\{h_j : T_j \to T\}$ と射
$a_j, b_j : T_j \to R$ で、
$a \circ h_j = \pi \circ a_j$ および
$b \circ h_j = \pi \circ b_j$ を満たすものが存在する。
$a_j, b_j$ は $U \times U$ で同じ像をもつので、
$g_j = c(a_j, i(b_j))$ は $c(g_j, b_j) = a_j$ を満たす
$G$ の $T_j$-値点である。言い換えれば、$a_j$ と $b_j$ は
$R'$ で同じ像をもち、主張が証明された。
$j : R \to U \times U$ は前同値関係であり（次を参照：
Groupoids in Spaces, Lemma
\ref{spaces-groupoids-lemma-groupoid-pre-equivalence})
かつ $R \to R'$ は層の写像として全射なので、
$j' : R' \to U \times U$ は同値関係である。したがって
Bootstrap, Theorem \ref{bootstrap-theorem-final-bootstrap}
により、$X = U/R'$ は代数空間である。
最後に、射
$$
\mathcal{X} = [U/R] \longrightarrow X = U/R'
$$
は $\mathcal{X}$ を $X$ 上のガーベにすると主張する。これは
Groupoids in Spaces, Lemma \ref{spaces-groupoids-lemma-when-gerbe}
から従う。実際、$R \to R'$ は全射、平坦、かつ局所有限表示である
（必要なら、これが要求される仮定を含意することを見るために
Bootstrap, Lemma
\ref{bootstrap-lemma-surjective-flat-locally-finite-presentation}
を用いよ）。
\end{proof}

\begin{lemma}
\label{stacks-morphisms-lemma-gerbe-isom-fppf}
$f : \mathcal{X} \to \mathcal{Y}$ を、$\mathcal{X}$ を
$\mathcal{Y}$ 上のガーベにする代数スタックの射とする。このとき
\begin{enumerate}
\item $\mathcal{I}_{\mathcal{X}/\mathcal{Y}} \to \mathcal{X}$
は平坦かつ局所有限表示である。
\item $\mathcal{X} \to \mathcal{X} \times_\mathcal{Y} \mathcal{X}$
は全射、平坦、かつ局所有限表示である。
\item 代数空間 $T_i$, $i = 1, 2$ と射
$x_i : T_i \to \mathcal{X}$ が与えられ、
$y_i = f \circ x_i$ とすると、射
$$
T_1 \times_{x_1, \mathcal{X}, x_2} T_2 \longrightarrow
T_1 \times_{y_1, \mathcal{Y}, y_2} T_2
$$
は全射、平坦、かつ局所有限表示である。
\item 代数空間 $T$ と射
$x_i : T \to \mathcal{X}$, $i = 1, 2$ が与えられ、
$y_i = f \circ x_i$ とすると、射
$$
\mathit{Isom}_\mathcal{X}(x_1, x_2) \longrightarrow
\mathit{Isom}_\mathcal{Y}(y_1, y_2)
$$
は全射、平坦、かつ局所有限表示である。
\end{enumerate}
\end{lemma}

\begin{proof}
(1) の証明。
スキーム $Y$ と全射滑らかな射 $Y \to \mathcal{Y}$ を選ぶ。
$\mathcal{X}' = \mathcal{X} \times_\mathcal{Y} Y$ と置く。
Lemma \ref{stacks-morphisms-lemma-cartesian-square-inertia} により、Cartesian な正方形
$$
\xymatrix{
\mathcal{I}_{\mathcal{X}'} \ar[r] \ar[d] &
\mathcal{X}' \ar[r] \ar[d] & Y \ar[d] \\
\mathcal{I}_{\mathcal{X}/\mathcal{Y}} \ar[r] &
\mathcal{X} \ar[r] & \mathcal{Y}
}
$$
を得る。Lemmas \ref{stacks-morphisms-lemma-descent-flat} と
\ref{stacks-morphisms-lemma-descent-finite-presentation}
により、$\mathcal{I}_{\mathcal{X}'} \to \mathcal{X}'$ が
平坦かつ局所有限表示であることを示せば十分である。
これは Proposition \ref{stacks-morphisms-proposition-when-gerbe} から従う
（Lemma \ref{stacks-morphisms-lemma-base-change-gerbe} により、
$\mathcal{X}'$ は $Y$ 上のガーベだからである）。

\medskip\noindent
(2) の証明。上の記法のもとで、ある $Y$ 上平坦かつ局所有限表示な
群代数空間 $G$ に対して $\mathcal{X}' = [Y/G]$ と仮定してよいことに
注意する。次を参照せよ：
Lemma \ref{stacks-morphisms-lemma-local-structure-gerbe}.
射
$\Delta : \mathcal{X} \to \mathcal{X} \times_\mathcal{Y} \mathcal{X}$
を $\mathcal{Y}$ 上で射 $Y \to \mathcal{Y}$ により基底変換したものは、射
$\Delta' : \mathcal{X}' \to \mathcal{X}' \times_Y \mathcal{X}'$.
である。したがって、$\Delta'$ が全射、平坦、かつ局所有限表示であることを
示せば十分である（Lemmas \ref{stacks-morphisms-lemma-descent-flat} と
\ref{stacks-morphisms-lemma-descent-finite-presentation}).
を参照）。言い換えれば、
$$
[Y/G] \longrightarrow [Y/G \times_Y G]
$$
が全射、平坦、かつ局所有限表示であることを示さなければならない。
これは、全射、平坦、かつ局所有限表示な射
$Y \to [Y/G \times_Y G]$ による基底変換が射 $G \to Y$ だからである。

\medskip\noindent
(3) の証明。図式
$$
\xymatrix{
T_1 \times_{x_1, \mathcal{X}, x_2} T_2 \ar[d] \ar[r] &
T_1 \times_{y_1, \mathcal{Y}, y_2} T_2 \ar[d] \\
\mathcal{X} \ar[r] & \mathcal{X} \times_\mathcal{Y} \mathcal{X}
}
$$
は Cartesian であることに注意する。したがって (3) は (2) から従う。

\medskip\noindent
(4) の証明。実際、
$$
\mathit{Isom}_\mathcal{X}(x_1, x_2) =
(T \times_{x_1, \mathcal{X}, x_2} T) \times_{T \times T, \Delta_T} T
$$
なので、(4) の射は (3) の射の基底変換である。
\end{proof}

\begin{proposition}
\label{stacks-morphisms-proposition-when-gerbe-over}
$f : \mathcal{X} \to \mathcal{Y}$ を代数スタックの射とする。
以下は同値である：
\begin{enumerate}
\item $\mathcal{X}$ は $\mathcal{Y}$ 上のガーベである。
\item $f : \mathcal{X} \to \mathcal{Y}$ および
$\Delta : \mathcal{X} \to \mathcal{X} \times_\mathcal{Y} \mathcal{X}$
は全射、平坦、かつ局所有限表示である。
\end{enumerate}
\end{proposition}

\begin{proof}
含意 (1) $\Rightarrow$ (2) は Lemmas \ref{stacks-morphisms-lemma-gerbe-fppf} および
\ref{stacks-morphisms-lemma-gerbe-isom-fppf} から従う。

\medskip\noindent
(2) を仮定する。全射、平坦、かつ局所有限表示な射
$\mathcal{Y}' \to \mathcal{Y}$ による $f$ の基底変換について
(1) を示せば十分である。Lemma \ref{stacks-morphisms-lemma-gerbe-descent} を参照せよ
（$f$ の対角射の基底変換は、基底変換した射の対角射であることに注意する）。
したがって $\mathcal{Y}$ はスキーム $Y$ であると仮定してよい。
この場合、$\mathcal{I}_\mathcal{X} \to \mathcal{X}$ は $\Delta$ の
基底変換であり、Proposition \ref{stacks-morphisms-proposition-when-gerbe} により
$\mathcal{X}$ はガーベである。
$\mathcal{X}$ が $Y$ 上のガーベであることをなお示さなければならない。
Lemma \ref{stacks-morphisms-lemma-gerbe-over-iso-classes}
の射 $\mathcal{X} \to X$ を取る。この射は $\mathcal{X}$ を、
$\mathcal{X}$ の対象の同型類を分類する代数空間 $X$ 上のガーベにする。
$f : \mathcal{X} \to Y$ が
$\mathcal{X} \to X \to Y$ と分解することは明らかである。
$f$ は全射、平坦、かつ局所有限表示なので、$X \to Y$ は
fppf 層の写像として全射である（例えば
Lemma \ref{stacks-morphisms-lemma-surjective-flat-locally-finite-presentation}).
を用いよ）。一方、$X \to Y$ は単射でもある。
実際、任意のスキーム $T$ と、$Y$ の同じ点へ写る $X$ の任意の二つの
$T$-値点 $x_1, x_2$ に対し、まず $T$ 上 fppf 局所的に
$x_1, x_2$ を $T$ 上の $\mathcal{X}$ の対象 $\xi_1, \xi_2$ へ持ち上げ、
次に
$\Delta : \mathcal{X} \to \mathcal{X} \times_Y \mathcal{X}$
が全射、平坦、かつ局所有限表示であるという仮定から、
$\xi_1$ と $\xi_2$ が fppf 局所的に同型であることを導ける。
したがって $X$ の構成により $x_1 = x_2$ である。
ゆえに $X = Y$ であり、証明が完了する。
\end{proof}

\noindent
ここまでで、剰余ガーベが存在するとき、それがガーベであることを
証明するのに十分な道具立てが整った。

\begin{lemma}
\label{stacks-morphisms-lemma-noetherian-singleton-stack-gerbe}
$\mathcal{Z}$ を、$|\mathcal{Z}|$ が一点集合であるような
被約局所 Noether 代数スタックとする。このとき $\mathcal{Z}$ は、
$|Z|$ が一点集合であるような被約局所 Noether 代数空間 $Z$ 上の
ガーベである。
\end{lemma}

\begin{proof}
Properties of Stacks, Lemma \ref{stacks-properties-lemma-unique-point-better}
により、$k$ を体として全射、平坦、かつ局所有限表示な射
$\Spec(k) \to \mathcal{Z}$ が存在する。
このとき $\mathcal{I}_Z \times_\mathcal{Z} \Spec(k) \to \Spec(k)$ は
代数空間で表現可能かつ局所有限型である
（$\mathcal{I}_\mathcal{Z} \to \mathcal{Z}$ の基底変換であることと、
Lemmas \ref{stacks-morphisms-lemma-inertia} および \ref{stacks-morphisms-lemma-base-change-finite-type}
を参照）。したがって局所有限表示である。次を参照せよ：
Morphisms of Spaces, Lemma
\ref{spaces-morphisms-lemma-noetherian-finite-type-finite-presentation}.
また $k$ は体なので、もちろん平坦でもある。したがって
Lemmas \ref{stacks-morphisms-lemma-descent-flat} と
\ref{stacks-morphisms-lemma-descent-finite-presentation}
を適用すると、$\mathcal{I}_\mathcal{Z} \to \mathcal{Z}$ は
平坦かつ局所有限表示であることが分かる。
Proposition \ref{stacks-morphisms-proposition-when-gerbe}.
により、$\mathcal{Z}$ はガーベであると結論する。
$\mathcal{Z}$ が $Z$ 上のガーベとなるような代数空間への射
$\pi : \mathcal{Z} \to Z$ を取る。このとき
Lemma \ref{stacks-morphisms-lemma-gerbe-fppf}.
により、$\pi$ は全射、平坦、かつ局所有限表示である。
したがって合成 $\Spec(k) \to Z$ も全射、平坦、かつ局所有限表示である。
次を参照せよ：
Properties of Stacks,
Lemma \ref{stacks-properties-lemma-composition-surjective}
および
Lemmas \ref{stacks-morphisms-lemma-composition-flat} と
\ref{stacks-morphisms-lemma-composition-finite-presentation}.
Properties of Stacks, Lemma \ref{stacks-properties-lemma-unique-point-better}
により、$|Z|$ は一点集合であり、$Z$ は局所 Noether かつ被約である。
\end{proof}

\begin{lemma}
\label{stacks-morphisms-lemma-gerbe-bijection-points}
$f : \mathcal{X} \to \mathcal{Y}$ を代数スタックの射とする。
$\mathcal{X}$ が $\mathcal{Y}$ 上のガーベならば、
$f$ は普遍同相射である。
\end{lemma}

\begin{proof}
Lemma \ref{stacks-morphisms-lemma-base-change-gerbe} により、$f$ に関する仮定は
基底変換で保たれる。したがって写像
$|\mathcal{X}| \to |\mathcal{Y}|$ が位相空間の同相写像であることを
示せば十分である。$k$ を体とし、$y$ を $\Spec(k)$ 上の
$\mathcal{Y}$ の対象とする。
Stacks, Lemma \ref{stacks-lemma-when-gerbe} の性質 (2)(a)
の性質 (2)(a) により、fppf 被覆 $\{T_i \to \Spec(k)\}$ と
$T_i$ 上の $\mathcal{X}$ の対象 $x_i$ で
$f(x_i) \cong y|_{T_i}$ を満たすものが存在する。
$T_i \not = \emptyset$ となる $i$ を選び、ある体 $K$ に対して
射 $\Spec(K) \to T_i$ を選ぶ。このとき $k \subset K$ であり、
$x_i|_K$ は $y|_K$ の上にある $\mathcal{X}$ の対象である。
したがって $|\mathcal{Y}| \to |\mathcal{X}|$ は全射である。
写像 $|\mathcal{Y}| \to |\mathcal{X}|$ は単射でもある。
実際、$x, x'$ を $\Spec(k)$ 上の $\mathcal{X}$ の対象とし、
その像 $f(x), f(x')$ が $\mathcal{Y}$ において拡大後に同型になるとする。
Stacks, Lemma \ref{stacks-lemma-when-gerbe} の性質 (2)(b)
の性質 (2)(b) により、$x$ と $x'$ が同型になるような $k$ の拡大が
存在する（細部は省略する）。
したがって $|\mathcal{X}| \to |\mathcal{Y}|$ は連続かつ全単射であり、
さらに開写像であることを示せば十分である。これは
Lemmas \ref{stacks-morphisms-lemma-gerbe-fppf} および \ref{stacks-morphisms-lemma-fppf-open} から従う。
\end{proof}

\begin{lemma}
\label{stacks-morphisms-lemma-gerbe-diagonal-quasi-compact}
$f : \mathcal{X} \to \mathcal{Y}$ を、
$\mathcal{X}$ が $\mathcal{Y}$ 上のガーベとなるような
代数スタックの射とする。
$\Delta_\mathcal{X}$ が準コンパクトならば、
$\Delta_\mathcal{Y}$ も準コンパクトである。
\end{lemma}

\begin{proof}
次の図式を考える：
$$
\xymatrix{
\mathcal{X} \ar[r] &
\mathcal{X} \times_\mathcal{Y} \mathcal{X} \ar[r] \ar[d] &
\mathcal{X} \times \mathcal{X} \ar[d] \\
&
\mathcal{Y} \ar[r] &
\mathcal{Y} \times \mathcal{Y}
}
$$
Proposition \ref{stacks-morphisms-proposition-when-gerbe-over} により、
左上の射は全射である。上の水平射の合成は準コンパクトなので、
Lemma \ref{stacks-morphisms-lemma-surjection-from-quasi-compact}.
により、右上の射は準コンパクトである。
正方形は Cartesian で、右の垂直射は全射、平坦、かつ局所有限表示である。
したがって Lemma \ref{stacks-morphisms-lemma-descent-quasi-compact} により結論を得る。
\end{proof}

\noindent
次の補題は、ガーベである任意の代数スタックのすべての点に対して
剰余ガーベが存在することを述べる。

\begin{lemma}
\label{stacks-morphisms-lemma-gerbe-residual-gerbe-exists}
$\mathcal{X}$ を代数スタックとする。$\mathcal{X}$ がガーベならば、
すべての $x \in |\mathcal{X}|$ に対して、
$x$ における $\mathcal{X}$ の剰余ガーベが存在する。
\end{lemma}

\begin{proof}
$\pi : \mathcal{X} \to X$ を、$\mathcal{X}$ から代数空間 $X$ への射であって、
$\mathcal{X}$ を $X$ 上のガーベにするものとする。
$Z_x \to X$ を $x$ における $X$ の剰余空間とする。
次を参照せよ：
Decent Spaces, Definition \ref{decent-spaces-definition-residual-space}.
$\mathcal{Z} = \mathcal{X} \times_X Z_x$ と置く。
Lemma \ref{stacks-morphisms-lemma-base-change-gerbe}
により、代数スタック $\mathcal{Z}$ は $Z_x$ 上のガーベである。
したがって $|\mathcal{Z}| = |Z_x|$
（Lemma \ref{stacks-morphisms-lemma-gerbe-bijection-points}）は一点集合である。
$\mathcal{Z} \to Z_x$ は $\pi$ の基底変換として局所有限表示なので
（Lemmas \ref{stacks-morphisms-lemma-gerbe-fppf} および
\ref{stacks-morphisms-lemma-base-change-finite-presentation} を参照）、
$\mathcal{Z}$ は局所 Noether である。次を参照せよ：
Lemma \ref{stacks-morphisms-lemma-locally-finite-type-locally-noetherian}.
したがって $x$ における $\mathcal{X}$ の剰余ガーベ
$\mathcal{Z}_x$ は存在し、代数スタック $\mathcal{Z}$ の被約化
$\mathcal{Z}_x = \mathcal{Z}_{red}$ に等しい。実際、上で
$|\mathcal{Z}_{red}|$ は $x \in |\mathcal{X}|$ に写る一点集合であることを見た。
これは構成により被約であり、また局所 Noether である
（局所 Noether 代数スタックの被約化は局所 Noether である。
細部は省略する）。
\end{proof}









\section{ガーベによる層化}
\label{stacks-morphisms-section-stratify}

\noindent
本節の目標は、多くの代数スタック $\mathcal{X}$ が、
各 $\mathcal{X}_i$ がガーベであるような局所閉部分スタック
$\mathcal{X}_i \subset \mathcal{X}$ による「層化」をもつことを示すことである。
これは、ある意味でガーベが任意の代数スタックを構成する基本部品であることを示す。
ここで層化とは、本節では
$$
|\mathcal{X}| = \bigcup\nolimits_i |\mathcal{X}_i|
$$
が $\mathcal{X}$ に付随する位相空間の層化であることだけを意味し、
それ以上ではないことに注意する。したがってこの層化を調べるために
$\mathcal{X}$ をその被約化で置き換えても差し支えない（
Properties of Stacks, Section \ref{stacks-properties-section-reduced})
を参照）。

\medskip\noindent
次の命題は、$\mathcal{X}$ の被約化に（ほとんど常に）
稠密な開部分スタックが存在することを述べる。

\begin{proposition}
\label{stacks-morphisms-proposition-open-stratum}
$\mathcal{X}$ を、$\mathcal{I}_\mathcal{X} \to \mathcal{X}$ が
準コンパクトであるような被約代数スタックとする。
このときガーベである稠密な開部分スタック
$\mathcal{U} \subset \mathcal{X}$ が存在する。
\end{proposition}

\begin{proof}
Proposition \ref{stacks-morphisms-proposition-when-gerbe}
により、$\mathcal{I}_\mathcal{U} \to \mathcal{U}$ が平坦かつ
局所有限表示となる稠密な開部分スタック $\mathcal{U}$ を
見つければ十分である。次に注意する：
$\mathcal{I}_\mathcal{U} =
\mathcal{I}_\mathcal{X} \times_\mathcal{X} \mathcal{U}$
。次を参照せよ：
Lemma \ref{stacks-morphisms-lemma-cartesian-square-inertia}.

\medskip\noindent
表示 $\mathcal{X} = [U/R]$ を選ぶ。$G \to U$ を亜群 $R$ の
安定化群代数空間とする。
Lemma \ref{stacks-morphisms-lemma-presentation-inertia}
により、$G \to U$ は $\mathcal{I}_\mathcal{X} \to \mathcal{X}$ の
基底変換であるから、準コンパクト（仮定による）かつ局所有限型
（Lemma \ref{stacks-morphisms-lemma-inertia} による）である。
制限 $G_W \to W$ が平坦かつ局所有限表示となる最大の開部分スキーム
$W \subset U$（空でもよい）を取る
（$W$ の存在証明は省略する。ヒント：これらの性質が局所的であることを用いよ）。
Morphisms of Spaces, Proposition
\ref{spaces-morphisms-proposition-generic-flatness-reduced}
により、$W \subset U$ は稠密である。また、
More on Groupoids in Spaces, Lemma
\ref{spaces-more-groupoids-lemma-property-G-invariant}.
により、$W \subset U$ は $R$-不変であることに注意する。したがって
Properties of Stacks, Lemma
\ref{stacks-properties-lemma-substacks-presentation}.
により、$W$ は開部分スタック $\mathcal{U} \subset \mathcal{X}$ に対応する。
$|U| \to |\mathcal{X}|$ は開写像で、$|W| \subset |U|$ は稠密なので、
$\mathcal{U}$ は $\mathcal{X}$ で稠密である。
最後に、全射滑らかな射 $W \to \mathcal{U}$ による
$\mathcal{I}_\mathcal{U} \to \mathcal{U}$ の基底変換は、
構成により平坦かつ局所有限表示な射 $G_W \to W$ である。
したがって元の射も平坦かつ局所有限表示である。次を参照せよ：
Lemmas \ref{stacks-morphisms-lemma-descent-flat} と
\ref{stacks-morphisms-lemma-descent-finite-presentation}.
\end{proof}

\noindent
上の命題から、準コンパクトな慣性をもつ代数スタックでは任意の点が
剰余ガーベをもつことが直ちに従う。これは
Lemma \ref{stacks-morphisms-lemma-every-point-residual-gerbe} で示す。
一方、ガーベによる有限層化が常に存在するわけではない。例を挙げる。

\begin{example}
\label{stacks-morphisms-example-infinite-stratification}
\leavevmode\par\noindent
$k$ を体とする。
$U = \Spec(k[x_0, x_1, x_2, \ldots])$ と取り、
$\mathbf{G}_m$ を
$t(x_0, x_1, x_2, \ldots) = (tx_0, t^p x_1, t^{p^2} x_2, \ldots)$
により作用させる。ここで $p$ は素数である。
$\mathcal{X} = [U/\mathbf{G}_m]$ と置く。これは代数スタックである。
$\mathcal{X}$ には次の層による層化がある：
\begin{enumerate}
\item $\mathcal{X}_0$ は $x_0$ が零でない部分である。
\item $\mathcal{X}_1$ は $x_0$ が零で $x_1$ が零でない部分である。
\item $\mathcal{X}_2$ は $x_0, x_1$ が零で $x_2$ が零でない部分である。
\item 以下同様である。
\item $\mathcal{X}_{\infty}$ はすべての $x_i$ が零である部分である。
\end{enumerate}
各層はスキーム上のガーベであり、その群は $\mathcal{X}_i$ に対して
$\mu_{p^i}$、$\mathcal{X}_{\infty}$ に対して $\mathbf{G}_m$ である。
これらの層は被約局所閉部分スタックである。同じ性質をもつ
より粗い層化は存在しない。
\end{example}

\noindent
それでも、超限帰納法と
Proposition \ref{stacks-morphisms-proposition-open-stratum}
を用いて、ガーベによる無限かもしれない層化を見つけることができる……！

\begin{lemma}
\label{stacks-morphisms-lemma-every-point-in-a-stratum}
$\mathcal{X}$ を、$\mathcal{I}_\mathcal{X} \to \mathcal{X}$ が
準コンパクトであるような代数スタックとする。
このとき整列された添字集合 $I$ と、各 $i \in I$ に対して
被約局所閉部分スタック $\mathcal{U}_i \subset \mathcal{X}$ で、
次を満たすものが存在する：
\begin{enumerate}
\item 各 $\mathcal{U}_i$ はガーベである。
\item $|\mathcal{X}| = \bigcup_{i \in I} |\mathcal{U}_i|$ である。
\item $T_i = |\mathcal{X}| \setminus \bigcup_{i' < i} |\mathcal{U}_{i'}|$
はすべての $i \in I$ に対して $|\mathcal{X}|$ で閉である。
\item $|\mathcal{U}_i|$ は $T_i$ で開である。
\end{enumerate}
さらに、(a) $|\mathcal{U}_i| \subset T_i$ が稠密であるか、
(b) $\mathcal{U}_i$ が準コンパクトであるかのいずれかになるように
選ぶことができる。(a) の場合、$\mathcal{U}_i$ を可能な限り大きく選べば
（細部は証明を参照）、層化は標準的である。
\end{lemma}

\begin{proof}
$T \subset |\mathcal{X}|$ を空でない閉部分集合とする。
このような各 $T$ に対して、$|\mathcal{U}(T)| \subset T$ が
開稠密（それぞれ空でなく準コンパクト）となる被約局所閉部分スタック
$\mathcal{U}(T) \subset \mathcal{X}$ を見つける（それぞれ選ぶ）。
実際、
Properties of Stacks, Lemma
\ref{stacks-properties-lemma-reduced-closed-substack}
により、$T = |\mathcal{X}'|$ を満たす一意な被約閉部分スタック
$\mathcal{X}' \subset \mathcal{X}$ が存在する。次に注意する：
$\mathcal{I}_{\mathcal{X}'} =
\mathcal{I}_\mathcal{X} \times_\mathcal{X} \mathcal{X}'$
。これは
Lemma \ref{stacks-morphisms-lemma-monomorphism-cartesian-square-inertia}.
による。したがって $\mathcal{I}_{\mathcal{X}'} \to \mathcal{X}'$ は
基底変換として準コンパクトである。次を参照せよ：
Lemma \ref{stacks-morphisms-lemma-base-change-quasi-compact}.
Proposition \ref{stacks-morphisms-proposition-open-stratum}
により、ガーベである稠密な最大開部分スタック
$\mathcal{U} \subset \mathcal{X}'$ が存在する
（命題の証明を参照）。場合 (a) では
$\mathcal{U}(T) = \mathcal{U}$ と置く（これは標準的である）。
場合 (b) では空でない準コンパクトな開
$\mathcal{U}(T) \subset \mathcal{U}$ を選ぶ。次を参照せよ：
Properties of Stacks, Lemma
\ref{stacks-properties-lemma-space-locally-quasi-compact}
（選択公理により、すべての $T$ に対して同時に選べる）。

\medskip\noindent
超限再帰により、各順序数 $\alpha$ に対して閉部分集合
$T_\alpha \subset |\mathcal{X}|$ を構成する。
$\alpha = 0$ に対して $T_0 = |\mathcal{X}|$ と置く。
$T_\alpha$ が与えられたとき、
$$
T_{\alpha + 1} = T_\alpha \setminus |\mathcal{U}(T_\alpha)|.
$$
と置く。$\beta$ が極限順序数ならば、
$$
T_\beta = \bigcap\nolimits_{\alpha < \beta} T_\alpha.
$$
と置く。十分大きいすべての $\alpha$ に対して
$T_\alpha = \emptyset$ であると主張する。実際、すべての $\alpha$ に対して
$T_\alpha \not = \emptyset$ であると仮定すると、順序数の類から
$|\mathcal{X}|$ の部分集合全体の集合への単射が得られ、矛盾する。

\medskip\noindent
この主張から補題が従う。実際、
$$
I = \{\alpha \mid \mathcal{U}_\alpha \not = \emptyset \}.
$$
と置く。この主張により、これは整列集合である。
$i = \alpha \in I$ に対して $\mathcal{U}_i = \mathcal{U}_\alpha$ と置く。
すると $\mathcal{U}_i$ は被約局所閉部分スタックかつガーベであり、(1) が成り立つ。
構成により、$i = \alpha \in I$ ならば $T_i = T\alpha$ なので、(3) が成り立つ。
また、$\mathcal{U}(T)$ の選び方により (4) と (a) または (b) も成り立つ。
最後に (2) を見るため、$x \in |\mathcal{X}|$ とする。
$x \not \in T_\beta$ を満たす最小の順序数 $\beta$ が存在する
（順序数は整列されているためである）。
極限順序数に対する $T_\beta$ の定義により、この $\beta$ は後続順序数で
なければならない。したがって $\beta = \alpha + 1$ かつ
$x \in |\mathcal{U}(T_\alpha)|$ であり、結論を得る。
\end{proof}

\begin{remark}
\label{stacks-morphisms-remark-order-type}
Lemma \ref{stacks-morphisms-lemma-every-point-in-a-stratum} で構成した型 (a) の層化として
得られる標準的層化の順序型について考えることができる。
自然な予想は、整列集合 $I$ の {\it 濃度} が高々 $\aleph_0$ であるという
ものである。これが真か偽かはまったく分からない。
分かる方は次へ電子メールを送っていただきたい：
\href{mailto:stacks.project@gmail.com}{stacks.project@gmail.com}.
\end{remark}







\section{代数スタックの位相空間}
\label{stacks-morphisms-section-topology}

\noindent
本節では、これまでの結果を代数スタックに付随する位相空間
$|\mathcal{X}|$ に適用する。

\begin{lemma}
\label{stacks-morphisms-lemma-spectral-qc-diagonal-qc}
$\mathcal{X}$ を、対角射 $\Delta$ が準コンパクトであるような
準コンパクト代数スタックとする。
このとき $|\mathcal{X}|$ はスペクトル位相空間である。
\end{lemma}

\begin{proof}
アフィンスキーム $U$ と全射滑らかな射 $U \to \mathcal{X}$ を選ぶ。
次を参照せよ：
Properties of Stacks, Lemma \ref{stacks-properties-lemma-quasi-compact-stack}.
このとき $|U| \to |\mathcal{X}|$ は連続、開、かつ全射である。次を参照せよ：
Properties of Stacks, Lemma \ref{stacks-properties-lemma-topology-points}.
したがって $|\mathcal{X}|$ の準コンパクト開集合は位相の基底をなす。
準コンパクト開集合 $W_1, W_2 \subset |\mathcal{X}|$ に対し、
$W_1$ と $W_2$ へ写る $U$ の準コンパクト開集合 $V_1, V_2$ を選べる。
$\Delta$ は準コンパクトなので、
$$
V_1 \times_\mathcal{X} V_2 =
(V_1 \times V_2) \times_{\mathcal{X} \times \mathcal{X}, \Delta} \mathcal{X}
$$
は準コンパクトである。また、
Properties of Stacks, Lemma \ref{stacks-properties-lemma-points-cartesian}.
により、$|\mathcal{X}|$ における
$|V_1 \times_\mathcal{X} V_2|$ の像は $W_1 \cap W_2$ である。
したがって $W_1 \cap W_2$ は準コンパクトである。
証明を終えるには、$|\mathcal{X}|$ が sober であることを示せば十分である。
Topology, Definition \ref{topology-definition-spectral-space} を参照せよ。

\medskip\noindent
$T \subset |\mathcal{X}|$ を既約閉部分集合とする。
$T$ が一意な生成点をもつことを示さなければならない。
$\mathcal{Z} \subset \mathcal{X}$ を $T$ に対応する被約誘導閉部分スタックとする。
次を参照せよ：
Properties of Stacks, Definition
\ref{stacks-properties-definition-reduced-induced-stack}.
$\mathcal{Z} \to \mathcal{X}$ は閉埋入なので、
$\Delta_\mathcal{Z}$ は準コンパクトである。実際、まず
$\mathcal{Z} \to \mathcal{X} \times \mathcal{X}$ が
$\mathcal{Z} \to \mathcal{X}$ と $\Delta$ の合成として
準コンパクトであることを示す。次に
$\mathcal{Z} \to \mathcal{X} \times \mathcal{X}$ を
$\Delta_\mathcal{Z}$ と
$\mathcal{Z} \times \mathcal{Z} \to \mathcal{X} \times \mathcal{X}$
の合成として書き、Lemma \ref{stacks-morphisms-lemma-quasi-compact-permanence} と、
射 $\mathcal{Z} \times \mathcal{Z} \to \mathcal{X} \times \mathcal{X}$
が分離であるという事実を用いる。
こうして次の段落で論じる場合に帰着する。

\medskip\noindent
$\mathcal{X}$ は準コンパクト、$\Delta$ は準コンパクト、
$\mathcal{X}$ は被約、かつ $|\mathcal{X}|$ は既約であると仮定する。
$|\mathcal{X}|$ が一意な生成点をもつことを示さなければならない。
$\mathcal{I}_\mathcal{X} \to \mathcal{X}$ は $\Delta$ の基底変換なので、
$\mathcal{I}_\mathcal{X} \to \mathcal{X}$ は準コンパクトである
（Lemma \ref{stacks-morphisms-lemma-base-change-quasi-compact}）。
したがって
Proposition \ref{stacks-morphisms-proposition-open-stratum}.
により、ガーベである稠密な開部分スタック
$\mathcal{U} \subset \mathcal{X}$ が存在する。
言い換えれば、$|\mathcal{U}| \subset |\mathcal{X}|$ は開稠密である。
したがって $\mathcal{X}$ がガーベであると仮定してよい。
$\mathcal{X} \to X$ が $\mathcal{X}$ を代数空間 $X$ 上の
ガーベにするとする。このとき
Lemma \ref{stacks-morphisms-lemma-gerbe-bijection-points}.
により $|\mathcal{X}| \cong |X|$ である。特に $X$ は準コンパクトである。
Lemma \ref{stacks-morphisms-lemma-gerbe-diagonal-quasi-compact} により、
$X$ は準コンパクトな対角射をもつ。すなわち $X$ は準分離代数空間である。
Properties of Spaces, Lemma
\ref{spaces-properties-lemma-quasi-compact-quasi-separated-spectral}
により $|X|$ はスペクトルであり、望む結論が従う。
\end{proof}

\begin{lemma}
\label{stacks-morphisms-lemma-spectral-qcqs}
$\mathcal{X}$ を準コンパクトかつ準分離な代数スタックとする。
このとき $|\mathcal{X}|$ はスペクトル位相空間である。
\end{lemma}

\begin{proof}
これは Lemma \ref{stacks-morphisms-lemma-spectral-qc-diagonal-qc} の特別な場合である。
\end{proof}

\begin{lemma}
\label{stacks-morphisms-lemma-sober-qs}
$\mathcal{X}$ を対角射が準コンパクトな代数スタックとする
（例えば $\mathcal{X}$ が準分離ならそうである）。
このとき $U_i$ がスペクトルであるような開被覆
$|\mathcal{X}| = \bigcup U_i$ が存在する。
特に $|\mathcal{X}|$ は sober 位相空間である。
\end{lemma}

\begin{proof}
Lemma \ref{stacks-morphisms-lemma-spectral-qc-diagonal-qc} から直ちに従う。
\end{proof}






\section{剰余ガーベの存在}
\label{stacks-morphisms-section-existence-residual-gerbes}

\noindent
代数スタック上の点の剰余ガーベの定義は
Properties of Stacks, Definition
\ref{stacks-properties-definition-residual-gerbe}.
である。有限型点（Lemma \ref{stacks-morphisms-lemma-point-finite-type-monomorphism}）および
ガーベの任意の点（Lemma \ref{stacks-morphisms-lemma-gerbe-residual-gerbe-exists}）について
剰余ガーベが存在することは既に示した。
本節では、多くの代数スタック上で剰余ガーベが存在することを証明する。
まず、予告した
Proposition \ref{stacks-morphisms-proposition-open-stratum}.
の応用を与える。

\begin{lemma}
\label{stacks-morphisms-lemma-every-point-residual-gerbe}
$\mathcal{X}$ を、$\mathcal{I}_\mathcal{X} \to \mathcal{X}$ が
準コンパクトであるような代数スタックとする。
このときすべての $x \in |\mathcal{X}|$ に対して、
$x$ における $\mathcal{X}$ の剰余ガーベが存在する。
\end{lemma}

\begin{proof}
$T = \overline{\{x\}} \subset |\mathcal{X}|$ を $x$ の閉包とする。
Properties of Stacks, Lemma
\ref{stacks-properties-lemma-reduced-closed-substack}
により、$T = |\mathcal{X}'|$ を満たす被約閉部分スタック
$\mathcal{X}' \subset \mathcal{X}$ が存在する。次に注意する：
$\mathcal{I}_{\mathcal{X}'} =
\mathcal{I}_\mathcal{X} \times_\mathcal{X} \mathcal{X}'$
Lemma \ref{stacks-morphisms-lemma-monomorphism-cartesian-square-inertia}.
による。したがって $\mathcal{I}_{\mathcal{X}'} \to \mathcal{X}'$ は
基底変換として準コンパクトである。次を参照せよ：
Lemma \ref{stacks-morphisms-lemma-base-change-quasi-compact}.
Proposition \ref{stacks-morphisms-proposition-open-stratum}
により、ガーベである稠密な開部分スタック
$\mathcal{U} \subset \mathcal{X}'$ が存在する。
$\{x\} \subset T$ も稠密なので、$x \in |\mathcal{U}|$ であることに
注意する。したがって
Lemma \ref{stacks-morphisms-lemma-gerbe-residual-gerbe-exists}.
により、$x$ における $\mathcal{U}$ の剰余ガーベ
$\mathcal{Z}_x \subset \mathcal{U}$ が存在する。
定義から直ちに、$\mathcal{Z}_x \to \mathcal{X}$ は
$x$ における $\mathcal{X}$ の剰余ガーベである。
\end{proof}

\noindent
スタックが準 DM ならば、剰余ガーベも存在する。
特に Deligne--Mumford スタックでは剰余ガーベが常に存在する。

\begin{lemma}
\label{stacks-morphisms-lemma-every-point-residual-gerbe-quasi-DM}
$\mathcal{X}$ を準 DM 代数スタックとする。
このときすべての $x \in |\mathcal{X}|$ に対して、
$x$ における $\mathcal{X}$ の剰余ガーベが存在する。
\end{lemma}

\begin{proof}
スキーム $U$ と、全射、平坦、局所有限表示、かつ局所準有限な射
$U \to \mathcal{X}$ を選ぶ。次を参照せよ：
Theorem \ref{stacks-morphisms-theorem-quasi-DM}.
$R = U \times_\mathcal{X} U$ と置く。射影 $s, t : R \to U$ は
射 $U \to \mathcal{X}$ の基底変換として、全射、平坦、
局所有限表示、かつ局所準有限である。標準射
$[U/R] \to \mathcal{X}$ が存在し（
Algebraic Stacks, Lemma \ref{algebraic-lemma-map-space-into-stack})
$U \to \mathcal{X}$ は全射、平坦、かつ局所有限表示なので、
これは同値である。次を参照せよ：
Algebraic Stacks, Remark \ref{algebraic-remark-flat-fp-presentation}.
したがって $\mathcal{X} = [U/R]$ と仮定してよい。ここで
$(U, R, s, t, c)$ は代数空間の亜群であり、
$s, t : R \to U$ は全射、平坦、局所有限表示、かつ局所準有限である。
次のように置く：
$$
U' = \coprod\nolimits_{u \in U\text{ lying over }x} \Spec(\kappa(u)).
$$
標準射 $U' \to U$ は単射である。また、
$$
R' = U' \times_\mathcal{X} U' =
R \times_{(U \times U)} (U' \times U')
$$
と置く。$U' \to U$ は単射なので、両射影
$s', t' : R' \to U'$ は単射と局所準有限射の合成として分解する。
したがって、$U'$ は体のスペクトルの非交和であるから、
Spaces over Fields, Lemma
\ref{spaces-over-fields-lemma-mono-towards-locally-quasi-finite-over-field}
を用いると、射 $s', t' : R' \to U'$ は局所準有限であると結論する。
再び $U'$ が体のスペクトルの非交和であることから、射
$s', t'$ は平坦でもある。最後に、$s', t'$ が局所準有限であることから
$s', t'$ は局所有限型であり、したがって $s', t'$ は局所有限表示である
（$U'$ は体のスペクトルの非交和で、特に局所 Noether なので、
Morphisms of Spaces, Lemma
\ref{spaces-morphisms-lemma-noetherian-finite-type-finite-presentation}
が適用できる）。したがって
Criteria for Representability, Theorem
\ref{criteria-theorem-flat-groupoid-gives-algebraic-stack}.
により、$\mathcal{Z} = [U'/R']$ は代数スタックである。
$R'$ は $U' \to U$ による $R$ の制限なので、
Groupoids in Spaces, Lemma
\ref{spaces-groupoids-lemma-quotient-stack-restrict}
および
Properties of Stacks, Lemma \ref{stacks-properties-lemma-monomorphism}.
により、$\mathcal{Z} \to \mathcal{X}$ は単射である。
$\mathcal{Z} \to \mathcal{X}$ は単射なので、
$|\mathcal{Z}| \to |\mathcal{X}|$ は単射である。次を参照せよ：
Properties of Stacks, Lemma
\ref{stacks-properties-lemma-monomorphism-injective-points}.
Properties of Stacks, Lemma \ref{stacks-properties-lemma-points-cartesian}
により、
$$
|U'| = |\mathcal{Z} \times_\mathcal{X} U'|
\longrightarrow
|\mathcal{Z}| \times_{|\mathcal{X}|} |U'|
$$
は全射である。これは、$U'$ の選び方により
$|\mathcal{Z}| \to |\mathcal{X}|$ の像が $\{x\}$ であることを意味する。
したがって $|\mathcal{Z}|$ は一点集合である。
最後に、構成により $U'$ は局所 Noether かつ被約、すなわち
$\mathcal{Z}$ は被約かつ局所 Noether である。
よって $\mathcal{Z} \to \mathcal{X}$ の本質像は
$x$ における $\mathcal{X}$ の剰余ガーベである。次を参照せよ：
Properties of Stacks, Lemma
\ref{stacks-properties-lemma-residual-gerbe-unique}.
\end{proof}

\begin{lemma}
\label{stacks-morphisms-lemma-every-point-residual-gerbe-locally-Noetherian}
$\mathcal{X}$ を局所 Noether 代数スタックとする。
このときすべての $x \in |\mathcal{X}|$ に対して、
$x$ における $\mathcal{X}$ の剰余ガーベが存在する。
\end{lemma}

\begin{proof}
アフィンスキーム $U$ と滑らかな射 $U \to \mathcal{X}$ で、
$x$ が開連続写像 $|U| \to |\mathcal{X}|$ の像に入るものを選ぶ。
$\mathcal{X}$ を $|U| \to |\mathcal{X}|$ の像に対応する
開部分スタックで置き換えることができ、実際そうする。次を参照せよ：
Properties of Stacks, Lemma \ref{stacks-properties-lemma-open-substacks}.
したがって、代数空間の滑らかな亜群 $(U, R, s, t, c)$ に対して
$\mathcal{X} = [U/R]$ であり、$U$ は Noether アフィンスキームであると
仮定してよい。次を参照せよ：
Algebraic Stacks, Section \ref{algebraic-section-stack-to-presentation}.

\medskip\noindent
$E \subset |U|$ を $\{x\} \subset |\mathcal{X}|$ の逆像とする。
もちろん $E \not = \emptyset$ である。
$|U|$ は Noether 位相空間なので、
$\overline{\{u\}} \cap E = \{u\}$ を満たす元 $u \in E$ を選べる。
いつものように、$u = \Spec(\kappa(u))$ をその剰余体のスペクトルとみなす。
次のように書く：
$$
F = u \times_{U, t} R = u \times_\mathcal{X} U
\quad\text{and}\quad
R' = (u \times u) \times_{(U \times U), (t, s)} R = u \times_\mathcal{X} u
$$
さらに、$Z = \overline{\{u\}} \subset U$ に被約誘導スキーム構造を
入れたものと記す。$p : F \to U$ を第二射影から誘導される射と記す
（$F$ の最初のファイバー積表示では $s : R \to U$ を用いる）。
このとき $E$ は $p$ の集合論的像である。
射 $R' \to F$ は単射で、$Z$ の逆像 $p^{-1}(Z)$ を経由する。
この逆像 $p^{-1}(Z) \subset F$ は閉部分スキームであり、
上で $u \in E$ を慎重に選んだことから、制限
$p|_{p^{-1}(Z)} : p^{-1}(Z) \to Z$ の像は集合論的に
$\{u\} \subset Z$ に含まれる。ここで $u = \lim W$ であり、
極限は既約被約スキーム $Z$ の空でないアフィン開部分スキーム全体にわたる。
したがって射 $p|_{p^{-1}(Z)} : p^{-1}(Z) \to Z$ は
射 $u \to Z$ を経由すると結論する。これは明らかに
$R' = p^{-1}(Z)$ を意味する。特に、射 $t' : R' \to u$ は、
局所 Noether 代数空間の閉埋入 $p^{-1}(Z) \to F$ と
滑らかな射 $\text{pr}_1 : F \to u$ の合成として局所有限表示である。
次を用いる：
Morphisms of Spaces, Lemmas
\ref{spaces-morphisms-lemma-locally-finite-type-locally-noetherian},
\ref{spaces-morphisms-lemma-closed-immersion-finite-presentation}, および
\ref{spaces-morphisms-lemma-composition-finite-presentation}.
したがって $u \to U$ による $(U, R, s, t, c)$ の制限
$(u, R', s', t', c')$ は代数空間の亜群であり、
$s'$ と $t'$ は平坦かつ局所有限表示である。ゆえに
Criteria for Representability, Theorem
\ref{criteria-theorem-flat-groupoid-gives-algebraic-stack}.
により、$\mathcal{Z} = [u/R']$ は代数スタックである。
$R'$ は $u \to U$ による $R$ の制限なので、
Groupoids in Spaces, Lemma
\ref{spaces-groupoids-lemma-quotient-stack-restrict}
および
Properties of Stacks, Lemma \ref{stacks-properties-lemma-monomorphism}.
により、$\mathcal{Z} \to \mathcal{X}$ は単射である。
そして $\mathcal{Z}$ は（同型を除いて）剰余ガーベである。これは
Properties of Stacks, Section
\ref{stacks-properties-section-residual-gerbe}.
の内容による。
\end{proof}













\section{エタール局所構造}
\label{stacks-morphisms-section-etale-local}

\noindent
本節では、代数スタックのエタール局所構造について何が言えるかを
論じ始める。

\begin{lemma}
\label{stacks-morphisms-lemma-quotient-etale}
$Y$ を代数空間とする。
$(U, R, s, t, c)$ を $Y$ 上の代数空間の亜群とする。
$U \to Y$ は平坦かつ局所有限表示であり、
$R \to U \times_Y U$ は開埋入であると仮定する。
このとき $X = [U/R] = U/R$ は代数空間であり、
$X \to Y$ はエタールである。
\end{lemma}

\begin{proof}
Criteria for Representability, Theorem
\ref{criteria-theorem-flat-groupoid-gives-algebraic-stack}.
により、商スタック $[U/R]$ は代数スタックである。
一方、$R \to U \times U$ は単射なので、これは代数空間であり
（我々の言葉の濫用と
Algebraic Stacks, Proposition
\ref{algebraic-proposition-algebraic-stack-no-automorphisms})
による）、もちろん代数空間 $U/R$
（Bootstrap, Theorem \ref{bootstrap-theorem-final-bootstrap} のもの）
に等しい。$U \to X$ は全射、平坦、かつ局所有限表示なので
（Bootstrap, Lemma \ref{bootstrap-lemma-covering-quotient}）、
Morphisms of Spaces, Lemma \ref{spaces-morphisms-lemma-flat-permanence}
および
Descent on Spaces, Lemma
\ref{spaces-descent-lemma-flat-finitely-presented-permanence}.
により、$X \to Y$ は平坦かつ局所有限表示である。
最後に、Cartesian 図式
$$
\xymatrix{
R \ar[d] \ar[r] & U \times_Y U \ar[d] \\
X \ar[r] & X \times_Y X
}
$$
を考える。右の垂直射は全射、平坦、かつ局所有限表示である
（細部は省略する）。上の水平射は仮定により開埋入なので、
$X \to X \times_Y X$ も開埋入である（
Properties of Stacks, Lemma
\ref{stacks-properties-lemma-check-property-covering} を用いよ）。
したがって
Morphisms of Spaces, Lemma
\ref{spaces-morphisms-lemma-diagonal-unramified-morphism}.
により、$X \to Y$ は非分岐である。結論は
Morphisms of Spaces, Lemma
\ref{spaces-morphisms-lemma-unramified-flat-lfp-etale}.
から従う。
\end{proof}

\begin{lemma}
\label{stacks-morphisms-lemma-quasi-splitting-etale}
$S$ をスキームとする。
$(U, R, s, t, c)$ を $S$ 上の代数空間の亜群とする。
$s, t$ は平坦かつ局所有限表示であると仮定する。
$P \subset R$ を、
$(U, P, s|_P, t|_P, c|_{P \times_{s, U, t} P})$ が
$S$ 上の代数空間の亜群となるような開部分空間とする。このとき
$$
[U/P] \longrightarrow [U/R]
$$
は代数空間で表現可能、全射、かつエタールな代数スタックの射である。
\end{lemma}

\begin{proof}
$P \subset R$ は開なので、$s|_P$ と $t|_P$ は平坦かつ局所有限表示である。
したがって
Criteria for Representability, Theorem
\ref{criteria-theorem-flat-groupoid-gives-algebraic-stack}.
により、$[U/R]$ と $[U/P]$ は代数スタックである。
この射が代数空間で表現可能であることを見るには、
$[U/P] \to [U/R]$ がファイバー圏上忠実であることを示せば十分である。
次を参照せよ：
Algebraic Stacks, Lemma
\ref{algebraic-lemma-characterize-representable-by-algebraic-spaces}.
これは $P \to R$ が単射であることと、
Groupoids in Spaces, Lemma
\ref{spaces-groupoids-lemma-quotient-stack-objects}.
で与えられる商スタックの明示的記述から直ちに従う。
以上により、
Algebraic Stacks, Definition
\ref{algebraic-definition-relative-representable-property}.
から $[U/P] \to [U/R]$ が全射かつエタールであることの意味が分かる。
全射性は、例えば
Criteria for Representability,
Lemma \ref{criteria-lemma-representable-by-spaces-surjective}
と、体のスペクトル上の商スタックの対象の記述
（上で引用した補題を参照）から従う。
残るのは、この射がエタールであることの証明である。

\medskip\noindent
そのためには、$U \times_{[U/R]} [U/P] \to U$ がエタールであることを
示せば十分である。次を参照せよ：Properties of Stacks, Lemma
\ref{stacks-properties-lemma-check-property-covering}.
Groupoids in Spaces, Lemma
\ref{spaces-groupoids-lemma-cartesian-square-of-morphism}
により、ファイバー積は $[R/P \times_{s, U, t} R]$ に等しく、
$U$ への射は $s : R \to U$ から誘導される。
写像 $s', t' : P \times_{s, U, t} R \to R$ は
$s' : (p, r) \mapsto r$ および
$t' : (p, r) \mapsto c(p, r)$ で与えられる。
$P \subset R$ は開なので、
$(t', s') : P \times_{s, U, t} R \to R \times_{s, U, s} R$
は開埋入である。したがって Lemma \ref{stacks-morphisms-lemma-quotient-etale} を
適用して結論を得る。
\end{proof}

\begin{lemma}
\label{stacks-morphisms-lemma-etale-local-quasi-DM}
$\mathcal{X}$ を代数スタックとする。$\mathcal{X}$ は分離対角射をもつ
準 DM である（同値なこととして、
$\mathcal{I}_\mathcal{X} \to \mathcal{X}$ は局所準有限かつ
分離である）と仮定する。$x \in |\mathcal{X}|$ とする。このとき、
代数スタックの射
$$
\mathcal{U} \longrightarrow \mathcal{X}
$$
で次の性質をもつものが存在する。
\begin{enumerate}
\item $x$ に写る点 $u \in |\mathcal{U}|$ が存在する。
\item $\mathcal{U} \to \mathcal{X}$ は代数空間で表現可能かつ
エタールである。
\item $\mathcal{U} = [U/R]$ であり、ここで $(U, R, s, t, c)$ は
$U$, $R$ がアフィンで、$s, t$ が有限、平坦、かつ
局所有限表示であるようなスキームの亜群である。
\end{enumerate}
\end{lemma}

\begin{proof}
（括弧内の主張は Lemma \ref{stacks-morphisms-lemma-diagonal-diagonal} の同値性から従う。）
アフィンスキーム $U$ と、ある点 $u \in U$ の像が $x$ となるような
平坦、局所有限表示、局所準有限な射 $U \to \mathcal{X}$ を選ぶ。
これは Theorem \ref{stacks-morphisms-theorem-quasi-DM} と $\mathcal{X}$ が準 DM であるという
仮定により可能である。$R = U \times_\mathcal{X} U$ とおいて得られる
代数空間の亜群を $(U, R, s, t, c)$ とする。次を参照せよ：
Algebraic Stacks, Lemma \ref{algebraic-lemma-map-space-into-stack}。
$|U| \to |\mathcal{X}|$ の開な像に対応する開部分スタックを
$\mathcal{X}' \subset \mathcal{X}$ とする
(Properties of Stacks, Lemmas
\ref{stacks-properties-lemma-topology-points} および
\ref{stacks-properties-lemma-open-substacks}).
明らかに、$\mathcal{X}$ を開部分スタック $\mathcal{X}'$ で置き換えてよい。
したがって $U \to \mathcal{X}$ は全射であると仮定してよく、このとき
Algebraic Stacks, Remark \ref{algebraic-remark-flat-fp-presentation}
から $\mathcal{X} = [U/R]$ を得る。
$s, t : R \to U$ は平坦、局所有限表示、かつ局所準有限であることに
注意する。
$R = U \times U \times_{(\mathcal{X} \times \mathcal{X})} \mathcal{X}$ であり、
$\mathcal{X}$ の対角射は分離だから、$R$ は分離である。
したがって $s, t : R \to U$ は分離である。ゆえに
Morphisms of Spaces, Proposition
\ref{spaces-morphisms-proposition-locally-quasi-finite-separated-over-scheme}
を $s : R \to U$ に適用すると、$R$ はスキームである。

\medskip\noindent
以上により、$(U, R, s, t, c)$ と $u$ について
More on Groupoids in Spaces, Lemma
\ref{spaces-more-groupoids-lemma-quasi-splitting-affine-scheme}
の仮定がすべて満たされることを確かめた。
したがって、$R$ の $U'$ への制限 $R'$ が $u$ 上で準分裂するような
初等エタール近傍 $(U', u') \to (U, u)$ を見いだせる。
$R' = U' \times_\mathcal{X} U'$ であることに注意する
（細部は省略する。ヒント：ファイバー積の推移性）。
$(U, R, s, t, c)$ を $(U', R', s', t', c')$ で置き換え、上と同様に
$\mathcal{X}$ を縮小すると、$(U, R, s, t, c)$ は $u$ 上で
準分裂をもつと仮定してよい（今後、点 $u$ は無関係であることは
More on Groupoids in Spaces, Definition
\ref{spaces-more-groupoids-definition-split-at-point} の脚注から分かる）。
$P \subset R$ を $u$ 上の $R$ の準分裂とする。
Lemma \ref{stacks-morphisms-lemma-quasi-splitting-etale} を適用すると、
$$
\mathcal{U} = [U/P] \longrightarrow [U/R] = \mathcal{X}
$$
が所望の性質をすべてもつことが分かる。
\end{proof}

\begin{lemma}
\label{stacks-morphisms-lemma-etale-local-quasi-DM-at-x}
$\mathcal{X}$ を代数スタックとする。$\mathcal{X}$ は分離対角射をもつ
準 DM である（同値なこととして、
$\mathcal{I}_\mathcal{X} \to \mathcal{X}$ は局所準有限かつ
分離である）と仮定する。$x \in |\mathcal{X}|$ とする。
$x$ における $\mathcal{X}$ の自己同型群が有限であると仮定する
（Remark \ref{stacks-morphisms-remark-property-automorphism-groups}）。
このとき、代数スタックの射
$$
g : \mathcal{U} \longrightarrow \mathcal{X}
$$
で次の性質をもつものが存在する。
\begin{enumerate}
\item $x$ に写る点 $u \in |\mathcal{U}|$ が存在し、$g$ は $u$ と $x$ に
おける自己同型群の間の同型を誘導する
（Remark \ref{stacks-morphisms-remark-identify-automorphism-groups}）。
\item $\mathcal{U} \to \mathcal{X}$ は代数空間で表現可能かつ
エタールである。
\item $\mathcal{U} = [U/R]$ であり、ここで $(U, R, s, t, c)$ は
$U$, $R$ がアフィンで、$s, t$ が有限、平坦、かつ
局所有限表示であるようなスキームの亜群である。
\end{enumerate}
\end{lemma}

\begin{proof}
$G_x$ は Lemma \ref{stacks-morphisms-lemma-automorphism-group-scheme} により
群スキームであることに注意する。
証明の前半は Lemma \ref{stacks-morphisms-lemma-etale-local-quasi-DM} の証明の前半と
{\bf まったく}同じである。
したがって、$x$ に写る $u \in U$ と $(U, R, s, t, c)$ が
More on Groupoids in Spaces, Lemma
\ref{spaces-more-groupoids-lemma-quasi-splitting-affine-scheme} の仮定を
すべて満たし、$\mathcal{X} = [U/R]$ であると仮定してよい。
$G_x$ に関する仮定から、$G_u$ は $u$ 上有限である。
したがって、
More on Groupoids in Spaces, Lemma
\ref{spaces-more-groupoids-lemma-splitting-affine-scheme}
の仮定がすべて満たされる。
ゆえに、$R$ の $U'$ への制限 $R'$ が $u$ 上で分裂するような
初等エタール近傍 $(U', u') \to (U, u)$ を見いだせる。
$R' = U' \times_\mathcal{X} U'$ であることに注意する
（細部は省略する。ヒント：ファイバー積の推移性）。
$(U, R, s, t, c)$ を $(U', R', s', t', c')$ で置き換え、上と同様に
$\mathcal{X}$ を縮小すると、$(U, R, s, t, c)$ は $u$ 上で分裂を
もつと仮定してよい。$P \subset R$ を $u$ 上の $R$ の分裂とする。
Lemma \ref{stacks-morphisms-lemma-quasi-splitting-etale} を適用すると、
$$
\mathcal{U} = [U/P] \longrightarrow [U/R] = \mathcal{X}
$$
は代数空間で表現可能かつエタールである。構成により $G_u$ は
$P$ に含まれるので、この射は所望のとおり $u$ における自己同型群上の
同型を定める。
\end{proof}

\begin{lemma}
\label{stacks-morphisms-lemma-etale-local-quasi-DM-at-x-inertia}
$\mathcal{X}$ を代数スタックとする。$\mathcal{X}$ は分離対角射をもつ
準 DM である（同値なこととして、
$\mathcal{I}_\mathcal{X} \to \mathcal{X}$ は局所準有限かつ
分離である）と仮定する。$x \in |\mathcal{X}|$ とする。$x$ は
準コンパクトな射 $\Spec(k) \to \mathcal{X}$ で表されると仮定する。
このとき、代数スタックの射
$$
g : \mathcal{U} \longrightarrow \mathcal{X}
$$
で次の性質をもつものが存在する。
\begin{enumerate}
\item $x$ に写る点 $u \in |\mathcal{U}|$ が存在し、$g$ は $u$ と $x$ に
おける剰余ガーベの間の同型を誘導する。
\item $\mathcal{U} \to \mathcal{X}$ は代数空間で表現可能かつ
エタールである。
\item $\mathcal{U} = [U/R]$ であり、ここで $(U, R, s, t, c)$ は
$U$, $R$ がアフィンで、$s, t$ が有限、平坦、かつ
局所有限表示であるようなスキームの亜群である。
\end{enumerate}
\end{lemma}

\begin{proof}
証明の前半は Lemma \ref{stacks-morphisms-lemma-etale-local-quasi-DM} の証明の前半と
{\bf まったく}同じである。
したがって、$x$ に写る $u \in U$ と $(U, R, s, t, c)$ が
More on Groupoids in Spaces, Lemma
\ref{spaces-more-groupoids-lemma-quasi-splitting-affine-scheme} の仮定を
すべて満たし、$\mathcal{X} = [U/R]$ であると仮定してよい。
$u = \Spec(\kappa(u)) \to \mathcal{X}$ は準コンパクトであることに注意する。
次を参照せよ：
Properties of Stacks, Lemma
\ref{stacks-properties-lemma-UR-quasi-compact-above-x}.
Cartesian 図式
$$
\xymatrix{
F \ar[d] \ar[r] & U \ar[d] \\
u \ar[r]^u & \mathcal{X}
}
$$
を考える。$U$ はアフィンスキームであり $F \to U$ は準コンパクトなので、
$F$ は準コンパクトである。$U \to \mathcal{X}$ は局所準有限なので、
$F \to u$ は局所準有限である。したがって $F \to u$ は準有限であり、
$F$ はその台位相空間が有限離散であるアフィンスキームである
（Spaces over Fields, Lemma
\ref{spaces-over-fields-lemma-locally-quasi-finite-over-field}).
単射 $u \times_\mathcal{X} u \to F$ があることに注意する。
特に、$|u \times_\mathcal{X} u| \to |R|$ の像である集合
$\{r \in R : s(r) = u, t(r) = u\}$ は有限である。
ゆえに、
More on Groupoids in Spaces, Lemma
\ref{spaces-more-groupoids-lemma-strong-splitting-affine-scheme}
の仮定がすべて成り立つ。

\medskip\noindent
したがって、$R$ の $U'$ への制限 $R'$ が $u'$ 上で強分裂するような
初等エタール近傍 $(U', u') \to (U, u)$ を見いだせる。
$R' = U' \times_\mathcal{X} U'$ であることに注意する
（細部は省略する。ヒント：ファイバー積の推移性）。
$(U, R, s, t, c)$ を $(U', R', s', t', c')$ で置き換え、上と同様に
$\mathcal{X}$ を縮小すると、$(U, R, s, t, c)$ は $u$ 上で強分裂を
もつと仮定してよい。$P \subset R$ を $u$ 上の $R$ の強分裂とする。
Lemma \ref{stacks-morphisms-lemma-quasi-splitting-etale} を適用すると、
$$
\mathcal{U} = [U/P] \longrightarrow [U/R] = \mathcal{X}
$$
は代数空間で表現可能かつエタールである。$P \subset R$ は開であり、
構成により $\{r \in R : s(r) = u, t(r) = u\}$ を含むので、
$u \times_\mathcal{U} u \to u \times_\mathcal{X} u$ は同型である。
したがって剰余ガーベに関する主張は
Properties of Stacks, Lemma
\ref{stacks-properties-lemma-residual-gerbe-isomorphic}
から従う（対象となる剰余ガーベが存在することは
Lemma \ref{stacks-morphisms-lemma-every-point-residual-gerbe-quasi-DM} による）。
\end{proof}







\section{滑らかな射}
\label{stacks-morphisms-section-smooth}

\noindent
代数空間の射が「滑らかである」という性質は、始域と終域上滑らか局所的である。
次を参照せよ：
Descent on Spaces, Remark \ref{spaces-descent-remark-list-local-source-target}.
これは基底変換の下でも安定であり、終域上 fpqc 局所的でもある。次を参照せよ：
Morphisms of Spaces,
Lemma \ref{spaces-morphisms-lemma-base-change-smooth}
および
Descent on Spaces, Lemma
\ref{spaces-descent-lemma-descending-property-smooth}.
したがって、上の Lemma \ref{stacks-morphisms-lemma-local-source-target} により、
代数スタックの射が滑らかであることを次のように定義できる。この定義は、射が
代数空間で表現可能な場合には、
Properties of Stacks,
Section \ref{stacks-properties-section-properties-morphisms}
で定義済みの概念と一致する。

\begin{definition}
\label{stacks-morphisms-definition-smooth}
$f : \mathcal{X} \to \mathcal{Y}$ を代数スタックの射とする。
Lemma \ref{stacks-morphisms-lemma-local-source-target}
の同値な条件が $\mathcal{P} = \text{smooth}$ について成り立つとき、
$f$ は {\it 滑らか}であるという。
\end{definition}

\begin{lemma}
\label{stacks-morphisms-lemma-composition-smooth}
滑らかな射の合成は滑らかである。
\end{lemma}

\begin{proof}
Remark \ref{stacks-morphisms-remark-composition} と
Morphisms of Spaces, Lemma
\ref{spaces-morphisms-lemma-composition-smooth} を組み合わせる。
\end{proof}

\begin{lemma}
\label{stacks-morphisms-lemma-base-change-smooth}
滑らかな射の基底変換は滑らかである。
\end{lemma}

\begin{proof}
Remark \ref{stacks-morphisms-remark-base-change} と
Morphisms of Spaces, Lemma
\ref{spaces-morphisms-lemma-base-change-smooth} を組み合わせる。
\end{proof}

\begin{lemma}
\label{stacks-morphisms-lemma-descent-smooth}
$f : \mathcal{X} \to \mathcal{Y}$ を代数スタックの射とする。
$\mathcal{Z} \to \mathcal{Y}$ を全射、平坦、かつ局所有限表示な
代数スタックの射とする。基底変換
$\mathcal{Z} \times_\mathcal{Y} \mathcal{X} \to \mathcal{Z}$
が滑らかならば、$f$ は滑らかである。
\end{lemma}

\begin{proof}
「滑らか」という性質は Lemma \ref{stacks-morphisms-lemma-descent-property} の条件を
満たす。始域と終域上滑らか局所的であることは本節の導入で見た。
終域上 fppf 局所的であることは
Descent on Spaces, Lemma
\ref{spaces-descent-lemma-descending-property-smooth} である。
\end{proof}

\begin{lemma}
\label{stacks-morphisms-lemma-smooth-locally-finite-presentation}
代数スタックの滑らかな射は局所有限表示である。
\end{lemma}

\begin{proof}
省略する。
\end{proof}

\begin{lemma}
\label{stacks-morphisms-lemma-where-smooth}
$f : \mathcal{X} \to \mathcal{Y}$ を代数スタックの射とする。
$f|_\mathcal{U} : \mathcal{U} \to \mathcal{Y}$ が滑らかとなるような
最大の開部分スタック $\mathcal{U} \subset \mathcal{X}$ が存在する。
さらに、この開部分スタックの形成は次と可換である。
\begin{enumerate}
\item 滑らかな射との前合成。
\item 平坦かつ局所有限表示な射による基底変換。
\item $f$ が局所有限表示である場合の、平坦な射による基底変換。
\end{enumerate}
\end{lemma}

\begin{proof}
可換図式
$$
\xymatrix{
U \ar[d]_a \ar[r]_h & V \ar[d]^b \\
\mathcal{X} \ar[r]^f & \mathcal{Y}
}
$$
を選ぶ。ここで $U$ と $V$ は代数空間であり、垂直射は滑らかで、
$a : U \to \mathcal{X}$ は全射である。
$h_{U'} : U' \to V$ が滑らかとなるような最大の開部分空間
$U' \subset U$ が存在する。次を参照せよ：
Morphisms of Spaces, Lemma \ref{spaces-morphisms-lemma-where-smooth}.
$|U'| \to |\mathcal{X}|$ の像に対応する開部分スタックを
$\mathcal{U} \subset \mathcal{X}$ とする
(Properties of Stacks, Lemmas \ref{stacks-properties-lemma-topology-points} および
\ref{stacks-properties-lemma-open-substacks}).
Lemma \ref{stacks-morphisms-lemma-local-source-target} の同値性により、
$f|_\mathcal{U} : \mathcal{U} \to \mathcal{Y}$ は滑らかであり、
$\mathcal{U}$ はこの性質をもつ最大の開部分スタックである。

\medskip\noindent
主張 (1) は、滑らかであることが始域上滑らか局所的であることから従う
（この性質は代数スタックの滑らかな射を定義するためにも用いた）。
主張 (2) と (3) は代数空間の場合から従う。次を参照せよ：
Morphisms of Spaces, Lemma \ref{spaces-morphisms-lemma-where-smooth}.
\end{proof}

\begin{lemma}
\label{stacks-morphisms-lemma-smooth-quotient-stack}
$X \to Y$ を代数空間の滑らかな射とする。
$G$ を $Y$ 上平坦かつ局所有限表示な、$Y$ 上の群代数空間とする。
$G$ は $Y$ 上 $X$ に作用するとする。このとき商スタック $[X/G]$ は
$Y$ 上滑らかである。
\end{lemma}

\noindent
$G$ が $Y$ 上滑らかでなくても、これは成り立つ！

\begin{proof}
$[X/G]$ は
Criteria for Representability, Theorem
\ref{criteria-theorem-flat-groupoid-gives-algebraic-stack} により
代数スタックである。$[X/G]$ が $Y$ 上滑らかであることは、滑らかさが
fppf 被覆の下で降下することから従う。
$U$ がスキームとなる全射滑らかな射 $U \to [X/G]$ を選ぶ。
$[X/G]$ が $Y$ 上滑らかであることは、$U$ が $Y$ 上滑らかであることと
同値である。$U \times_{[X/G]} X$ は $X$ 上滑らかであり、したがって
$Y$ 上滑らかであることに注意する（滑らかな射の合成は滑らかだからである）。
一方、$U \times_{[X/G]} X \to U$ は局所有限表示、平坦、かつ全射である
（これは $X \to [X/G]$ の基底変換であり、この射がそれらの性質をもつことは、
例えば Criteria for Representability, Lemma
\ref{criteria-lemma-flat-quotient-flat-presentation}).
による）。したがって Descent on Spaces,
Lemma \ref{spaces-descent-lemma-smooth-permanence}.
を適用できる。
\end{proof}

\begin{lemma}
\label{stacks-morphisms-lemma-gerbe-smooth}
$\pi : \mathcal{X} \to \mathcal{Y}$ を代数スタックの射とする。
$\mathcal{X}$ が $\mathcal{Y}$ 上のガーベならば、$\pi$ は全射かつ
滑らかである。
\end{lemma}

\begin{proof}
全射性は Lemma \ref{stacks-morphisms-lemma-gerbe-fppf} で見た。
Lemma \ref{stacks-morphisms-lemma-descent-smooth} により、$\pi$ を全射、平坦、かつ
局所有限表示な射 $\mathcal{Y}' \to \mathcal{Y}$ による基底変換で
置き換えた後に補題を証明すれば十分である。
Lemma \ref{stacks-morphisms-lemma-local-structure-gerbe} により、$\mathcal{Y} = U$ は
代数空間であり、$\mathcal{X} = [U/G]$ は $U$ 上の商で、
$G \to U$ は平坦かつ局所有限表示であると仮定してよい。
このとき Lemma \ref{stacks-morphisms-lemma-smooth-quotient-stack} から結論を得る。
\end{proof}






\section{始域と終域上エタール・滑らか局所的な射の型}
\label{stacks-morphisms-section-etale-smooth-local-source-target}

\noindent
代数空間の射の性質で、{\it 始域と終域上エタール・滑らか局所的}なものが
与えられたとする。次を参照せよ：
Descent on Spaces,
Definition \ref{spaces-descent-definition-etale-smooth-local-source-target}
この性質を用いて代数スタックの DM 射に対応する性質を定義できる。
すなわち、次の補題の同値な条件のいずれかを課せばよい。

\begin{lemma}
\label{stacks-morphisms-lemma-etale-smooth-local-source-target}
$\mathcal{P}$ を、始域と終域上エタール・滑らか局所的な
代数空間の射の性質とする。
$f : \mathcal{X} \to \mathcal{Y}$ を代数スタックの DM 射とする。
可換図式
$$
\xymatrix{
U \ar[d]_a \ar[r]_h & V \ar[d]^b \\
\mathcal{X} \ar[r]^f & \mathcal{Y}
}
$$
で、$U$ と $V$ は代数空間、$V \to \mathcal{Y}$ は滑らか、かつ
$U \to \mathcal{X} \times_\mathcal{Y} V$ はエタールであるものを考える。
次は同値である。
\begin{enumerate}
\item 上のような任意の図式について、射 $h$ は性質 $\mathcal{P}$ をもつ。
\item $a : U \to \mathcal{X}$ が全射となる上のようなある図式について、
射 $h$ は性質 $\mathcal{P}$ をもつ。
\end{enumerate}
$\mathcal{X}$ と $\mathcal{Y}$ が代数空間で表現可能ならば、これは
$f$ が（代数空間の射として）性質 $\mathcal{P}$ をもつこととも同値である。
さらに $\mathcal{P}$ が任意の基底変換の下で保たれ、基底上 fppf 局所的ならば、
代数空間で表現可能な射 $f$ について、これは $f$ が
Properties of Stacks,
Section \ref{stacks-properties-section-properties-morphisms}.
の意味で性質 $\mathcal{P}$ をもつこととも同値である。
\end{lemma}

\begin{proof}
含意 (1) $\Rightarrow$ (2) を証明しよう。代数空間 $V$ と全射滑らかな射
$V \to \mathcal{Y}$ を選ぶ。$f$ は DM なので、スキーム $U$ と
全射エタールな射 $U \to V \times_\mathcal{Y} \mathcal{X}$ が存在する。
Lemma \ref{stacks-morphisms-lemma-DM} を参照せよ。したがって (2) が成り立つことが分かる。
$U \to \mathcal{X}$ も全射かつ滑らかであることに注意する。実際、これは
基底変換 $\mathcal{X} \times_\mathcal{Y} V \to \mathcal{X}$ と、選んだ写像
$U \to \mathcal{X} \times_\mathcal{Y} V$ の合成である。ゆえに (1) のような
図式を得る。したがって (1) が成り立つならば、$h : U \to V$ は性質
$\mathcal{P}$ をもつ。$U \to \mathcal{X}$ は全射なので、これは (2) が
成り立つことを意味する。

\medskip\noindent
逆に (2) が成り立つと仮定し、$U, V, a, b, h$ を (2) のとおりとする。
また、$U', V', a', b', h'$ を (1) のような任意の図式とする。
図示すると
$$
\xymatrix{
U \ar[d] \ar[r]_h & V \ar[d] \\
\mathcal{X} \ar[r]^f & \mathcal{Y}
}
\quad\quad
\xymatrix{
U' \ar[d] \ar[r]_{h'} & V' \ar[d] \\
\mathcal{X} \ar[r]^f & \mathcal{Y}
}
$$
となる。(2) が (1) を含意することを示すには、$h'$ が $\mathcal{P}$ を
もつことを証明しなければならない。そのため可換図式
$$
\xymatrix{
U \ar[d]^h &
U \times_\mathcal{X} U' \ar[l] \ar[d]^{(h, h')} \ar[r] &
U' \ar[d]^{h'} \\
V &
V \times_\mathcal{Y} V' \ar[l] \ar[r] &
V'
}
$$
を代数空間について考える。水平射は滑らかな射
$V \to \mathcal{Y}$、$V' \to \mathcal{Y}$、$U \to \mathcal{X}$、および
$U' \to \mathcal{X}$ の基底変換なので滑らかである。正方形
$$
\xymatrix{
U \ar[d] & U \times_\mathcal{X} U' \ar[l] \ar[d] &
U \times_\mathcal{X} U' \ar[d] \ar[r] & U' \ar[d] \\
V \times_\mathcal{Y} \mathcal{X} &
V \times_\mathcal{Y} U' \ar[l] &
U \times_\mathcal{Y} V' \ar[r] &
\mathcal{X} \times_\mathcal{Y} V'
}
$$
は Cartesian であることに注意する。したがって $U', V', a', b', h'$ と
$U, V, a, b, h$ に関する仮定により、垂直射はエタールである。
$\mathcal{P}$ は
Descent on Spaces, Lemma
\ref{spaces-descent-lemma-etale-smooth-local-source-target-implies} の (2)
により終域上滑らか局所的なので、$h$ の基底変換
$t : U \times_\mathcal{Y} V' \to V \times_\mathcal{Y} V'$ は
$\mathcal{P}$ をもつ。また $\mathcal{P}$ は
Descent on Spaces, Lemma
\ref{spaces-descent-lemma-etale-smooth-local-source-target-implies} の (1)
により始域上エタール局所的であり、
$s : U \times_\mathcal{X} U' \to U \times_\mathcal{Y} V'$ はエタールだから、
射 $(h, h') = t \circ s$ は $\mathcal{P}$ をもつ。図式
$$
\xymatrix{
U \times_\mathcal{X} U' \ar[r]_{(h, h')} \ar[d] &
V \times_\mathcal{Y} V' \ar[d] \\
U' \ar[r]^{h'} & V'
}
$$
を考える。左の垂直射は全射、右の垂直射は滑らかであり、誘導射
$$
U \times_\mathcal{X} U'
\longrightarrow
U' \times_{V'} (V \times_\mathcal{Y} V') = V \times_\mathcal{Y} U'
$$
は上で見たようにエタールである。したがって
Descent on Spaces, Definition
\ref{spaces-descent-definition-etale-smooth-local-source-target} の (3)
により、$h'$ は $\mathcal{P}$ をもつと結論できる。これで (1) と (2) の
同値性の証明が完了する。

\medskip\noindent
$\mathcal{X}$ と $\mathcal{Y}$ が表現可能ならば、
Descent on Spaces,
Lemma \ref{spaces-descent-lemma-etale-smooth-local-source-target-characterize}
を適用でき、(1) と (2) は $f$ が $\mathcal{P}$ をもつことと同値であると
分かる。

\medskip\noindent
最後に、$f$ は表現可能であり、$U, V, a, b, h$ は補題の (2) のとおりで、
$\mathcal{P}$ は任意の基底変換の下で保たれると仮定する。任意のスキーム
$Z$ と射 $Z \to \mathcal{X}$ について、基底変換
$Z \times_\mathcal{Y} \mathcal{X} \to Z$
が性質 $\mathcal{P}$ をもつことを示さなければならない。図式
$$
\xymatrix{
Z \times_\mathcal{Y} U \ar[d] \ar[r] &
Z \times_\mathcal{Y} V \ar[d] \\
Z \times_\mathcal{Y} \mathcal{X} \ar[r] &
Z
}
$$
を考える。上の水平射は $h$ の基底変換なので、性質 $\mathcal{P}$ をもつ。
左の垂直射は全射であり、誘導射
$$
Z \times_\mathcal{Y} U \longrightarrow
(Z \times_\mathcal{Y} \mathcal{X}) \times_Z (Z \times_\mathcal{Y} V)
$$
はエタールで、右の垂直射は滑らかである。したがって
Descent on Spaces,
Lemma \ref{spaces-descent-lemma-etale-smooth-local-source-target-characterize}
を適用すると、$Z \times_\mathcal{Y} \mathcal{X} \to Z$ は
性質 $\mathcal{P}$ をもつと分かる。
\end{proof}

\begin{definition}
\label{stacks-morphisms-definition-etale-smooth-P}
$\mathcal{P}$ を、始域と終域上エタール・滑らか局所的な
代数空間の射の性質とする。代数スタックの DM 射
$f : \mathcal{X} \to \mathcal{Y}$ が
Lemma \ref{stacks-morphisms-lemma-local-source-target}
の同値な条件を満たすとき、これは {\it 性質 $\mathcal{P}$ をもつ}という。
\end{definition}

\begin{remark}
\label{stacks-morphisms-remark-etale-smooth-composition}
$\mathcal{P}$ を、始域と終域上エタール・滑らか局所的で、合成の下で
安定な代数空間の射の性質とする。このとき
Definition \ref{stacks-morphisms-definition-etale-smooth-P} で定義した代数スタックの
DM 射の性質は、合成の下で安定である。実際、
$f : \mathcal{X} \to \mathcal{Y}$ と $g : \mathcal{Y} \to \mathcal{Z}$ を
性質 $\mathcal{P}$ をもつ代数スタックの DM 射とする。
Lemma \ref{stacks-morphisms-lemma-composition-separated} により、合成 $g \circ f$ は DM である。
代数空間 $W$ と全射滑らかな射 $W \to \mathcal{Z}$ を選ぶ。
代数空間 $V$ と全射エタールな射
$V \to \mathcal{Y} \times_\mathcal{Z} W$ を選ぶ
（Lemma \ref{stacks-morphisms-lemma-DM}）。代数空間 $U$ と全射エタールな射
$U \to \mathcal{X} \times_\mathcal{Y} V$ を選ぶ。このとき定義により、
射 $V \to W$ と $U \to V$ は性質 $\mathcal{P}$ をもつ。
$\mathcal{P}$ は合成の下で安定であると仮定したので、$U \to W$ は
性質 $\mathcal{P}$ をもつ。したがって再び定義により、
$g \circ f : \mathcal{X} \to \mathcal{Z}$ は性質 $\mathcal{P}$ をもつ。
\end{remark}

\begin{remark}
\label{stacks-morphisms-remark-etale-smooth-base-change}
$\mathcal{P}$ を、始域と終域上エタール・滑らか局所的で、基底変換の下で
安定な代数空間の射の性質とする。このとき
Definition \ref{stacks-morphisms-definition-etale-smooth-P} で定義した代数スタックの
DM 射の性質は、任意の基底変換の下で安定である。
実際、$f : \mathcal{X} \to \mathcal{Y}$ を代数スタックの DM 射、
$g : \mathcal{Y}' \to \mathcal{Y}$ を代数スタックの射とし、
$f$ は性質 $\mathcal{P}$ をもつと仮定する。このとき基底変換
$\mathcal{Y}' \times_\mathcal{Y} \mathcal{X} \to \mathcal{Y}'$
は Lemma \ref{stacks-morphisms-lemma-base-change-separated} により DM 射である。
代数空間 $V$ と全射滑らかな射 $V \to \mathcal{Y}$ を選ぶ。
代数空間 $U$ と全射エタールな射
$U \to \mathcal{X} \times_\mathcal{Y} V$ を選ぶ（Lemma \ref{stacks-morphisms-lemma-DM}）。
最後に、代数空間 $V'$ と全射滑らかな射
$V' \to \mathcal{Y}' \times_\mathcal{Y} V$ を選ぶ。このとき定義により
射 $U \to V$ は性質 $\mathcal{P}$ をもつ。$\mathcal{P}$ は基底変換の下で
安定であると仮定したので、$V' \times_V U \to V'$ は
性質 $\mathcal{P}$ をもつ。図式
$$
\xymatrix{
V' \times_V U \ar[r] \ar[d] &
\mathcal{Y}' \times_\mathcal{Y} \mathcal{X} \ar[r] \ar[d] &
\mathcal{X} \ar[d] \\
V' \ar[r] & \mathcal{Y}' \ar[r] & \mathcal{Y}
}
$$
を考えると、左上の水平射は全射であり、
$$
V' \times_V U \to V' \times_\mathcal{Y}
(\mathcal{Y}' \times_{\mathcal{Y}'} \mathcal{X}) =
V' \times_V (\mathcal{X} \times_\mathcal{Y} V)
$$
は $U \to \mathcal{X} \times_\mathcal{Y} V$ の基底変換としてエタールである。
したがって定義により、射影
$\mathcal{Y}' \times_\mathcal{Y} \mathcal{X} \to \mathcal{Y}'$ は
性質 $\mathcal{P}$ をもつ。
\end{remark}

\begin{remark}
\label{stacks-morphisms-remark-etale-smooth-implication}
$\mathcal{P}, \mathcal{P}'$ を、始域と終域上エタール・滑らか局所的な
代数空間の射の性質とする。代数空間の射について
$\mathcal{P} \Rightarrow \mathcal{P}'$ が成り立つと仮定する。
このとき、$\mathcal{P}$ と $\mathcal{P}'$ を用いて
Definition \ref{stacks-morphisms-definition-etale-smooth-P} で定義した代数スタックの射の
性質についても $\mathcal{P} \Rightarrow \mathcal{P}'$ が成り立つ。
これは定義から明らかである。
\end{remark}








\section{エタール射}
\label{stacks-morphisms-section-etale}

\noindent
代数スタックのエタール射を、相対次元 $0$ の滑らかな射として
定義すべきではない。実際、射
$$
[\mathbf{A}^1_k/\mathbf{G}_{m, k}] \longrightarrow \Spec(k)
$$
は、群スキーム $\mathbf{G}_{m, k}$ の $\mathbf{A}^1_k$ への作用を
どのように選んでも相対次元 $0$ で滑らかである。これは、エタール射は
接空間を同一視すべきだという通常の考え方に対応しない。上の例は
準有限ではないが、射
$$
\mathcal{X} = [\Spec(k)/\mu_{p, k}] \longrightarrow \Spec(k)
$$
は滑らかかつ準有限である（Section \ref{stacks-morphisms-section-locally-quasi-finite}）。
しかし $k$ の標数が $p > 0$ ならば、表現可能な射
$\Spec(k) \to \mathcal{X}$ はエタールではない。実際、その基底変換
$\mu_{p, k} = \Spec(k) \times_\mathcal{X} \Spec(k) \to \Spec(k)$
は非被約スキームから体のスペクトルへの射だからである。
したがって、エタール射を滑らかかつ局所準有限として定義すると、
Morphisms of Spaces, Lemma
\ref{spaces-morphisms-lemma-etale-permanence} の類似は成り立たない。

\medskip\noindent
代わりに、エタール性は基底変換の下で保たれる性質であるべきこと、また
$\mathcal{X} \to X$ が代数スタックからスキームへのエタール射ならば
$\mathcal{X}$ は Deligne--Mumford であるべきことを要請するところから
始める。言い換えると、エタール射は DM であることを要求し、
Section \ref{stacks-morphisms-section-etale-smooth-local-source-target} の内容を用いて
代数スタックのエタール射を定義する。

\medskip\noindent
Lemma \ref{stacks-morphisms-lemma-characterize-etale} では、代数スタックのエタール射を、
(a) 局所有限表示、(b) 平坦、かつ (c) エタールな対角射をもつ射として
特徴づける。

\medskip\noindent
代数空間の射の「エタール」という性質は、始域と終域上
エタール・滑らか局所的である。次を参照せよ：
Descent on Spaces, Remark
\ref{spaces-descent-remark-list-etale-smooth-local-source-target}.
これは基底変換の下でも安定であり、終域上 fpqc 局所的でもある。
次を参照せよ：
Morphisms of Spaces,
Lemma \ref{spaces-morphisms-lemma-base-change-etale}
および
Descent on Spaces,
Lemma \ref{spaces-descent-lemma-descending-property-etale}.
したがって、上の Lemma \ref{stacks-morphisms-lemma-etale-smooth-local-source-target} により、
代数スタックの射がエタールであることを次のように定義できる。
この定義は、射が代数空間で表現可能な場合には Properties of Stacks,
Section
\ref{stacks-properties-section-properties-morphisms}
で定義済みの概念と一致する。そのような射は
Lemma \ref{stacks-morphisms-lemma-trivial-implications} により自動的に DM だからである。

\begin{definition}
\label{stacks-morphisms-definition-etale}
$f : \mathcal{X} \to \mathcal{Y}$ を代数スタックの射とする。$f$ が DM で、
Lemma \ref{stacks-morphisms-lemma-etale-smooth-local-source-target}
の同値な条件が $\mathcal{P} = \etale$ について成り立つとき、
$f$ は {\it エタール}であるという。
\end{definition}

\noindent
$f$ が代数空間で表現可能である場合、または $\mathcal{X}$ と $\mathcal{Y}$ が
代数空間で表現可能である場合には、これが既存のエタール射の概念と一致する
ことを、今後断りなく用いる。

\begin{lemma}
\label{stacks-morphisms-lemma-composition-etale}
エタール射の合成はエタールである。
\end{lemma}

\begin{proof}
Remark \ref{stacks-morphisms-remark-etale-smooth-composition} と Morphisms of Spaces,
Lemma \ref{spaces-morphisms-lemma-composition-etale} を組み合わせる。
\end{proof}

\begin{lemma}
\label{stacks-morphisms-lemma-base-change-etale}
エタール射の基底変換はエタールである。
\end{lemma}

\begin{proof}
Remark \ref{stacks-morphisms-remark-etale-smooth-base-change} と Morphisms of Spaces,
Lemma \ref{spaces-morphisms-lemma-base-change-etale} を組み合わせる。
\end{proof}

\begin{lemma}
\label{stacks-morphisms-lemma-open-immersion-etale}
開埋入はエタールである。
\end{lemma}

\begin{proof}
$j : \mathcal{U} \to \mathcal{X}$ を代数スタックの開埋入とする。
$j$ は表現可能なので、Lemma \ref{stacks-morphisms-lemma-trivial-implications} により DM である。
一方、$X$ がスキームで $X \to \mathcal{X}$ が滑らかかつ全射ならば、
$U = \mathcal{U} \times_\mathcal{X} X$ は $X$ の開部分スキームである。
したがって $U \to X$ はエタールであり
（Morphisms, Lemma \ref{morphisms-lemma-open-immersion-etale}）、
定義から $j$ はエタールであると結論できる。
\end{proof}

\begin{lemma}
\label{stacks-morphisms-lemma-etale}
$f : \mathcal{X} \to \mathcal{Y}$ を代数スタックの射とする。
次は同値である。
\begin{enumerate}
\item $f$ はエタールである。
\item $f$ は DM であり、$V$ が代数空間である任意の射
$V \to \mathcal{Y}$ と、$U$ が代数空間である任意のエタール射
$U \to V \times_\mathcal{Y} \mathcal{X}$ について、射 $U \to V$ は
エタールである。
\item $W$ が代数空間である全射、局所有限表示、かつ平坦な射
$W \to \mathcal{Y}$ と、$T$ が代数空間である全射エタールな射
$T \to W \times_\mathcal{Y} \mathcal{X}$ で、射 $T \to W$ が
エタールとなるものが存在する。
\end{enumerate}
\end{lemma}

\begin{proof}
(1) を仮定する。このとき $f$ は DM であり、エタール性は基底変換の下で
保たれるので、(2) が成り立つ。

\medskip\noindent
(2) を仮定する。スキーム $V$ と全射エタールな射
$V \to \mathcal{Y}$ を選ぶ。スキーム $U$ と全射エタールな射
$U \to V \times_\mathcal{Y} \mathcal{X}$ を選ぶ
（Lemma \ref{stacks-morphisms-lemma-DM}）。したがって (3) が成り立つ。

\medskip\noindent
$W \to \mathcal{Y}$ と $T \to W \times_\mathcal{Y} \mathcal{X}$ は
(3) のとおりと仮定する。まず $f$ が DM であることを確かめる。
すなわち、$W \times_\mathcal{Y} \mathcal{X} \to W$ が DM であることを
確かめれば十分である。Lemma \ref{stacks-morphisms-lemma-check-separated-covering} を参照せよ。
Lemma \ref{stacks-morphisms-lemma-compose-after-separated} により、
$W \times_\mathcal{Y} \mathcal{X}$ が DM であることを確かめれば十分である。
これは $T \to W \times_\mathcal{Y} \mathcal{X}$ の存在と
Theorem \ref{stacks-morphisms-theorem-DM} の容易な向きから従う。

\medskip\noindent
$f$ は DM であり、$W \to \mathcal{Y}$ と
$T \to W \times_\mathcal{Y} \mathcal{X}$ は (3) のとおりと仮定する。
$V$ を代数空間、$V \to \mathcal{Y}$ を全射滑らかな射、$U$ を代数空間、
$U \to V \times_\mathcal{Y} \mathcal{X}$ を全射エタールな射とする
（Lemma \ref{stacks-morphisms-lemma-DM}）。$U \to V$ がエタールであることを
確かめなければならない。Descent on Spaces, Lemma
\ref{spaces-descent-lemma-descending-property-etale}.
により、$U \times_\mathcal{Y} W \to V \times_\mathcal{Y} W$ が
エタールであることを証明すれば十分である。
$\mathcal{X}, \mathcal{Y}, W, T, U, V$ を
$\mathcal{X} \times_\mathcal{Y} W, W, W, T, U \times_\mathcal{Y} W,
V \times_\mathcal{Y} W$ で置き換えてよい（細部は省略する）。
したがって $Y = \mathcal{Y}$ は代数空間であり、代数空間 $T$ と
全射エタールな射 $T \to \mathcal{X}$ で $T \to Y$ がエタールとなるものが
存在し、$U$ と $V$ は上のとおりであると仮定してよい。この場合、
$$
U \to V\text{ is \'etale}
\Leftrightarrow
\mathcal{X} \to Y\text{ is \'etale}
\Leftrightarrow
T \to Y\text{ is \'etale}
$$
であることを、Lemma \ref{stacks-morphisms-lemma-etale-smooth-local-source-target} の
性質 (1) と (2) の同値性、および Definition \ref{stacks-morphisms-definition-etale} から得る。
これで証明が完了する。
\end{proof}

\begin{lemma}
\label{stacks-morphisms-lemma-etale-permanence}
$\mathcal{X}, \mathcal{Y}$ を代数スタック $\mathcal{Z}$ 上エタールな
代数スタックとする。$\mathcal{Z}$ 上の任意の射
$\mathcal{X} \to \mathcal{Y}$ はエタールである。
\end{lemma}

\begin{proof}
射 $\mathcal{X} \to \mathcal{Y}$ は
Lemma \ref{stacks-morphisms-lemma-compose-after-separated} により DM である。
$W \to \mathcal{Z}$ を、始域が代数空間である全射滑らかな射とする。
$V \to \mathcal{Y} \times_\mathcal{Z} W$ を、始域が代数空間である
全射エタールな射とする（Lemma \ref{stacks-morphisms-lemma-DM}）。
$U \to \mathcal{X} \times_\mathcal{Y} V$ を、始域が代数空間である
全射エタールな射とする（Lemma \ref{stacks-morphisms-lemma-DM}）。このとき
$$
U \longrightarrow \mathcal{X} \times_\mathcal{Z} W
$$
は、$U \to \mathcal{X} \times_\mathcal{Y} V$ と、
$\mathcal{X} \times_\mathcal{Z} W \to \mathcal{Y} \times_\mathcal{Z} W$
による $V \to \mathcal{Y} \times_\mathcal{Z} W$ の基底変換との合成なので、
全射エタールである。したがって $U \to W$ がエタールであることを
示せば十分である。$\mathcal{X} \to \mathcal{Z}$ と
$\mathcal{Y} \to \mathcal{Z}$ はエタールなので、$U \to W$ と $V \to W$ は
エタールである。ゆえにこれは代数空間についての補題、すなわち
Morphisms of Spaces, Lemma \ref{spaces-morphisms-lemma-etale-permanence}
から従う。
\end{proof}








\section{非分岐射}
\label{stacks-morphisms-section-unramified}

\noindent
非分岐射の定義をこのように選ぶ理由については、エタール射を扱う
Section \ref{stacks-morphisms-section-etale} の議論を参照されたい。

\medskip\noindent
Lemma \ref{stacks-morphisms-lemma-characterize-unramified} では、代数スタックの非分岐射を、
局所有限型でエタールな対角射をもつ射として特徴づける。

\medskip\noindent
代数空間の射の「非分岐」という性質は、始域と終域上
エタール・滑らか局所的である。次を参照せよ：
Descent on Spaces, Remark
\ref{spaces-descent-remark-list-etale-smooth-local-source-target}.
これは基底変換の下でも安定であり、終域上 fpqc 局所的でもある。
次を参照せよ：
Morphisms of Spaces,
Lemma \ref{spaces-morphisms-lemma-base-change-unramified}
および
Descent on Spaces,
Lemma \ref{spaces-descent-lemma-descending-property-unramified}.
したがって、上の Lemma \ref{stacks-morphisms-lemma-etale-smooth-local-source-target} により、
代数スタックの射が非分岐であることを次のように定義できる。
この定義は、射が代数空間で表現可能な場合には Properties of Stacks,
Section
\ref{stacks-properties-section-properties-morphisms}
で定義済みの概念と一致する。そのような射は
Lemma \ref{stacks-morphisms-lemma-trivial-implications} により自動的に DM だからである。

\begin{definition}
\label{stacks-morphisms-definition-unramified}
$f : \mathcal{X} \to \mathcal{Y}$ を代数スタックの射とする。$f$ が DM で、
Lemma \ref{stacks-morphisms-lemma-etale-smooth-local-source-target}
の同値な条件が $\mathcal{P} =$「非分岐」について成り立つとき、
$f$ は {\it 非分岐}であるという。
\end{definition}

\noindent
$f$ が代数空間で表現可能である場合、または $\mathcal{X}$ と $\mathcal{Y}$ が
代数空間で表現可能である場合には、これが既存の非分岐射の概念と一致する
ことを、今後断りなく用いる。

\begin{lemma}
\label{stacks-morphisms-lemma-composition-unramified}
非分岐射の合成は非分岐である。
\end{lemma}

\begin{proof}
Remark \ref{stacks-morphisms-remark-etale-smooth-composition} と Morphisms of Spaces,
Lemma \ref{spaces-morphisms-lemma-composition-unramified} を組み合わせる。
\end{proof}

\begin{lemma}
\label{stacks-morphisms-lemma-base-change-unramified}
非分岐射の基底変換は非分岐である。
\end{lemma}

\begin{proof}
Remark \ref{stacks-morphisms-remark-etale-smooth-base-change} と Morphisms of Spaces,
Lemma \ref{spaces-morphisms-lemma-base-change-unramified} を組み合わせる。
\end{proof}

\begin{lemma}
\label{stacks-morphisms-lemma-etale-unramified}
エタール射は非分岐である。
\end{lemma}

\begin{proof}
Remark \ref{stacks-morphisms-remark-etale-smooth-implication} と Morphisms of Spaces,
Lemma \ref{spaces-morphisms-lemma-etale-unramified} から従う。
\end{proof}

\begin{lemma}
\label{stacks-morphisms-lemma-immersion-unramified}
埋入は非分岐である。
\end{lemma}

\begin{proof}
$j : \mathcal{Z} \to \mathcal{X}$ を代数スタックの埋入とする。
$j$ は表現可能なので、Lemma \ref{stacks-morphisms-lemma-trivial-implications} により DM である。
一方、$X$ がスキームで $X \to \mathcal{X}$ が滑らかかつ全射ならば、
$Z = \mathcal{Z} \times_\mathcal{X} X$ は $X$ の局所閉部分スキームである。
したがって $Z \to X$ は非分岐であり
（Morphisms, Lemmas \ref{morphisms-lemma-open-immersion-unramified} および
\ref{morphisms-lemma-closed-immersion-unramified})
定義から $j$ は非分岐であると結論できる。
\end{proof}

\begin{lemma}
\label{stacks-morphisms-lemma-unramified}
$f : \mathcal{X} \to \mathcal{Y}$ を代数スタックの射とする。
次は同値である。
\begin{enumerate}
\item $f$ は非分岐である。
\item $f$ は DM であり、$V$ が代数空間である任意の射
$V \to \mathcal{Y}$ と、$U$ が代数空間である任意のエタール射
$U \to V \times_\mathcal{Y} \mathcal{X}$ について、射 $U \to V$ は
非分岐である。
\item $W$ が代数空間である全射、局所有限表示、かつ平坦な射
$W \to \mathcal{Y}$ と、$T$ が代数空間である全射エタールな射
$T \to W \times_\mathcal{Y} \mathcal{X}$ で、射 $T \to W$ が
非分岐となるものが存在する。
\end{enumerate}
\end{lemma}

\begin{proof}
(1) を仮定する。このとき $f$ は DM であり、非分岐性は基底変換の下で
保たれるので、(2) が成り立つ。

\medskip\noindent
(2) を仮定する。スキーム $V$ と全射エタールな射
$V \to \mathcal{Y}$ を選ぶ。スキーム $U$ と全射エタールな射
$U \to V \times_\mathcal{Y} \mathcal{X}$ を選ぶ
（Lemma \ref{stacks-morphisms-lemma-DM}）。したがって (3) が成り立つ。

\medskip\noindent
$W \to \mathcal{Y}$ と $T \to W \times_\mathcal{Y} \mathcal{X}$ は
(3) のとおりと仮定する。まず $f$ が DM であることを確かめる。
すなわち、$W \times_\mathcal{Y} \mathcal{X} \to W$ が DM であることを
確かめれば十分である。Lemma \ref{stacks-morphisms-lemma-check-separated-covering} を参照せよ。
Lemma \ref{stacks-morphisms-lemma-compose-after-separated} により、
$W \times_\mathcal{Y} \mathcal{X}$ が DM であることを確かめれば十分である。
これは $T \to W \times_\mathcal{Y} \mathcal{X}$ の存在と
Theorem \ref{stacks-morphisms-theorem-DM} の容易な向きから従う。

\medskip\noindent
$f$ は DM であり、$W \to \mathcal{Y}$ と
$T \to W \times_\mathcal{Y} \mathcal{X}$ は (3) のとおりと仮定する。
$V$ を代数空間、$V \to \mathcal{Y}$ を全射滑らかな射、$U$ を代数空間、
$U \to V \times_\mathcal{Y} \mathcal{X}$ を全射エタールな射とする
（Lemma \ref{stacks-morphisms-lemma-DM}）。$U \to V$ が非分岐であることを
確かめなければならない。Descent on Spaces, Lemma
\ref{spaces-descent-lemma-descending-property-unramified}.
により、$U \times_\mathcal{Y} W \to V \times_\mathcal{Y} W$ が
非分岐であることを証明すれば十分である。
$\mathcal{X}, \mathcal{Y}, W, T, U, V$ を
$\mathcal{X} \times_\mathcal{Y} W, W, W, T, U \times_\mathcal{Y} W,
V \times_\mathcal{Y} W$ で置き換えてよい（細部は省略する）。
したがって $Y = \mathcal{Y}$ は代数空間であり、代数空間 $T$ と
全射エタールな射 $T \to \mathcal{X}$ で $T \to Y$ が非分岐となるものが
存在し、$U$ と $V$ は上のとおりであると仮定してよい。この場合、
$$
U \to V\text{ is unramified}
\Leftrightarrow
\mathcal{X} \to Y\text{ is unramified}
\Leftrightarrow
T \to Y\text{ is unramified}
$$
であることを、Lemma \ref{stacks-morphisms-lemma-etale-smooth-local-source-target} の
性質 (1) と (2) の同値性、および Definition \ref{stacks-morphisms-definition-unramified}
から得る。これで証明が完了する。
\end{proof}

\begin{lemma}
\label{stacks-morphisms-lemma-unramified-quasi-finite}
代数スタックの非分岐射は局所準有限である。
\end{lemma}

\begin{proof}
これは Lemma \ref{stacks-morphisms-lemma-unramified}（非分岐射の特徴づけ）、
Lemma \ref{stacks-morphisms-lemma-characterize-locally-quasi-finite}
（局所準有限射の特徴づけ）、および
Morphisms of Spaces, Lemma
\ref{spaces-morphisms-lemma-unramified-quasi-finite}
（代数空間についての対応する結果）から従う。
\end{proof}

\begin{lemma}
\label{stacks-morphisms-lemma-permanence-unramified}
$\mathcal{X} \to \mathcal{Y} \to \mathcal{Z}$ を代数スタックの射とする。
$\mathcal{X} \to \mathcal{Z}$ が非分岐で
$\mathcal{Y} \to \mathcal{Z}$ が DM ならば、
$\mathcal{X} \to \mathcal{Y}$ は非分岐である。
\end{lemma}

\begin{proof}
$\mathcal{X} \to \mathcal{Z}$ は非分岐であると仮定する。
Lemma \ref{stacks-morphisms-lemma-compose-after-separated} により、射
$\mathcal{X} \to \mathcal{Y}$ は DM である。可換図式
$$
\xymatrix{
U \ar[d] \ar[r] & V \ar[d] \ar[r] & W \ar[d] \\
\mathcal{X} \ar[r] & \mathcal{Y} \ar[r] & \mathcal{Z}
}
$$
を選ぶ。ここで $U, V, W$ は代数空間、$W \to \mathcal{Z}$ は全射滑らか、
$V \to \mathcal{Y} \times_\mathcal{Z} W$ は全射エタール、かつ
$U \to \mathcal{X} \times_\mathcal{Y} V$ は全射エタールである
（Lemma \ref{stacks-morphisms-lemma-DM} を参照）。このとき
$U \to \mathcal{X} \times_\mathcal{Z} W$ も全射エタールである。
したがって $U \to W$ は非分岐であることが分かり、$U \to V$ が
非分岐であることを示せばよい。これは Morphisms of Spaces, Lemma
\ref{spaces-morphisms-lemma-permanence-unramified} から従う。
\end{proof}

\begin{lemma}
\label{stacks-morphisms-lemma-characterize-unramified}
$f : \mathcal{X} \to \mathcal{Y}$ を代数スタックの射とする。
次は同値である。
\begin{enumerate}
\item $f$ は非分岐である。
\item $f$ は局所有限型であり、その対角射はエタールである。
\end{enumerate}
\end{lemma}

\begin{proof}
$f$ は非分岐であると仮定する。このとき $f$ は DM なので、代数空間
$U$, $V$、全射滑らかな射 $V \to \mathcal{Y}$、および全射エタールな射
$U \to \mathcal{X} \times_\mathcal{Y} V$ を選べる
（Lemma \ref{stacks-morphisms-lemma-DM}）。$f$ は非分岐なので、誘導射 $U \to V$ は
非分岐である。したがって $U \to V$ は局所有限型である
(Morphisms of Spaces, Lemma
\ref{spaces-morphisms-lemma-unramified-locally-finite-type})
から、$f$ は局所有限型であると結論できる。対角射
$\Delta : \mathcal{X} \to \mathcal{X} \times_\mathcal{Y} \mathcal{X}$ は
$\mathcal{Y}$ 上の代数スタックの射である。全射滑らかな射
$V \to \mathcal{Y}$ による $\Delta$ の基底変換は、$f$ の基底変換、
すなわち $\mathcal{X}_V = \mathcal{X} \times_\mathcal{Y} V \to V$ の
対角射である。言い換えると、図式
$$
\xymatrix{
\mathcal{X}_V \ar[r] \ar[d] &  \mathcal{X}_V \times_V \mathcal{X}_V \ar[d] \\
\mathcal{X} \ar[r] & \mathcal{X} \times_\mathcal{Y} \mathcal{X}
}
$$
は Cartesian である。右の垂直射は全射かつ滑らかなので、
Properties of Stacks, Lemma
\ref{stacks-properties-lemma-check-property-weak-covering}.
により、上の水平射がエタールであることを示せば十分である。
可換図式
$$
\xymatrix{
U \ar[d] \ar[r] & U \times_V U \ar[d] \\
\mathcal{X}_V \ar[r] & \mathcal{X}_V \times_V \mathcal{X}_V
}
$$
を考える。すべての射は代数空間で表現可能であり、垂直射はエタール、
左の垂直射は全射で、上の水平射は
Morphisms of Spaces, Lemma
\ref{spaces-morphisms-lemma-diagonal-unramified-morphism} により開埋入である。
これから所望の結論を得る。まず、エタール射の合成として
$U \to \mathcal{X}_V \times_V \mathcal{X}_V$ はエタールである。
次に
Properties of Stacks, Lemma
\ref{stacks-properties-lemma-check-property-after-precomposing}
を用いると、$\mathcal{X}_V \to \mathcal{X}_V \times_V \mathcal{X}_V$ は
エタールであると分かる。実際、（代数空間の射について）エタールであることは
エタール位相で始域上局所的である
（Descent on Spaces, Lemma \ref{spaces-descent-lemma-etale-etale-local-source}）。

\medskip\noindent
$f$ は局所有限型であり、その対角射はエタールであると仮定する。
このとき $f$ は定義により DM である（代数空間のエタール射は
非分岐だからである）。上と同様に、代数空間 $U$, $V$、全射滑らかな射
$V \to \mathcal{Y}$、および全射エタールな射
$U \to \mathcal{X} \times_\mathcal{Y} V$ を選べる
（Lemma \ref{stacks-morphisms-lemma-DM}）。証明を終えるには $U \to V$ が非分岐であることを
示さなければならない。$U \to V$ は局所有限型であることが既に分かっている。
上と同様に議論すると、可換図式
$$
\xymatrix{
U \ar[d] \ar[r] & U \times_V U \ar[d] \\
\mathcal{X}_V \ar[r] & \mathcal{X}_V \times_V \mathcal{X}_V
}
$$
を得る。ここですべての射は代数空間で表現可能であり、垂直射はエタール、
下の水平射は $\Delta$ の基底変換としてエタールである。したがって、例えば
Lemma \ref{stacks-morphisms-lemma-etale-permanence}\footnote{これは Morphisms of Spaces,
Lemma \ref{spaces-morphisms-lemma-etale-permanence} から直接導くことも
容易である。} により、$U \to U \times_V U$ はエタールである。
ゆえに $U \to U \times_V U$ はエタールな単射、したがって開埋入である
（Morphisms of Spaces, Lemma
\ref{spaces-morphisms-lemma-etale-universally-injective-open}).
よって $U \to V$ は
Morphisms of Spaces, Lemma
\ref{spaces-morphisms-lemma-diagonal-unramified-morphism} により非分岐である。
\end{proof}

\begin{lemma}
\label{stacks-morphisms-lemma-characterize-etale}
$f : \mathcal{X} \to \mathcal{Y}$ を代数スタックの射とする。
次は同値である。
\begin{enumerate}
\item $f$ はエタールである。
\item $f$ は局所有限表示、平坦、かつ非分岐である。
\item $f$ は局所有限表示かつ平坦であり、その対角射はエタールである。
\end{enumerate}
\end{lemma}

\begin{proof}
(2) と (3) の同値性は Lemma \ref{stacks-morphisms-lemma-characterize-unramified} から
直ちに従う。したがって、いずれの場合にも射 $f$ は DM である。
このとき代数空間 $U$, $V$、全射滑らかな射 $V \to \mathcal{Y}$、および
全射エタールな射 $U \to \mathcal{X} \times_\mathcal{Y} V$ を選べる
（Lemma \ref{stacks-morphisms-lemma-DM}）。証明を終えるには、$U \to V$ がエタールである
ことと、局所有限表示、平坦、かつ非分岐であることが同値であると
示さなければならない。これは Morphisms of Spaces, Lemma
\ref{spaces-morphisms-lemma-unramified-flat-lfp-etale}
（および、より自明な Morphisms of Spaces, Lemmas
\ref{spaces-morphisms-lemma-etale-unramified},
\ref{spaces-morphisms-lemma-etale-locally-finite-presentation}、および
\ref{spaces-morphisms-lemma-etale-flat}）から従う。
\end{proof}








\section{固有射}
\label{stacks-morphisms-section-proper}

\noindent
固有射の概念は代数幾何において重要な役割を果たす。
代数スタックの固有射を次のように定義する。

\begin{definition}
\label{stacks-morphisms-definition-proper}
$f : \mathcal{X} \to \mathcal{Y}$ を代数スタックの射とする。
$f$ が分離、有限型、かつ普遍閉であるとき、$f$ は {\it 固有}であるという。
\end{definition}

\noindent
これは代数空間の固有射という既存の概念と矛盾しない。代数空間の射が
固有であることと、分離、有限型、かつ普遍閉であることは同値であり
(Morphisms of Spaces, Definition \ref{spaces-morphisms-definition-proper})
これらの概念の整合性は既に
Lemma \ref{stacks-morphisms-lemma-representable-separated-diagonal-closed},
Section \ref{stacks-morphisms-section-finite-type}、および
Lemmas \ref{stacks-morphisms-lemma-characterize-representable-universally-closed} で
確かめた。同様に、$f : \mathcal{X} \to \mathcal{Y}$ が代数空間で
表現可能な代数スタックの射ならば、$f$ が固有であることは
Properties of Stacks, Section
\ref{stacks-properties-section-properties-morphisms} で定義した。
しかし同節の議論により、これは $f$ が分離、有限型、かつ普遍閉であることと
同値であり、上と同じ参照により整合性が得られる。

\begin{lemma}
\label{stacks-morphisms-lemma-base-change-proper}
固有射の基底変換は固有である。
\end{lemma}

\begin{proof}
Lemmas \ref{stacks-morphisms-lemma-base-change-separated}、
\ref{stacks-morphisms-lemma-base-change-finite-type}、および
\ref{stacks-morphisms-lemma-base-change-universally-closed} を参照せよ。
\end{proof}

\begin{lemma}
\label{stacks-morphisms-lemma-composition-proper}
固有射の合成は固有である。
\end{lemma}

\begin{proof}
Lemmas \ref{stacks-morphisms-lemma-composition-separated}、
\ref{stacks-morphisms-lemma-composition-finite-type}、および
\ref{stacks-morphisms-lemma-composition-universally-closed} を参照せよ。
\end{proof}

\begin{lemma}
\label{stacks-morphisms-lemma-closed-immersion-proper}
代数スタックの閉埋入は代数スタックの固有射である。
\end{lemma}

\begin{proof}
閉埋入は定義により表現可能である
(Properties of Stacks, Definition
\ref{stacks-properties-definition-immersion}).
したがって、これは
Properties of Stacks, Section
\ref{stacks-properties-section-properties-morphisms}
の議論と、代数空間の射についての対応する結果から従う。次を参照せよ：
Morphisms of Spaces, Lemma
\ref{spaces-morphisms-lemma-closed-immersion-proper}.
\end{proof}

\begin{lemma}
\label{stacks-morphisms-lemma-universally-closed-permanence}
代数スタックの可換図式
$$
\xymatrix{
\mathcal{X} \ar[rr] \ar[rd] & &
\mathcal{Y} \ar[ld] \\
& \mathcal{Z} &
}
$$
を考える。
\begin{enumerate}
\item $\mathcal{X} \to \mathcal{Z}$ が普遍閉で
$\mathcal{Y} \to \mathcal{Z}$ が分離ならば、射
$\mathcal{X} \to \mathcal{Y}$ は普遍閉である。特に、$|\mathcal{X}|$ の
$|\mathcal{Y}|$ における像は閉である。
\item $\mathcal{X} \to \mathcal{Z}$ が固有で
$\mathcal{Y} \to \mathcal{Z}$ が分離ならば、射
$\mathcal{X} \to \mathcal{Y}$ は固有である。
\end{enumerate}
\end{lemma}

\begin{proof}
$\mathcal{X} \to \mathcal{Z}$ は普遍閉で、
$\mathcal{Y} \to \mathcal{Z}$ は分離であると仮定する。射を
$\mathcal{X} \to \mathcal{X} \times_\mathcal{Z} \mathcal{Y} \to \mathcal{Y}$.
と分解する。最初の射は固有であり（Lemma \ref{stacks-morphisms-lemma-semi-diagonal}）、
したがって普遍閉である。射影
$\mathcal{X} \times_\mathcal{Z} \mathcal{Y} \to \mathcal{Y}$ は普遍閉射の
基底変換なので普遍閉である。次を参照せよ：
Lemma \ref{stacks-morphisms-lemma-base-change-universally-closed}。
したがって $\mathcal{X} \to \mathcal{Y}$ は普遍閉射の合成として普遍閉である。
次を参照せよ：Lemma \ref{stacks-morphisms-lemma-composition-universally-closed}。
これで (1) が証明された。(2) を導くには、(1) と
Lemmas \ref{stacks-morphisms-lemma-compose-after-separated}、
\ref{stacks-morphisms-lemma-quasi-compact-permanence}、および
\ref{stacks-morphisms-lemma-finite-type-permanence} を組み合わせる。
\end{proof}

\begin{lemma}
\label{stacks-morphisms-lemma-image-proper-is-proper}
$\mathcal{Z}$ を代数スタックとする。
$f : \mathcal{X} \to \mathcal{Y}$ を $\mathcal{Z}$ 上の代数スタックの射とする。
$\mathcal{X}$ が $\mathcal{Z}$ 上普遍閉で $f$ が全射ならば、
$\mathcal{Y}$ は $\mathcal{Z}$ 上普遍閉である。特に、さらに
$\mathcal{Y}$ が $\mathcal{Z}$ 上分離かつ有限型ならば、
$\mathcal{Y}$ は $\mathcal{Z}$ 上固有である。
\end{lemma}

\begin{proof}
$\mathcal{X}$ は普遍閉で $f$ は全射であると仮定する。構造射を
$p : \mathcal{X} \to \mathcal{Z}$、$q : \mathcal{Y} \to \mathcal{Z}$ と書く。
$\mathcal{Z}' \to \mathcal{Z}$ を代数スタックの射とする。
$\mathcal{Z}' \to \mathcal{Z}$ による $f$ の基底変換
$f' : \mathcal{X}' \to \mathcal{Y}'$ は全射であり
(Properties of Stacks, Lemma
\ref{stacks-properties-lemma-base-change-surjective})、$p$ の基底変換
$p' : \mathcal{X}' \to \mathcal{Z}'$ は閉である。
$T \subset |\mathcal{Y}'|$ が閉ならば、
$(f')^{-1}(T) \subset |\mathcal{X}'|$ は閉であり、したがって
$p'((f')^{-1}(T)) = q'(T)$ は閉である。ゆえに $q'$ は閉である。
\end{proof}








\section{スキーム論的像}
\label{stacks-morphisms-section-scheme-theoretic-image}

\noindent
定義は次のとおりである。

\begin{definition}
\label{stacks-morphisms-definition-scheme-theoretic-image}
$f : \mathcal{X} \to \mathcal{Y}$ を代数スタックの射とする。
$f$ の {\it スキーム論的像}とは、$f$ が経由する最小の閉部分スタック
$\mathcal{Z} \subset \mathcal{Y}$ である\footnote{
Lemma \ref{stacks-morphisms-lemma-scheme-theoretic-image-existence}
で、スキーム論的像が常に存在することを見る。}。
\end{definition}

\noindent
$f$ のこの分解をしばしば $f : \mathcal{X} \to \mathcal{Z}$ と書く。
射 $f$ が準コンパクトでなければ、一般にスキーム論的像の構成は
$\mathcal{Y}$ の開部分スタックへの制限と可換でない。しかし $f$ が
準コンパクトならば、スキーム論的像は平坦基底変換と可換である
（Lemma \ref{stacks-morphisms-lemma-existence-plus-flat-base-change}）。

\begin{lemma}
\label{stacks-morphisms-lemma-cover-upstairs}
$f : \mathcal{X} \to \mathcal{Y}$ を代数スタックの射とする。
$g : \mathcal{W} \to \mathcal{X}$ を全射、平坦、かつ局所有限表示な
代数スタックの射とする。このとき $f$ のスキーム論的像が存在することと
$f \circ g$ のスキーム論的像が存在することは同値であり、存在する場合、
両者のスキーム論的像は同じである。
\end{lemma}

\begin{proof}
$\mathcal{Z} \subset \mathcal{Y}$ を閉部分スタックとし、$f \circ g$ は
$\mathcal{Z}$ を経由すると仮定する。補題を証明するには、$f$ が
$\mathcal{Z}$ を経由することを示せば十分である。スキーム $T$ と、
$T$ 上の $\mathcal{X}$ のファイバー圏の対象 $x$ で与えられる射
$T \to \mathcal{X}$ を考える。$f(x)$ が実際に $T$ 上の $\mathcal{Z}$ の
ファイバー圏に属することを示す。射影
$T \times_\mathcal{X} \mathcal{W} \to T$ は全射、平坦、かつ
局所有限表示である。したがって、すべての $i$ について $T_i \to T$ が
$T \times_\mathcal{X} \mathcal{W} \to T$ を経由するような fppf 被覆
$\{T_i \to T\}$ が存在する。このとき $T_i \to \mathcal{X}$ は
$\mathcal{W}$ を経由し、したがって $T_i \to \mathcal{Y}$ は
$\mathcal{Z}$ を経由する。ゆえに $f(x)|_{T_i}$ は $\mathcal{Z}$ の対象である。
$\mathcal{Z}$ は厳密充満部分スタックなので、所望のとおり $f(x)$ は
$\mathcal{Z}$ の対象であると結論できる。
\end{proof}

\begin{lemma}
\label{stacks-morphisms-lemma-scheme-theoretic-image-existence}
$f : \mathcal{Y} \to \mathcal{X}$ を代数スタックの射とする。
このとき $f$ のスキーム論的像は存在する。
\end{lemma}

\begin{proof}
スキーム $V$ と全射滑らかな射 $V \to \mathcal{Y}$ を選ぶ。
Lemma \ref{stacks-morphisms-lemma-cover-upstairs} により、$\mathcal{Y}$ を $V$ で置き換えてよい。
したがって、スキームから代数スタックへの射 $X \to \mathcal{X}$ の
スキーム論的像が存在することを示せば十分である。スキーム $U$ と
全射滑らかな射 $U \to \mathcal{X}$ を選び、$R = U \times_\mathcal{X} U$ と
おく。
Algebraic Stacks, Lemma \ref{algebraic-lemma-stack-presentation}.
により $\mathcal{X} = [U/R]$ である。Properties of Stacks, Lemma
\ref{stacks-properties-lemma-substacks-presentation}
により、$\mathcal{X}$ の閉部分スタック $\mathcal{Z}$ と、
$R$ 不変な閉部分スキーム $Z \subset U$ は $1$ 対 $1$ に対応する。
$Z_1 \subset U$ を $X \times_\mathcal{X} U \to U$ のスキーム論的像とする。
$X \to \mathcal{X}$ が $\mathcal{Z}$ を経由することと、
$X \times_\mathcal{X} U \to U$ が対応する $R$ 不変な閉部分スキーム $Z$ を
経由することは同値であることに注意する（細部は省略する。ヒント：これは
$X \times_\mathcal{X} U \to X$ が全射かつ滑らかであることから従う）。
したがって、$Z_1$ を含む最小の $R$ 不変な閉部分スキーム
$Z \subset U$ が存在することを示さなければならない。

\medskip\noindent
$\mathcal{I}_1 \subset \mathcal{O}_U$ を閉部分スキーム
$Z_1 \subset U$ に対応する準連接イデアル層とする。
$Z_\alpha$, $\alpha \in A$ を、$Z_1$ を含む $U$ のすべての
$R$ 不変な閉部分スキームの集合とする。$\alpha \in A$ について、
$\mathcal{I}_\alpha \subset \mathcal{O}_U$ を閉部分スキーム
$Z_\alpha \subset U$ に対応する準連接イデアル層とする。
包含 $Z_1 \subset Z_\alpha$ は
$\mathcal{I}_\alpha \subset \mathcal{I}_1$ を意味する。
$Z_\alpha$ が $R$ 不変であることは、
$$
s^{-1}\mathcal{I}_\alpha \cdot \mathcal{O}_R =
t^{-1}\mathcal{I}_\alpha \cdot \mathcal{O}_R
$$
が（代数空間）$R$ 上の（準連接）イデアル層として成り立つことを意味する。
像
$$
\mathcal{I} =
\Im\left(
\bigoplus\nolimits_{\alpha \in A} \mathcal{I}_\alpha \to \mathcal{I}_1
\right) =
\Im\left(
\bigoplus\nolimits_{\alpha \in A} \mathcal{I}_\alpha \to \mathcal{O}_X
\right)
$$
を考える。準連接層の直和は準連接であり、準連接層間の写像の像も
準連接なので、$\mathcal{I}$ は準連接である。引き戻しは完全で直和と
可換なので、
$$
s^{-1}\mathcal{I} \cdot \mathcal{O}_R =
t^{-1}\mathcal{I} \cdot \mathcal{O}_R
$$
を得る。したがって $\mathcal{I}$ は、すべての $Z_\alpha$ に含まれ、
$Z_1$ を含む $R$ 不変な閉部分スキーム $Z \subset U$ を定める。
これは所望のものである。
\end{proof}

\begin{lemma}
\label{stacks-morphisms-lemma-factor-factor}
代数スタックの可換図式
$$
\xymatrix{
\mathcal{X}_1 \ar[d] \ar[r]_{f_1} & \mathcal{Y}_1 \ar[d] \\
\mathcal{X}_2 \ar[r]^{f_2} & \mathcal{Y}_2
}
$$
を考える。$\mathcal{Z}_i \subset \mathcal{Y}_i$, $i = 1, 2$ を
$f_i$ のスキーム論的像とする。このとき射
$\mathcal{Y}_1 \to \mathcal{Y}_2$ は射
$\mathcal{Z}_1 \to \mathcal{Z}_2$ と可換図式
$$
\xymatrix{
\mathcal{X}_1 \ar[r] \ar[d] &
\mathcal{Z}_1 \ar[d] \ar[r] &
\mathcal{Y}_1 \ar[d] \\
\mathcal{X}_2 \ar[r] &
\mathcal{Z}_2 \ar[r] &
\mathcal{Y}_2
}
$$
を誘導する。
\end{lemma}

\begin{proof}
$\mathcal{Y}_1$ における $\mathcal{Z}_2$ のスキーム論的逆像は、
$f_1$ が経由する $\mathcal{Y}_1$ の閉部分スタックである。
したがって $\mathcal{Z}_1$ はこれに含まれる。これで補題が証明された。
\end{proof}

\begin{lemma}
\label{stacks-morphisms-lemma-existence-plus-flat-base-change}
$f : \mathcal{X} \to \mathcal{Y}$ を代数スタックの準コンパクトな射とする。
このときスキーム論的像の形成は平坦基底変換と可換である。
\end{lemma}

\begin{proof}
$\mathcal{Y}' \to \mathcal{Y}$ を代数スタックの平坦な射とする。
スキーム $V$ と全射滑らかな射 $V \to \mathcal{Y}$ を選ぶ。
スキーム $V'$ と全射滑らかな射
$V' \to \mathcal{Y}' \times_\mathcal{Y} V$ を選ぶ。
$V = \coprod_{i \in I} V_i$ はアフィンスキームの非交和であり、
$V' = \coprod_{i \in I} \coprod_{j \in J_i} V_{i, j}$
は各 $V_{i, j}$ が $V_i$ に写るアフィンスキームの非交和であると仮定してよく、
実際そのように仮定する。次のようにおく。
\begin{enumerate}
\item $\mathcal{Z} \subset \mathcal{Y}$ は $f$ のスキーム論的像。
\item $\mathcal{Z}' \subset \mathcal{Y}'$ は
$\mathcal{Y}' \to \mathcal{Y}$ による $f$ の基底変換のスキーム論的像。
\item $Z \subset V$ は $V \to \mathcal{Y}$ による $f$ の
基底変換のスキーム論的像。
\item $Z' \subset V'$ は $V' \to \mathcal{Y}$ による $f$ の
基底変換のスキーム論的像。
\end{enumerate}
(a) $Z = V \times_\mathcal{Y} \mathcal{Z}$、
(b) $Z' = V' \times_{\mathcal{Y}'} \mathcal{Z}'$、および
(c) $Z' = V' \times_V Z$ を示せれば、補題は従う。実際、包含
$\mathcal{Z}' \to \mathcal{Z} \times_\mathcal{Y} \mathcal{Y}'$
(Lemma \ref{stacks-morphisms-lemma-factor-factor})
は、全射滑らかな射 $V' \to \mathcal{Y}'$ による基底変換後に同型なので、
同型でなければならない。

\medskip\noindent
(a) の証明。$R = V \times_\mathcal{Y} V$ とおく。
Properties of Stacks, Lemma
\ref{stacks-properties-lemma-substacks-presentation}
により、規則 $\mathcal{Z} \mapsto \mathcal{Z} \times_\mathcal{Y} V$ は、
$\mathcal{Y}$ の閉部分スタックと $V$ の $R$ 不変な閉部分空間との間の
$1$ 対 $1$ の対応を定める。さらに、
$f : \mathcal{X} \to \mathcal{Y}$ が $\mathcal{Z}$ を経由することと、基底変換
$g : \mathcal{X} \times_\mathcal{Y} V \to V$ が
$\mathcal{Z} \times_\mathcal{Y} V$ を経由することは同値である。
$g$ のスキーム論的像 $Z \subset V$ は $R$ 不変であると主張する。
この主張と直前の議論から (a) が従う。

\medskip\noindent
各 $i$ について射 $\mathcal{X} \times_\mathcal{Y} V_i \to V_i$ は
準コンパクトであり、したがって $\mathcal{X} \times_\mathcal{Y} V_i$ は
準コンパクトである。ゆえにアフィンスキーム $W_i$ と全射滑らかな射
$W_i \to \mathcal{X} \times_\mathcal{Y} V_i$ を選べる。
$W = \coprod W_i$ は、$g$ との合成 $W \to V$ が準コンパクトとなる
滑らかかつ全射な射 $W \to \mathcal{X} \times_\mathcal{Y} V$ を備えた
スキームであることに注意する。$Z \to V$ を $W \to V$ の
スキーム論的像とする。次を参照せよ：
Morphisms, Section
\ref{morphisms-section-scheme-theoretic-image} および
Morphisms of Spaces, Section
\ref{spaces-morphisms-section-scheme-theoretic-image}.
Lemma \ref{stacks-morphisms-lemma-cover-upstairs} から、$Z \subset V$ は $g$ の
スキーム論的像である。$Z$ が $R$ 不変であることを示すため、
$$
\text{pr}_0^{-1}(Z), \text{pr}_1^{-1}(Z) \subset R = V \times_\mathcal{Y} V
$$
の双方が $\mathcal{X} \times_\mathcal{Y} R \to R$ のスキーム論的像であると
主張する。実際、まず Morphisms of Spaces, Lemma
\ref{spaces-morphisms-lemma-flat-base-change-scheme-theoretic-image}
を用いると、$\text{pr}_0^{-1}(Z)$ は合成
$$
W \times_{V, \text{pr}_0} R = W \times_\mathcal{Y} V \to
\mathcal{X} \times_\mathcal{Y} R \to R
$$
のスキーム論的像であることが分かる。ここで最初の射は全射かつ滑らかなので、
$\text{pr}_0^{-1}(Z)$ は $\mathcal{X} \times_\mathcal{Y} R \to R$ の
スキーム論的像である。同じ議論を $\text{pr}_1^{-1}(Z)$ に適用できる。
したがって $Z$ は $R$ 不変である。

\medskip\noindent
主張 (b) は (a) とまったく同じ方法で証明される。

\medskip\noindent
(c) の証明。$Z_i \subset V_i$ を
$\mathcal{X} \times_\mathcal{Y} V_i \to V_i$ のスキーム論的像とし、
$Z_{i, j} \subset V_{i, j}$ を
$\mathcal{X} \times_\mathcal{Y} V_{i, j} \to V_{i, j}$ の
スキーム論的像とする。$V_{i, j}$ における $Z_i$ の逆像が
$Z_{i, j}$ であることを示せば十分である。上で $Z_i$ は
$W_i \to V_i$ のスキーム論的像であることを見た。同じ理由により、
$Z_{i, j}$ は $W_i \times_{V_i} V_{i, j} \to V_{i, j}$ の
スキーム論的像である。したがって等式はスキームの場合
(Morphisms, Lemma
\ref{morphisms-lemma-flat-base-change-scheme-theoretic-image})
と、$V_{i, j} \to V_i$ が平坦であることから従う。
\end{proof}

\begin{lemma}
\label{stacks-morphisms-lemma-topology-scheme-theoretic-image}
$f : \mathcal{X} \to \mathcal{Y}$ を代数スタックの準コンパクトな射とする。
$\mathcal{Z} \subset \mathcal{Y}$ を $f$ のスキーム論的像とする。
このとき $|\mathcal{Z}|$ は $|f|$ の像の閉包である。
\end{lemma}

\begin{proof}
$z \in |\mathcal{Z}|$ を点とする。アフィンスキーム $V$、点 $v \in V$、
および $v$ を $z$ に写す滑らかな射 $V \to \mathcal{Y}$ を選ぶ。
このとき $\mathcal{X} \times_\mathcal{Y} V$ は準コンパクトな
代数スタックである。したがって、アフィンスキーム $W$ と全射滑らかな射
$W \to \mathcal{X} \times_\mathcal{Y} V$ を見いだせる。
Lemma \ref{stacks-morphisms-lemma-existence-plus-flat-base-change} により、
$\mathcal{X} \times_\mathcal{Y} V \to V$ のスキーム論的像は
$Z = \mathcal{Z} \times_\mathcal{Y} V$ である。
したがって、$|V|$ における $|\mathcal{Z}|$ の逆像は
Properties of Stacks, Lemma \ref{stacks-properties-lemma-points-cartesian}.
$|Z|$ である。Lemma \ref{stacks-morphisms-lemma-cover-upstairs} により、$Z$ は
$W \to V$ のスキーム論的像である。Morphisms of Spaces, Lemma
\ref{spaces-morphisms-lemma-quasi-compact-scheme-theoretic-image}
により、$|W| \to |Z|$ の像は稠密である。したがって
$|\mathcal{X} \times_\mathcal{Y} V| \to |Z|$ の像は稠密である。
$v \in Z$ に注意する。$|V| \to |\mathcal{Y}|$ は開なので、位相的議論から
$z$ は所望のとおり $|f|$ の像の閉包に属することが分かる。
\end{proof}

\begin{lemma}
\label{stacks-morphisms-lemma-scheme-theoretic-image-of-partial-section}
$f : \mathcal{X} \to \mathcal{Y}$ を、代数空間で表現可能かつ分離な
代数スタックの射とする。$\mathcal{V} \subset \mathcal{Y}$ を、
$\mathcal{V} \to \mathcal{Y}$ が準コンパクトとなる開部分スタックとする。
$s : \mathcal{V} \to \mathcal{X}$ を
$f \circ s = \text{id}_\mathcal{V}$ となる射とする。
$\mathcal{Y}'$ を $s$ のスキーム論的像とする。このとき
$\mathcal{Y}' \to \mathcal{Y}$ は $\mathcal{V}$ 上同型である。
\end{lemma}

\begin{proof}
Lemma \ref{stacks-morphisms-lemma-quasi-compact-permanence} により、射
$s : \mathcal{V} \to \mathcal{Y}$ は準コンパクトである。
したがって Lemma \ref{stacks-morphisms-lemma-existence-plus-flat-base-change} により、
$s$ のスキーム論的像 $\mathcal{Y}'$ の構成は平坦基底変換と可換である。
ゆえに補題を証明するには、$\mathcal{Y}$ は代数空間で表現可能であると
仮定してよく、代数空間の場合、すなわち
Morphisms of Spaces, Lemma
\ref{spaces-morphisms-lemma-scheme-theoretic-image-of-partial-section}
に帰着する。
\end{proof}









\section{付値判定法}
\label{stacks-morphisms-section-valuative}

\noindent
付値判定法を述べる際には注意が必要である。すなわち、その定式化では
可換図式について語る必要があるが、我々は $2$-圏で作業しているので、
$2$-射も正しく合成されることを確かめなければならない！

\begin{definition}
\label{stacks-morphisms-definition-fill-in-diagram}
$f : \mathcal{X} \to \mathcal{Y}$ を代数スタックの射とする。
$2$-可換な実線図式
\begin{equation}
\label{stacks-morphisms-equation-diagram}
\vcenter{
\xymatrix{
\Spec(K) \ar[r]_-x \ar[d]_j & \mathcal{X} \ar[d]^f \\
\Spec(A) \ar[r]^-y \ar@{..>}[ru] & \mathcal{Y}
}
}
\end{equation}
を考える。ここで $A$ は分数体 $K$ をもつ付値環である。
$$
\gamma : y \circ j \longrightarrow f \circ x
$$
を図式の $2$-可換性を示す $2$-射とする
（記法は Categories, Sections \ref{categories-section-formal-cat-cat} および
\ref{categories-section-2-categories} に従う）。
(\ref{stacks-morphisms-equation-diagram}) と $\gamma$ が与えられたとき、
{\it 点線矢印}とは、射 $a : \Spec(A) \to \mathcal{X}$ と $2$-矢印
$\alpha : a \circ j \to x$, $\beta : y \to f \circ a$ からなる三つ組
$ (a, \alpha, \beta) $ で、
$\gamma = (\text{id}_f \star \alpha) \circ (\beta \star \text{id}_j)$
を満たすもの、言い換えると
$$
\xymatrix{
& f \circ a \circ j \ar[rd]^{\text{id}_f \star \alpha} \\
y \circ j \ar[ru]^{\beta \star \text{id}_j} \ar[rr]^\gamma & &
f \circ x
}
$$
が可換となるものである。{\it 点線矢印の射}
$(a, \alpha, \beta) \to (a', \alpha', \beta')$ とは、
$\alpha = \alpha' \circ (\theta \star \text{id}_j)$ および
$\beta' = (\text{id}_f \star \theta) \circ \beta$ を満たす $2$-矢印
$\theta : a \to a'$ である。
\end{definition}

\noindent
上の定義は Categories, Definition \ref{categories-definition-dotted-arrows}
の特別な場合である。点線矢印の圏は一般に $\gamma$ に依存する。
$\mathcal{Y}$ が代数空間で表現可能ならば（または体上の対象の自己同型群が
自明ならば）、もちろん $\gamma$ は高々一つであり、三角形の可換性を
確かめる必要はない。より一般には
Lemma \ref{stacks-morphisms-lemma-cat-dotted-arrows-independent} がある。
三角形の可換性は基底変換との整合性の証明で重要である。
Lemma \ref{stacks-morphisms-lemma-cat-dotted-arrows-base-change} の証明を参照せよ。

\begin{lemma}
\label{stacks-morphisms-lemma-cat-dotted-arrows}
Definition \ref{stacks-morphisms-definition-fill-in-diagram} の状況で、点線矢印の圏は
亜群である。$\Delta_f$ が分離ならば、これは集合亜群である。
\end{lemma}

\begin{proof}
$2$-矢印は可逆なので、点線矢印の圏が亜群であることは明らかである。
点線矢印 $(a, \alpha, \beta)$ が与えられたとき、$(a, \alpha, \beta)$ の自己同型は二つの条件を
満たす $2$-射 $\theta : a \to a$ である。第一の条件
$\beta = (\text{id}_f \star \theta) \circ \beta$ は、$\theta$ が射
$(a, \theta) : \Spec(A) \to \mathcal{I}_{\mathcal{X}/\mathcal{Y}}$.
を定めることを意味する。第二の条件
$\alpha = \alpha \circ (\theta \star \text{id}_j)$ は、$(a, \theta)$ の
$\Spec(K)$ への制限が恒等射であることを含意する。図示すると
$$
\xymatrix{
\mathcal{I}_{\mathcal{X}/\mathcal{Y}} \ar[d] & &
\Spec(K) \ar[d]^j \ar[ll]_{(a \circ j, \text{id})} \\
\mathcal{X} & & \Spec(A) \ar[ll]_a \ar[llu]_{(a, \theta)}
}
$$
となる。言い換えると、$G \to \Spec(A)$ を相対慣性を $a$ で引き戻して
得られる群代数空間とすると、$\theta$ は $G(K)$ における像が自明な点
$\theta \in G(A)$ を定める。恒等射 $e : \Spec(A) \to G$ が閉埋入ならば、
これは $\theta$ が恒等射である場合にしか起こりえない。
Lemma \ref{stacks-morphisms-lemma-diagonal-diagonal} を見ると、所望の結果を得る。
\end{proof}

\begin{lemma}
\label{stacks-morphisms-lemma-cat-dotted-arrows-independent}
Definition \ref{stacks-morphisms-definition-fill-in-diagram} において、
$\mathcal{I}_\mathcal{Y} \to \mathcal{Y}$ が固有であると仮定する
（例えば $\mathcal{Y}$ が分離であるか、または $\mathcal{Y}$ が
代数空間上分離である場合）。このとき、点線矢印の圏は
（非標準的な圏同値を除いて）$\gamma$ の選択に依存しない。また、
点線矢印が存在すること
（あるものに対し、したがって同値にすべての $\gamma$ に対し）は、
図式
$$
\xymatrix{
\Spec(K) \ar[r]_-x \ar[d]_j & \mathcal{X} \ar[d]^f \\
\Spec(A) \ar[r]^-y \ar[ru]_a & \mathcal{Y}
}
$$
が存在し、その二つの三角形が $2$-可換であること
（$2$-射が正しく合成されるかは確認しない）と同値である。
\end{lemma}

\begin{proof}
$\gamma, \gamma' : y \circ j \longrightarrow f \circ x$ を二つの
$2$-射とする。すると $\gamma^{-1} \circ \gamma'$ は
$\Spec(K)$ 上の $y$ の自己同型である。
したがって、$\mathit{Isom}_\mathcal{Y}(y, y) \to \Spec(A)$ が
固有ならば、固有性の付値判定法
（Morphisms of Spaces, Lemma \ref{spaces-morphisms-lemma-characterize-proper}）により、
$\Spec(K)$ への制限が $\gamma^{-1} \circ \gamma'$ である
$\delta : y \to y$ を見つけられる。
このとき $\delta$ を用い、$(a, \alpha, \beta)$ を
$(a, \alpha, \beta \circ \delta)$ に送ることで、$\gamma$ に対する
点線矢印の圏から $\gamma'$ に対する点線矢印の圏への同値を
定義できる。最後の主張は明らかである。
\end{proof}

\begin{lemma}
\label{stacks-morphisms-lemma-cat-dotted-arrows-base-change}
$2$-可換図式
$$
\xymatrix{
\Spec(K) \ar[r]_-{x'} \ar[d]_j &
\mathcal{X}' \ar[d]^p \ar[r]_q &
\mathcal{X} \ar[d]^f \\
\Spec(A) \ar[r]^-{y'} &
\mathcal{Y}' \ar[r]^g &
\mathcal{Y}
}
$$
が与えられ、右の正方形が $2$-Cartesian であると仮定する。$2$-矢印
$\gamma' : y' \circ j \to p \circ x'$ を選ぶ。
$x = q \circ x'$, $y = g \circ y'$ とおき、
$\gamma : y \circ j \to f \circ x$ を、$\gamma'$ と右の正方形の
$2$-可換性に暗黙に含まれる $2$-矢印との合成とする。
このとき、左の正方形と $\gamma'$ に対する点線矢印の圏は、
外側の長方形と $\gamma$ に対する点線矢印の圏と同値である。
\end{lemma}

\begin{proof}
（点線矢印の圏の代わりに、Lemma
\ref{stacks-morphisms-lemma-cat-dotted-arrows-independent} におけるような二つの
$2$-可換三角形を生じる射の同型類の集合を考える場合、
この補題の類似をどのように証明すればよいか我々は知らない。
実際、その類似は誤りである可能性も十分にある。）
第一の証明：この補題は Categories, Lemma
\ref{categories-lemma-cat-dotted-arrows-base-change} の特別な場合である。
第二の証明：Categories, Lemma
\ref{categories-lemma-2-product-categories-over-C} に述べられているように、
$\mathcal{X}'$ を $2$-ファイバー積
$\mathcal{Y}' \times_\mathcal{Y} \mathcal{X}$ で置き換えてよい。
すると対象 $x'$ は三つ組 $(y' \circ j, x, \gamma)$ となる。
外側の長方形に対する点線矢印 $(a, \alpha, \beta)$ から左の
正方形に対する点線矢印 $(a', \alpha', \beta')$ へ、
$a' = (y', a, \beta)$, $\alpha' = (\text{id}_{y' \circ j}, \alpha)$,
$\beta' = \text{id}_{y'}$ とおくことにより移れる。詳細は省略する。
\end{proof}

\begin{lemma}
\label{stacks-morphisms-lemma-cat-dotted-arrows-composition}
$2$-可換図式
$$
\xymatrix{
\Spec(K) \ar[r]_-x \ar[dd]_j & \mathcal{X} \ar[d]^f \\
& \mathcal{Y} \ar[d]^g \\
\Spec(A) \ar[r]^-z & \mathcal{Z}
}
$$
が与えられたと仮定する。$2$-矢印
$\gamma : z \circ j \to g \circ f \circ x$ を選ぶ。
$\mathcal{C}$ を外側の長方形と $\gamma$ に対する点線矢印の圏とする。
$\mathcal{C}'$ を正方形
$$
\xymatrix{
\Spec(K) \ar[r]_-{f \circ x} \ar[d]_j & \mathcal{Y} \ar[d]^g \\
\Spec(A) \ar[r]^-z & \mathcal{Z}
}
$$
と $\gamma$ に対する点線矢印の圏とする。このとき $\mathcal{C}$ は、
次の性質をもつ圏 $\mathcal{C}''$ と同値である：関手
$\mathcal{C}'' \to \mathcal{C}'$ が存在し、これによって $\mathcal{C}''$ は
$\mathcal{C}'$ 上の亜群をファイバーとする圏となり、そのファイバー圏は
形
$$
\xymatrix{
\Spec(K) \ar[r]_-x \ar[d]_j & \mathcal{X} \ar[d]^f \\
\Spec(A) \ar[r]^-y & \mathcal{Y}
}
$$
のある正方形と、$y \circ j \to f \circ x$ のある選択に対する
点線矢印の圏である。
\end{lemma}

\begin{proof}
この補題は Categories, Lemma
\ref{categories-lemma-cat-dotted-arrows-composition} の特別な場合である。
\end{proof}

\begin{definition}
\label{stacks-morphisms-definition-uniqueness}
$f : \mathcal{X} \to \mathcal{Y}$ を代数スタックの射とする。
すべての図式 (\ref{stacks-morphisms-equation-diagram}) と、Definition
\ref{stacks-morphisms-definition-fill-in-diagram} におけるような $\gamma$ に対し、
点線矢印の圏が空であるか、
または同型類をちょうど一つもつ集合亜群であるとき、
$f$ は{\it 付値判定法の一意性部分}を満たすという。
\end{definition}

\begin{lemma}
\label{stacks-morphisms-lemma-base-change-uniqueness}
付値判定法の一意性部分を満たす代数スタックの射を、
代数スタックの任意の射によって基底変換して得られる射も、
付値判定法の一意性部分を満たす代数スタックの射である。
\end{lemma}

\begin{proof}
Lemma \ref{stacks-morphisms-lemma-cat-dotted-arrows-base-change} と定義から従う。
\end{proof}

\begin{lemma}
\label{stacks-morphisms-lemma-composition-uniqueness}
付値判定法の一意性部分を満たす代数スタックの射の合成も、
付値判定法の一意性部分を満たす代数スタックの射である。
\end{lemma}

\begin{proof}
Lemma \ref{stacks-morphisms-lemma-cat-dotted-arrows-composition} と定義から従う。
\end{proof}

\begin{lemma}
\label{stacks-morphisms-lemma-uniqueness-representable}
$f : \mathcal{X} \to \mathcal{Y}$ を、代数空間により表現可能な
代数スタックの射とする。このとき次は同値である。
\begin{enumerate}
\item $f$ は付値判定法の一意性部分を満たす。
\item すべてのスキーム $T$ と射 $T \to \mathcal{Y}$ に対し、射
$\mathcal{X} \times_\mathcal{Y} T \to T$ は代数空間の射として
付値判定法の一意性部分を満たす。
\end{enumerate}
\end{lemma}

\begin{proof}
Lemma \ref{stacks-morphisms-lemma-cat-dotted-arrows-base-change} と定義から従う。
\end{proof}

\begin{definition}
\label{stacks-morphisms-definition-existence}
$f : \mathcal{X} \to \mathcal{Y}$ を代数スタックの射とする。
すべての図式 (\ref{stacks-morphisms-equation-diagram}) と、Definition
\ref{stacks-morphisms-definition-fill-in-diagram} におけるような $\gamma$ に対し、
体の拡大 $K'/K$ と、$A$ を支配する付値環 $A' \subset K'$ が存在し、
図式
$$
\xymatrix{
\Spec(K') \ar[r] \ar@/^2em/[rr]_{x'} \ar[d]_{j'} &
\Spec(K) \ar[d]_j \ar[r]_-x &
\mathcal{X} \ar[d]^f \\
\Spec(A') \ar[r] \ar@/_2em/[rr]^{y'} &
\Spec(A) \ar[r]^-y &
\mathcal{Y}
}
$$
の外側の長方形と、誘導される $2$-矢印
$\gamma' : y' \circ j' \to f \circ x'$ に対する点線矢印の圏が
空でないとき、$f$ は{\it 付値判定法の存在性部分}を満たすという。
\end{definition}

\noindent
代数空間の射の場合には、付値判定法の正しい概念を得るために、
分数体の拡大を許す必要があることをすでに見た。Morphisms of Spaces,
Example \ref{spaces-morphisms-example-finite-separable-needed} を参照せよ。
それでも、分離代数空間の間の射については、このような拡大は必要ない
（Morphisms of Spaces, Lemma \ref{spaces-morphisms-lemma-usual-enough}）。
しかし、代数スタックの間の射については、$\mathcal{X}$ と $\mathcal{Y}$ が
ともに分離であっても拡大が必要となることがある。
例えば、代数スタックの射
$$
[\Spec(\mathbf{C})/G] \to \Spec(\mathbf{C})
$$
を考えよ。これは基底スキーム $\Spec(\mathbf{C})$ 上の射で、$G$ は位数 $2$ の群である。
始域と終域はどちらも分離代数スタックであり、この射は固有である。
したがって、この射は付値判定法の一意性部分と存在性部分を満たす
（Lemma \ref{stacks-morphisms-lemma-criterion-proper} を参照せよ）。
しかし一方で、図式
$$
\xymatrix{
\Spec(K) \ar[r] \ar[d] & [\Spec(\mathbf{C})/G] \ar[d] \\
\Spec(A) \ar[r] & \Spec(\mathbf{C})
}
$$
が存在し、$A = \mathbf{C}[[t]]$ および $K = \mathbf{C}((t))$ とすると、
点線矢印は存在しない。実際、上の水平矢印は、$K = \mathbf{C}((t))$ の
スペクトル上の $G$-トーサーによって与えられる。厳密 Hensel 局所環
$A = \mathbf{C}[[t]]$ 上の任意の $G$-トーサーは自明であるから、
もし点線矢印が常に存在するなら、$K$ 上のすべての $G$-トーサーが
自明であることが分かる。これは真でない。なぜなら $G = \{+1, -1\}$ であり、
Kummer 理論により $K$ 上の $G$-トーサーは非自明な
$K^*/(K^*)^2$ で分類されるからである。

\begin{lemma}
\label{stacks-morphisms-lemma-base-change-existence}
付値判定法の存在性部分を満たす代数スタックの射を、
代数スタックの任意の射によって基底変換して得られる射も、
付値判定法の存在性部分を満たす代数スタックの射である。
\end{lemma}

\begin{proof}
Lemma \ref{stacks-morphisms-lemma-cat-dotted-arrows-base-change} と定義から従う。
\end{proof}

\begin{lemma}
\label{stacks-morphisms-lemma-composition-existence}
付値判定法の存在性部分を満たす代数スタックの射の合成も、
付値判定法の存在性部分を満たす代数スタックの射である。
\end{lemma}

\begin{proof}
Lemma \ref{stacks-morphisms-lemma-cat-dotted-arrows-composition} と定義から従う。
\end{proof}

\begin{lemma}
\label{stacks-morphisms-lemma-existence-representable}
$f : \mathcal{X} \to \mathcal{Y}$ を、代数空間により表現可能な
代数スタックの射とする。このとき次は同値である。
\begin{enumerate}
\item $f$ は付値判定法の存在性部分を満たす。
\item すべてのスキーム $T$ と射 $T \to \mathcal{Y}$ に対し、射
$\mathcal{X} \times_\mathcal{Y} T \to T$ は代数空間の射として
付値判定法の存在性部分を満たす。
\end{enumerate}
\end{lemma}

\begin{proof}
Lemma \ref{stacks-morphisms-lemma-cat-dotted-arrows-base-change} と定義から従う。
\end{proof}

\begin{lemma}
\label{stacks-morphisms-lemma-closed-immersion-valuative-criteria}
代数スタックの閉埋入は、付値判定法の存在性部分と
一意性部分の両方を満たす。
\end{lemma}

\begin{proof}
省略する。ヒント：Lemmas \ref{stacks-morphisms-lemma-uniqueness-representable} および
\ref{stacks-morphisms-lemma-existence-representable} により、スキームの閉埋入の場合に帰着させよ。
\end{proof}





\section{第二対角射の付値判定法}
\label{stacks-morphisms-section-valuative-second}

\noindent
逆向きの主張は Lemma \ref{stacks-morphisms-lemma-cat-dotted-arrows} ですでに証明した。
判定法そのものは次のとおりである。

\begin{lemma}
\label{stacks-morphisms-lemma-setoids-and-diagonal}
$f : \mathcal{X} \to \mathcal{Y}$ を代数スタックの射とする。
$\Delta_f$ が準分離であり、すべての図式 (\ref{stacks-morphisms-equation-diagram}) と、
Definition \ref{stacks-morphisms-definition-fill-in-diagram} における $\gamma$ の選択に対して
点線矢印の圏が集合亜群であるならば、$\Delta_f$ は分離である。
\end{lemma}

\begin{proof}
以下で詳細な証明を書くが、読者には Lemma
\ref{stacks-morphisms-lemma-cat-dotted-arrows} の証明で与えた議論を読み、
それを手がかりに自分で証明を見つけることを強く勧めたい。

\medskip\noindent
$\Delta_f$ が準分離であり、すべての図式 (\ref{stacks-morphisms-equation-diagram}) と、
Definition \ref{stacks-morphisms-definition-fill-in-diagram} における $\gamma$ の選択に対して
点線矢印の圏が集合亜群であると仮定する。
Lemma \ref{stacks-morphisms-lemma-diagonal-diagonal} により、
$e : \mathcal{X} \to \mathcal{I}_{\mathcal{X}/\mathcal{Y}}$ が閉埋入であることを
示せば十分である。実際、Lemma
\ref{stacks-morphisms-lemma-first-diagonal-separated-second-diagonal-closed} により、
$e = \Delta_{f, 2}$ が普遍閉であることを示せば十分である。
これらの補題のいずれからも、$\Delta_f$ が準分離であるという仮定により、
$e = \Delta_{f, 2}$ が準コンパクトであることが分かる。

\medskip\noindent
この段落では、$e$ が付値判定法の存在性部分を満たすことを示す。
$2$-可換な実線図式
$$
\xymatrix{
\Spec(K) \ar[r]_x \ar[d]_j & \mathcal{X} \ar[d]^e \\
\Spec(A) \ar[r]^{(a, \theta)} & \mathcal{I}_{\mathcal{X}/\mathcal{Y}}
}
$$
を考え、図式の $2$-可換性を与える任意の $2$-射
$\alpha : (a, \theta) \circ j \to e \circ x$ をとる
（Definition \ref{stacks-morphisms-definition-fill-in-diagram} で用いた文字 $\gamma$ の代わりに $\alpha$ を用いる）。
$f \circ \theta = \text{id}$ であることに注意せよ。これを以下で用いる。
$e \circ x = (x, \text{id}_x)$ および
$(a, \theta) \circ j = (a \circ j, \theta \star \text{id}_j)$ である。
したがって $\alpha$ は、$\theta \star \text{id}_j$ および $\text{id}_x$ と
整合する $2$-矢印 $\alpha : a \circ j \to x$ である。
$y = f \circ x$ および $\beta = \text{id}_{f \circ a}$ とおく。
Lemma \ref{stacks-morphisms-lemma-cat-dotted-arrows} の証明で与えた議論を逆に読むと、
$\theta$ は、次のような点線矢印 $(a, \alpha, \beta)$ の自己同型である。
$$
\gamma : y \circ j \to f \circ x
\quad\text{equal to}\quad
\text{id}_f \star \alpha : f \circ a \circ j \to f \circ x
$$
一方、$\text{id}_a$ も自己同型であるから、$f$ に対する仮定より
$\theta = \text{id}_a$ を得る。そこで、上の表示図式に対する点線矢印として、
射 $a : \Spec(A) \to \mathcal{X}$ と、$\alpha$ で与えられる $2$-射
$(a, \text{id}_a) \circ j \to (x, \text{id}_x)$ および $\text{id}_a$ で与えられる
$(a, \theta) \to e \circ a$ をとれる。

\medskip\noindent
Lemma \ref{stacks-morphisms-lemma-base-change-existence} により、$e$ の任意の基底変換は
付値判定法の存在性部分を満たす。$e$ は代数空間により
表現可能であるから、全射かつ滑らかな射
$I \to \mathcal{I}_{\mathcal{X}/\mathcal{Y}}$ で、$I$ が代数空間であるものによる基底変換の後で、
$e$ が普遍閉であることを示せば十分である
（Properties of Stacks, Section
\ref{stacks-properties-section-properties-morphisms} を参照せよ）。
この基底変換 $e' : X' \to I'$ は代数空間の準コンパクトな射で、
付値判定法の存在性部分を満たす。したがって、Morphisms of Spaces,
Lemma \ref{spaces-morphisms-lemma-quasi-compact-existence-universally-closed} により
普遍閉である。
\end{proof}






\section{対角射の付値判定法}
\label{stacks-morphisms-section-valuative-diagonal}

\noindent
結果は Lemma \ref{stacks-morphisms-lemma-uniqueness-and-diagonal} である。
まず、形式的な補助補題を述べて証明する。

\begin{lemma}
\label{stacks-morphisms-lemma-helper-diagonal}
$f : \mathcal{X} \to \mathcal{Y}$ を代数スタックの射とする。
$2$-可換な実線図式
$$
\xymatrix{
\Spec(K) \ar[rr]_-x \ar[d]_j & &
\mathcal{X} \ar[d]^{\Delta_f} \\
\Spec(A) \ar[rr]^{(a_1, a_2, \varphi)} \ar@{..>}[rru] & &
\mathcal{X} \times_\mathcal{Y} \mathcal{X}
}
$$
を考える。ここで $A$ は分数体 $K$ をもつ付値環である。
$\gamma : (a_1, a_2, \varphi) \circ j \longrightarrow \Delta_f \circ x$ を、
図式の $2$-可換性を与える $2$-射とする。このとき次が成り立つ。
\begin{enumerate}
\item $\gamma = (\alpha_1, \alpha_2)$ と書き、
$\alpha_i : a_i \circ j \to x$ とすると、図式
$$
\xymatrix{
\Spec(K) \ar[r]_-x \ar[d]_j & \mathcal{X} \ar[d]^f \\
\Spec(A) \ar[r]^-{f \circ a_1} \ar@{..>}[ru] & \mathcal{Y}
}
$$
において二つの点線矢印 $(a_1, \alpha_1, \text{id})$ および
$(a_2, \alpha_2, \varphi)$ を得る。
\item 元の図式と $\gamma$ に対する点線矢印の圏は集合亜群であり、
その対象の同型類の集合は、点線矢印の圏における射の集合
$(a_1, \alpha_1, \text{id}) \to (a_2, \alpha_2, \varphi)$ に等しい。
\end{enumerate}
\end{lemma}

\begin{proof}
$\Delta_f$ は代数空間により表現可能である（したがって $\Delta_f$ の
対角射は分離である）から、Lemma \ref{stacks-morphisms-lemma-cat-dotted-arrows} により、
補題の最初の可換図式に対する点線矢印の圏は集合亜群である。
補題の他の主張は、ここでは省略する $2$-図式的計算から従う。
\end{proof}

\begin{lemma}
\label{stacks-morphisms-lemma-uniqueness-and-diagonal}
$f : \mathcal{X} \to \mathcal{Y}$ を代数スタックの射とする。
$f$ が準分離であると仮定する。$f$ が付値判定法の
一意性部分を満たすならば、$f$ は分離である。
\end{lemma}

\begin{proof}
$f$ に対する仮定は、$\Delta_f$ が準コンパクトかつ準分離であることを
意味する（Definition \ref{stacks-morphisms-definition-separated}）。$\Delta_f$ が固有であることを
示さねばならない。Lemma \ref{stacks-morphisms-lemma-setoids-and-diagonal} により $\Delta_f$ は分離である。
Lemma \ref{stacks-morphisms-lemma-properties-diagonal} により $\Delta_f$ が局所有限型であることも分かる。
証明を完了するには、$\Delta_f$ が普遍閉であることを示さねばならない。
形式的な議論（Lemma \ref{stacks-morphisms-lemma-helper-diagonal} を参照せよ）により、
付値判定法の一意性部分から、次のような任意の実線図式において
点線矢印が存在することが従う。
$$
\xymatrix{
\Spec(K) \ar[d] \ar[r] & \mathcal{X} \ar[d]^{\Delta_f} \\
\Spec(A) \ar[r] \ar@{..>}[ru] & \mathcal{X} \times_\mathcal{Y} \mathcal{X}
}
$$
この性質は任意の基底変換によって保たれるから、
$\mathcal{X} \times_\mathcal{Y} \mathcal{X}$ へ写る代数空間による $\Delta_f$ の任意の基底変換は、
付値判定法の存在性部分を満たす。したがって、代数空間の射の
付値判定法により、それは普遍閉である。Morphisms of Spaces,
Lemma \ref{spaces-morphisms-lemma-quasi-compact-existence-universally-closed} を参照せよ。
\end{proof}

\noindent
次が逆である。

\begin{lemma}
\label{stacks-morphisms-lemma-converse-uniqueness-and-diagonal}
$f : \mathcal{X} \to \mathcal{Y}$ を代数スタックの射とする。
$f$ が分離であるならば、$f$ は付値判定法の一意性部分を満たす。
\end{lemma}

\begin{proof}
$f$ は分離であるから、Lemma \ref{stacks-morphisms-lemma-cat-dotted-arrows} により、
すべての点線矢印の圏が集合亜群であることが分かる。図式
$$
\xymatrix{
\Spec(K) \ar[r]_-x \ar[d]_j & \mathcal{X} \ar[d]^f \\
\Spec(A) \ar[r]^-y \ar@{..>}[ru] & \mathcal{Y}
}
$$
と、Definition \ref{stacks-morphisms-definition-fill-in-diagram} におけるような $2$-射
$\gamma : y \circ j \to f \circ x$ を考える。点線矢印の圏の二つの対象
$(a, \alpha, \beta)$ および $(a', \beta', \alpha')$ を考える。
証明を完了するには、これらの対象が同型であることを示さねばならない。
同型
$$
f \circ a \xrightarrow{\beta^{-1}} y \xrightarrow{\beta'} f \circ a'
$$
は、$(a, a', \beta' \circ \beta^{-1})$ が射
$\Spec(A) \to \mathcal{X} \times_\mathcal{Y} \mathcal{X}$ であることを意味する。
一方、$\alpha$ と $\alpha'$ は $2$-矢印
$$
(a, a', \beta' \circ \beta^{-1}) \circ j =
(a \circ j, a' \circ j,
(\beta' \star \text{id}_j) \circ (\beta \star \text{id}_j)^{-1})
\xrightarrow{(\alpha, \alpha')} (x, x, \text{id}) = \Delta_f \circ x
$$
を定める。ここで、$(a, \alpha, \beta)$ と $(a', \alpha', \beta')$ の両方が
$\gamma$ に関する点線矢印であることを用いた。可換図式
$$
\xymatrix{
\Spec(K) \ar[d]_j \ar[rr]_x & & \mathcal{X} \ar[d]^{\Delta_f} \\
\Spec(A) \ar[rr]^{(a, a', \beta' \circ \beta^{-1})} & &
\mathcal{X} \times_\mathcal{Y} \mathcal{X}
}
$$
を得る。その $2$-可換性は $(\alpha, \alpha')$ で与えられる。
さて、$\Delta_f$ は代数空間により表現可能であり
（Lemma \ref{stacks-morphisms-lemma-properties-diagonal}）、$f$ が分離であるから固有である。
したがって、Lemma \ref{stacks-morphisms-lemma-existence-representable} と代数空間の射に対する
固有性の付値判定法（Morphisms of Spaces, Lemma
\ref{spaces-morphisms-lemma-characterize-proper}）により、点線矢印が存在することが分かる。
その構成を展開すると、$(a, \alpha, \beta)$ と $(a', \alpha', \beta')$ が
望まれたとおり点線矢印の圏で同型であることを意味する。
\end{proof}






\section{普遍閉性の付値判定法}
\label{stacks-morphisms-section-valutive-criterion}

\noindent
次の主張がある。

\begin{lemma}
\label{stacks-morphisms-lemma-quasi-compact-existence-universally-closed}
$f : \mathcal{X} \to \mathcal{Y}$ を代数スタックの射とする。次を仮定する。
\begin{enumerate}
\item $f$ は準コンパクトである。
\item $f$ は付値判定法の存在性部分を満たす。
\end{enumerate}
このとき $f$ は普遍閉である。
\end{lemma}

\begin{proof}
Lemmas \ref{stacks-morphisms-lemma-base-change-quasi-compact} および
\ref{stacks-morphisms-lemma-base-change-existence} により、性質 (1), (2) は任意の基底変換で保たれる。
Lemma \ref{stacks-morphisms-lemma-universally-closed-local} により、$\mathcal{Y}$ へ写る任意のアフィンスキーム $T$
に対して $|T \times_\mathcal{Y} \mathcal{X}| \to |T|$ が閉写像であることを
示せばよい。したがって次を示せば十分である：
$f : \mathcal{X} \to Y$ が代数スタックからアフィンスキームへの準コンパクトな射で、
付値判定法の存在性部分を満たすならば、$|f|$ は閉写像である。
$T \subset |\mathcal{X}|$ を閉部分集合とする。証明を完了するには、$f(T)$ が閉であることを示さねばならない。

\medskip\noindent
$\mathcal{Z} \subset \mathcal{X}$ を、$T$ 上の誘導される既約代数スタック構造とする
（Properties of Stacks, Definition
\ref{stacks-properties-definition-reduced-induced-stack}）。すると
$i : \mathcal{Z} \to \mathcal{X}$ は閉埋入であり、$|\mathcal{Z}| \to |Y|$ の像が閉であることを
示さねばならない。閉埋入は準コンパクトであり
（Lemma \ref{stacks-morphisms-lemma-closed-immersion-quasi-compact}）、付値判定法の存在性部分を満たす
（Lemma \ref{stacks-morphisms-lemma-closed-immersion-valuative-criteria}）。また、準コンパクトな射の合成は
準コンパクトであり（Lemma \ref{stacks-morphisms-lemma-composition-quasi-compact}）、
付値判定法の存在性部分を満たすという性質も合成で保たれる
（Lemma \ref{stacks-morphisms-lemma-composition-existence}）。したがって、次を示せば十分である：
$f : \mathcal{X} \to Y$ が代数スタックからアフィンスキームへの準コンパクトな射で、
付値判定法の存在性部分を満たすならば、$|f|(|\mathcal{X}|)$ は閉である。

\medskip\noindent
$\mathcal{X}$ は準コンパクトである（アフィンスキーム $Y$ 上で準コンパクトであるから）ので、
アフィンスキーム $U$ と全射滑らかな射 $U \to \mathcal{X}$ を選べる
（Properties of Stacks, Lemma
\ref{stacks-properties-lemma-quasi-compact-stack}）。
$y \in Y$ が $U \to Y$ の像の閉包に属すると仮定する
（言い換えれば、$|f|$ の像の閉包に属すると仮定する）。
すると Morphisms, Lemma
\ref{morphisms-lemma-reach-points-scheme-theoretic-image} により、分数体 $K$ をもつ
付値環 $A$ と可換図式
$$
\xymatrix{
\Spec(K) \ar[r] \ar[d] & U \ar[d] \\
\Spec(A) \ar[r] & Y
}
$$
で、$\Spec(A)$ の閉点が $y$ に写るものを見つけられる。
仮定により、拡大 $K'/K$ と、$A$ を支配する付値環 $A' \subset K'$ および
次の図式の点線矢印を得る。
$$
\xymatrix{
\Spec(K') \ar[r] \ar[d] &
\Spec(K) \ar[r] \ar[d] &
U \ar[d] \ar[r] &
\mathcal{X} \ar[d]^f \\
\Spec(A') \ar[r] \ar@{..>}[rrru] &
\Spec(A) \ar[r] &
Y  \ar@{=}[r] &
Y
}
$$
したがって $y$ は $|f|$ の像に属し、証明が完了する。
\end{proof}

\noindent
次が逆である。

\begin{lemma}
\label{stacks-morphisms-lemma-converse-existence-universally-closed}
$f : \mathcal{X} \to \mathcal{Y}$ を代数スタックの射とする。次を仮定する。
\begin{enumerate}
\item $f$ は準分離である。
\item $f$ は普遍閉である。
\end{enumerate}
このとき $f$ は付値判定法の存在性部分を満たす。
\end{lemma}

\begin{proof}
実線図式
$$
\xymatrix{
\Spec(K) \ar[r]_-x \ar[d]_j & \mathcal{X} \ar[d]^f \\
\Spec(A) \ar[r]^-y \ar@{..>}[ru] & \mathcal{Y}
}
$$
を考える。ここで $A$ は分数体 $K$ をもつ付値環であり、
$\gamma : y \circ j \longrightarrow f \circ x$ は Definition
\ref{stacks-morphisms-definition-fill-in-diagram} におけるようなものである。
Lemma \ref{stacks-morphisms-lemma-cat-dotted-arrows-base-change} により、
（$K$ を拡大で、$A$ をそれを支配する付値環で置き換えることを必要に応じて許して）
点線矢印を見つけるには、$\mathcal{Y}$ を $\Spec(A)$ で、$\mathcal{X}$ を
$\Spec(A) \times_\mathcal{Y} \mathcal{X}$ で置き換えてよい。
もちろんここで、準分離であることと普遍閉であることが基底変換で
保たれることを用いている。したがって次の段落で論じる場合に帰着する。

\medskip\noindent
実線図式
$$
\xymatrix{
\Spec(K) \ar[r]_-x \ar[d]_j & \mathcal{X} \ar[d]^f \\
\Spec(A) \ar@{=}[r] \ar@{..>}[ru] & \Spec(A)
}
$$
を考える。ここで $A$ は、Definition \ref{stacks-morphisms-definition-fill-in-diagram} におけるような、
分数体 $K$ をもつ付値環である。Lemma \ref{stacks-morphisms-lemma-quasi-compact-permanence} と
$f$ が準分離であることにより、射 $x$ は準コンパクトである。
$f$ は普遍閉であるから、特に $|f|(\overline{\{x\}})$ は $\Spec(A)$ の閉集合である。
この像は $\Spec(A)$ の一般点を含むから、$x$ の閉包に含まれ、
$\Spec(A)$ の閉点に写る点 $x' \in |\mathcal{X}|$ が存在する。
Lemma \ref{stacks-morphisms-lemma-reach-points-scheme-theoretic-image} により、可換図式
$$
\xymatrix{
\Spec(K') \ar[r] \ar[d] & \Spec(K) \ar[d] \\
\Spec(A') \ar[r] & \mathcal{X}
}
$$
で、$\Spec(A')$ の閉点が $x' \in |\mathcal{X}|$ に写るものを見つけられる。
したがって $\Spec(A') \to \Spec(A)$ は閉点を閉点に写す。
すなわち $A'$ は $A$ を支配し、証明が完了する。
\end{proof}





\section{固有性の付値判定法}
\label{stacks-morphisms-section-valutive-criterion-properness}

\noindent
主張は次のとおりである。

\begin{lemma}
\label{stacks-morphisms-lemma-criterion-proper}
$f : \mathcal{X} \to \mathcal{Y}$ を代数スタックの射とする。
$f$ が有限型かつ準分離であると仮定する。このとき次は同値である。
\begin{enumerate}
\item $f$ は固有である。
\item $f$ は付値判定法の一意性部分と存在性部分の両方を満たす。
\end{enumerate}
\end{lemma}

\begin{proof}
固有射とは、分離、有限型かつ普遍閉な射にほかならない。
したがって、この補題は Lemmas
\ref{stacks-morphisms-lemma-uniqueness-and-diagonal},
\ref{stacks-morphisms-lemma-converse-uniqueness-and-diagonal},
\ref{stacks-morphisms-lemma-quasi-compact-existence-universally-closed}, および
\ref{stacks-morphisms-lemma-converse-existence-universally-closed} から従う。
\end{proof}







\section{局所完全交叉射}
\label{stacks-morphisms-section-lci}

\noindent
代数空間の射の性質「局所完全交叉射である」は、
始域と終域上滑らか局所的である。Descent on Spaces, Lemma
\ref{spaces-descent-lemma-smooth-local-source-target} および More on Morphisms of Spaces,
Lemmas \ref{spaces-more-morphisms-lemma-descending-property-lci} および
\ref{spaces-more-morphisms-lemma-lci-syntomic-local-source} を参照せよ。
上の Lemma \ref{stacks-morphisms-lemma-local-source-target} により、代数空間の射が
局所完全交叉射であるということを次のように定義できる。
この定義は、始域と終域の両方が代数空間であるとき、
More on Morphisms of Spaces, Section
\ref{spaces-more-morphisms-section-lci} ですでに定義されている概念と一致する。

\begin{definition}
\label{stacks-morphisms-definition-lci}
$f : \mathcal{X} \to \mathcal{Y}$ を代数スタックの射とする。
Lemma \ref{stacks-morphisms-lemma-local-source-target} の同値な条件が
$\mathcal{P} = \text{local complete intersection}$ に対して成り立つとき、
$f$ を{\it 局所完全交叉射}、または{\it Koszul} であるという。
\end{definition}

\begin{lemma}
\label{stacks-morphisms-lemma-composition-lci}
局所完全交叉射の合成は局所完全交叉射である。
\end{lemma}

\begin{proof}
Remark \ref{stacks-morphisms-remark-composition} と More on Morphisms of Spaces, Lemma
\ref{spaces-more-morphisms-lemma-composition-lci} を組み合わせよ。
\end{proof}

\begin{lemma}
\label{stacks-morphisms-lemma-flat-base-change-lci}
局所完全交叉射の平坦基底変換は局所完全交叉射である。
\end{lemma}

\begin{proof}
省略する。ヒント：Remark \ref{stacks-morphisms-remark-base-change} とまったく同様に議論せよ
（ただし、平坦な $\mathcal{Y}' \to \mathcal{Y}$ に対してのみ）。その際、
More on Morphisms of Spaces, Lemma
\ref{spaces-more-morphisms-lemma-flat-base-change-lci} を用いよ。
\end{proof}

\begin{lemma}
\label{stacks-morphisms-lemma-lci-permanence}
代数スタックの射の可換図式
$$
\xymatrix{
\mathcal{X} \ar[rr]_f \ar[rd] & & \mathcal{Y} \ar[ld] \\
& \mathcal{Z}
}
$$
を考える。$\mathcal{Y} \to \mathcal{Z}$ が滑らかで、
$\mathcal{X} \to \mathcal{Z}$ が局所完全交叉射であると仮定する。
このとき $f : \mathcal{X} \to \mathcal{Y}$ は局所完全交叉射である。
\end{lemma}

\begin{proof}
スキーム $W$ と全射滑らかな射 $W \to \mathcal{Z}$ を選ぶ。
スキーム $V$ と全射滑らかな射
$V \to W \times_\mathcal{Z} \mathcal{Y}$ を選ぶ。
スキーム $U$ と全射滑らかな射
$U \to V \times_\mathcal{Y} \mathcal{X}$ を選ぶ。
すると $U \to W$ はスキームの局所完全交叉射であり、
$V \to W$ はスキームの滑らかな射である。スキームに対する結果
（More on Morphisms, Lemma \ref{more-morphisms-lemma-lci-permanence}）により、
$U \to V$ は局所完全交叉射である。定義により、これは $f$ が
局所完全交叉射であることを意味する。
\end{proof}





\section{安定化群を保存する射}
\label{stacks-morphisms-section-stabilizer-preserving}

\noindent
文献では、代数スタックの射 $f : \mathcal{X} \to \mathcal{Y}$ により誘導される射
$\mathcal{I}_\mathcal{X} \to
\mathcal{X} \times_\mathcal{Y} \mathcal{I}_\mathcal{Y}$
が同型であるとき、この射は{\it 安定化群を保存する}、または
{\it 固定点を反映する}といわれる。このような射は、$\mathcal{X}$ のすべての点で
自己同型群の間の同型を誘導する（Remark
\ref{stacks-morphisms-remark-identify-automorphism-groups}）。この節では、この概念に関するいくつかの
簡単な補題を証明する。

\begin{lemma}
\label{stacks-morphisms-lemma-stabilizer-preserving}
$f : \mathcal{X} \to \mathcal{Y}$ を代数スタックの射とする。射
$\mathcal{I}_\mathcal{X} \to
\mathcal{X} \times_\mathcal{Y} \mathcal{I}_\mathcal{Y}$ が同型ならば、$f$ は代数空間により表現可能である。
\end{lemma}

\begin{proof}
Lemma \ref{stacks-morphisms-lemma-second-diagonal} から直ちに従う。
\end{proof}

\begin{remark}
\label{stacks-morphisms-remark-get-property-auts-from-diagonal}
$f : \mathcal{X} \to \mathcal{Y}$ を代数スタックの射とする。
$U \to \mathcal{X}$ を、始域が代数空間である射とする。
$G \to H$ を、射
$\mathcal{I}_\mathcal{X} \to
\mathcal{X} \times_\mathcal{Y} \mathcal{I}_\mathcal{Y}$
の $U$ への引き戻しとする。$\Delta_f$ が非分岐、エタールなどであれば、
$G \to H$ もそうである。これは、図式
$$
\xymatrix{
U \times_\mathcal{X} U \ar[r] \ar[d] & \mathcal{X} \ar[d]^{\Delta_f} \\
U \times_\mathcal{Y} U \ar[r] & \mathcal{X} \times_\mathcal{Y} \mathcal{X}
}
$$
が Cartesian であり、射 $G \to H$ が左の垂直矢印を対角射
$U \to U \times U$ によって基底変換したものだからである。
Lemma \ref{stacks-morphisms-lemma-separated-implies-isom} の証明と比較せよ。
\end{remark}

\begin{lemma}
\label{stacks-morphisms-lemma-aut-iso-unramified}
$f : \mathcal{X} \to \mathcal{Y}$ を代数スタックの非分岐射とする。
次は同値である。
\begin{enumerate}
\item $\mathcal{I}_\mathcal{X} \to
\mathcal{X} \times_\mathcal{Y} \mathcal{I}_\mathcal{Y}$
は同型である。
\item すべての $x \in |\mathcal{X}|$ に対して、$f$ は $x$ と $f(x)$ における
自己同型群の間の同型を誘導する
（Remark \ref{stacks-morphisms-remark-identify-automorphism-groups}）。
\end{enumerate}
\end{lemma}

\begin{proof}
スキーム $U$ と全射滑らかな射 $U \to \mathcal{X}$ を選ぶ。
$G \to H$ を、射
$\mathcal{I}_\mathcal{X} \to
\mathcal{X} \times_\mathcal{Y} \mathcal{I}_\mathcal{Y}$
の $U$ への引き戻しとする。Remark
\ref{stacks-morphisms-remark-get-property-auts-from-diagonal} および Lemma
\ref{stacks-morphisms-lemma-characterize-unramified} により、射 $G \to H$ はエタールである。
条件 (1) は、$G \to H$ が同型であるという条件と同値である
（例えば Properties of Stacks, Lemma
\ref{stacks-properties-lemma-check-property-covering} を適用すれば従う）。
条件 (2) は、すべての $u \in U$ に対して、ファイバーの射
$G_u \to H_u$ が同型であるという条件と同値である。
したがって (1) $\Rightarrow$ (2) は明らかである。
(2) が成り立つならば、$G \to H$ は代数空間の全射、普遍単射かつエタールな射である。
このような射は Morphisms of Spaces, Lemma
\ref{spaces-morphisms-lemma-etale-universally-injective-open} により同型である。
\end{proof}

\begin{lemma}
\label{stacks-morphisms-lemma-stabilizer-preserving-unramified}
\begin{reference}
\cite[Proposition 3.5]{rydh_quotients} および
\cite[Proposition 2.5]{alper_quotient}
\end{reference}
$f : \mathcal{X} \to \mathcal{Y}$ を代数スタックの射とする。次を仮定する。
\begin{enumerate}
\item $f$ は代数空間により表現可能かつ非分岐である。
\item $\mathcal{I}_\mathcal{Y} \to \mathcal{Y}$ は固有である。
\end{enumerate}
このとき、$f$ が $x$ と $f(x)$ における自己同型群の間の同型を
誘導するような $x \in |\mathcal{X}|$ の集合は開である
（Remark \ref{stacks-morphisms-remark-identify-automorphism-groups}）。
$\mathcal{U} \subset \mathcal{X}$ を対応する開部分スタックとすると、射
$\mathcal{I}_\mathcal{U} \to
\mathcal{U} \times_\mathcal{Y} \mathcal{I}_\mathcal{Y}$
は同型である。
\end{lemma}

\begin{proof}
スキーム $U$ と全射滑らかな射 $U \to \mathcal{X}$ を選ぶ。
$G \to H$ を、射
$\mathcal{I}_\mathcal{X} \to
\mathcal{X} \times_\mathcal{Y} \mathcal{I}_\mathcal{Y}$
の $U$ への引き戻しとする。Remark
\ref{stacks-morphisms-remark-get-property-auts-from-diagonal} および Lemma
\ref{stacks-morphisms-lemma-characterize-unramified} により、射 $G \to H$ はエタールである。
$f$ は代数空間により表現可能であるから、$G \to H$ は単射である。
したがって $G \to H$ は開埋入である。Morphisms of Spaces, Lemma
\ref{spaces-morphisms-lemma-etale-universally-injective-open} を参照せよ。
仮定により $H \to U$ は固有である。

\medskip\noindent
以上の準備のもとで、次のように補題を証明できる。
補題における $|\mathcal{X}|$ の部分集合の逆像は、$G_u \to H_u$ が同型となる
$u \in U$ の集合にほかならない
（というのも、$G_u$ は $H_u$ の開部分群代数空間だからである）。
これは開部分集合である。実際、その補集合は閉部分集合
$|H| \setminus |G|$ の像であり、$|H| \to |U|$ は閉写像である。
Properties of Stacks, Lemma \ref{stacks-properties-lemma-open-substacks} により、
対応する $\mathcal{X}$ の開部分スタック $\mathcal{U}$ を考えられる。
補題の最後の主張は、$\mathcal{U} \to \mathcal{Y}$ に Lemma
\ref{stacks-morphisms-lemma-aut-iso-unramified} を適用すれば従う。
\end{proof}

\begin{lemma}
\label{stacks-morphisms-lemma-base-change-stabilizer-preserving}
代数スタックの Cartesian 図式
$$
\xymatrix{
\mathcal{X}' \ar[r] \ar[d]_{f'} & \mathcal{X} \ar[d]^f \\
\mathcal{Y}' \ar[r] & \mathcal{Y}
}
$$
\begin{enumerate}
\item $x' \in |\mathcal{X}'|$ の像を $x \in |\mathcal{X}|$ とする。
$f$ が $x$ と $f(x)$ における自己同型群の間の同型を誘導するならば
（Remark \ref{stacks-morphisms-remark-identify-automorphism-groups}）、$f'$ は $x'$ と $f(x')$ における
自己同型群の間の同型を誘導する。
\item $\mathcal{I}_\mathcal{X} \to
\mathcal{X} \times_\mathcal{Y} \mathcal{I}_\mathcal{Y}$ が同型ならば、
$\mathcal{I}_{\mathcal{X}'} \to
\mathcal{X}' \times_{\mathcal{Y}'} \mathcal{I}_{\mathcal{Y}'}$ も同型である。
\end{enumerate}
\end{lemma}

\begin{proof}
省略する。
\end{proof}

\begin{lemma}
\label{stacks-morphisms-lemma-stabilizer-preserving-points-cartesian}
代数スタックの Cartesian 図式
$$
\xymatrix{
\mathcal{X}' \ar[r] \ar[d]_{f'} & \mathcal{X} \ar[d]^f \\
\mathcal{Y}' \ar[r]^g & \mathcal{Y}
}
$$
を考える。$f$ が点における自己同型群の間の同型を誘導するならば
（Remark \ref{stacks-morphisms-remark-identify-automorphism-groups}）、
$$
\Mor(\Spec(k), \mathcal{X}')
\longrightarrow
\Mor(\Spec(k), \mathcal{Y}') \times \Mor(\Spec(k), \mathcal{X})
$$
は任意の体 $k$ に対して同型類上で単射である。
\end{lemma}

\begin{proof}
与えられた $(y', x)$ に写る $x'$ が高々一つであることを示さねばならない。
$2$-ファイバー積の我々の構成によれば、射 $x'$ は三つ組
$(x, y', \alpha)$ で与えられる。ここで
$\alpha : g \circ y' \to f \circ x$ は $2$-射である。
もう一つのこのような三つ組 $(x, y', \beta)$ があると仮定する。
すると $\alpha$ と $\beta$ は、自己同型群代数空間 $G_{f(x)}$ の
$k$-値点 $\epsilon$ だけ異なる。仮定により $f$ は同型
$G_x \to G_{f(x)}$ を誘導するから、$\epsilon$ を $G_x$ の $k$-値点
$\gamma$ に持ち上げられる。このとき $(\gamma, \text{id}) : (x, y', \alpha) \to
(x, y', \beta)$ は望まれた同型である。
\end{proof}

\begin{lemma}
\label{stacks-morphisms-lemma-etale-iso}
$f : \mathcal{X} \to \mathcal{Y}$ を代数スタックの射とする。
$f$ はエタールで、$f$ は点における自己同型群の間の同型を誘導し
（Remark \ref{stacks-morphisms-remark-identify-automorphism-groups}）、かつすべての代数閉体 $k$ に対して関手
$$
f : \Mor(\Spec(k), \mathcal{X}) \longrightarrow \Mor(\Spec(k), \mathcal{Y})
$$
が圏同値であると仮定する。このとき $f$ は同型である。
\end{lemma}

\begin{proof}
Lemma \ref{stacks-morphisms-lemma-universally-injective} により $f$ は普遍単射である。
Lemmas \ref{stacks-morphisms-lemma-stabilizer-preserving} および
\ref{stacks-morphisms-lemma-aut-iso-unramified} を組み合わせると、$f$ が代数空間により
表現可能であることが分かる。したがって Morphisms of Spaces, Lemma
\ref{spaces-morphisms-lemma-etale-universally-injective-open} により $f$ は開埋入である。
最後に、補題の条件から $f$ が全射であることも保証されることに注意すれば、
証明が完了する。
\end{proof}







\section{正規化}
\label{stacks-morphisms-section-normalization}

\noindent
この節は Morphisms of Spaces, Section
\ref{spaces-morphisms-section-normalization} の類似である。

\begin{lemma}
\label{stacks-morphisms-lemma-prepare-normalization}
$\mathcal{X}$ を代数スタックとする。次は同値である。
\begin{enumerate}
\item スキーム $U$ を始域とする全射滑らかな射 $U \to \mathcal{X}$ が存在し、
$U$ のすべての準コンパクト開部分は有限個の既約成分をもつ。
\item すべてのスキーム $U$ と滑らかな射 $U \to \mathcal{X}$ に対し、
$U$ のすべての準コンパクト開部分は有限個の既約成分をもつ。
\item すべての代数空間 $Y$ と滑らかな射 $Y \to \mathcal{X}$ に対し、空間 $Y$ は
Morphisms of Spaces, Lemma \ref{spaces-morphisms-lemma-prepare-normalization}
の同値な条件を満たす。
\item $\mathcal{X}$ 上滑らかなすべての準コンパクト代数スタック $\mathcal{Y}$ に対し、
空間 $|\mathcal{Y}|$ は有限個の既約成分をもつ。
\end{enumerate}
\end{lemma}

\begin{proof}
(1), (2), (3) の同値性は、Descent, Lemma
\ref{descent-lemma-locally-finite-nr-irred-local-fppf}, Properties of Stacks, Lemma
\ref{stacks-properties-lemma-type-property}, および Morphisms of Spaces, Lemma
\ref{spaces-morphisms-lemma-prepare-normalization} から従う。
これらの参照先から、条件 (4) が条件 (1) を含意することも明らかである。
逆に、同値な条件 (1), (2), (3) が成り立つと仮定し、
$\mathcal{Y}$ が準コンパクトであるような代数スタックの滑らかな射
$\mathcal{Y} \to \mathcal{X}$ をとる。Properties of Stacks, Lemma
\ref{stacks-properties-lemma-quasi-compact-stack} により、アフィンスキーム $V$ と
全射滑らかな射 $V \to \mathcal{Y}$ を選べる。
(2) により $V$ は有限個の既約成分をもち、$|V| \to |\mathcal{Y}|$ は全射かつ
連続であるから、Topology, Lemma
\ref{topology-lemma-surjective-continuous-irreducible-components} により、
$|\mathcal{Y}|$ は有限個の既約成分をもつ。
\end{proof}

\begin{lemma}
\label{stacks-morphisms-lemma-normalization}
$\mathcal{X}$ を Lemma \ref{stacks-morphisms-lemma-prepare-normalization} の同値な条件を満たす
代数スタックとする。このとき、代数スタックの整射
$$
\mathcal{X}^\nu \longrightarrow \mathcal{X}
$$
が存在し、すべてのスキーム $U$ と滑らかな射 $U \to \mathcal{X}$ に対して、
ファイバー積 $\mathcal{X}^\nu \times_\mathcal{X} U$ は $U$ の正規化である。
\end{lemma}

\begin{proof}
$U$ がスキームであるような全射滑らかな射 $U \to \mathcal{X}$ をとる。
$R = U \times_\mathcal{X} U$ とおく。代数空間における滑らかな亜群
$(U, R, s, t, c)$ と、$\mathcal{X}$ の表示 $\mathcal{X} = [U/R]$ を得ることを思い出そう。
Algebraic Stacks, Lemmas \ref{algebraic-lemma-map-space-into-stack} および
\ref{algebraic-lemma-stack-presentation} と Definition
\ref{algebraic-definition-presentation} を参照せよ。
$\mathcal{X}$ に対する仮定は、$U$ の正規化 $U^\nu$ が定義されることを意味する。
Morphisms, Definition \ref{morphisms-definition-normalization} を参照せよ。
Morphisms of Spaces, Lemma \ref{spaces-morphisms-lemma-normalization} により、
正規化をとる操作は代数空間の滑らかな射と可換である。
したがって、$R$ の正規化 $R^\nu$ は $R \times_{s, U} U^\nu$ および
$U^\nu \times_{U, t} R$ のどちらとも同型である。よって、代数空間の二つの
滑らかな射 $s^\nu : R^\nu \to U^\nu$ および $t^\nu : R^\nu \to U^\nu$ を得る。
ファイバー積に関する形式的計算から、
$R^\nu \times_{s^\nu, U^\nu, t^\nu} R^\nu$ は $R \times_{s, U, t} R$ の正規化であることが分かる。
したがって、滑らかな射 $c : R \times_{s, U, t} R \to R$ も $c^\nu$ に延長される。
同様に、逆写像 $i : R \to R$（同型である）は同型
$i^\nu : R^\nu \to R^\nu$ を誘導する。最後に、例えば $e$ が $s$ の切断であり、
$R^\nu = R \times_{U, s} U^\nu$ であることから、単位写像 $e : U \to R$ は
$e^\nu : U^\nu \to R^\nu$ に持ち上がる。
$(U^\nu, R^\nu, s^\nu, t^\nu, c^\nu)$ は代数空間における滑らかな亜群であると主張する。
これを示すには、$(U^\nu, R^\nu, s^\nu, t^\nu, c^\nu, e^\nu, i^\nu)$ に対して
Groupoids, Section \ref{groupoids-section-groupoids} の公理 (1), (2)(a), (2)(b), (3)(a), (3)(b) を
確認する必要がある。例えば (1) に対しては、得られる二つの射
$a, b : R^\nu \times_{s^\nu, U^\nu, t^\nu} R^\nu
\times_{s^\nu, U^\nu, t^\nu} R^\nu \to R^\nu$ が一致することを示さねばならない。
これは、対応する射の組 $R \times_{s, U, t} R \times_{s, U, t} R \to R$ が
一致することが分かっており、射 $a$ と $b$ はこの射の正規化への
一意的な延長だから成り立つ。他の公理についても同様である。

\medskip\noindent
代数スタック $\mathcal{X}^\nu = [U^\nu/R^\nu]$ を考える
（Algebraic Stacks, Theorem
\ref{algebraic-theorem-smooth-groupoid-gives-algebraic-stack}）。
代数空間における亜群の射
$(U^\nu, R^\nu, s^\nu, t^\nu, c^\nu) \to (U, R, s, t, c)$ があるから、
代数スタックの射 $\nu : \mathcal{X}^\nu \to \mathcal{X}$ を得る。
$R^\nu = R \times_{s, U} U^\nu$ であるから、Groupoids in Spaces, Lemma
\ref{spaces-groupoids-lemma-criterion-fibre-product} により
$U^\nu = \mathcal{X}^\nu \times_\mathcal{X} U$ であることが分かる。
特に、$U^\nu \to U$ は整であるから、$\nu$ も整である。
補題で述べた基底変換性質が、スキームから $\mathcal{X}$ へのすべての
滑らかな射に対して成り立つことの確認は省略する。
\end{proof}

\noindent
これにより次の定義に至る。

\begin{definition}
\label{stacks-morphisms-definition-normalization}
$\mathcal{X}$ を Lemma \ref{stacks-morphisms-lemma-prepare-normalization} の同値な条件を満たす
代数スタックとする。$\mathcal{X}$ の{\it 正規化}を、射
$$
\nu : \mathcal{X}^\nu \longrightarrow \mathcal{X}
$$
で、Lemma \ref{stacks-morphisms-lemma-normalization} において構成されたものと定義する。
\end{definition}










\section{点と特殊化}
\label{stacks-morphisms-section-specializations}

\noindent
この節は Decent Spaces, Section
\ref{decent-spaces-section-specializations} の類似である。

\begin{lemma}
\label{stacks-morphisms-lemma-specialization}
$\mathcal{X}$ を代数スタックとする。$U$ を代数空間とし、
$f : U \to \mathcal{X}$ を滑らかな射とする。
$x' \leadsto x$ を $|\mathcal{X}|$ の点の特殊化とし、$f(u) = x$ となる
$u \in |U|$ をとる。$(\mathcal{X}, x')$ が Properties of Stacks, Lemma
\ref{stacks-properties-lemma-UR-quasi-compact-above-x} の同値な条件を満たすならば、
$f(u') = x'$ となる $|U|$ 内の特殊化 $u' \leadsto u$ が存在する。
\end{lemma}

\begin{proof}
$U_1$ がアフィンスキームであるようなエタール射
$(U_1, u_1) \to (U, u)$ を選ぶ。そこで $U$ を $U_1$ で置き換えてよいし、実際そうする。
したがって $U$ がアフィンスキームであると仮定してよい。
滑らかな射影 $t, s : R \to U$ をもつ代数空間
$R = U \times_\mathcal{X} U$ を考える。$x'$ に写る点 $w \in U$ を選ぶ。
$f : |U| \to |\mathcal{X}|$ は開写像なので、これは可能である。
$x'$ に対する仮定により、$t : R \to U$ の $w$ 上のファイバー
$F' = t^{-1}(w) = R \times_{t, U} w$ は準コンパクト代数空間である。
アフィンスキーム $T$ と全射エタール射 $T \to F'$ を選ぶ。
$x' \leadsto x$ であることは、$u$ が射
$$
T \to F' \to R \xrightarrow{s} U
$$
の像の閉包に属することを意味する。実際、この像は
$x'$ 上の $|U| \to |\mathcal{X}'|$ のファイバーである。
もしある開部分集合 $u \in V \subset |U|$ がこのファイバーと交わらないならば、
$f(V)$ は $x'$ を含まない $x$ の開近傍となり、矛盾する。
したがって Morphisms, Lemma
\ref{morphisms-lemma-reach-points-scheme-theoretic-image} により、
$x'$ 上の $|U| \to |\mathcal{X}|$ のファイバーの中に、$u$ に特殊化する
$u' \in |U|$ が存在することが分かる。
\end{proof}



\section{適正な代数スタック}
\label{stacks-morphisms-section-reasonable-decent}

\noindent
この節は Decent Spaces, Section
\ref{decent-spaces-section-reasonable-decent} の類似である。
特に、次の定義はそこで定義された適正な代数空間の概念と整合する。

\begin{definition}
\label{stacks-morphisms-definition-decent}
$\mathcal{X}$ を代数スタックとする。すべての $x \in |\mathcal{X}|$ に対して、
Properties of Stacks, Lemma
\ref{stacks-properties-lemma-UR-quasi-compact-above-x} の同値な条件が成り立つとき、
$\mathcal{X}$ は{\it 適正}であるという。
\end{definition}

\noindent
この定義を、$\mathcal{X}$ のすべての点が準コンパクトであると
言い換える人もいるだろう。やや強い条件は、$x$ の同値類に属する任意の射
$\Spec(k) \to \mathcal{X}$ が準コンパクトであるだけでなく準分離でもあることを要求するものである。

\begin{lemma}
\label{stacks-morphisms-lemma-quasi-separated-decent}
準分離代数スタック $\mathcal{X}$ は適正である。
より一般に、$\Delta : \mathcal{X} \to \mathcal{X} \times \mathcal{X}$ が準コンパクトならば、
$\mathcal{X}$ は適正である。
\end{lemma}

\begin{proof}
実際、$\mathcal{X}$ が準分離であるならば、始域が準コンパクトスキーム $T$ である
任意の射 $f : T \to \mathcal{X}$ は準コンパクトである。
Lemma \ref{stacks-morphisms-lemma-quasi-compact-permanence} を参照せよ。
$\Delta$ が準コンパクトであることだけが分かっている場合には、記述
$$
T \times_{f, \mathcal{X}, f'} T' =
(T \times T')
\times_{(f, f'), \mathcal{X} \times \mathcal{X}, \Delta}
\mathcal{X}
$$
を用いて証明する。詳細は省略する。
\end{proof}

\begin{lemma}
\label{stacks-morphisms-lemma-representable-decent}
$f : \mathcal{X} \to \mathcal{Y}$ を代数スタックの射とする。
$Y$ が適正であり、$f$ が（スキームにより）表現可能であるか、または $f$ が
代数空間により表現可能かつ準分離であると仮定する。
このとき $\mathcal{X}$ は適正である。
\end{lemma}

\begin{proof}
$x \in |\mathcal{X}|$ の像を $y \in |\mathcal{Y}|$ とする。
$y$ を定める同値類の中から射 $y : \Spec(k) \to \mathcal{Y}$ を選ぶ。
$\mathcal{X}_y = \Spec(k) \times_{y, \mathcal{Y}} \mathcal{X}$ とおく。
$x$ に写る点 $x' \in |\mathcal{X}_y|$ を選ぶ。Properties of Stacks, Lemma
\ref{stacks-properties-lemma-points-cartesian} を参照せよ。
$x'$ の同値類の中から射 $x' : \Spec(k') \to \mathcal{X}_y$ を選ぶ。図式
$$
\xymatrix{
\Spec(k') \ar[r]_{x'} &
\mathcal{X}_y \ar[r] \ar[d] &
\mathcal{X} \ar[d] \\
&
\Spec(k) \ar[r]^y &
\mathcal{Y}
}
$$
$\mathcal{Y}$ が適正ならば、射 $y$ は準コンパクトである。
したがって、基底変換である $\mathcal{X}_y \to \mathcal{X}$ は準コンパクトである
（Lemma \ref{stacks-morphisms-lemma-base-change-quasi-compact}）。よって結論を得るには、$x'$ が準コンパクトであることを
示せば十分である（Lemma \ref{stacks-morphisms-lemma-composition-quasi-compact}）。
$f$ が表現可能ならば、$\mathcal{X}_y$ はスキームであり、$x'$ は準コンパクトである。
$f$ が代数空間により表現可能かつ準分離ならば、$\mathcal{X}_y$ は
準分離代数空間であり、したがって適正である（Decent Spaces, Lemma
\ref{decent-spaces-lemma-properties-trivial-implications}）。
\end{proof}

\begin{lemma}
\label{stacks-morphisms-lemma-qc-decent}
$f : \mathcal{X} \to \mathcal{Y}$ を代数スタックの射とする。
$f$ が準コンパクトかつ全射で、$\mathcal{X}$ が適正ならば、$\mathcal{Y}$ は適正である。
\end{lemma}

\begin{proof}
$k$ を体とし、$x : \Spec(k) \to \mathcal{X}$ を射として、$y = f \circ x$ と書く。
$f$ は全射であるから、$|\mathcal{Y}|$ のすべての点はこのようにして得られる。
Properties of Stacks, Lemma \ref{stacks-properties-lemma-characterize-surjective} を参照せよ。
アフィンスキーム $T$ と射 $T \to \mathcal{Y}$ を考える。
$T \times_{\mathcal{Y}, y} \Spec(k)$ が準コンパクトであることを示せば十分である。
Lemma \ref{stacks-morphisms-lemma-check-quasi-compact-covering} を参照せよ。等式
$$
(T \times_{\mathcal{Y}} \mathcal{X}) \times_{\mathcal{X}, x} \Spec(k) =
T \times_{\mathcal{Y}, y} \Spec(k)
$$
が成り立つ。射 $T \times_{\mathcal{Y}} \mathcal{X} \to T$ は準コンパクトだから、
$T \times_\mathcal{Y} \mathcal{X}$ は準コンパクトである。$\mathcal{X}$ は適正なので、$x$ は
準コンパクトな射である。したがって、表示したファイバー積は準コンパクトである。
\end{proof}

\begin{lemma}
\label{stacks-morphisms-lemma-gerbe-decent}
$f : \mathcal{X} \to \mathcal{Y}$ を代数スタックの射とする。
$\mathcal{X}$ が $\mathcal{Y}$ 上のガーベであり、$\mathcal{X}$ が適正ならば、$\mathcal{Y}$ は適正である。
\end{lemma}

\begin{proof}
$\mathcal{X}$ が $\mathcal{Y}$ 上のガーベであり、$\mathcal{X}$ が適正であると仮定する。
Lemma \ref{stacks-morphisms-lemma-gerbe-bijection-points} により $f$ は普遍同相写像であることに注意せよ。
したがって補題は Lemma \ref{stacks-morphisms-lemma-qc-decent} から従う。
\end{proof}



\section{適正なスタック上の点}
\label{stacks-morphisms-section-points-decent}

\noindent
この節は Decent Spaces, Section \ref{decent-spaces-section-points} の類似である。
適正な代数スタックに付随する位相空間が常に sober であるかどうかは分からない。
やや弱い結果について Proposition \ref{stacks-morphisms-proposition-decent-sober} を参照せよ。

\begin{lemma}
\label{stacks-morphisms-lemma-kolmogorov}
$\mathcal{X}$ を適正な代数スタックとする。このとき $|\mathcal{X}|$ は Kolmogorov である
（Topology, Definition \ref{topology-definition-generic-point} を参照せよ）。
\end{lemma}

\begin{proof}
$x_1 \leadsto x_2$ および $x_2 \leadsto x_1$ を満たす $x_1, x_2 \in |\mathcal{X}|$ をとる。
$x_1 = x_2$ を示さねばならない。$|\mathcal{Z}|$ が
$\overline{\{x_1\}} = \overline{\{x_2\}}$ に等しいような既約閉部分スタック
$\mathcal{Z} \subset \mathcal{X}$ をとる。Lemma \ref{stacks-morphisms-lemma-representable-decent} により
$\mathcal{Z}$ は適正である。$\mathcal{X}$ を $\mathcal{Z}$ で置き換えることにより、
次の段落で論じる場合に帰着する。

\medskip\noindent
$|\mathcal{X}|$ が既約で、$x_1$ と $x_2$ が一般点であると仮定する。
アフィンスキーム $U$、$u_1, u_2 \in U$ および、$f(u_i) = x_i$ を満たす滑らかな射
$f : U \to \mathcal{X}$ を選ぶ。すると、$U$ のある既約成分の一般点である第三の点
$u_3 \in U$ で、その像 $x_3 \in |\mathcal{X}|$ {\it もまた} $|\mathcal{X}|$ の一般点であるものを見つけられる。
実際、$u_1$ を通る既約成分（望むなら $u_2$ を通るもの）の任意の一般点を
$u_3$ として選べばよい。次の段落で $x_1 = x_3$ および $x_2 = x_3$ を示し、
それによって望む結果を証明する。

\medskip\noindent
対称性により、$x_1 = x_3$ を証明すれば十分である。$x_1$ は $|\mathcal{X}|$ の一般点なので、
特殊化 $x_1 \leadsto x_3$ がある。Lemma \ref{stacks-morphisms-lemma-specialization} により、
$x_1 \leadsto x_3$ に写る $U$ 内の特殊化 $u'_1 \leadsto u_3$（！）を見つけられる。
しかし $u_3$ は既約成分の一般点だから、望まれたとおり $u'_1 = u_3$ である。
\end{proof}

\begin{lemma}
\label{stacks-morphisms-lemma-decent-locally-noetherian-sober}
$\mathcal{X}$ を適正で局所 Noether な代数スタックとする。
このとき $|\mathcal{X}|$ は sober かつ局所 Noether な位相空間である。
\end{lemma}

\begin{proof}
Lemma \ref{stacks-morphisms-lemma-Noetherian-topology} により、位相空間 $|\mathcal{X}|$ は局所 Noether である。
Lemma \ref{stacks-morphisms-lemma-kolmogorov} により、位相空間 $|\mathcal{X}|$ は Kolmogorov である。
Lemma \ref{stacks-morphisms-lemma-locally-Noetherian-quasi-sober} により、位相空間 $|\mathcal{X}|$ は準 sober である。
これで証明が完了する。Topology, Definition
\ref{topology-definition-generic-point} を参照せよ。
\end{proof}

\begin{proposition}
\label{stacks-morphisms-proposition-decent-sober}
$\mathcal{X}$ を適正な代数スタックとし、
$\mathcal{I}_\mathcal{X} \to \mathcal{X}$ が準コンパクトであると仮定する。
このとき $|\mathcal{X}|$ は sober である。
\end{proposition}

\begin{proof}
Lemma \ref{stacks-morphisms-lemma-kolmogorov} により $|\mathcal{X}|$ が Kolmogorov であることが分かる
（実際には、これを改めて証明する）。$T \subset |\mathcal{X}|$ を既約閉部分集合とする。
$T$ が一般点をもつことを示さねばならない。$\mathcal{Z} \subset \mathcal{X}$ を、$T$ に対応する
誘導される既約閉部分スタックとする。Properties of Stacks, Definition
\ref{stacks-properties-definition-reduced-induced-stack} を参照せよ。
$\mathcal{Z} \to \mathcal{X}$ は閉埋入なので、$\mathcal{Z}$ は適正な代数スタックである。
Lemma \ref{stacks-morphisms-lemma-representable-decent} を参照せよ。また、射
$\mathcal{I}_\mathcal{Z} \to \mathcal{Z}$ は $\mathcal{I}_\mathcal{X} \to \mathcal{X}$ の基底変換である
（Lemma \ref{stacks-morphisms-lemma-monomorphism-cartesian-square-inertia}）。したがって
$\mathcal{I}_\mathcal{Z} \to \mathcal{Z}$ は準コンパクトである
（Lemma \ref{stacks-morphisms-lemma-base-change-quasi-compact}）。よって、次の段落で論じる場合に帰着する。

\medskip\noindent
$\mathcal{X}$ が適正で、$\mathcal{I}_\mathcal{X} \to \mathcal{X}$ が準コンパクトで、
$\mathcal{X}$ が既約で、$|\mathcal{X}|$ が既約であると仮定する。
$|\mathcal{X}|$ が一般点をもつことを示さねばならない。
Proposition \ref{stacks-morphisms-proposition-open-stratum} により、ガーベである稠密開部分スタック
$\mathcal{U} \subset \mathcal{X}$ が存在する。言い換えれば、$|\mathcal{U}| \subset |\mathcal{X}|$ は開かつ稠密である。
したがって、他のすべての性質に加えて $\mathcal{X}$ がガーベであると仮定してよい。
$\mathcal{X} \to X$ によって $\mathcal{X}$ が代数空間 $X$ 上のガーベになるとしよう。
すると Lemma \ref{stacks-morphisms-lemma-gerbe-bijection-points} により
$|\mathcal{X}| \cong |X|$ である。特に $X$ は準コンパクトで、$|X|$ は既約である。
また Lemma \ref{stacks-morphisms-lemma-gerbe-decent} により、$X$ は適正な代数空間である。
したがって Decent Spaces, Proposition
\ref{decent-spaces-proposition-reasonable-sober} により $|\mathcal{X}| = |X|$ は sober であり、
よって（一意的な）一般点をもつ。
\end{proof}






\section{整代数スタック}
\label{stacks-morphisms-section-integral-stacks}

\noindent
この節は Spaces over Fields, Section
\ref{spaces-over-fields-section-integral-spaces} の類似である。
その節での考察と Proposition \ref{stacks-morphisms-proposition-decent-sober} の結果に動機づけられ、
整代数スタックを次のように定義する
（これは、すでに存在する整スキームと整代数空間の定義と抵触しない）。

\begin{definition}
\label{stacks-morphisms-definition-integral-algebraic-stack}
代数スタック $\mathcal{X}$ が既約かつ適正であり、
$\mathcal{I}_\mathcal{X} \to \mathcal{X}$ が準コンパクトで、$|\mathcal{X}|$ が既約であるとき、
このとき{\it 整}であるという。
\end{definition}

\noindent
$\mathcal{X}$ が準分離ならば、それが整であるためには、$\mathcal{X}$ が既約で
$|\mathcal{X}|$ が既約であれば十分であることに注意せよ。
Lemma \ref{stacks-morphisms-lemma-totaro-case} を参照せよ。

\begin{lemma}
\label{stacks-morphisms-lemma-structure-integral-stack}
$\mathcal{X}$ を整代数スタックとする。このとき次が成り立つ。
\begin{enumerate}
\item $|\mathcal{X}|$ は sober かつ既約で、一意的な一般点をもつ。
\item ある整スキーム $U$ 上のガーベである開部分スタック
$\mathcal{U} \subset \mathcal{X}$ が存在する。
\end{enumerate}
\end{lemma}

\begin{proof}
Proposition \ref{stacks-morphisms-proposition-decent-sober} により $|\mathcal{X}|$ は sober である。
もちろんこれは既約でもあり、したがって（sober 性の定義により）一意的な一般点 $x$ をもつ。
Proposition \ref{stacks-morphisms-proposition-open-stratum} により、代数空間 $U$ 上のガーベである
稠密開部分 $\mathcal{U} \subset \mathcal{X}$ が存在する。Lemma \ref{stacks-morphisms-lemma-gerbe-decent} により
$U$ は適正な代数空間である
（Lemma \ref{stacks-morphisms-lemma-representable-decent} により $\mathcal{U}$ が適正であることも用いる）。
$|U| = |\mathcal{U}|$ なので、$|U|$ は既約である。
最後に、$\mathcal{U}$ は既約なので、射 $\mathcal{U} \to U$ は $U_{red}$ を経由する。
Properties of Stacks, Lemma \ref{stacks-properties-lemma-map-into-reduction} を参照せよ。
さて、$\mathcal{U} \to U$ は平坦、局所有限表示かつ全射であるから
（Lemma \ref{stacks-morphisms-lemma-gerbe-fppf}）、これから $U = U_{red}$、すなわち $U$ が既約であることが従う
（小さな詳細は省略する）。よって $U$ は整代数空間である。Spaces over Fields,
Definition \ref{spaces-over-fields-definition-integral-algebraic-space} を参照せよ。
最後に $U$（それに応じて $\mathcal{U}$）を開部分空間で置き換え、$U$ が整スキームであると
仮定してよい。Spaces over Fields, Section
\ref{spaces-over-fields-section-integral-spaces} の議論を参照せよ。
\end{proof}

\begin{lemma}
\label{stacks-morphisms-lemma-totaro-case}
$\mathcal{X}$ を既約かつ準分離な代数スタックとし、その付随位相空間
$|\mathcal{X}|$ が既約であるとする。このとき $\mathcal{X}$ は整である。
\end{lemma}

\begin{proof}
$\mathcal{X}$ が準分離ならば、Lemma \ref{stacks-morphisms-lemma-quasi-separated-decent} により
$\mathcal{X}$ は適正である。$\mathcal{X}$ が準分離ならば、
$\Delta : \mathcal{X} \to \mathcal{X} \times \mathcal{X}$ は準コンパクトである。
したがって、$\Delta$ を $\Delta$ によって基底変換したものとして
$\mathcal{I}_\mathcal{X} \to \mathcal{X}$ は準コンパクトである。
Lemma \ref{stacks-morphisms-lemma-base-change-quasi-compact} を参照せよ。
よって Definition \ref{stacks-morphisms-definition-integral-algebraic-stack} のすべての仮定が成り立つ
（また、「準分離」を「$\Delta_\mathcal{X}$ が準コンパクト」で置き換えてよいことも分かる）。
\end{proof}

\begin{lemma}
\label{stacks-morphisms-lemma-decent-irreducible-closed}
$\mathcal{X}$ を適正な代数スタックとし、
$\mathcal{I}_\mathcal{X} \to \mathcal{X}$ が準コンパクトであると仮定する。
次の集合の間に標準的な全単射がある。
\begin{enumerate}
\item $\mathcal{X}$ の点の集合、すなわち $|\mathcal{X}|$。
\item $|\mathcal{X}|$ の既約閉部分集合の集合。
\item $\mathcal{X}$ の整閉部分スタックの集合。
\end{enumerate}
(1) から (2) への全単射は $x$ を $\overline{\{x\}}$ に送る。
(3) から (2) への全単射は $\mathcal{Z}$ を $|\mathcal{Z}|$ に送る。
\end{lemma}

\begin{proof}
Proposition \ref{stacks-morphisms-proposition-decent-sober} により $|\mathcal{X}|$ は sober であるから、
上の写像は (1) と (2) の間の全単射を定める。
$T \subset |\mathcal{X}|$ が閉かつ既約ならば、$|\mathcal{Z}| = T$ となる一意的な
既約閉部分スタック $\mathcal{Z} \subset \mathcal{X}$ が存在する。すなわち、$\mathcal{Z}$ は
$T$ 上の誘導される既約部分空間構造である。Properties of Stacks, Definition
\ref{stacks-properties-definition-reduced-induced-stack} を参照せよ。
このとき $\mathcal{Z}$ は整代数スタックである。実際、これは適正であり
（Lemma \ref{stacks-morphisms-lemma-representable-decent}）、射
$\mathcal{I}_\mathcal{Z} \to \mathcal{Z}$ は準コンパクトである
（$\mathcal{I}_\mathcal{X} \to \mathcal{X}$ の基底変換であることによる。Lemma
\ref{stacks-morphisms-lemma-monomorphism-cartesian-square-inertia} を参照せよ）。また $\mathcal{Z}$ は既約で、
$|\mathcal{Z}|$ は既約である。
\end{proof}










\section{剰余ガーベ}
\label{stacks-morphisms-section-residual-gerbe}

\noindent
この節は Properties of Stacks, Section
\ref{stacks-properties-section-residual-gerbe} の続きである。

\begin{lemma}
\label{stacks-morphisms-lemma-gerbe-residual-gerbe}
$\pi : \mathcal{X} \to \mathcal{Y}$ を代数スタックの射とする。
$x \in |\mathcal{X}|$ の像を $y \in |\mathcal{Y}|$ とする。
$y$ における $\mathcal{Y}$ の剰余ガーベ $\mathcal{Z}_y \subset \mathcal{Y}$ が存在し、
$\mathcal{X}$ が $\mathcal{Y}$ 上のガーベであると仮定する。このとき
$\mathcal{Z}_x = \mathcal{Z}_y \times_\mathcal{Y} \mathcal{X}$ は $x$ における $\mathcal{X}$ の剰余ガーベである。
\end{lemma}

\begin{proof}
射 $\mathcal{Z}_x \to \mathcal{X}$ は、単射 $\mathcal{Z}_y \to \mathcal{Y}$ の基底変換として単射である。
Lemma \ref{stacks-morphisms-lemma-gerbe-bijection-points} により射 $\pi$ は普遍同相写像であり、
したがって $|\mathcal{Z}_x| = \{x\}$ である。最後に、射
$\mathcal{Z}_x \to \mathcal{Z}_y$ は滑らかな射 $\pi$ の基底変換として滑らかである。
Lemma \ref{stacks-morphisms-lemma-gerbe-smooth} を参照せよ。よって $\mathcal{Z}_y$ が既約かつ局所 Noether であることから、
$\mathcal{Z}_x$ もそうである（詳細は省略する）。したがって Properties of Stacks, Definition
\ref{stacks-properties-definition-residual-gerbe} により、$\mathcal{Z}_x$ は $x$ における
$\mathcal{X}$ の剰余ガーベである。
\end{proof}

\begin{lemma}
\label{stacks-morphisms-lemma-factor-through-residual-space}
$f : \mathcal{Y} \to \mathcal{X}$ を代数スタックの射とし、$x \in |\mathcal{X}|$ を点とする。
次を仮定する。
\begin{enumerate}
\item $\mathcal{X}$ は適正または局所 Noether である（またはその両方）。
\item $\mathcal{I}_\mathcal{X} \to \mathcal{X}$ は準コンパクトである。
\item $|f|(|\mathcal{Y}|)$ は $\{x\} \subset |\mathcal{X}|$ に含まれる。
\item $\mathcal{Y}$ は既約である。
\end{enumerate}
このとき $f$ は、$x$ における $\mathcal{X}$ の剰余ガーベ $\mathcal{Z}_x$ を経由する
（その存在は Lemma \ref{stacks-morphisms-lemma-every-point-residual-gerbe} または
\ref{stacks-morphisms-lemma-every-point-residual-gerbe-locally-Noetherian} により保証される）。
\end{lemma}

\begin{proof}
$T = \overline{\{x\}} \subset |\mathcal{X}|$ を $x$ の閉包とする。Properties of Stacks, Lemma
\ref{stacks-properties-lemma-reduced-closed-substack} により、
$T = |\mathcal{X}'|$ となる既約閉部分スタック $\mathcal{X}' \subset \mathcal{X}$ が存在する。
Properties of Stacks, Lemma \ref{stacks-properties-lemma-map-into-reduction} により、
射 $f$ は $\mathcal{X}'$ を経由する。$\mathcal{X}$ が適正ならば、Lemma
\ref{stacks-morphisms-lemma-representable-decent} によりスタック $\mathcal{X}'$ は適正である。
$\mathcal{X}$ が局所 Noether ならば、$\mathcal{X}'$ は局所 Noether である
（詳細は省略する）。Lemma \ref{stacks-morphisms-lemma-monomorphism-cartesian-square-inertia} により、
$\mathcal{I}_{\mathcal{X}'} \to \mathcal{X}'$ は $\mathcal{I}_\mathcal{X} \to \mathcal{X}$ の基底変換であることに注意せよ。
Lemma \ref{stacks-morphisms-lemma-base-change-quasi-compact} により
$\mathcal{I}_{\mathcal{X}'} \to \mathcal{X}'$ は準コンパクトである。
これにより次の段落で論じる場合に帰着する。

\medskip\noindent
$\mathcal{X}$ が既約で、$x \in |\mathcal{X}|$ が一般点であると仮定する。
Proposition \ref{stacks-morphisms-proposition-open-stratum} により、ガーベである稠密開部分スタック
$\mathcal{U} \subset \mathcal{X}'$ が存在することが従う。$x \in |\mathcal{U}|$ であることに注意せよ。
上の議論を繰り返すことにより、次の段落で論じる場合に帰着する。

\medskip\noindent
$\mathcal{X} \to X$ によってこの代数スタックが代数空間 $X$ 上のガーベになると仮定する。
$\mathcal{X}$ が適正ならば、Lemmas \ref{stacks-morphisms-lemma-gerbe-bijection-points} および
\ref{stacks-morphisms-lemma-qc-decent} により、空間 $X$ は適正である。$\mathcal{X}$ が局所 Noether ならば、
fppf 降下により $X$ は局所 Noether である（詳細は省略する）。
したがって対応する結果が $X$ に対して成り立つ。Decent Spaces, Lemma
\ref{decent-spaces-lemma-factor-through-residual-space-decent} または
\ref{decent-spaces-lemma-factor-through-residual-space-Noetherian} を参照せよ
（小さな詳細は省略する）。Lemma \ref{stacks-morphisms-lemma-gerbe-residual-gerbe} を適用すると、
この結果が $\mathcal{X}$ に対しても成り立つことが従う。
\end{proof}

\begin{remark}
\label{stacks-morphisms-lemma-factor-through-residual-space-Noetherian}
$\mathcal{X}$ が局所 Noether であることのみを仮定する場合、すなわち慣性が準コンパクトであるという
仮定を外す場合に、Lemma \ref{stacks-morphisms-lemma-factor-through-residual-space} が成り立つかどうかは分からない。
この場合でも、$x$ が閉点ならば、次のはるかに容易な補題からそれが真であることは
確実である。
\end{remark}

\begin{lemma}
\label{stacks-morphisms-lemma-residual-space-closed}
$\mathcal{X}$ を局所 Noether 代数スタックとする。$x \in |\mathcal{X}|$ を点とし、
$\mathcal{Z}_x \subset \mathcal{X}$ をその剰余ガーベとする
（Lemma \ref{stacks-morphisms-lemma-every-point-residual-gerbe-locally-Noetherian}）。
このとき、$x$ が $|\mathcal{X}|$ の閉点であることと、射
$\mathcal{Z}_x \to \mathcal{X}$ が閉埋入であることは同値である。
\end{lemma}

\begin{proof}
$\mathcal{Z}_x \to \mathcal{X}$ が閉埋入ならば、$x$ は $|\mathcal{X}|$ の閉点である。
例えば Lemma \ref{stacks-morphisms-lemma-closed-immersion-proper} を参照せよ。
逆に、$x$ が $|\mathcal{X}|$ の閉点であると仮定する。$\mathcal{Z} \subset \mathcal{X}$ を、
$|Z| = \{x\}$ となる既約閉部分スタックとする（Properties of Stacks, Lemma
\ref{stacks-properties-lemma-reduced-closed-substack}）。
Lemmas \ref{stacks-morphisms-lemma-immersion-locally-finite-type} および
\ref{stacks-morphisms-lemma-locally-finite-type-locally-noetherian} により、$\mathcal{Z}$ は局所 Noether 代数スタックである。
また $\mathcal{Z}$ は既約で $|\mathcal{Z}| = \{x\}$ であるから、定義により
$\mathcal{Z} = \mathcal{Z}_x$ が剰余ガーベであることが従う。
\end{proof}
